\section{E. The fourfold and transverse fibre products}

In this section $X=X_c$ has the canonical sphere twists
$v_1=(1,2,-4,0)^t$, $v_2=(-1,-3,3,0)^t$, the untwisted cusp collar,
and an arbitrary admissible parameter $c$.  Orient all complex manifolds by
their complex structures.  Cohomology is integral unless a coefficient
field is displayed.  Put
\[
 L=\mathcal O_X(S_1-S_2).
\]
The divisor comparisons $3S_1\sim4S_2\sim W$ give
$L^4=\mathcal O_X(S_1)$, $L^6=\mathcal O_X(2S_2)$ and
$L^{12}=\mathcal O_X(W)$.  The comparisons use the meromorphic functions
$t$, $t-1$ and their quotient; thus these are comparisons of actual
holomorphic lines, not just of Chern classes.

\paragraph{Proposition E.1.1 (Refuted: the normalization in Lead 6).}
The local space $(F\times_{\widetilde\Delta}E)/\mu_m$ in Lead 6 is not
the normalization of the literal fibre product $X\times_B S$ at an
additive fibre.  It is a different proper model of its punctured family.

\begin{proof}
The reduced fibres of $S$ of types IV* and III have respectively seven
and two irreducible components.  In the literal fibre product the inverse
image of each such component contains an irreducible component which
dominates it: at a general smooth point of the component, the base-change
equation is $s^m=u^a$, times smooth factors.  Each branch of its
normalization still dominates that component.  A finite normalization
cannot contract any of these components.  Consequently the normalized
central fibre has at least seven, respectively two, components.  In the
proposed replacement the reduced central fibre is
$(F_0\times E_0)/\mu_m$, which is irreducible.  The component counts
disprove the asserted identification.  Good reduction of the elliptic
surface uses contraction as well as normalization; omitting that
contraction is the exact error in the reader's argument.
\end{proof}

\paragraph{Theorem E.1.2 (Proved: a corrected compact construction).}
There is a smooth compact complex fourfold $Z$ with proper holomorphic
maps $g:Z\to B$ and $\Pi:Z\to X$, with $g=f\Pi$, satisfying the local
models and the marked punctured-fibre product in Lead 6.  The map $\Pi$
has a holomorphic section.  The construction is a small crepant
resolution of a Weierstrass model over $X$, and is bimeromorphic to
$X\times_B S$.

\begin{proof}
Let $a$ be a nonzero scalar multiple of the divisor section of $L^4$
and let $b$ be a nonzero scalar multiple of the square of the divisor
section of $\mathcal O_X(S_2)=L^3$.  Choose the scalars so that
$4a^3+27b^2$ is the pullback of the section of $\mathcal O_B(1)$
vanishing at $p_0$.  This is possible because $a^3$ and $b^2$ have
divisors $3S_1$ and $4S_2$, which correspond to a basis of
$H^0(B,\mathcal O_B(1))$.  The Weierstrass equation
\[
 y^2z=x^3+axz^2+bz^3
 \tag{E.1}
\]
in the projective bundle associated to
$\mathcal O_X\oplus L^2\oplus L^3$ defines a proper complex space
$\mathscr W\to X$.  Its section at infinity is disjoint from its
singular locus.  Its $j$-function is $1728t$.

Over $B^\circ$, equation (E.1) is the product with Engel's elliptic
curve: both have that marked Weierstrass family.  At $p_1$, pass to the
actual root atlas $t=s^3$.  The coefficient $a$ has a simple zero in $s$
and $b$ is a unit.  At $p_2$, on $t-1=s^4$, $a$ is a unit and $b$ has
order two.  These are the smooth good-reduction families $E$.
The deck action on $F\times E$ is diagonal, with the specified
$A_j\oplus C_j$ linear part and translation $(v_j,0)/m_j$.
It is free because its projection to $F$ is free.  Its quotient is
therefore smooth and is exactly the local model asserted in the theorem.
The section at infinity is invariant.  Notice that the full fibres here
have multiplicities three and four; only their reductions are smooth.

At the cusp, use the angular--deck marking of the supplied toroidal
bridge, and write the two-torus period as
$Z_2=sB_0+C(t_c)$, where
\[
 B_0=\begin{pmatrix}0&1\\-1&0\end{pmatrix},\quad
 P=\begin{pmatrix}1&-1&0\\0&1&-1\end{pmatrix},\quad
 Q=B_0^{-1}P,\quad
 S=\begin{pmatrix}1&1\\0&1\\0&0\end{pmatrix}.
\]
Here $PS=I_2$.  Put $k=(1,1,1)^t$, $r=(0,0,1)$, and write the Tate
period of the elliptic surface as $s+c_E(t_c)$, with $c_E$ holomorphic
and bounded.  Define
\[
 Z_3=sI_3+C_3(t_c),\qquad
 C_3=SCQ+kc_Er.
 \tag{E.2}
\]
For $A=\binom{P}{r}$ and $B=\binom{Q}{r}$ we have
\[
 \det A=\det B=1,\qquad
 AZ_3B^{-1}=\operatorname{diag}(Z_2,s+c_E).
 \tag{E.3}
\]
This proves an integral, holomorphic product identification on the
punctured overlap, including the bounded correction which the untwisted
Tate curve alone would miss.

Take the fan over the periodic Kuhn tetrahedra
\[
 [m,m+e_{\pi_1},m+e_{\pi_1}+e_{\pi_2},m+(1,1,1)],
 \quad m\in\mathbb Z^3,\quad\pi\in S_3.
 \tag{E.4}
\]
Every cone is unimodular, since subtracting its first vertex reduces
its determinant to that of a permuted triangular integral matrix with
diagonal entries one.  Every maximal chart is therefore $\mathbb C^4$,
with $t_c=z_0z_1z_2z_3$.  The action
\[
 (x,t_c)\longmapsto
 \bigl(e^{2\pi iC_3(t_c)\lambda}t_c^\lambda x,t_c\bigr),
 \qquad\lambda\in\mathbb Z^3,
 \tag{E.5}
\]
extends to the fan by translation of its cells.  On the boundary a
fixed point would require a bounded cell to equal its nonzero translate.
On a punctured fibre, invertibility of $\operatorname{Im}Z_3$ excludes
a fixed point.  If $y=\operatorname{Im}s$ and
$K=\sup\|\operatorname{Im}C_3\|$, shrinking the disc gives
\[
 \tfrac12\|\lambda\|\leq
 \|(I+y^{-1}\operatorname{Im}C_3)\lambda\|
 \leq\tfrac32\|\lambda\|.
\]
This proves proper discontinuity by tropical coordinates.  Rounding
the three coordinates of
$(I+y^{-1}\operatorname{Im}C_3)^{-1}\operatorname{trop}(x)$ puts every
orbit in finitely many toric charts with bounded coordinates.  Over a
smaller closed disc their closures are compact.  Thus the quotient
$N_3^C$ is smooth and proper over that disc.

The map $(P,r)$ takes (E.4) to a subdivision of the product of an
$A_2$ triangle and an interval.  There are three tetrahedra in each
triangular prism.  No ray is added.  Consequently the induced map is
a small toric resolution of the local equation
$xy=w_0w_1w_2$ in (E.1); over a double curve it resolves
$xy=w_0w_1$ times a smooth factor.  All rays have height one, so the
resolution is crepant.  The formula $PZ_3=Z_2Q$ proves equivariance
of its map to the cusp of $X$.  The slice $x_3=1$ gives the zero
section: its lattice map is $S$, and $S$ takes each $A_2$ triangle
to a face of a Kuhn tetrahedron.  Thus the section extends over every
double curve and both triple points, rather than only over the
punctured cusp.

We glue $N_3^C$ and the two finite models to the punctured product by
(E.3) and their root-atlas identifications.  These overlap maps are
holomorphic isomorphisms and satisfy the cocycle condition because
they are identifications with the same marked product.  Properness
over the three closed discs and compactness over the remaining compact
base give compactness of the glued Hausdorff space.  Smoothness has
been checked in all its charts.  The maps to $X$ and the zero section
agree on their overlaps.  Since both this model and the literal
fibre product agree on a dense open set, they are bimeromorphic.
\end{proof}

\paragraph{Theorem E.1.3 (Proved: the integral cusp calculation).}
For the central fibre $W_3$ of (E.5),
\[
 H_q(W_3;\mathbb Z)=\mathbb Z^{b_q},\qquad
 (b_0,\ldots,b_6)=(1,3,9,9,10,3,1).
 \tag{E.6}
\]
Integral specialization is onto the monodromy invariant lattice in
every degree.  Its only kernel is a copy of $\mathbb Z$ in degree four.

\begin{proof}
The normalization is the three-dimensional permutohedral toric variety.
Its paired boundary surfaces consist of four copies of $\mathrm{dP}_6$
and three copies of $\mathbb P^1\times\mathbb P^1$.  There are twelve
normalized triple curves and six quadruple points.  The normalization
complex for the constant sheaf is the alternating branch complex.
Local exactness follows from the augmented simplex on the branches
of $z_0z_1z_2z_3=0$.  Its nonzero rows after taking cohomology are
\[
\begin{array}{c|rrrr|rr}
q& p=0&p=1&p=2&p=3&\operatorname{rank}d^0&
 \operatorname{rank}d^1\ \text{or }d^2\\ \hline
0&1&7&12&6&0&4,5\\
2&11&22&12&0&5&9\\
4&11&7&0&0&4&0\\
6&1&0&0&0&0&0
\end{array}
\tag{E.7}
\]
Every nonzero Smith invariant in these matrices is one.

Here is a reproducible specification of the integral matrices, in
addition to their complete arrays in the accompanying code output.
The fourteen rays are $\pm v$, with
$v\in\{0,1\}^3\setminus\{0\}$.  The twenty-four maximal cones
are the translated Kuhn tetrahedra containing zero.  This enumeration
requires exactly the eight cubes with lower vertex in $\{-1,0\}^3$;
no omitted cube can contain zero.  The Picard group of the normalization
is the group of ray coefficients modulo the three character relations.
We set the coefficients at $e_1,e_2,e_3$ to zero, giving an integral
basis of rank eleven.  On a boundary surface $D_v$, choose adjacent
rays $u_1,u_2$ for which $(v,u_1,u_2)$ is unimodular and set the two
remaining character coefficients to zero.  Restriction is evaluated
by subtracting the unique integral linear function with the prescribed
values at these three rays.  The paired restriction is
$\operatorname{res}_{D_v}-\operatorname{res}_{D_{-v}}$, with the
second adjacent ray corresponding to $u-v$.

For a triangle $(0,v,u)$ with $0<v<u$ coordinatewise, the next map
is the alternating sum of the degrees on the three boundary curves
of $D_{u-v},D_u,D_v$.  In a surface fan, if the neighbours of a ray
$w$ are $h_1,h_2$ and $h_1+h_2=bw$ modulo $v$, that degree is
$a_{h_1}+a_{h_2}-ba_w$.  In degree four, the restriction map is dual
to the seven divisor differences $[D_v]-[D_{-v}]$.  In degree zero
the matrices are the cellular incidence matrices of the one-vertex
Kuhn triangulation of $T^3$.  These rules define (E.7) over
$\mathbb Z$, and the program checks both compositions are zero.
Its exact Smith reductions give the ranks and primitivity just stated.

For completeness, the possible higher differential of this
normalization spectral sequence vanishes integrally.  Replace $C_3$
by $hC_3$, $0\leq h\leq1$.  The preceding uniform properness estimate
applies simultaneously.  Local parameter vector fields preserving
$t_c$ and each coordinate stratum patch by a partition of unity;
their flow identifies the central fibres and their specialization
maps.  We may therefore use $C_3=0$.  This toroidal degeneration is
projective after passing to the finite cover with deck lattice
$2\mathbb Z^3$.  An explicit coherent polarization is supplied by
\[
 Q(n)=4\sum n_i^2-(\sum n_i)^2.
\]
For the tetrahedron ordered $1,2,3$, its affine interpolant is
$3n_1+n_2-n_3$, and their difference is
\[
 \sum_i n_i(n_i-1)+
 \sum_{i<j}(n_i-n_j)(n_i-n_j-1)\geq0
 \quad(n\in\mathbb Z^3).
\]
Equality holds exactly at its four vertices.  Translation and
permutation give strict coherence on every face.  The finite cover
has simple normal crossings, and its descent complex is a complex
of pure toric cohomology groups.  Over $\mathbb Q$ the higher
differentials vanish by their different weights.  Taking finite
deck invariants returns the original fibre.  The integral $E_2$
groups of (E.7) are free, so a higher differential which is zero
over $\mathbb Q$ is zero integrally.  Summing its diagonals proves
(E.6), with no extension torsion.

The same projective comparison allows the local invariant cycle
theorem, in the form of \cite[Theorem B]{ffChen}, followed by exact
finite-group invariants over $\mathbb Q$.  The cusp monodromy on
$H_1$ is $\left(\begin{smallmatrix}I_3&I_3\\0&I_3\end{smallmatrix}\right)$.
Its exterior-power invariant dimensions are
$(1,3,9,9,9,3,1)$, as checked by exact exterior matrices in the
attached program.  Comparison with (E.6) proves the last assertion.
We now strengthen this to an integral statement.  Write $b_i$ for the
three degree-one invariant deck classes and $a_i$ for the angular
classes, so the infinitesimal monodromy sends $a_i$ to $b_i$.
At $C_3=0$, every integral symmetric matrix $S$ gives a line bundle
on the toroidal family.  Here is its descent data, so no integral
invariant-cycle theorem is being assumed.  Assign to a fan vertex $n$
the integer
\[
 q_S(n)=\frac{n^tSn-\sum_iS_{ii}n_i}{2}.
\]
Interpolate on each unimodular Kuhn tetrahedron.  This is an integral
piecewise linear support function.  Translation by $\lambda$ changes
it by the integral affine function
$n^tS\lambda+q_S(\lambda)$, giving the deck linearization
$t_c^{-q_S(\lambda)}x^{-S\lambda}$.
The identity
$q_S(\lambda+\mu)=q_S(\lambda)+q_S(\mu)+\lambda^tS\mu$
is exactly the cocycle condition.  On a smooth fibre its Chern class
is $\sum S_{ij}a_i b_j$.  Restriction of this line bundle to the
central fibre therefore supplies every class
\[
 a_i b_i,\qquad a_i b_j+a_j b_i\quad(i<j)
\]
in the specialization image, as well as the $b_i$ and their products.
These elements generate the full integral invariant exterior algebra
for three pairs.  To check saturation explicitly, the accompanying
code enumerates all monomials of weighted degree at most six in these
nine generators and reduces their coordinates in an integral kernel
basis of $\bigwedge^q U-I$.  In every degree the nonzero Smith factors
are all one.  This finite check has only $2^3$ possible products of
the odd generators and at most three even factors.  It can also be
seen directly: in degree three the mixed kernel is the trace-zero
rank-eight lattice, generated by six off-diagonal terms and two
primitive diagonal differences; the complementary degrees are
generated by the corresponding products of two or three $a_i b_i$.
The parameter isotopy already constructed identifies the integral
specialization maps for general $C_3$.  Thus the images are primitive
and surjective in every degree.  Comparing their ranks with (E.6)
gives exactly the asserted free degree-four kernel.
\end{proof}

\paragraph{Theorem E.1.4 (Proved: fundamental group and rational homology).}
The fourfold in Theorem E.1.2 has
\[
 \pi_1(Z)=1,\qquad \chi(Z)=6,\qquad
 (b_0,\ldots,b_8)=(1,0,2,1,2,1,2,0,1).
 \tag{E.8}
\]
The first homology groups of its two finite reduced fibres are,
respectively, $\mathbb Z^2\oplus\mathbb Z/3$ and
$\mathbb Z^2\oplus\mathbb Z/2$.

\begin{proof}
The section of $\Pi$ splits the fundamental-group sequence over its
smooth locus $X\setminus W$.  Every loop in $Z$ can be moved off
the discriminant inverse image and is a product of a loop in the
elliptic fibre and a section loop.  A section loop becomes trivial
in $Z$ because it bounds in the embedded simply connected $X$.
Thus $\pi_1(Z)$ is a quotient of the coinvariants of
$\mathbb Z^2$ under $C_1,C_2$.  Loops realizing these two
monodromies are the root-circle paths closed by the affine
identifications in $X\setminus W$.  The Smith form of
$(C_1-I\mid C_2-I)$ is $(1,1)$, so these coinvariants vanish.

Euler characteristic localizes on the exceptional fibres: all
regular torus fibres and the finite free torus quotients have Euler
characteristic zero.  The cusp has six zero-dimensional torus
strata, and every positive-dimensional torus stratum has Euler
characteristic zero.  Hence $\chi(Z)=6$.

Write $V=R^1g_*\mathbb Q$ on the punctured base, and $j$ for its
inclusion.  Finite quotient specialization is an isomorphism onto
invariants over $\mathbb Q$.  Theorem E.1.3 gives
$R^qg_*\mathbb Q=j_*\bigwedge^qV$ except for one kernel
$\mathbb Q_{p_0}$ in degree four.  For any local system $L'$ on
the punctured sphere,
\[
 H^2(B,j_*L')=(L')_{\pi_1},\qquad
 \chi(B,j_*L')=-\operatorname{rank}L'+
       \sum_{j=0}^2\dim (L')^{M_j}.
 \tag{E.9}
\]
The first equality is compactly supported duality and the second is
Euler additivity on the open stratum and the three point stalks.
The exact exterior-matrix computation gives the following rows of
the actual rational Leray $E_2$ page:
\[
\begin{array}{c|rrrrrrr}
q&0&1&2&3&4&5&6\\ \hline
\dim E_2^{0,q}&1&1&2&3&3&1&1\\
\dim E_2^{1,q}&0&1&0&1&0&1&0\\
\dim E_2^{2,q}&1&1&2&3&2&1&1.
\end{array}
\tag{E.10}
\]
This is not an assumed degeneration.  The differential ranks
$r_q=\operatorname{rank}(E_2^{0,q}\to E_2^{2,q-1})$ are
\[
 (r_1,\ldots,r_6)=(1,1,2,2,1,1).
 \tag{E.11}
\]
Indeed $\Pi^*$ transports the four nonzero transgressions of the
sphere fibration to the corresponding direct summands here.
These four transgressions are forced over $\mathbb Q$ by the
one-dimensional rows and $H^*(X;\mathbb Q)=H^*(S^6;\mathbb Q)$;
this also checks the particular reduction of Engel's Lemma 7.6
which we use.  The divisor class of the zero section gives a
permanent class $u$ restricting to the elliptic orientation class.
Multiplication by $u$ gives another independent image in degree
three, and proves $r_3=2$, $r_5=1$, $r_6=1$, and $r_4\geq2$.
The target summands are independent by the marked decomposition
$\Lambda_X\oplus\Lambda_E$.  The resulting formulas for $b_3$
and $b_5$ are $b_3=1$ and $b_5=3-r_4$.  Poincare duality of
the compact smooth fourfold gives $r_4=2$.  No later differential
exists because the base has real dimension two.  Taking diagonal
sums of (E.10) with (E.11) proves (E.8).

Finally, for a finite quotient its affine fundamental group has
generators $\Lambda\oplus\mathbb Z a$ and abelianized relations
$(A_j\oplus C_j-I)\Lambda$ and $m_ja-(v_j,0)$.  The complete
relation matrix is printed by \texttt{fourfolds\_checks.py}.
Its nonzero Smith factors are $(1,1,1,1,3)$ and
$(1,1,1,1,2)$, respectively.  This proves the integral assertions
about the finite fibres.
\end{proof}

\paragraph{Theorem E.1.5 (Proved: integral homology of the fourfold).}
The fourfold of Theorem E.1.2 has torsion-free integral homology:
\[
 H_i(Z;\mathbb Z)=\mathbb Z^{b_i},\qquad
 (b_0,\ldots,b_8)=(1,0,2,1,2,1,2,0,1).
\]
In particular the finite-fibre torsion in E.1.4 does not survive
as torsion in the homology of the total space.

\begin{proof}
We compute the integral covering cokernels, rather than substituting
integral invariants in the rational calculation. Write
\(\mathcal V=R^1g_*\mathbb Z|_{B^\circ}\), and put
\(\mathcal J_q=j_*\bigwedge^q\mathcal V\).
The cohomology monodromies are the inverse transposes of
\(A_1\oplus C_1,A_2\oplus C_2,M_0\oplus C_0\).
All lattices below are in this marking.

First we describe a finite filling without making any equivariant
formality assumption. Since the first coordinate of each twist
\(v_j\) is \(\pm1\), the integral change of basis with first
column \(v_j/(v_j)_\gamma\) and remaining columns the last five
standard basis vectors changes its linear monodromy to
\(\operatorname{diag}(1,D_j)\) and its translation to
\(\pm e_0/m_j\). The reduced finite fibre is the mapping torus
of the actual linear automorphism \(D_j:T^5\to T^5\).
The covering \(T^6\to B_j\) is the \(m_j\)-fold cover of its
suspension circle. Set \(T_j=D_j^{-t}\). The two vertical matrices
on homology, in coordinates \((u,w,\delta,e_1,e_2)\), are
\[
D_1=\begin{pmatrix}
0&1&0&0&0\\-1&-1&0&0&0\\1&0&1&0&0\\
0&0&0&-1&-1\\0&0&0&1&0
\end{pmatrix},\qquad
D_2=\begin{pmatrix}
0&-1&0&0&0\\1&0&0&0&0\\0&1&1&0&0\\
0&0&0&0&1\\0&0&0&-1&0
\end{pmatrix}.
\]
The Wang sequences of the mapping torus and its cover give the
following exact description of the cokernel of covering pullback:
\[
 C_j^q\simeq
 \frac{\ker(\bigwedge^{q-1}T_j-I)}
 {\operatorname{im}\mathcal N_j^{q-1}},\qquad
 \mathcal N_j^{q-1}=\sum_{a=0}^{m_j-1}
                  (\bigwedge^{q-1}T_j)^a .
 \tag{E.I1}
\]
Here the target of pullback is the invariant lattice in
\(H^q(T^6;\mathbb Z)\). To verify (E.I1), its restriction to
the vertical torus surjects onto
\(\ker(\bigwedge^qT_j-I)\), and restriction of a mapping-torus
class has exactly the same image, by Wang exactness. Thus every
cokernel class has a representative containing the suspension
covector. On this kernel the covering map is the displayed norm:
cut the covered circle into its \(m_j\) consecutive intervals
and transport the vertical cochain through them. This proves both
surjectivity and the precise kernel of the horizontal description.
It does not choose a splitting of either Wang sequence.

For \(q=2\), the two quotients in (E.I1) are respectively zero
and \(\mathbb Z/2\). For \(q=3\), they are
\[
 C_1^3=\mathbb Z/3,\qquad
 C_2^3=(\mathbb Z/2)^2\oplus\mathbb Z/4.
 \tag{E.I2}
\]
The same Wang sequences show that the kernels of finite
specialization in degree two are \(\mathbb Z/3\) and
\(\mathbb Z/2\), and that the degree-three and degree-four
kernels vanish. The cusp specialization is integrally onto,
with the sole free degree-four kernel of E.1.3.

We give the integral global maps needed to use these local groups.
The global invariant two-form
\(q_X=u^*w^*+6\gamma\delta^*\) maps onto \(C_2^2\).
Indeed, if its pullback descended to the six-dimensional finite
fibre, restriction to the elliptic zero section would make it
descend to the original bielliptic surface \(S_2\). Its square
there would be \(\pm12/4=\pm3\), which is impossible in the
even intersection form of \(S_2\). That form is even because
\(H^2(S_2;\mathbb Z)\) is torsion free and its canonical class
is torsion, as checked in D.1.4. Thus the global-to-local map in
degree two is onto its only nonzero covering cokernel.

For degree three all three global invariant classes contain
\(\gamma\), and remain purely horizontal in both adapted
finite markings. They are
\[
 \gamma e_1^*e_2^*,\qquad
 \gamma(-u^*e_1^*+w^*e_2^*),\qquad
 \gamma u^*w^* .
 \tag{E.I3}
\]
Juxtaposition denotes exterior product. Consequently their images
in (E.I2) are exactly the vertical two-forms in (E.I3) modulo the
norm lattices; there is no unspecified lift in this computation.
In the integral invariant bases printed by the accompanying program,
the two norm matrices and this common image matrix are
\[
\begin{split}
N_3&=\begin{pmatrix}
0&0&0&0&0&0&0&0&0&3\\
0&0&-2&1&0&1&1&1&-1&0\\
0&0&-1&2&0&2&-1&0&-1&0\\
3&-1&0&0&2&0&0&0&0&0
\end{pmatrix},\\
N_4&=\begin{pmatrix}
0&0&0&0&0&0&0&0&0&4\\
0&0&-2&0&0&0&2&1&-1&0\\
0&0&0&2&0&2&0&-1&-1&0\\
4&-2&0&0&2&0&0&0&0&0
\end{pmatrix},\qquad
G=\begin{pmatrix}1&0&0\\0&1&0\\0&0&0\\0&0&1\end{pmatrix}.
\end{split}
\]
Every nonzero Smith factor of
\(\bigl(\operatorname{diag}(N_3,N_4)\mid(G,G)^t\bigr)\)
is one; there are eight. Thus the global classes surject onto
the direct sum of the two degree-three covering cokernels.

We also compute the integral parabolic groups, retaining torsion.
The degree-two global coboundary matrix, obtained by stacking
the two monodromy-minus-identity matrices, has a single nontrivial
Smith factor six. Its torsion generator restricts with orders
three, two and one at the three boundary circles. Hence no
nonzero torsion class is parabolic. Its rational parabolic rank
is zero, so \(H^1(B,\mathcal J_2)=0\). In degree three the
global coboundary image is primitive and the rational parabolic
rank is one; hence \(H^1(B,\mathcal J_3)=\mathbb Z\).
These assertions use the cocycle description as the kernel of
restriction from the free-group cohomology to the three
peripheral circle groups. The matrices and the restriction
orders are computed exactly in
\texttt{fourfolds\_integral\_audit.py}.

Passing from \(\mathcal J_q\) to \(R^qg_*\mathbb Z\)
adds the point-supported kernels just described and replaces
the finite stalks by the pullback sublattices. A point-supported
sheaf has no positive cohomology. The long exact sequences,
together with the two surjections onto covering cokernels,
therefore give
\[
 H^1(B,R^2g_*\mathbb Z)=0,\qquad
 H^1(B,R^3g_*\mathbb Z)=\mathbb Z.
 \tag{E.I4}
\]
They also identify degree-two base cohomology with the global
coinvariants. In the relevant degrees these are torsion-free,
with
\[
 H^2(B,R^1g_*\mathbb Z)=\mathbb Z,
 \qquad H^2(B,R^2g_*\mathbb Z)=\mathbb Z^2.
 \tag{E.I5}
\]
The first is the corresponding group from \(X\). The second
is the direct sum of that from \(\bigwedge^2\mathcal V_X\)
and the elliptic orientation line; the mixed coinvariants are
zero. These identifications, including primitivity, are also
verified by Smith forms, so no finite index is suppressed.

The integral differential onto the first group in (E.I5) is
surjective by naturality from the sphere fibration: its own
corresponding Leray quotient must vanish because \(H^3(X)=0\).
The differential onto the second group is surjective as well.
The \(X\) summand is killed by the differential transported
from \(H^4(X)=0\); the elliptic orientation summand is killed
by multiplying the primitive degree-one transgression of \(X\)
by the permanent class of the elliptic zero-section divisor.
That divisor has degree one on each elliptic fibre. Thus these
are two primitive independent target generators, not merely two
independent rational images.

Now \(H^0(B,R^3g_*\mathbb Z)\) and
\(H^0(B,R^4g_*\mathbb Z)\) are free: in the first degree
specialization has no kernel, and in the second its only
kernel is the free cusp class. There is no later Leray
differential, since the base has cohomological dimension two.
Equation (E.I4) and the first surjectivity in (E.I5) show that
\(H^3(Z;\mathbb Z)\) is a free group, of rank one by E.1.4.
Equation (E.I4) and the second surjectivity show that
\(H^4(Z;\mathbb Z)\) is an extension of two free groups of
rank one, and hence is \(\mathbb Z^2\).
Simple connectivity makes \(H^2(Z;\mathbb Z)\) torsion free.
The universal coefficient theorem now proves successively that
\(H_2(Z;\mathbb Z)=\mathbb Z^2\) and
\(H_3(Z;\mathbb Z)=\mathbb Z\). Poincare duality determines
the other degrees and gives exactly the displayed groups.
\end{proof}


\paragraph{Theorem E.2.1 (Proved: complex invariants and middle form).}
For the fourfold of Theorem E.1.2,
\[
 K_Z=\Pi^*L^{-5},\qquad a(Z)=2,\qquad
 (h^{0,0},\ldots,h^{0,4})=(1,1,0,0,0).
 \tag{E.12}
\]
It is not in Fujiki's class $\mathcal C$.  Moreover
$c_1(Z)=0$ integrally, $\int_Zc_2^2=2$, $\sigma(Z)=2$, and
the intersection form on $H^4(Z;\mathbb Z)/\mathrm{torsion}$ is
$\langle1\rangle\oplus\langle1\rangle$.

\begin{proof}
We first check the structure-sheaf input from the supplied family
calculation.  The orbifold base has roots of orders three and four.
Write its line generator as $\mathcal O(1)$; it pulls back to $L$.
For the coarse map $\rho$,
\[
 \rho_*\mathcal O(k)=\mathcal O_B(d(k)),\qquad
 d(k)=\frac{k}{12}-\frac{k\bmod3}{3}
                    -\frac{-k\bmod4}{4}.
 \tag{E.13}
\]
This follows by dividing a local invariant section by its necessary
root-coordinate power.  Finite-group invariants over $\mathbb C$
are exact.  In particular $\mathcal O(-1)$ and $\mathcal O(-2)$
are acyclic, since their coarse degree is $-1$.

For $\bar f:X\to\mathfrak B$, the actual direct images are
\[
 R^0\bar f_*\mathcal O_X=\mathcal O,\quad
 R^1\bar f_*\mathcal O_X=\mathcal O\oplus\mathcal O(-1),
 \quad R^2\bar f_*\mathcal O_X=\mathcal O(-1).
 \tag{E.14}
\]
Here is the step of the supplied proof that verifies (E.14).
The cusp branch resolution of $\mathcal O_W$ is the cellular
complex of its dual two-torus, and has cohomology dimensions
$(1,2,1)$ with nonzero exterior cup pairing.  Proper flat base
change therefore gives vector bundles with those ranks.  The
relative logarithmic vertical bundle has a fixed line generated
by $\widehat\delta$, and its determinant has twelfth power
$\mathcal O(-p_0)$, by the modular discriminant.  Thus the
determinant is $\mathcal O(-1)$.  Relative duality gives the last
term of (E.14).  The invariant covector $\gamma$ gives a nowhere
zero section of $R^1\bar f_*\mathcal O_X$: on a torus its class
in $\Lambda^*_{\mathbb C}/\operatorname{rowspan}\Pi_c$ is
nonzero by the last two columns of $\Pi_c$; at the cusp it is
the nonzero deck cocycle on the dual two-torus.  Its determinant
and the nonzero cusp cup pairing then give the extension
$0\to\mathcal O\to R^1\bar f_*\mathcal O_X
\to\mathcal O(-1)\to0$.  It splits because
$H^1(\mathfrak B,\mathcal O(1))=0$, also from (E.13).
The same logarithmic determinant gives $K_X=L^{-6}$.
These checks use the actual root atlases and reduced cusp; they
do not replace the multiple fibres by an unramified base disc.

The Weierstrass adjunction calculation gives
$K_{\mathscr W}=\Pi^*(K_X\otimes L)$, and the small toric
resolution is crepant.  This proves the canonical formula in
(E.12).  Every fibre of $\Pi$ is a connected genus-one curve;
over the deeper strata its toric resolution is a cycle of rational
curves.  Thus $\Pi$ is flat, by equidimensionality between smooth
spaces, and relative duality gives
$R^1\Pi_*\mathcal O_Z=L^{-1}$ and
$R^0\Pi_*\mathcal O_Z=\mathcal O_X$.
Restriction to the section splits the unit map in the derived
direct image, so
\[
 R\Pi_*\mathcal O_Z\simeq
 \mathcal O_X\oplus L^{-1}[-1].
\]
Projection formula, (E.13), and (E.14) show that
$H^*(X,L^{-1})=0$ and $h^*(X,\mathcal O_X)=(1,1,0,0)$.
This proves the asserted Hodge numbers.

A meromorphic trivialization of $L$ makes the coordinate $x$ in
(E.1) a meromorphic function over $X$, of degree two on each
elliptic fibre.  It gives a generically finite meromorphic map to
$X\times\mathbb P^1$.  Traces and elementary symmetric functions
over its finite generic fibre show that the meromorphic function
field is an algebraic extension of that of the target.  Since
$a(X)=1$, this proves $a(Z)=2$.  If $Z$ were in class
$\mathcal C$, a sequence of blowups along smooth centres would
give a compact Kahler model, by the form of the model theorem
in \cite[Theorem 1.1 and Remark 1.2]{ffJiaMeng}.
Both $b_1$ and $h^{0,1}$ are unchanged by such blowups: the
exceptional fibres have no first homology and no positive-degree
structure-sheaf cohomology.  The Kahler identity
$b_1=2h^{0,1}$ would give $0=2$, a contradiction.

Since $L$ is pulled from a line on $S^6$, its integral first Chern
class vanishes, so $c_1(Z)=0$.  Riemann--Roch and the signature
formula, in the explicit normalization of
\cite[(1.2)--(1.4)]{ffEsoleKang}, yield
\[
 0=\chi(\mathcal O_Z)=\frac{3c_2^2-c_4}{720},\qquad
 \sigma(Z)=\frac{3c_2^2+14c_4}{45}.
\]
Here $c_4[Z]=6$, hence $c_2^2[Z]=2$ and $\sigma(Z)=2$.
The integral middle pairing modulo torsion is unimodular by
Poincare duality, has rank two by (E.8), and is positive definite
by its signature.  The rank-at-most-four lattice classification
\cite[Theorem 2.2.1]{ffPetersSterk} therefore gives precisely
$\langle1\rangle\oplus\langle1\rangle$.
\end{proof}

\paragraph{Theorem E.2.2 (Proved: transverse pullback and the trivial case).}
Let $Y$ be a compact complex manifold and $h:Y\to B$ a holomorphic
map which is a submersion at every point over $p_0,p_1,p_2$.
Then $P=X\times_B Y$ is a compact complex manifold, of dimension
$\dim_{\mathbb C}Y+2$.  If $Y=B\times N$ with its projection,
then $P\simeq X\times N$ holomorphically, is diffeomorphic to
$S^6\times N$, and
\[
 \pi_1(P)=\pi_1(N),\quad
 H_k(P;\mathbb Z)=H_k(N;\mathbb Z)\oplus H_{k-6}(N;\mathbb Z),
 \quad \chi(P)=2\chi(N).
 \tag{E.15}
\]

\begin{proof}
The pullback is the inverse image of the diagonal under
$(f,h):X\times Y\to B\times B$.  At any point, either $df$ is
onto or its base value is exceptional and $dh$ is onto.  Thus
the difference of these two differentials is onto, proving
transversality and smoothness.  It is a closed subset of a
compact product, hence compact.  The maps
$(x,(t,n))\mapsto(x,n)$ and $(x,n)\mapsto(x,(f(x),n))$ give
the asserted holomorphic identification.  The homology statement
is the integral Kunneth theorem with $H_*(S^6)$, which is free;
there are no Tor terms.  Fundamental groups and Euler
characteristics of products give the remaining formulas.
\end{proof}

\paragraph{Theorem E.2.3 (Proved: a shifted second sphere).}
Let $X'$ be a second canonical member and choose $\varphi\in
\operatorname{PGL}_2(\mathbb C)$ so that its three exceptional
values are disjoint from those of $X$.  The smooth fivefold
$P=X\times_B\varphi^*X'$ has Euler characteristic zero and
an infinite fundamental group with presentation
\[
\begin{split}
G=\langle a_1,a_2,b_1,b_2,h_1,h_2\mid\ &
 h_1,h_2\ \text{central},\quad
 a_1^3=h_1,\ a_2^4=h_1^{-1},\\
 &b_1^3=h_2,\ b_2^4=h_2^{-1},\quad
 a_1a_2b_1b_2=1\rangle .
\end{split}
\tag{E.16}
\]
Its abelianization is $\mathbb Z$, and its rational Betti numbers are
\[
 (1,1,1,20,63,88,63,20,1,1,1).
 \tag{E.17}
\]
For fixed $X,X'$, all admissible $\varphi$ give the same oriented
diffeomorphism type.

\begin{proof}
Theorem E.2.2 proves smoothness.  At each special value the other
factor is a smooth four-torus.  The Euler characteristic of every
fibre is therefore zero, and Euler additivity gives $\chi(P)=0$.
We give van Kampen with the actual cusp relations, rather than
deducing simple connectivity from that of the factors.
Over the six-punctured sphere the fibre group is
$\Lambda_X\oplus\Lambda_{X'}$ and the base is free on five
peripheral loops.  Each monodromy acts on its own direct summand.
Filling the first cusp kills its angular generators
$\widehat w,\widehat\delta$ and its cusp meridian.  Their normal
closure under $A_1$ also kills $\widehat u$, since
$A_1\widehat w=\widehat u-\widehat w$.  The same calculation
in the other summand leaves only $h_1=\widehat\gamma$ and
$h_2=\widehat\gamma'$.  The finite monodromies now act trivially
on them.  The four finite filling relations are precisely those
in (E.16), since $\gamma(v_1)=1$ and $\gamma(v_2)=-1$.
Both cusp collars are untwisted; their meridians are killed.
The product of the remaining four meridians gives the last
relation.  These are all relations introduced by the six
fillings, so van Kampen gives exactly (E.16).

Quotienting by $h_1,h_2$ gives the orbifold group of the sphere
with cone orders $3,4,3,4$.  Its orbifold Euler characteristic
is $-5/6$, so its universal orbifold cover is the disc and the
group is infinite.  Equivalently one can build the Seifert
$T^2$-bundle with these four primitive translation vectors; over
the disc it is a trivial torus bundle.  Its exact sequence shows
that the central subgroup in (E.16) is $\mathbb Z^2$.
The abelianized five-by-six relation matrix has Smith factors
$(1,1,1,1,1)$, as printed in the code, proving $G_{\rm ab}=\mathbb Z$.

For the rational homology there is no extra cusp kernel: each
cusp is the original rank-two cusp times a smooth torus, and its
specialization is an isomorphism onto invariants over $\mathbb Q$.
The two monodromy groups act on independent summands.  In exterior
degrees $q=0,\ldots,8$, the $H^0$ and $H^2$ rows are both
\[
 (1,2,3,4,5,4,3,2,1),
\]
and the $H^1$ row is
\[
 (0,0,18,60,86,60,18,0,0).
\]
To verify this without a matrix-size guess, the local invariant
polynomials of the first factor are
$A(z)=1+2z+2z^2+2z^3+z^4$ at a finite point and
$C(z)=A(z)+2z^2$ at its cusp.  Hence the Euler polynomial of
the direct-image rows is
\[
 -4(1+z)^8+(4A(z)+2C(z))(1+z)^4.
\]
Formula (E.9), with six point strata, gives the displayed rows.

Each invariant algebra of a factor has one generator in degrees
zero through four.  Its invariant-to-coinvariant map is nonzero
in even degrees and zero in odd degrees.  In the marking used
above, this follows from the invariant form
$6\gamma\wedge\delta+u\wedge w$ and its exterior products:
the odd invariant classes pair to zero with every invariant
complementary class, and the even ones have nonzero square or
top pairing.  Its four transgressions are all nonzero, as the
factor is a sphere.  The derivation rule therefore sends a tensor
of bidegree $(i,j)$ into bidegrees $(i-1,j)$ and $(i,j-1)$,
with the first term nonzero exactly when $i>0$ and $j$ is even,
and the second exactly when $j>0$ and $i$ is even.  The nonzero
column supports are disjoint except for pairs in a single row,
so their scalar normalizations do not affect ranks.  This gives
\[
 (r_1,\ldots,r_8)=(1,2,2,3,3,2,2,1).
\]
The exact program forms these matrices and takes their ranks.
There are no later differentials.  Their diagonal sums give
(E.17), which also satisfies Poincare duality and Euler zero.

Finally, the admissible set of $\varphi$ is the complement of
finitely many proper complex hypersurfaces in the connected
complex manifold $\operatorname{PGL}_2(\mathbb C)$, so it is
path connected.  The universal pullback over this set is a
proper smooth submersion: transversality was checked pointwise
above.  Parallel transport along a path gives an oriented
diffeomorphism of its fibres.
\end{proof}

\paragraph{Conclusion E.2.4 (Proved: the shifted product's smooth type).}
The shifted product in (E.16) is the closed oriented ten-manifold
$\mathcal H_{10}$ specified by the $2432$ ordinary framed handles
of E.2.21. Its regular-fibre complement is, relative to its entire
marked boundary, the $76$-handle six-manifold $\mathcal H_6$ of
that theorem. The full attaching embeddings, their ordered normal
coordinates, and every dual belt embedding are prescribed by
finite coordinate formulas. The identification follows by an
actual relative diffeomorphism of the cusp, core and finite caps,
and then by the handle-by-handle comparison of E.2.21 using
Milnor's Theorems~1.4, 3.13 and~5.8. Top attaching maps are
retained, so there is no uncomputed sphere-capping ambiguity.

Together with E.2.1--3, E.2.5--20, this supplies the complex
invariants, the general smooth compact pullback construction,
the product case $S^6\times N$, and the shifted product's
fundamental group and full integral homology. The smooth type
is given by the explicit framed-handle presentation, not by an
inference from its homology or by an asserted connected-sum
normal form. E.2.18's actual quadratic obstruction is consistent
with this presentation. The former open paragraph is retained
in the research records.

\paragraph{Proposition E.2.5 (Proved: the first integral torsion of the shifted product).}
Let \(Z_\varphi=X\times_B\varphi^*X\) be the transverse
fivefold in E.2.3, with the two sets of three critical values
disjoint. Its second integral homology is
\[
 H_2(Z_\varphi;\mathbb Z)
   =\mathbb Z\oplus(\mathbb Z/12)^7.
\]
In particular the rational Betti vector in E.2.3 cannot be
interpreted as integral torsion-free homology.
\begin{proof}
Use subscripts one and two for the two rank-four cohomology
local systems. At a critical value of one factor the other
factor is a smooth torus bundle over a disc, hence a product
there. Each local specialization is therefore the tensor
product of the original specialization with the identity on
the exterior algebra of \(\mathbb Z^4\). The original finite
reduced surfaces and the original cusp have free cohomology.
Their specialization maps are injective; at the cusp they
are integral isomorphisms onto invariants. For the latter,
apply the support-function construction of E.1.3 in two
coordinates. It supplies the invariant deck classes
\(b_1,b_2\), the three classes
\(a_1b_1,a_2b_2,a_1b_2+a_2b_1\), and their products.
Together with \(b_1b_2\) these give an integral kernel
basis in degree two; their products give integral bases
in degrees three and four. Both the central-fibre and
invariant lattices have ranks \((1,2,4,2,1)\), as follows
also from the normalization matrix in D.1.2 at \(d=1\).
Thus the specialization is onto between free groups of
equal ranks in every degree, proving the claimed isomorphism.
Thus the only difference between \(R^qg_*\mathbb Z\) and
\(\mathcal J_q=j_*\bigwedge^q(\mathcal V_1\oplus\mathcal V_2)\)
is their finite covering cokernel, with no point-supported kernel.

The finite Wang calculation (E.I1), now with the original
three-dimensional vertical lattice, gives
\[
 C_3^1=\mathbb Z/3,\quad C_4^1=\mathbb Z/4,
 \qquad C_3^2=0,\quad C_4^2=\mathbb Z/2.
\]
The invariant covector \(\gamma\) generates the first two
cokernels, and the invariant two-form \(q_X\) generates the
last, by the same even-intersection argument as in E.1.5.
Consequently the degree-two covering cokernel at all four
finite values, combined, is
\[
 C^2=\left((\mathbb Z/3)^4\oplus
                 (\mathbb Z/4)^4\oplus\mathbb Z/2\right)^2.
\]
The two monodromy groups act independently on the two factors.
The three global invariant two-forms have integral basis
\(q_1,\gamma_1\gamma_2,q_2\). The first and third map
onto the respective \(\mathbb Z/2\) summands and vanish
in the other summands. The middle form maps to the primitive
\(\gamma\) coordinate of each of the four order-three or
order-four blocks. After the two order-two summands have
been removed, it therefore generates a direct cyclic summand
of order twelve in \((\mathbb Z/12)^8\). Hence
\[
 \operatorname{coker}\{H^0(B,\mathcal J_2)\to C^2\}
      =(\mathbb Z/12)^7.
 \tag{E.20}
\]
The program \texttt{fourfolds\_shifted\_integral.py} realizes
this as an eighteen-row integer presentation; its Smith form
has eleven unit factors and seven factors twelve.

In degree two the global coboundary image is primitive.
Thus the parabolic group, a subgroup of the free-group
cohomology on the punctured base, is torsion free. Its rank
is eighteen: the six local invariant ranks are
\((16,16,18,16,16,18)\), and its global invariant and
coinvariant ranks are both three. The Euler formula on the
six-punctured sphere gives \(3+3+4\cdot28-100=18\).
The primitivity and all these ranks are checked by the exact
exterior matrices in the same program. The sheaf sequence
\(0\to R^2g_*\mathbb Z\to\mathcal J_2\to C^2\to0\)
and (E.20) now give
\[
 H^1(B,R^2g_*\mathbb Z)
     =\mathbb Z^{18}\oplus(\mathbb Z/12)^7.
\]
Here the quotient by its finite subgroup is free, so this
abelian-group extension splits.

The group \(H^2(B,R^1g_*\mathbb Z)\) is the free
rank-two group of coinvariants, one primitive summand from
each factor. Pulling back the integral degree-two
transgression of the original sphere fibration through each
projection kills these two summands separately and
surjectively. Thus \(E_\infty^{2,1}=0\) in the Leray
sequence of \(g\). Its other middle group
\(E_\infty^{0,3}\) is free, since \(R^3g_*\mathbb Z\)
injects into \(\mathcal J_3\). There are no higher
differentials. Consequently
\(\operatorname{Tor}H^3(Z_\varphi;\mathbb Z)
=(\mathbb Z/12)^7\). The universal coefficient theorem
identifies this with \(\operatorname{Tor}H_2(Z_\varphi;\mathbb Z)\).
Finally E.2.3 gives \(b_2=1\), proving the assertion.
This argument does not determine the extensions in the
higher integral Leray filtrations.
\end{proof}


\paragraph{Theorem E.2.6 (Proved: full integral homology of the shifted product).}
For the transverse shifted product $P$ of E.2.3, write
\[
 A=(\mathbb Z/12)^7,\qquad
 B=(\mathbb Z/12)^{10}\oplus(\mathbb Z/2)^7,\qquad
 C=(\mathbb Z/12)^6\oplus\mathbb Z/24\oplus(\mathbb Z/2)^{17}.
\]
Its complete integral homology is
\[
 \begin{array}{c|l}
 i&H_i(P;\mathbb Z)\\\hline
 0,1,8,9,10&\mathbb Z\\
 2&\mathbb Z\oplus A\\
 3&\mathbb Z^{20}\oplus B\\
 4&\mathbb Z^{63}\oplus C\\
 5&\mathbb Z^{88}\oplus C\\
 6&\mathbb Z^{63}\oplus B\\
 7&\mathbb Z^{20}\oplus A.
 \end{array}
\]
In particular the higher integral filtration extensions are nontrivial:
the middle torsion includes a cyclic factor of order $24$.
\begin{proof}
\emph{A decomposition with free groups on every piece.}
Choose a separating circle in $B=\mathbb P^1$ so that one closed
disc contains the three critical values of $X$ and the other
contains the three critical values of $X'$. The smooth torus
family of the other factor is trivial over each disc. Set
\[
 F=T^4,\qquad E=X\setminus\operatorname{int}(F\times D^2),
 \qquad E'=X'\setminus\operatorname{int}(F'\times D^2).
\]
The product trivializations give the actual decomposition
\[
 P=(E\times F')\ \cup_{F\times F'\times S^1}\ (F\times E').
 \tag{E.MV1}
\]
The boundary identification preserves the two marked torus factors
and reverses the normal circle, because the two base discs have
opposite boundary orientations. There is no additional torus
clutching: the same regular-disc trivializations define both the
boundary identifications of $E,E'$ and (E.MV1).

Since $X,X'$ are six-spheres, the Mayer--Vietoris sequence of
$S^6=E\cup(F\times D^2)$ gives, integrally,
\[
 H_0(E)=\mathbb Z,\qquad H_k(E)=H_{k-1}(F)\quad(1\leq k\leq4),
 \qquad H_k(E)=0\quad(k\geq5).
 \tag{E.MV2}
\]
The isomorphism is the normal-circle sweep
$\alpha\mapsto[S^1]\times\alpha$. Indeed, in degrees one
through four the boundary map to $H_k(E)\oplus H_k(F)$ is
an isomorphism, and its second component is the identity on
the unswept fibre summand and zero on the swept summand.
In degree five the boundary fundamental class is the image
of the sphere fundamental class, so $H_5(E)=0$.
Thus all homology groups of the two pieces and their common
boundary in (E.MV1) are free abelian. The same statements hold
for $E'$.

\emph{Determining every boundary map integrally.}
Let $V=H^1(F;\mathbb Z)$, in the basis $(\gamma,u,w,\delta)$,
and use the positively ordered volume
$v=\gamma u w\delta$. The primitive global invariant classes
of the original torus fibration are
\[
 I_0=1,\quad I_1=\gamma,\quad
 I_2=q=uw+6\gamma\delta,\quad
 I_3=\gamma uw,\quad I_4=v.
\]
The images of $H^k(E)\to H^k(F)$ are exactly
\[
 n_k\mathbb Z I_k,\qquad
 (n_1,n_2,n_3,n_4)=(12,2,2,12).
 \tag{E.MV3}
\]
Here is an integral verification, rather than descent of invariant
real forms. On the disc the only local restrictions on a global
invariant class are the actual finite-cover pullback lattices;
the cusp specialization is an integral isomorphism onto its
invariant lattice by E.2.5. In degree one the finite images
contain $3\gamma$ and $4\gamma$, giving the intersection
$12\mathbb Z\gamma$. In degree two the order-three image
contains $q$, while the order-four image requires an even
coefficient of $q$, by the finite integral lattices of B.5.13.
In degree three use the perfect integral pairing downstairs:
an invariant class is a pullback precisely when its pairings
with every pulled-back degree-one class are divisible by the
covering degree. The other degree-one generators are
$2u+w+3\delta$ and $u+w+2\delta$. Their pairings with
$t\gamma uw$ are, up to signs, $3t$ and $2t$.
Thus the first cover imposes no further condition, and the
second requires $t$ even. In degree four the pullback degrees
are three and four, giving $12\mathbb Z v$.

These local restrictions are also sufficient for the image
from $E$. For a local system on a disc with punctures, the
restriction from its free-group $H^1$ to the direct sum of
the local coinvariants is onto: each puncture has its own
independent free generator. Consequently
$H^2(D,j_*\mathcal L)=0$. Replacing $j_*\mathcal L$ by the
actual direct-image sheaf changes it only by point-supported
cokernels, and so does not change this vanishing.
The Leray sequence over the disc therefore makes
$H^k(E)\to H^0(D,R^kf_*\mathbb Z)$ onto. This proves
(E.MV3) without an integral formality assumption.

Write $\theta$ for the normal-circle covector. The dual of
(E.MV2) identifies $H^k(E)$ with $\bigwedge^{k-1}V$, and
its boundary restriction has the graph form
\[
 r_k(\alpha)=\phi_k(\alpha)+\theta\alpha,
 \qquad
 \phi_k:\bigwedge^{k-1}V\longrightarrow\bigwedge^kV.
 \tag{E.MV4}
\]
The image of $\phi_k$ is exactly $n_k\mathbb Z I_k$ by
(E.MV3). If $k+l=5$, classes extending to $E$ have zero
product evaluated on $\partial E$. Therefore
\[
 (-1)^k\phi_k(\alpha)\beta+
                    \alpha\phi_l(\beta)=0
 \quad\text{in }\bigwedge^4V.
 \tag{E.MV5}
\]
Each $I_k$ is primitive, and wedge pairing between the two
complementary exterior powers is unimodular. Writing
$\phi_k=f_k I_k$, equation (E.MV5), tested against a class
pairing to one with $I_k$, proves that $f_k$ is an integral
multiple of $\alpha\mapsto\langle\alpha I_{5-k},[F]\rangle$.
The exact image (E.MV3) fixes the absolute multiple. Thus
\[
 \begin{split}
 \phi_k(\alpha)&=c_k
       \langle\alpha I_{5-k},[F]\rangle I_k,\\
 (c_1,c_2,c_3,c_4)&=(12\varepsilon,2\eta,-2\eta,-12\varepsilon),
 \qquad \varepsilon,\eta\in\{1,-1\}.
 \end{split}
 \tag{E.MV6}
\]
This proves every boundary map up to two signs. We allow those
signs independently for $E$ and $E'$ below; the answer is the
same for all sixteen possibilities.

\emph{The full integral Mayer--Vietoris matrices.}
Let $i_p:H_p(F)\to H_p(E)$ be the transpose of $\phi_p$ in
the sweep bases, with $i_0=1$; let $j_{p+1}$ denote the sweep
isomorphism for $0\leq p\leq3$, and zero for $p=4$.
The analogous maps of the second copy are $i'_q,j'_{q+1}$.
For a fibre tensor $\alpha\otimes\beta$, with degrees $p,q$,
the matrix $D_k$ of (E.MV1) is
\[
 \alpha\otimes\beta\longmapsto
  (i_p\alpha\otimes\beta,-\alpha\otimes i'_q\beta),
 \qquad p+q=k.
 \tag{E.MV7}
\]
For a swept tensor, place the normal circle first. Then
\[
 \theta\times\alpha\times\beta\longmapsto
 \bigl(j_{p+1}\alpha\otimes\beta,
                 (-1)^p\alpha\otimes j'_{q+1}\beta\bigr),
 \qquad p+q=k-1.
 \tag{E.MV8}
\]
The second sign contains the Koszul sign from moving the circle
past $\alpha$, the reversal of the two boundary-circle
orientations, and the second minus in Mayer--Vietoris. The last
two cancel. This specifies every integer entry of $D_k$.

The executable \texttt{code/manifolds\_shifted\_mv.py} constructs
the exterior bases, the four adjoint maps (E.MV6), and the full
matrices (E.MV7)--(E.MV8). Exact Smith reduction gives the table
below. The source ranks are $\binom9k$; a zero binomial is
understood outside $0\leq k\leq9$.
\[
 \begin{array}{c|rrrr|l}
 k&\text{target rank}&\text{source rank}&\operatorname{rank}D_k&
       \text{unit factors}&\text{nonunit factors}\\\hline
 0&2&1&1&1&--\\
 1&10&9&9&9&--\\
 2&28&36&27&20&12^{[7]}\\
 3&64&84&53&36&2^{[7]},12^{[10]}\\
 4&114&126&82&58&2^{[17]},12^{[6]},24\\
 5&138&126&94&70&2^{[17]},12^{[6]},24\\
 6&104&84&73&56&2^{[7]},12^{[10]}\\
 7&44&36&35&28&12^{[7]}\\
 8&8&9&8&8&--\\
 9&0&1&0&0&--\\
 10&0&0&0&0&--
 \end{array}
\]
An exponent $[r]$ records repetition, not an integer power.
All sixteen choices of the two signs in each copy give precisely
this table; the retained output and JSON contain every case.

Finally the exact homology sequence gives
\[
 0\longrightarrow\operatorname{coker}D_k
 \longrightarrow H_k(P)
 \longrightarrow\ker D_{k-1}\longrightarrow0.
\]
The last group is free, since the boundary homology is free.
This sequence therefore splits as abelian groups, so the Smith
factors in the table are the actual torsion of $H_k(P)$.
Adding the free kernel ranks gives exactly the rational Betti
vector in E.2.3 and the integral groups asserted here.
The resulting torsion also satisfies the required ten-dimensional
duality $\operatorname{Tor}H_k(P)\cong
\operatorname{Tor}H_{9-k}(P)$.
\end{proof}

\paragraph{Proposition E.2.7 (Proved: a unique transverse family type and its stable tangent bundle).}
For fixed marked sphere fibrations, all M\"obius transformations
whose critical triples are disjoint give orientation-preservingly
diffeomorphic shifted products. The same conclusion holds while
either sphere varies in its proper smooth marked period family.
For every such fivefold $P$ there is an isomorphism of real vector
bundles
\[
 TP\oplus\mathbb R^4\simeq\mathbb R^{14}.
 \tag{E.MV9}
\]
Thus all its positive-degree Stiefel--Whitney and Pontryagin
classes vanish. For its given complex tangent bundle,
$c_1(P)=c_2(P)=0$ integrally.
These assertions and E.2.6 do not identify the complete smooth
structure from its ordinary homology and fundamental group.
\begin{proof}
The allowed transformations form the complement in
$\operatorname{PGL}_2(\mathbb C)$ of the nine analytic
hypersurfaces imposing an equality of one critical value from
each triple. This complement is path connected: a piecewise
smooth path between two points can be perturbed relative to
its endpoints to avoid these sets, which have real codimension
two. Equivalently, use local complex coordinate balls and
perturb each segment there. Choose a smooth such path.

The corresponding fibre products form a compact smooth family
over its parameter interval. Indeed the vertical derivative of
$f_1(x)-f_2(y)$ is surjective at every point, since its two
factors cannot be critical at the same base value. The same
surjectivity solves the equation for a lift of the interval
coordinate vector field. A partition of unity supplies a
global smooth lift; compactness makes its flow exist for the
whole interval. This flow gives the claimed diffeomorphisms.
The complex orientations vary continuously and hence are
preserved. The identical argument applies to smooth proper
period families of either factor, keeping the critical triples
disjoint throughout.

View $P$ as the transverse inverse image of the diagonal in
$B\times B$ under $f_1\times f_2$. Its real oriented normal
two-plane bundle is the underlying bundle of $g^*T_B$,
where $g:P\to B$ is the common base map. It is trivial:
its first Chern class is
\[
 c_1(g^*T_B)=2g^*[B]^\vee
    =2\operatorname{pr}_1^*f_1^*[B]^\vee=0,
\]
because $H^2(X;\mathbb Z)=0$, and complex line bundles are
classified by their first Chern class. The tangent bundles
of the two six-spheres become trivial after adding one real
line each. Restriction of their product tangent bundle gives
\[
 (TP\oplus\mathbb R^2)\oplus\mathbb R^2
       \simeq\mathbb R^{14},
\]
proving (E.MV9) and the characteristic-class assertions for
the real tangent bundle. Finally the holomorphic tangent
sequence is
\[
 0\longrightarrow TP\longrightarrow
 \operatorname{pr}_1^*TX\oplus\operatorname{pr}_2^*TX'
 \longrightarrow g^*T_B\longrightarrow0.
\]
Both sphere factors have integral $c_1=c_2=0$, since their
cohomology in degrees two and four vanishes. The last line
has zero first Chern class, as just proved. Multiplying total
Chern classes in this exact sequence gives
$c_1(P)=c_2(P)=0$ integrally.
\end{proof}

\paragraph{Proposition E.2.8 (Proved: the shifted product bounds with its natural stable framing).}
Equip the two oriented six-spheres with their standard bounding
stable tangent framings. Trivialize the oriented two-plane bundle
$f_1^*T\mathbb P^1$ on the first sphere, and equip the transverse
shifted product $P$ with the induced stable framing in (E.MV9).
Then there is a compact smooth oriented eleven-manifold $V$ with
\[
 \partial V=P,\qquad TV\oplus\mathbb R^3\simeq\mathbb R^{14},
 \qquad [P]=0\quad\hbox{in }\Omega^{\mathrm{fr}}_{10}.
 \tag{E.B1}
\]
The boundary framing is the stated one. In particular its induced
spin structure bounds. This supplies a framed bordism invariant;
it does not identify the smooth type from ordinary homology.
\begin{proof}
Write $f_2:S^6\to\mathbb P^1$ for the shifted map of the second
factor. By B.2.1 the first map $f_1:S^6\to\mathbb P^1$ is
null-homotopic. It therefore extends to a smooth map
$F:D^7\to\mathbb P^1$, which can be chosen to equal $f_1$ on
each sphere of a collar of the boundary. To see this with the
boundary fixed, first use the null-homotopy to obtain a continuous
extension constant in the collar direction. Smooth its interior
in coordinate charts, with cutoffs supported away from the collar;
on the overlap the extension is already smooth.

Let $C_2$ be the three critical values of $f_2$. The restriction
$f_1$ is transverse to each point of $C_2$, since the two critical
triples are disjoint. Perturb $F$ only in the interior to be
transverse to all these points. The applicable finite simultaneous
transversality statement is
\cite[Chapter 3, Section 2, Theorems 2.1 and 2.5]{rootHirsch76}.
For the relative assertion, the proof of that theorem uses local
perturbations multiplied by cutoff functions. Their supports may
be restricted to the complement of the collar, where
transversality has not already been secured; the open condition
on the collar persists. Compactness permits finitely many such
perturbations. Thus no assumption that $F$ is a submersion
everywhere is needed.

Define the smooth fibre product
\[
 V=\{(x,y)\in D^7\times S^6:F(x)=f_2(y)\}.
 \tag{E.B2}
\]
At a point whose common value is outside $C_2$, the differential
$df_2$ is onto. At a point over $C_2$, the differential $dF$ is
onto. Consequently $dF-df_2$ is onto everywhere on (E.B2).
The same argument holds on its boundary, using $df_1$ there.
The implicit-function theorem, including its boundary version,
makes $V$ a smooth eleven-manifold with boundary exactly $P$.
It is compact as a closed subset of $D^7\times S^6$.
The inverse-image orientation using the complex orientation of
$T\mathbb P^1$ induces the complex orientation of $P$ at the
boundary, because the collar keeps the first map fixed.

Let $g:V\to\mathbb P^1$ be the common map. The real normal
two-plane bundle of (E.B2) is $g^*T\mathbb P^1$, identified by
$dF-df_2$. This is the pullback from $D^7$ of $F^*T\mathbb P^1$,
which is trivial because the disc is contractible. A choice of
its oriented framing restricts on the boundary to the chosen
framing of $f_1^*T\mathbb P^1$ up to homotopy: two oriented
framings of this two-plane bundle differ by a map
$S^6\to SO(2)$, and $H^1(S^6;\mathbb Z)=0$.
The tangent exact sequence and the bounding stable framing of
the second sphere now give
\[
 \begin{split}
 TV\oplus\mathbb R^2&\simeq
       T(D^7\times S^6)|_V,\\
 TV\oplus\mathbb R^3&\simeq
       \mathbb R^7\oplus(TS^6\oplus\mathbb R)|_V
       \simeq\mathbb R^{14}.
 \end{split}
\]
At the boundary, the outward normal of $D^7$ supplies the
additional real line in $TV|_P=TP\oplus\mathbb R$.
The resulting framing is precisely the stable tangent framing
obtained from the two sphere framings and the normal framing
in E.2.7. This proves (E.B1) by the definition of framed bordism.
It also proves the assertion about the induced spin structure,
by forgetting the framing.

This disc-extension construction does not by itself extend a
classifying map $P\to B\pi_1(P)$ over $V$.
The different filling in E.2.9 below does extend that map.
Both are bordism statements; the higher equivariant surgery
data required for the smooth classification remain separate.
\end{proof}


\paragraph{Theorem E.2.9 (Proved: framed null-bordism over the fundamental group).}
Let $G$ be the group in (E.16). There is a closed smooth
aspherical four-manifold $N$ with
\[
 \pi_1(N)=G,\qquad \widetilde N\simeq\mathbb H^2\times\mathbb R^2,
 \qquad H_*(N;\mathbb Z)\cong H_*(S^1\times S^3;\mathbb Z),
 \tag{E.B3}
\]
and a smooth classifying map $c:P\to N$ inducing the displayed
identification of fundamental groups. With the natural stable
framing of E.2.8, one has
\[
                 [P,c]=0\quad\hbox{in }\Omega^{\mathrm{fr}}_{10}(N).
 \tag{E.B4}
\]
Thus $P$ has a compact stably framed eleven-dimensional filling
over $BG$. This is stronger than the filling of E.2.8; it is
still a bordism statement, not a diffeomorphism classification.
\begin{proof}
\emph{The aspherical target.}
The repaired quotient maps of A.2.1 are smooth maps
$e_i:S^6\to S^3$ over the base, and B.2.1 gives the exact
factorizations $f_i=h_i e_i$. Here $h_i:S^3\to\mathbb P^1$
is the weighted Hopf map of orders $3,4$, with the second
map shifted by the chosen M\"obius transformation. Define
\[
 N=\{(u,v)\in S^3\times S^3:h_1(u)=h_2(v)\},\qquad
 c(x,y)=(e_1(x),e_2(y)).
 \tag{E.B5}
\]
The critical values of each $h_i$ are only its two exceptional
circle values. They are disjoint, so the same transversality
argument as in E.2.2 makes $N$ a compact smooth four-manifold.
The two circle actions on the weighted Hopf spheres give a
principal $T^2$ bundle in the orbifold sense over
\[
                  \mathcal B=S^2(3,4,3,4).
\]
At a cone of order $m$, its uniformizing chart is
$(D^2\times T^2)/C_m$, with the generator rotating the disc
and translating one torus coordinate by a primitive fraction
of denominator $m$. This local description is smooth and is
the same finite-filling description as A.1.3, with the unused
circle retained.

The orbifold has Euler characteristic
$2-2(1-1/3)-2(1-1/4)=-5/6$. It is good and has hyperbolic
orbifold universal cover, by
\cite[Theorems 2.3--2.4 and the Euler-number discussion in
Section 2]{rootScott83}. Pulling the principal torus bundle to
$\mathbb H^2$ gives an ordinary principal $T^2$ bundle over a
contractible base, hence the trivial bundle. Its universal
cover is $\mathbb H^2\times\mathbb R^2$. The local chart above
shows that these orbifold pullbacks are ordinary covering maps
of total spaces. This proves asphericity and injectivity of
the fibre subgroup $\mathbb Z^2$.

Using the marked fibre generators $h_1,h_2$ and the four
positive peripheral generators gives exactly
\[
 \begin{split}
 \pi_1(N)=\langle a_1,a_2,b_1,b_2,h_1,h_2\mid {}&h_1,h_2
       \text{ central},\quad a_1^3=h_1,\quad a_2^4=h_1^{-1},\\
 &b_1^3=h_2,\quad b_2^4=h_2^{-1},\quad
 a_1a_2b_1b_2=1\rangle.
 \end{split}
\]
The map $c$ sends the surviving regular-fibre generators
$\widehat\gamma,\widehat\gamma'$ to $h_1,h_2$, respectively,
and each peripheral generator to the one with the same name.
This is because $e_i$ is the literal circle quotient on these
fibres and finite collars. Its cusp replacement extends over
a solid torus and kills precisely the cusp meridian in this
presentation. Comparing with the full van Kampen calculation
in E.2.3 therefore shows that $c_*$ is the identity on the
displayed presentation, and in particular is an isomorphism.
As $N$ is aspherical, $c$ is a classifying map.

For completeness the abelian relation matrix has Smith
diagonal $(1,1,1,1,1)$; a primitive integer kernel vector in
the order $(a_1,a_2,b_1,b_2,h_1,h_2)$ is
$(4,-3,-4,3,12,-12)$. Thus $H_1(N)=\mathbb Z$.
Every stratum of its base has a torus fibre, so Euler
additivity gives $\chi(N)=0$. Integral Poincare duality and
the universal coefficient theorem give $H_3(N)=\mathbb Z$
and no torsion in $H_2(N)$. Euler characteristic then forces
$H_2(N)=0$, proving the integral homology assertion in (E.B3).

\emph{The bordism computation before taking the inverse image.}
For a based sphere, the stable splitting
$(S^3)_+\simeq S^0\vee S^3$ identifies
\[
 \Omega^{\mathrm{fr}}_6(S^3)
       =\pi^{\mathrm S}_6\oplus\pi^{\mathrm S}_3.
 \tag{E.B6}
\]
We use Pontryagin--Thom in its form for maps to a space,
\cite[Definition 8.16 and Theorem 8.17]{rootDavisKirk01}.
The first component of $[S^6,e_i]$ is obtained by forgetting
the map, so it is zero for the bounding sphere framing.
The second component is the stable class of $e_i$, equivalently
the framed inverse image of a regular point. It is $\eta^3$
by B.1.2. One can also check the splitting directly: the
Pontryagin--Thom collapse of the framed sphere into its
disjoint-basepoint suspension has degree one on the sphere
summand; composing with $e_i$ gives its stable homotopy class.
Thus
\[
                         [S^6,e_i]=(0,\eta^3).
\]
Smashing the two sphere splittings gives
\[
 (S^3\times S^3)_+\simeq
       S^0\vee S^3\vee S^3\vee S^6.
\]
Under this splitting the external product
$[S^6\times S^6,e_1\times e_2]$ has four components.
The first three each contain a factor $[S^6]=0$; the last
is $\eta^3\eta^3=\eta^6=0$. Indeed $\eta^4$ belongs to
$\pi^{\mathrm S}_4=0$, as recorded in
\cite[Theorem 2.1 and Table 1]{rootIWX20}. Therefore
\[
 [S^6\times S^6,e_1\times e_2]=0
      \quad\hbox{in }\Omega^{\mathrm{fr}}_{12}(S^3\times S^3).
 \tag{E.B7}
\]
This calculation uses the product of the specific bounding
sphere framings, not just vanishing of ordinary characteristic
classes.

\emph{The comparison with the requested classifying map.}
By (E.B7) there is a compact stably framed thirteen-manifold
$Z$ with boundary $S^6\times S^6$ and a map
$E:Z\to S^3\times S^3$ extending $e_1\times e_2$.
Smooth it relative to the boundary and make it transverse
to $N$ there. Transversality on the boundary already holds:
the derivative of its defining equation is $df_1-df_2$,
which is onto. Relative transversality therefore supplies
the interior perturbation without changing that boundary.

The normal bundle of $N\subset S^3\times S^3$ is the
restriction of $h_1^*T\mathbb P^1$ pulled back from the first
sphere. This oriented two-plane bundle is trivial because
$H^2(S^3;\mathbb Z)=0$. Choose a framing on it before
restriction to $N$. Its pullback under $e_1\times e_2$
is a framing of $f_1^*T\mathbb P^1$, agreeing with that
in E.2.8 up to homotopy since $H^1(S^6;\mathbb Z)=0$.
It follows that
\[
                  V=E^{-1}(N),\qquad E|_V:V\longrightarrow N
\]
is a compact stably framed eleven-manifold with boundary
$(P,c)$ and the required boundary framing. This proves (E.B4).

The exact group-matrix and Euler calculation, together with
the two equivalent stable-product reductions, is recorded in
\texttt{code/shifted\_group\_bordism.py}.
The earlier disc-extension filling need not extend the
classifying map; (E.B7) supplies a different filling that does.
Neither construction computes the higher equivariant surgery
data needed to distinguish smooth ten-manifolds over $N$.
\end{proof}


\paragraph{Proposition E.2.10 (Proved: the fibre complement and a spinning obstruction).}
Let $F=T^4$ be a regular fibre of the original sphere fibration and
$E=S^6\setminus\operatorname{int}(F\times D^2)$, with the marking of
E.2.6. Its group and peripheral homomorphism are
\[
 \begin{split}
 K=\pi_1(E)&=\langle a,b,h\mid h\text{ central},\ a^3=h,
                                      \ b^4=h^{-1}\rangle,\\
 g&\longmapsto h,\qquad u,w,d\longmapsto1,
       \qquad \mu\longmapsto ab.
 \end{split}
 \tag{E.S1}
\]
Here $\mu$ is the normal meridian, with the peripheral orientation
chosen so that $[a]=4[\mu]$, $[b]=-3[\mu]$ and $[h]=12[\mu]$.
The image of the boundary group is the embedded copy
$\langle h,\mu\rangle\simeq\mathbb Z^2$. In the unique homotopy
class of normal framings for which every pushed-off loop of $F$
has linking number zero, the image of $g$ is
\[
             \lambda=h\mu^{-12},
 \qquad \ker\{\pi_1(F)\longrightarrow K\}
                 =\mathbb Z\langle u,w,d\rangle.
 \tag{E.S2}
\]
In particular the peripheral structure is that of the $(3,4)$
torus knot, with its regular Seifert fibre rather than its preferred
longitude as the original $g$ coordinate. Nevertheless the embedded
four-torus $F\subset S^6$ is not the ordinary threefold closed spin
of that classical knot, even allowing ambient orientation reversal
and arbitrary changes of its normal framing. More strongly, it is
not a single ordinary closed spin of any embedded $T^3\subset S^5$.

\begin{proof}
Apply the van Kampen presentation in B.5.2 with the regular base disc
removed. The cusp kills $w,d$ and its peripheral loop. Conjugating
$w$ by the order-three generator also kills $u$. The two finite caps
impose $a^3=h$ and $b^4=h^{-1}$; there is now no relation $ab=1$,
since the outer boundary has not been filled. The outer meridian
is $ab$ in the stated orientation. This proves (E.S1), including
the actual peripheral map, rather than only an abstract group
isomorphism. Eliminating $h$ gives $\langle a,b\mid a^3b^4=1\rangle$.
Its quotient by $h$ is $C_3*C_4$. The word $ab$ has infinite order
in that free product, so a relation $h^r\mu^s=1$ forces $s=0$.
Abelianization then forces $r=0$. This proves peripheral injectivity.
It also shows that $\lambda$ has infinite order: its image in
$C_3*C_4$ is $(ab)^{-12}$.

Write $\theta\in H^1(\partial E;\mathbb Z)$ for the normal-circle
covector and retain $(\gamma,u,w,\delta)$ for $H^1(F;\mathbb Z)$.
The generator $\xi\in H^1(E;\mathbb Z)$ restricts to
$\theta+12\gamma$. Homotopy classes of normal framings differ by
$H^1(F;\mathbb Z)$. Consequently there is a unique framing for
which the associated section $s:F\to\partial E$ satisfies
$s^*\xi=0$: it has $s^*\theta=-12\gamma$. Its induced peripheral
map is exactly (E.S2). Simultaneously reversing the meridian
conventions changes the displayed signs and preserves all the
following assertions.

For this section set
\[
 \mathcal R=\operatorname{im}\{H^2(E;\mathbb Z)
                         \longrightarrow H^2(F;\mathbb Z)\},
 \qquad
 \operatorname{rad}(\mathcal R)=
 \{v\in H_1(F;\mathbb Z):\iota_v\omega=0
                                    \text{ for all }\omega\in\mathcal R\}.
 \tag{E.S3}
\]
This uses the canonical identification of the cohomology of a torus
with the exterior algebra of its first cohomology. Equations
(E.MV4)--(E.MV6) give classes whose boundary restrictions are
$\theta u$ and $\theta w$: their unswept terms are zero because
$u\gamma uw=w\gamma uw=0$. Restricting them along $s$ proves
\[
                -12\gamma u,\ -12\gamma w\ \in\mathcal R.
 \tag{E.S4}
\]
The simultaneous contraction kernel of these two forms is
$\mathbb Z d$. Indeed their contractions force successively
$\gamma(v)=u(v)=w(v)=0$. Thus
\[
       \operatorname{rad}(\mathcal R)\subseteq\ker\gamma
        =\ker\{\pi_1(F)\longrightarrow K\}.
 \tag{E.S5}
\]
Both the zero-linking framing and this inclusion are invariants of
the embedded pair, with no choice of the original torus coordinates.

We define the comparison construction precisely. For a closed
submanifold $L\subset\operatorname{int}D^n$, its ordinary closed
spin is $L\times S^1$ in
$D^n\times S^1\subset\partial(D^n\times D^2)=S^{n+1}$.
Start with the closed classical knot $K_0\subset\operatorname{int}D^3$
and repeat this operation three times, each time placing the
previous compact pair in the interior of a disc. This is the
ordinary torus spinning construction reviewed in
\cite[Section 3, immediately before Theorem 3.4]{mapsWinter08};
no turning, twisting, or family of deformations is inserted.
If $\Sigma$ is a Seifert surface for $K_0$, the same three spins
give the embedded Seifert hypersurface
$\Sigma\times T^3$ with boundary $K_0\times T^3$.
Pushing the latter off using the normal of this hypersurface
gives its zero-linking framing.

Let $v$ be the original knot-circle class and let $v_1,v_2,v_3$
be the spinning circles. Each pushed-off two-cycle $v\wedge v_i$
bounds the corresponding copy of $\Sigma\times S^1$ in a parallel
copy of that hypersurface, disjoint from the knot. Hence every
restriction of a complement two-class vanishes on $v\wedge H_1(F)$:
$v$ belongs to $\operatorname{rad}(\mathcal R)$.
On the other hand its pushed-off loop is the original preferred
longitude and is nontrivial for $K_0=T(3,4)$, whereas the spinning
circles are trivial in the complement group. These peripheral
claims also follow directly at each spin: its complement is
\[
 ((D^n\setminus\operatorname{int}\nu L)\times S^1)
       \ \cup_{S^{n-1}\times S^1}\ (S^{n-1}\times D^2),
\]
so van Kampen kills precisely the new spinning circle and preserves
the old complement group. They agree with the primary peripheral
calculation \cite[Theorem 4.5]{mapsWinter08} in the first spin.
Thus the ordinary spun torus violates (E.S5). An identification
of embedded pairs would preserve its two sides, giving a
contradiction.

For the stronger assertion, let $L=T^3\subset S^5$ be arbitrary
and suppose that its ordinary closed spin were our four-torus.
Write $E_L$ for its exterior. Deleting a ball disjoint from
$L$ changes cohomology only in degrees at least four. In the
displayed spun-complement decomposition, with $n=5$, the cap
and intersection are $S^4\times D^2$ and $S^4\times S^1$.
Mayer--Vietoris and the surjectivity onto the latter first
cohomology therefore give, in degrees two and three, precisely
the cohomology of $(E_L\setminus D^5)\times S^1$.
Let $s$ be the spinning-circle covector and $v_s$ its primitive
homology generator. Put $\mathcal R_L^j=\operatorname{im}
\{H^j(E_L)\to H^j(L)\}$, using the zero-linking framing.
One has $\mathcal R_L^1=0$. Also $\mathcal R_L^3=0$:
the pushed-off $L$ bounds a parallel Seifert hypersurface in
its exterior, so every top-degree restriction evaluates to
zero. The product restrictions consequently satisfy
\[
       \mathcal R^2=\mathcal R_L^2,\qquad
       \mathcal R^3=\mathcal R_L^2\wedge s.
 \tag{E.S12}
\]
In particular $v_s\in\operatorname{rad}(\mathcal R^2)$,
and contraction by $v_s$ is injective on $\mathcal R^3$.
For our actual four-torus, (E.S4) forces $v_s=\pm d$.
But (E.MV4)--(E.MV6), now applied to $\alpha=\gamma\delta$
in degree three, show that $\mathcal R^3$ contains
$\pm2\gamma uw$. Its swept term vanishes in the zero-linking
framing, and this is a nonzero class annihilated by $d$.
This contradicts that injectivity. No ordinary closed spin
is possible. Turned or deform-spun constructions insert
additional data and are not excluded by this argument.
\end{proof}

\paragraph{Theorem E.2.11 (Proved: exact second-homotopy coinvariants).}
For the classifying map $c:P\to N=K(G,1)$ of E.2.9, one has
\[
 c_*:H_3(P;\mathbb Z)\longrightarrow H_3(N;\mathbb Z)=\mathbb Z
       \quad\text{equal to zero},
 \qquad
 \bigl(\pi_2(P)\bigr)_G
       \simeq\mathbb Z^2\oplus(\mathbb Z/12)^6.
 \tag{E.S6}
\]
The pullback of a generator of $H^3(N;\mathbb Z)$ has exact order
twelve in $H^3(P;\mathbb Z)$. In particular $\pi_2(P)$ is not a
projective $\mathbb ZG$-module. The fivefold cannot be diffeomorphic,
or homotopy equivalent, to $N\times S^6$ connected-summed with any
finite collection of $S^r\times S^{10-r}$, $2\leq r\leq8$, and
any homotopy ten-sphere. This computes coinvariants and a definite
module obstruction; it does not assert a presentation or a
classification of the full $\mathbb ZG$-module.

\begin{proof}
Use the separating base circle in E.2.6 also in the target of
E.2.9. There $M$ denotes the exterior of a regular orbit of the
weighted Hopf map, namely the torus-knot exterior of A.1.3. The
two decompositions, and the restriction of $c$, are
\[
 \begin{array}{ccc}
 (E\times F')\cup(F\times E')&\longrightarrow&
                  (M\times S^1)\cup(S^1\times M')=N,\\
 F\times F'\times S^1&\longrightarrow&S^1\times S^1\times S^1.
 \end{array}
 \tag{E.S7}
\]
On the common boundary the map retains exactly the covectors
$\gamma_1,\gamma_2,\theta$. It is the literal quotient there.
Each knot exterior has the integral homology of a circle, so the
third homology and cohomology of each target piece in (E.S7) vanish.

We first check the homological statement, retaining the full
integer boundary maps. In the degree-two source of (E.MV7)--(E.MV8)
write $t_1,t_2,t_{12}$ for the coefficients of
$\theta\times g_1$, $\theta\times g_2$, and $g_1\times g_2$.
The image of $i_2:H_2(F)\to H_2(E)$ lies entirely in the
$d$-sweep coordinate, by (E.MV6). Consequently the $g_1$-sweep
coordinate in the first $H_2(E)$ summand of the target is just
$t_1$, and the corresponding second coordinate is just $t_2$.
Every vector in the kernel of the boundary matrix therefore has
$t_1=t_2=0$. Its first $H_1(E)\otimes\mathbb Zg_2$ coordinate
then forces $12t_{12}=0$, up to an immaterial orientation sign.
These three coordinates are the only ones retained in
$H_2(T^9)\to H_2(T^3)$. The target Mayer--Vietoris sequence
injects $H_3(N)$ into $H_2(T^3)$; its source-piece third groups
are zero. Naturality now proves $c_*H_3(P)=0$.

The cohomological calculation records the extension which this
vanishing alone would miss. On the target boundary write
$v=\gamma_1\gamma_2$. The generators of the two degree-two
piece cohomology groups restrict, with consistent meridian
orientations, to
\[
       \theta\gamma_2+12v,
       \qquad \theta\gamma_1+12v.
 \tag{E.S8}
\]
Their image is a primitive rank-two subgroup of $H^2(T^3)$:
the minor in the last two coordinates is a unit. Its quotient
is freely generated by $v$. Since the degree-three piece groups
are zero, the connecting class $\partial v$ is a generator of
$H^3(N;\mathbb Z)$.

In the source decomposition both $\theta\gamma_1$ and
$\theta\gamma_2$ already belong to the image of the degree-two
piece restrictions. This is (E.MV4) with $\alpha=\gamma$,
since $\phi_2(\gamma)=0$; the second normal circle is reversed,
which only changes its sign. Subtracting one of these classes
from the corresponding product of a degree-one class of $E$
with $\gamma_2$ proves that $12v$ is a source boundary image.
Conversely the coefficient of $v$ in every such boundary image
is divisible by twelve. The pure degree-two terms from a single
factor have no $v$ coefficient, and the mixed terms use the
restriction $\theta+12\varepsilon\gamma$ of the degree-one
generator, with $\varepsilon=\pm1$. It follows that
\[
                  c^*(\partial v)=\partial v
        \quad\text{has exact order }12.
 \tag{E.S9}
\]
This uses the actual Mayer--Vietoris image; no splitting of a
torsion filtration has been assumed.

Put $A=(\mathbb Z/12)^7$ and $L=(\pi_2(P))_G$.
The low-degree homology sequence of the universal-cover
fibration, together with E.2.6 and E.2.9, gives
\[
             0\longrightarrow\mathbb Z z\longrightarrow L
                \longrightarrow\mathbb Z\oplus A\longrightarrow0.
 \tag{E.S10}
\]
The zero on the left uses $c_*H_3(P)=0$, just proved.
To determine this extension, replace $c$ by its mapping cylinder
and use relative groups for $(N,P)$. The induced map on fundamental
groups is an isomorphism and the target has zero second homotopy
group. On universal covers the relative pair is two-connected;
relative Hurewicz gives
$H_3(\widetilde N,\widetilde P)=\pi_2(P)$ and all lower relative
groups vanish. The relative covering homology sequence therefore
identifies $H_3(N,P)=L$ and $H_2(N,P)=0$.
The universal coefficient theorem and the relative cohomology
sequence identify the map
\[
       H^3(N,P)=\operatorname{Hom}(L,\mathbb Z)
         \longrightarrow H^3(N)=\mathbb Z
\]
with evaluation on $z$. Its image is $12\mathbb Z$ by (E.S9).

The free summand of the quotient in (E.S10) splits. Lift its
seven order-twelve generators to $x_1,\ldots,x_7$ in the
remaining extension. There are integers $e_i$, defined modulo
twelve, with the complete abelian presentation
\[
        \langle z,x_1,\ldots,x_7\mid
                      12x_i=e_i z\ (1\leq i\leq7)\rangle.
 \tag{E.S11}
\]
Evaluation of homomorphisms to $\mathbb Z$ on $z$ has image
$\bigl(12/\gcd(12,e_1,\ldots,e_7)\bigr)\mathbb Z$.
Thus that gcd is one. The relation matrix in (E.S11) has
Smith factors $1,12,12,12,12,12,12$, giving
$\mathbb Z\oplus(\mathbb Z/12)^6$. Adding the split free
summand proves (E.S6).

If $\pi_2(P)$ were projective over $\mathbb ZG$, it would be
a direct summand of a free module. Taking coinvariants preserves
that splitting, so $L$ would be a projective, hence free,
abelian group. The torsion in (E.S6) contradicts this. Finally
$N\times S^6$ has zero second homotopy group. Connected sum in
dimension ten with a simply connected sphere product adds its
second homotopy as a free induced $\mathbb ZG$-module, as seen
by lifting the connected-sum necks and using their simply
connected boundary spheres. Only $r=2,8$ contribute, with one
free generator each; a homotopy sphere contributes none.
Their second homotopy modules are therefore free, so none is
homotopy equivalent to $P$.

The program \texttt{maps\_shifted\_pi2.py} checks the simultaneous
contraction kernel, the three homological boundary coordinates,
and the exact-order-twelve image certificate for all sixteen
independent boundary-adjoint sign choices in E.2.6.
\end{proof}

\paragraph{Theorem E.2.12 (Proved: central localizations of the homotopy module).}
Retain the central generators $h_1,h_2$ of $G$ from E.2.9 and put
\[
 R=\mathbb ZG,\qquad R_i=R[(h_i-1)^{-1}],\qquad
              R_{12}=R[(h_1-1)^{-1},(h_2-1)^{-1}].
 \tag{E.S13}
\]
These are flat central localizations. As modules over the indicated
localized group rings, the actual second homotopy group satisfies
\[
 \begin{split}
 R_1\otimes_R\pi_2(P)&\simeq R_1/(h_2-1),\\
 R_2\otimes_R\pi_2(P)&\simeq R_2/(h_1-1),\\
 R_{12}\otimes_R\pi_2(P)&=0.
 \end{split}
 \tag{E.S14}
\]
More generally, $H_j(P;R_1)$ is zero outside $2\leq j\leq7$;
in those six degrees it is respectively a direct sum of
\[
                  1,\ 3,\ 4,\ 4,\ 3,\ 1
 \tag{E.S15}
\]
copies of $R_1/(h_2-1)$. The analogous statement holds with the
indices interchanged. All homology with coefficients in $R_{12}$
vanishes. For the fibre complement of E.2.10, with
$S=\mathbb ZK[(h-1)^{-1}]$, one likewise has
\[
             S\otimes_{\mathbb ZK}\pi_2(E)\simeq S,
       \qquad H_4(E;S)\simeq S,
 \tag{E.S16}
\]
and its other localized homology groups vanish. These are
localizations of the actual group-ring modules, not a replacement
of $G$ or $K$ by their abelianizations. Moreover, if
$q=(h_1-1)(h_2-1)$, then
\[
                  q^2\pi_2(P)=0,
        \qquad \dim_{\mathbb Q}\bigl(\pi_2(P)\otimes\mathbb Q\bigr)
                     =\infty.
 \tag{E.S20}
\]
Consequently the universal cover has no finite homotopy type.
There is no fibration of $P$ over $N$ with compact manifold fibre
and inducing an isomorphism on fundamental groups.

\begin{proof}
\emph{The cusp module with its marked deck transformation.}
In the original cusp, the angular generators $w,d$ vanish and
$\pi_1(N_0)=\mathbb Z\langle g,u\rangle$. In $E$ their images
are respectively $h,1$; in the first factor of $P$ they are
$h_1,1$. Thus the connected cover of this cusp occurring inside
the universal cover is the cover $\widehat W$ of its central
fibre in which the $g$ deck transformation is $z$, whereas $u$
is retained as a loop. Write $A=\mathbb Z[z,z^{-1}]$.
Its ordinary homology modules are
\[
 H_0(\widehat W)=H_1(\widehat W)=A/(z-1),\quad
 H_2(\widehat W)=A\oplus(A/(z-1))^2,
 \quad H_3(\widehat W)=0,\quad H_4(\widehat W)=A.
 \tag{E.S17}
\]
Here is an integral normalization calculation of this input.
The dual complex of the cover is an infinite cylinder, whose
ordinary homology is $A/(z-1)$ in degrees zero and one and zero
in degree two. The normalization degree-four row is a free
copy of $A$. In degree two, the homology map from the three
double curves to the four component classes is the transpose,
up to the unimodular intersection marking and possible inversion
of $z$, of
\[
 B(z)=\begin{pmatrix}
 -1&-1&-1&-1\\ 1&1&1&z\\-z&-z&-1&-z
 \end{pmatrix}
 \ \sim\
 \begin{pmatrix}1&0&0&0\\0&z-1&0&0\\0&0&z-1&0\end{pmatrix}.
 \tag{E.S18}
\]
For example multiply on the left by
$\left(\begin{smallmatrix}-1&0&0\\1&1&0\\-z&0&1\end{smallmatrix}\right)$,
then use as new columns $e_1,e_4-e_1,e_3-e_1,e_2-e_1$.
Both operations are unimodular over $A$. The transpose is
injective and its cokernel is the displayed degree-two module.
All strata have only even homology, and the normalization
direction has length two, so no higher differential is possible.
There is no degree-two extension from the dual complex, whose
second homology is zero. This proves (E.S17). Ordinary chains
in this locally finite normalization have compact support and
use direct sums, so no compactly supported cohomology has been
substituted for ordinary homology. The exact matrix certificate
is \texttt{manifolds\_cusp\_alexander.py}, contributed by the
six-manifolds worker and independently rerun here.

\emph{Localization and the vanishing of every other piece.}
The central subgroup $\langle h_1,h_2\rangle$ is an injected
$\mathbb Z^2$ by E.2.9. The group ring is free over its central
Laurent subring, by choosing coset representatives. In particular
$h_i-1$ and their products are non-zero-divisors. Inverting them
is central localization and is flat; it commutes with homology
of the universal-cover chain complex. The analogous assertions
hold for $h$ in $K$.

Use the regular torus bundle over a graph and the six closed
cap neighborhoods in E.2.3. For $R_1$ coefficients, the first
regular torus has its $g_1$ loop acting by $h_1$. The torus
chain complex is the Koszul complex of its commuting loop
actions. Its $h_1-1$ term is invertible, providing a chain
contraction. Thus every regular fibre has zero homology with
these coefficients. The homology sequence of a bundle over
a graph, and the same calculation over each boundary circle,
show that the whole regular piece and every cap boundary are
acyclic. Each cap of the second factor is a product with a
first-factor regular torus over a disc, so it too is acyclic.

At a finite cap of the first factor, the affine generator
$a_j$ has $a_j^{m_j}=h_1^{\pm1}$ in its image in $G$.
The connected cover of the reduced flat surface corresponding
to its map onto $\langle a_j\rangle$ is $T^3\times\mathbb R$.
The $m_j$-th power acts by an integral translation with linear
part $A_j^{m_j}=I$, and therefore acts identically on all of
its homology. Thus $h_1-1$ annihilates every such homology
module before localization. Retaining the other regular torus
does not change that assertion, and induction to $G$ commutes
with the central action. Both first-factor finite caps are
therefore acyclic after localization as well. This argument
keeps their actual affine translation, rather than treating
the quotient as an untwisted torus product.

\emph{The one surviving cap and induction.}
The first cusp is $N_{0,1}\times F_2$, since the two critical
triples are disjoint. Its image in $G$ is precisely
$H=\langle h_1,h_2\rangle$. A connected component of its
preimage in the universal cover is
\[
                \widehat W\times\mathbb R\times T^3
\]
up to deformation retraction: $g_1$ acts on $\widehat W$ by
$z=h_1$, and $g_2=h_2$ translates the real factor. It acts
trivially on the homology of the retained three-torus. All
components together give induction from $\mathbb ZH$ to
$R$. By (E.S17), inverting $h_1-1$ leaves one free
$\mathbb Z[h_1^{\pm1},(h_1-1)^{-1}]$ summand in cusp degrees
two and four and nothing else. The integral product homology
with $T^3$ has no torsion correction. Inducing the result gives
\[
 H_j(N_{0,1}\times F_2;R_1)
   \simeq
 \left(R_1/(h_2-1)\right)^{
           \binom3{j-2}+\binom3{j-4}},
 \tag{E.S19}
\]
where a binomial coefficient is zero outside its usual range.
All other pieces and their intersections are acyclic, as
proved above, so Mayer--Vietoris identifies this with
$H_j(P;R_1)$. This proves (E.S15) and (E.S19), and interchanging
the factors proves the second version. Inverting $h_2-1$
also kills (E.S19), proving the double-localization assertion.
Since $\widetilde P$ is simply connected, Hurewicz identifies
its second homology with $\pi_2(P)$; flat localization now
gives (E.S14).

For $E$ the same decomposition has only one cusp and two
finite caps. Every other piece is acyclic after $h-1$ is
inverted by the identical argument. Its cusp image is
$\langle h\rangle$, so inducing (E.S17) after localization
gives one free $S$ summand in degrees two and four. This
proves (E.S16).

\emph{A uniform annihilator without localization.}
Put $I=(h_1-1,h_2-1)\subset R$. Before localization, the homology
of a regular eight-torus with $R$ coefficients is
\[
                 (R/I)\otimes_{\mathbb Z}\bigwedge^j\mathbb Z^6
                      \quad(0\leq j\leq6).
 \tag{E.S21}
\]
Indeed the two nontrivial loop images are the independent central
generators $h_1,h_2$. Their differences from one form a regular
sequence in the central Laurent ring, and $R$ is free over that
ring. The two-variable Koszul complex therefore has only its
degree-zero quotient as homology; the other six loop differentials
are zero. This proves (E.S21).
The homology sequence of the regular bundle over a graph has
only base degrees zero and one. Both filtration quotients are
annihilated by $I$, by (E.S21), so $I^2$ annihilates the homology
of the regular piece. The same two-stage Wang argument applies
to every boundary mapping torus. It does not require splitting
either filtration or treating the monodromy as trivial.

At a finite cap its connected lift is a product of six torus
directions and two real directions up to deformation retraction.
Each central deck generator acts trivially on homology: the
one belonging to that cap is the finite affine power already
checked above, and the other translates the regular real
direction. Hence $I$ annihilates its induced homology. At the
first cusp, $h_2$ translates the real factor of the other
regular torus, so $h_2-1$ annihilates all of its homology;
the second cusp is annihilated by $h_1-1$. Thus $q$ annihilates
the homology of every piece and every intersection, since
$q\in I^2$. The decomposition has a regular piece and disjoint
caps with disjoint boundary collars, so its Mayer--Vietoris
sequence gives in every degree
\[
 0\longrightarrow\operatorname{coker}D_j
   \longrightarrow H_j(P;R)
   \longrightarrow\ker D_{j-1}\longrightarrow0.
\]
Both end groups are killed by $q$, and therefore the middle
group is killed by $q^2$. Since $H_2(P;R)=\pi_2(P)$, this
proves the first assertion of (E.S20) integrally.

\emph{Infinite rational rank.}
If $\pi_2(P)\otimes\mathbb Q$ were finite dimensional,
localizing its linear automorphism $h_1$ at $h_1-1$ would
remain finite dimensional. Explicitly, this removes the
generalized eigenspace with eigenvalue one and inverts an
automorphism on the remaining finite-dimensional quotient.
But (E.S14), tensored with $\mathbb Q$, contains a copy of
$\mathbb Q[h_1^{\pm1},(h_1-1)^{-1}]$, as follows by taking
the identity coset in the free module over the central
Laurent subring and setting $h_2=1$. That vector space is
infinite dimensional, a contradiction. This proves the
second assertion and rules out a finite homotopy type of
the universal cover. A fibration over the aspherical $N$
inducing a fundamental-group isomorphism would have simply
connected fibre, and pulling it to $\widetilde N$ would
identify that fibre up to homotopy with $\widetilde P$.
A compact manifold fibre has finite homotopy type, so such
a fibration is impossible.

The quotient $R_1/(h_2-1)$ is nonzero, since the two central
generators are independent, and it is annihilated by the
non-zero-divisor $h_2-1$. A projective $R_1$-module is a
direct summand of a free module and has no such nonzero
annihilated element. Thus (E.S14) also gives an independent
proof of the nonprojectivity in E.2.11. It does not determine
the part of $\pi_2(P)$ supported where both central
elements specialize to one, or its remaining extension data.
\end{proof}

\paragraph{Theorem E.2.13 (Proved: an actual circle fibration and infinite-order return).}
The shifted fivefold admits a smooth fibre bundle
\[
                  Y^9\longrightarrow P\stackrel{\rho}{\longrightarrow}S^1
 \tag{E.S22}
\]
with connected closed fibre, inducing the primitive character
\[
 \chi(a_1,a_2,b_1,b_2,h_1,h_2)
                    =(4,-3,-4,3,12,-12)
 \tag{E.S23}
\]
on the presentation of E.2.9. Its fibre group is the fundamental
group of the oriented principal circle bundle of Euler number
$-1$ over a surface of genus six. Here the base has its complex
orientation and the positive circle is $h_1h_2$; reversing that
circle changes the Euler number to $+1$. In particular
\[
 \pi_1(Y)=\ker\chi,\qquad H_1(Y;\mathbb Z)=\mathbb Z^{12}.
 \tag{E.S24}
\]
The return diffeomorphism of (E.S22) has infinite order even
among unbased homotopy classes of self-equivalences of $Y$,
although its induced outer automorphism of $\pi_1(Y)$ has
exact order twelve. No connected finite cover of $P$ is
homotopy equivalent to a product $S^1\times Z$.
These assertions give a geometric mapping-torus description;
they do not identify the smooth nine-dimensional fibre or
the full return diffeomorphism up to classification.

\begin{proof}
\emph{A closed one-form with a nonzero derivative at every point.}
Use the actual smooth map $c=(e_1,e_2):P\to N$ in E.2.9,
not an arbitrary representative of its classifying homotopy
class. The manifold $N$ has the locally free $T^2$ action
induced by the two weighted Hopf actions. Take a closed
one-form representing the primitive integral class $\chi$
on $N$ and average it over this compact torus. Its average
$\alpha$ is closed, represents the same integral class,
and is invariant. If $V_i$ is the fundamental vector field
with time-one orbit $h_i$, Cartan's formula gives
$d(\alpha(V_i))=0$. Integrating along an ordinary orbit
therefore gives
\[
                  \alpha(V_1)=12,\qquad \alpha(V_2)=-12.
 \tag{E.S25}
\]
The integrality of the periods integrates $\alpha$ to a
smooth map $\rho_N:N\to\mathbb R/\mathbb Z$.

Choose the two cusp replacements defining $e_i$ in disjoint
base discs, themselves disjoint from the other factor's
critical values. At every point of $P$, at least one factor
is a regular torus at which its $e_i$ is the literal character
quotient of A.1.1. In a local flat marking of that factor,
the constant vertical vector $g_i$ maps to $V_i$. Holding
the other factor fixed gives a tangent vector to $P$, and
\[
                     (c^*\alpha)(g_i)=12\quad\hbox{or }-12.
\]
Thus $c^*\alpha$ is nowhere zero. This is a local derivative
calculation on the actual map, not an inference from its
fundamental-group character. The map $\rho=\rho_Nc$ is a
proper submersion. Equivalently, choose a vector field $W$
with $(c^*\alpha)(W)=1$; its complete flow supplies local
trivializations and a return diffeomorphism. The fibre is
connected because $\chi$ is primitive: the image in
$\pi_1(S^1)$ of a connected component of a bundle's total
space has index equal to the number of fibre components.
The homotopy sequence now proves $\pi_1(Y)=\ker\chi$.

\emph{The exact surface cover and its circle Euler class.}
Write $\Gamma=G/\langle h_1,h_2\rangle$ for the oriented
orbifold group of $S^2(3,4,3,4)$. The character modulo
twelve descends to the surjection
\[
 \overline\chi:\Gamma\longrightarrow\mathbb Z/12,
       \qquad(a_1,a_2,b_1,b_2)\longmapsto(4,-3,-4,3).
 \tag{E.S26}
\]
Its restriction to every cone group is injective. Hence its
kernel $\Gamma_0$ is torsion free and defines a degree-twelve
ordinary surface cover of the orbifold. Its Euler characteristic
is $12(-5/6)=-10$, so that surface is $\Sigma_6$. The image
of $\ker\chi$ in $\Gamma$ is exactly $\Gamma_0$: a lift of
an element of $\Gamma_0$ has character a multiple of twelve,
which can be cancelled by a power of $h_1$. Its central
kernel is the primitive subgroup
\[
 \ker\chi\cap\langle h_1,h_2\rangle
                =\mathbb Z\langle h_1h_2\rangle.
 \tag{E.S27}
\]

The circle Euler number can be fixed without an ambiguous
Seifert sign convention. On either weighted Hopf sphere put
$z_1=\cos r\,e^{i\varphi_1}$ and
$z_2=\sin r\,e^{i\varphi_2}$. For the positive action with
weights $(4,3)$ the smooth invariant connection is
\[
 \vartheta=\tfrac14\cos^2r\,d\varphi_1
                    +\tfrac13\sin^2r\,d\varphi_2.
\]
It evaluates to one on the generator with period $2\pi$.
With the boundary orientation of the unit ball in $\mathbb C^2$,
\[
             \frac1{(2\pi)^2}\int_{S^3}\vartheta\wedge d\vartheta
                         =\frac1{12}.
 \tag{E.S28}
\]
For the associated complex line whose unit circle is this
positive circle bundle, the convention agreeing with the
tautological Hopf line is $c_1=-[d\vartheta/(2\pi)]$.
Thus each orbifold circle Euler number is $-1/12$.
Both shifted base identifications preserve the complex
orientation. Pullback to $\Sigma_6$ multiplies the two Euler
numbers by twelve, giving $(-1,-1)$; this is also the
covering formula of \cite[Theorem 3.6]{rootScott83}.
In the integral fibre basis
\[
                  v=(1,1),\qquad w=(0,1)
\]
the determinant is one, the Euler vector is $-v$, and
$\chi(v)=0$, $\chi(w)=-12$. Consequently the circle kernel
has Euler number $-1$ with generator $h_1h_2$.

One can see this bundle directly in $N$. Invariance of
$\alpha$ implies
$\rho_N((s_1,s_2)x)=\rho_N(x)+12(s_1-s_2)$, where
$s_i\in\mathbb R/\mathbb Z$. A level contains twelve
diagonal circles over an ordinary base point. The peripheral
permutations (E.S26) give exactly the connected cover
$\Sigma_6$; they have no stabilizers at the cone points.
The diagonal circle acts freely, since at every cone only
one of the two Hopf actions has a nontrivial stabilizer.
The level is therefore the circle bundle just computed.
This identifies its fundamental group with $\ker\chi$
and gives an actual three-dimensional aspherical model.
The homology sequence of a circle bundle of Euler number
$\pm1$ gives its first homology $\mathbb Z^{12}$, proving
the second assertion in (E.S24).

\emph{Faithful central action and the return obstruction.}
The centre of $G$ is exactly $\langle h_1,h_2\rangle$:
the quotient is a nonelementary oriented hyperbolic orbifold
group with trivial centre. This centre acts faithfully on
$\pi_2(P)$. Indeed if $h_1^a h_2^b$ acts as the identity,
localize first at $h_1-1$ and apply E.2.12. On
$R_1/(h_2-1)$ its action is multiplication by $h_1^a$.
This module is free over the central ring
$\mathbb Z[h_1^{\pm1},(h_1-1)^{-1}]$ after choosing coset
representatives, so this is the identity only if $a=0$.
The symmetric localization then gives $b=0$.

If $H<G$ has finite index, its image in $\Gamma$ is again
nonelementary. Its centralizer in $\Gamma$ is trivial:
commuting with two hyperbolic elements with distinct axes
fixes at least three boundary points and hence is the
identity. It follows that
$Z(H)=H\cap Z(G)$. A homotopy product $S^1\times Z$ for
the corresponding finite cover would supply an infinite
central element acting trivially on $\pi_2$, namely the
circle factor; this contradicts the faithful action just
proved. If a positive power of the return of (E.S22) were
homotopic to the identity, its mapping torus would be
homotopy equivalent to $S^1\times Y$. That mapping torus
is the corresponding finite cyclic cover of $P$, so this
too is impossible.

Finally take $t=a_1a_2$, for which $\chi(t)=1$, as a
stable letter. The image of $\ker\chi$ in $\Gamma$ is
$\Gamma_0$, so the same centralizer argument gives
$C_G(\ker\chi)=Z(G)$. The automorphism
$\operatorname{Ad}(t^n)|_{\ker\chi}$ is inner precisely
when $t^n$ differs from an element of $\ker\chi$ by a
central element. Since $\chi(Z(G))=12\mathbb Z$, this
holds precisely for $12\mid n$. Thus the outer action has
order twelve, while the actual return has infinite homotopy
order. This return is the circle fibration of the shifted
ten-manifold; it is distinct from the six-sphere parameter
returns computed in B.6.24.

The exact integral character, cone orders, primitive kernel
basis and Euler calculation are recorded in
\texttt{code/maps\_shifted\_circle.py}. No identification
of $Y$ with a product or connected sum is used here.
\end{proof}

\paragraph{Theorem E.2.14 (Proved: the cyclic-cover splice and its separating torus twist).}
Let $E_{12}\to E$ be the connected cyclic cover defined by
reduction modulo twelve of the positive meridian character.
The circle fibration of E.2.13 may be chosen with fibre
\[
 Y=(E_{12}\times T^3)\ \cup_{T^8}\ (T^3\times E_{12}),
 \tag{E.S29}
\]
with the following exact boundary marking. Write
$\theta_i'$ for the lifted normal circle and retain the
two regular-torus markings. Then
\[
       \theta_1'=\gamma_2-\gamma_1,\qquad
       \theta_2'=\gamma_1-\gamma_2.
 \tag{E.S30}
\]
There is a return representative $\psi$ whose twelfth power
is the separating torus twist
\[
 \begin{gathered}
 T(x,s)=(x+\eta(s)(g_1+g_2),s)\quad\hbox{on }T^8\times[-1,1],
 \\
 T=\operatorname{id}\text{ elsewhere}.
 \end{gathered}
 \tag{E.S31}
\]
Here $\eta$ is zero near $-1$ and one near $1$, and circle
coordinates have period one. In particular (E.S31) is the
identity near both ends of its collar. Its rational
cohomological variation satisfies
\[
 (T^*-I)^2=0,\qquad
 \operatorname{rank}(T^*-I)|_{H^j(Y;\mathbb Q)}=
 \begin{cases}1&j=4,5,\\0&j\ne4,5.\end{cases}
 \tag{E.S32}
\]
The rational Betti numbers of the actual fibre are
\[
       (b_0,\ldots,b_9)(Y)
          =(1,12,121,407,693,693,407,121,12,1).
 \tag{E.S33}
\]
Every eigenvalue of $\psi^*$ is a twelfth root of unity.
At eigenvalue one there is exactly one Jordan block of
size two in each of degrees four and five; all its other
blocks, and all blocks at other eigenvalues, have size one.
Thus the infinite order is already visible on ordinary
rational cohomology. The smooth complement $E_{12}$ in
(E.S29) remains the actual finite cover of $E$, not a
replacement inferred from its group or homology.

\begin{proof}
\emph{Adapting the fibration to the actual splice.}
Cut the base along the circle used in E.2.6, with the two
critical triples on opposite sides. Over the first disc the
second factor is a regular marked $T^4$, and over the second
disc the first factor is regular. Let
$\xi_E:E\to\mathbb R/\mathbb Z$ represent the positive
meridian class. Its restriction to the boundary can be
chosen exactly
\[
                         \xi_E=\theta+12\gamma.
 \tag{E.S34}
\]
Indeed this is its integral restriction by E.2.10; homotopy
extension makes the representative equal to that linear
map on a collar. The two maps
\[
        \rho_A=\xi_{E_1}-12\gamma_2,
             \qquad \rho_B=12\gamma_1-\xi_{E_2}
 \tag{E.S35}
\]
agree exactly on the common boundary, where
$\theta_2=-\theta_1$. Smooth them with this common product
collar. The regular-factor circle derivative is $-12$ on
the first piece and $12$ on the second, so they give a
proper submersion throughout, including both cusp pieces.
These are also permissible choices in E.2.13: on the
quotient $N$, an invariant closed one-form differs from
these representatives only by the differential of a
single-valued function on each of the two pieces. Absorb
that function by translation in the unused regular circle.
Thus this construction uses the same actual quotient and
primitive character.

A level of $\xi_E-12\gamma_2$ is the pullback of the
degree-twelve circle cover along $\xi_E$, times the three
remaining regular circles. The pullback is $E_{12}$.
Its boundary is connected because the normal meridian has
character one. The covering of that boundary has
$\theta=12\theta'$ and unchanged $\gamma,u,w,\delta$.
Equation (E.S34) in the level condition gives precisely
(E.S30). This proves (E.S29) with its actual smooth
boundary maps.

Choose the vector field integrating (E.S35) to be
$-g_2/12$ on the first piece off its collar and $g_1/12$
on the second. On the common regular collar interpolate by
\[
                     W_s=\frac{\eta(s)g_1-(1-\eta(s))g_2}{12}.
 \tag{E.S36}
\]
It has $d\rho(W_s)=1$, and it preserves the collar
coordinate $s$ and the original base point. Its flow
therefore defines a return $\psi$. At time twelve it is
the identity off the collar, since the two indicated
regular circles have period one. On the collar it
translates by $\eta(s)g_1-(1-\eta(s))g_2$; removing the
integral translation $-g_2$ gives (E.S31). This comparison
uses a globally regular factor on each piece, so there
is no continuation through a singular cap to justify.

\emph{The variation map and the invariant boundary calculation.}
Put $D=T^8$ and let $r:H^j(Y;\mathbb Q)\to H^j(D;\mathbb Q)$
be restriction. Denote the Mayer--Vietoris connecting map
by $\partial:H^{j-1}(D;\mathbb Q)\to H^j(Y;\mathbb Q)$.
Writing $v=g_1+g_2$, the collar pullback of a translation
invariant form is
\[
       T^*\omega=\omega+d\eta\wedge\iota_v\omega.
\]
Average its restriction to the torus and change the global
representative by an exact form to make this equality
applicable on the collar. The difference represents the
connecting class, and hence, with the connecting sign
chosen by the increasing collar coordinate,
\[
                          T^*-I=\partial\,\iota_v\,r.
 \tag{E.S37}
\]
As $r\partial=0$, its square is zero.

No unknown part of $H^*(E_{12})$ is needed to compute
(E.S37). Its cyclic deck action is a translation on the
boundary torus, hence acts trivially on boundary cohomology.
Every noninvariant deck eigenspace therefore restricts to
zero. Rational transfer identifies the invariant part
with $H^*(E)$, whose marked restriction is E.2.6. Its
pullback to (E.S30) is
\[
 R_{k,1}(a)=\phi_k(a)+12(\gamma_2-\gamma_1)a,
 \qquad
 R_{k,2}(b)=\phi_k(b)+12(\gamma_1-\gamma_2)b.
 \tag{E.S38}
\]
Retain the other factor's three covectors $u,w,\delta$,
and take exterior products to obtain the two subspaces
$U_j^k\subset\bigwedge^k\mathbb Q^8$.
The image of $r$ is exactly $U_1^k\cap U_2^k$, and the
kernel of $\partial$ is $U_1^{k-1}+U_2^{k-1}$.
Consequently the rank in (E.S32) is exactly the rank of
\[
 U_1^k\cap U_2^k\ \stackrel{\iota_v}{\longrightarrow}
 \bigwedge^{k-1}\mathbb Q^8/(U_1^{k-1}+U_2^{k-1}).
 \tag{E.S39}
\]

Here are the full dimensions for this small rational
calculation. Let $d_k$ be twice the degree-$k$ dimension
of $H^*(E)\otimes H^*(T^3)$ and let
$r_k=\dim(U_1^k+U_2^k)$:
\[
 \begin{array}{c|rrrrrrrrr}
 k&0&1&2&3&4&5&6&7&8\\ \hline
 d_k&2&8&20&44&70&68&36&8&0\\
 r_k&1&8&19&34&47&46&27&8&0
 \end{array}
 \tag{E.S40}
\]
Contraction in (E.S39) has rank one only for
$k=4,5$ and is zero in every other degree. All entries
are computed by exterior products of (E.S38) in
\texttt{maps\_shifted\_return.py}; its output retains an
explicit nonzero boundary witness in each of degrees
four and five. One can normalize the remaining sign in
$\phi_2$ by the ring structure: with
$r_1(1)=\theta+12\gamma$, the equality
$r_1(1)r_2(\delta)=r_3(a)$ forces the coefficient
$c_2=-2$, and then $c_3=2$, in the notation of E.2.6.
This is also the normalization used by that executable.
No rank calculation treats a rational cokernel as an
integral lattice.

\emph{The finite-cover cohomology needed for the full Betti numbers.}
For clarity, the full rational cohomology of $E_{12}$
can also be obtained without an uncomputed extension.
Let $\mathcal L_\lambda$ be the rank-one system on $E$
whose meridian acts by $\lambda$, with $\lambda^{12}=1$.
On a regular torus it is trivial. Its three base transports
in vertical degree $q$ are
\[
 \lambda^4\bigwedge^q A_1^{-t},\qquad
 \lambda^{-3}\bigwedge^q A_2^{-t},\qquad
                         \bigwedge^q M_0^{-t}.
 \tag{E.S41}
\]
The exponents retain the actual affine meridian character;
using only the finite linear matrices would give a different
cover.

At a finite cap, averaging on its finite root cover
identifies the cap cohomology with the invariant part of
(E.S41). At the unchanged cusp the system is trivial,
because $g$ has character twelve and $u$ has character zero.
Its rational specialization is an isomorphism onto the
monodromy invariants in every degree. Here is an explicit
check of that last input. The cusp Betti numbers are
$(1,2,4,2,1)$ by its normalization calculation. In degree
two, the integral surjectivity of the specialization on
homology in (B.5.16) gives injectivity on cohomology;
the monodromy-invariant subspace has the same dimension
four. Its image therefore includes $uw$ and $\gamma\delta$.
In degree one the image is $\langle\gamma,u\rangle$;
their products with those two degree-two classes give the
two independent invariant classes in degree three and
the top class in degree four. Thus all the specialization
maps have the asserted ranks, with no additional
skyscraper summand. The paragraph following (B.5.16)
also proves these specialization assertions integrally
in all degrees by the actual honeycomb quotient map.

It follows that $R^q f_*\mathcal L_\lambda=j_*L_q$ over
the disc containing the three critical values. On this
disc $H^p(j_*L_q)=0$ for $p\geq2$. To check this directly,
retract the punctured disc to a free graph on three loops.
Restriction of a graph cocycle to the three local
coinvariant spaces is onto: its values on the three free
generators can be chosen independently. The exact
restriction sequence therefore kills $H^2(j_*L_q)$.
If $d_q=\binom4q$, $g_q$ is the common fixed dimension of
(E.S41), and $k_{i,q}$ are their three individual fixed
dimensions, this same graph calculation gives
\[
 \dim H^0(j_*L_q)=g_q,\qquad
 \dim H^1(j_*L_q)=2d_q+g_q-\sum_i k_{i,q}.
 \tag{E.S42}
\]
The Leray sequence has only these two base degrees and
degenerates. Exact exterior-matrix ranks give the following
table; a row applies to each character of the indicated
order:
\[
 \begin{array}{c|rrrrrr}
 \operatorname{ord}(\lambda)&b_0&b_1&b_2&b_3&b_4&b_5\\ \hline
 1 &1&1&4&6&4&0\\
 2 &0&0&4&6&4&0\\
 3 &0&0&3&4&3&0\\
 4 &0&0&3&4&3&0\\
 6 &0&1&5&6&5&1\\
 12&0&1&4&4&4&1
 \end{array}
 \tag{E.S43}
\]
Here $b_j=\dim H^j(E;\mathcal L_\lambda)$ and $b_6=0$.
For the nontrivial characters all $g_q$ are zero; for the
trivial one they are all one. The program
\texttt{maps\_shifted\_cover.py} computes (E.S41)--(E.S43)
over the exact cyclotomic fields, represented by rational
companion matrices. It prints every local fixed dimension.
Summing the rows with multiplicities $1,1,2,2,2,4$ gives
\[
             (b_0,\ldots,b_6)(E_{12})=(1,7,46,56,46,6,0).
 \tag{E.S44}
\]
Both sides have Euler characteristic $24$, as required by
the degree-twelve covering of $E$.

\emph{The Betti and Jordan calculations on the actual fibre.}
The return on each of the two pieces of (E.S29), up to
isotopy allowing its boundary to move, is its cyclic deck
transformation times the identity of $T^3$. This follows
also by moving the interpolation in (E.S36) entirely into
the other piece's collar. It acts trivially on the
cohomology of $D$. The generalized eigenvalue-one part of
the Mayer--Vietoris sequence therefore uses just the
invariant part $H^*(E)$ in (E.S40). Its dimensions are
\[
       (1,0,1,19,45,45,19,1,0,1).
 \tag{E.S45}
\]
Every other eigenspace is the direct sum of the two piece
eigenspaces, because the boundary has none. In particular
those eigenspaces are semisimple. Their combined Poincare
polynomial is
\[
     2(1+t)^3(6t+42t^2+50t^3+42t^4+6t^5),
 \tag{E.S46}
\]
by (E.S43)--(E.S44). Adding (E.S45) proves (E.S33).

Finally the Wang sequence for $P=M_\psi$ and the rational
Betti numbers in E.2.7 give the numbers of Jordan blocks
at eigenvalue one as
\[
          (1,0,1,19,44,44,19,1,0,1).
 \tag{E.S47}
\]
The only possible block sizes are one and two: the
twelfth power is (E.S31), and (E.S37) squares to zero.
Subtracting (E.S47) from (E.S45) gives exactly one block
of size two in each of degrees four and five. This also
independently checks the direct variation ranks (E.S39).
All other eigenvalues and blocks are already described
by (E.S43) and (E.S46). None of these rational calculations
claims an integral torsion classification or replaces the
smooth $E_{12}$ by a manifold with the same invariants.
\end{proof}

\paragraph{Theorem E.2.15 (Proved: integral fibre groups and the coefficient of the torus twist).}
The groups $H^*(E_{12};\mathbb Z)$ are torsion free, with
the ranks in (E.S44). For the actual fibre $Y$, the only
torsion in cohomology is
\[
       H^3(Y;\mathbb Z)_{\rm tors}
       =H^7(Y;\mathbb Z)_{\rm tors}=(\mathbb Z/6)^2.
 \tag{E.S48}
\]
Equivalently, its integral homology has the ranks (E.S33)
and torsion $(\mathbb Z/6)^2$ exactly in degrees two and six.
The twelfth return $T=\psi^{12}$ acts as the identity on
all integral cohomology except in degrees four and five.
In each of those two free groups it has, in an integral
basis, the form
\[
             \begin{pmatrix}1&72\\0&1\end{pmatrix}
                                  \oplus I_{691}.
 \tag{E.S49}
\]

For $n\geq1$, let $P_{12n}$ be the connected cyclic cover
of $P$ of degree $12n$ defined by $\chi$. Its Betti numbers,
independent of $n$, are
\[
       (1,13,133,528,1099,1384,1099,528,133,13,1).
 \tag{E.S50}
\]
Its integral homology torsion is exactly
\[
 H_j(P_{12n};\mathbb Z)_{\rm tors}=
 \begin{cases}
       (\mathbb Z/6)^2,&j=2,3,6,7,\\
       \mathbb Z/(72n),&j=4,5,\\
       0,&\text{otherwise}.
 \end{cases}
 \tag{E.S51}
\]
In particular these covers are pairwise nonhomotopy-equivalent.
These integral conclusions still do not recognize the smooth
$E_{12}$ or replace its boundary marking by homological data.

\begin{proof}
\emph{The integral cohomology of the actual cover.}
Apply the proof of (E.S42) with the coefficient lattice
$\mathbb Z[C_{12}]$, using its regular permutation matrix
$P_{12}$. In vertical degree $q$ set
\[
 \begin{aligned}
 T_{1,q}=P_{12}^4\otimes\bigwedge^q A_1^{-t},\quad
 &T_{2,q}=P_{12}^{-3}\otimes\bigwedge^q A_2^{-t},\\
 T_{0,q}=I_{12}\otimes\bigwedge^q M_0^{-t},
 &V_q=\mathbb Z^{12\binom4q}.
 \end{aligned}
 \tag{E.S52}
\]
This coefficient system computes $H^*(E_{12};\mathbb Z)$.
There is no integral averaging step. At the first finite
cap the cover has four components, each the ordinary
degree-three torus root cover; at the second there are
three ordinary degree-four torus root covers. An invariant
vector is determined by its value on one sheet of such
an orbit, with the other values its integral monodromy
translates. This identifies cap cohomology with the full
primitive invariant lattice. The cusp has twelve unchanged
components and the integral specialization of (B.5.16).
Consequently its direct-image sheaf is exactly the integral
$j_*L_q$, including at every critical point.

The three graph-cocycle values again surject onto the
three local coinvariant groups, integrally. Thus $H^2(j_*L_q)=0$,
and its remaining cohomology is the kernel and cokernel
of the exact lattice map
\[
 \begin{gathered}
 V_q\longrightarrow
        \bigoplus_{i=0}^2\operatorname{im}(T_{i,q}-I),\\
           v\longmapsto((T_{i,q}-I)v)_i.
 \end{gathered}
 \tag{E.S53}
\]
Use Hermite column bases for the three image lattices;
this retains their integral indices rather than replacing
them by rational spans. The source and target ranks,
and the rank of (E.S53), are respectively
\[
 \begin{array}{c|rrrrr}
 q&0&1&2&3&4\\ \hline
 \operatorname{rank}V_q&12&48&72&48&12\\
 \operatorname{rank}\bigoplus\operatorname{im}(T_{i,q}-I)
                      &17&92&126&92&17\\
 \operatorname{rank}(\text{E.S53})&11&47&71&47&11
 \end{array}
 \tag{E.S54}
\]
Every nonzero Smith entry is one. Hence $H^0(j_*L_q)$
has rank one and $H^1(j_*L_q)$ is free with ranks
$(6,45,55,45,6)$. The two-column Leray filtration has
free end groups and therefore splits as abelian groups.
This proves the torsion-free assertion and (E.S44)
integrally. The full integer matrices and the exact
Hermite/Smith check are in
\texttt{maps\_shifted\_cover\_integral.py}.

\emph{Why the marked rational images now determine the integral images.}
Each piece $A=E_{12}\times T^3$ in (E.S29) has free
integral homology. By Poincare--Lefschetz duality,
$H^{k+1}(A,\partial A;\mathbb Z)$ is free. The cokernel
of $H^k(A;\mathbb Z)\to H^k(\partial A;\mathbb Z)$
embeds in that free group, so the boundary image is a
primitive lattice. Its rational span has already been
computed in (E.S38). Therefore the actual integral
image is precisely the saturation of that displayed
image in $\bigwedge^k\mathbb Z^8$.

Choose saturated column bases $\overline A_k,\overline B_k$
for the two images. The combined matrix
$[\overline A_k\ \overline B_k]$ has the ranks in the
second row of (E.S40). All its nonzero Smith entries are
one except that, in each of degrees two and six, exactly
two entries are six. The Mayer--Vietoris exact sequence
has a free kernel term on its right, since all piece
cohomology is free. Thus its abelian extension splits,
and these Smith entries give exactly (E.S48). The free
ranks remain (E.S33).

\emph{The integral variation, including possible torsion components.}
Formula (E.S37) holds integrally. One can use the singular
prism of the loop of translations instead of differential
forms: the action map $S^1\times T^8\to T^8$ has
cohomology formula
\[
 a\longmapsto 1\otimes a+[S^1]^*\otimes\iota_v a.
\]
The difference between the two collar homotopies is
exactly this loop, so its connecting class is
$\partial\iota_v r(a)$. This proves the integral formula
also on torsion and still gives $(T^*-I)^2=0$.

For complete lattice control, take an integral kernel
basis of $[\overline A_k\ -\overline B_k]$. Its first
block, multiplied by $\overline A_k$, is a basis of the
common boundary lattice. Contract it by $v$ and reduce
modulo $[\overline A_{k-1}\ \overline B_{k-1}]$ using
the integral Smith transforms. In degrees four and five
the result has rank one and unique nonzero Smith entry
$72$. In degree three its two possible components
modulo six are both zero; in degree seven the common
boundary image is zero. All the remaining reductions
are zero. These computations are performed, with the
unimodular transforms verified, in
\texttt{maps\_shifted\_fibre\_integral.py}.
In degrees four and five the connecting cokernel is
free and embeds as a primitive summand of $H^k(Y)$,
because the following Mayer--Vietoris kernel is free.
Thus $72$ is the divisibility in the whole cohomology
group, not just in a chosen sublattice.

A rank-one integral map of Smith entry $72$ has the
form $72u\ell$, where $u$ and $\ell$ are primitive.
Its square is zero precisely when $\ell(u)=0$.
Extend $u$ to a basis of $\ker\ell$ and choose a vector
on which $\ell$ is one. This gives (E.S49). The same
calculation proves identity on every other integral
cohomology group, including its torsion.

\emph{The finite-cover groups and the two potentially nontrivial extensions.}
The cover $P_{12n}$ is the mapping torus of $T^n$.
Since $(T^*-I)^2=0$, its variation is $n(T^*-I)$,
with Smith entry $72n$ in degrees four and five.
For every $k$ the integral Wang sequence gives
\[
 0\longrightarrow\operatorname{coker}(T^{n*}-I)_{k-1}
 \longrightarrow H^k(P_{12n};\mathbb Z)
 \longrightarrow\ker(T^{n*}-I)_k\longrightarrow0.
 \tag{E.S55}
\]
The ranks immediately give (E.S50). Whenever the right
group is free, (E.S55) splits as an abelian sequence.
In particular degrees four and eight have torsion
$(\mathbb Z/6)^2$, contributed by $H^3(Y)$ and $H^7(Y)$
respectively, while degrees five and six have torsion
$\mathbb Z/(72n)$. Every other such degree has no torsion.

The only right groups in (E.S55) that are not free
occur in degrees three and seven. Their extensions
must not be split by assumption. Instead use the
nondegenerate integral torsion linking pairing of the
closed oriented ten-manifold $P_{12n}$: cohomological
torsion in degree $k$ is Pontryagin dual to that in
degree $11-k$. The already fixed torsion in degree
eight therefore forces the degree-three torsion to
be $(\mathbb Z/6)^2$, and the fixed degree-four torsion
forces the same answer in degree seven. The free
ranks are already known. This resolves precisely
the two possibly absorbing extensions. The universal
coefficient theorem now translates these cohomology
groups to (E.S51). The order $72n$ in the middle
groups distinguishes the covers for distinct $n$.
\end{proof}

\paragraph{Theorem E.2.16 (Proved: the full complement module and its peripheral map).}
Let $K=\pi_1(E)$ have the marking (E.S1), and put
\[
 R=\mathbb ZK,\qquad Q=K/\langle h\rangle=C_3*C_4,
 \qquad \overline R=\mathbb ZQ.
\]
Use right deck modules; conversion to left modules applies the
involution $g\mapsto g^{-1}$ to every coefficient. Then
\[
 \pi_2(E)\simeq R\oplus\overline R^{\,2}\oplus L,
 \qquad
 L=\langle x,t\mid t(a-1)=t(b-1)=0,
                         \ x(1+b)+3t=0\rangle_{\overline R}.
 \tag{E.C1}
\]
Equivalently, replacing $x$ by $x+t$ and eliminating $t$ gives
\[
 L\simeq\overline R\big/
       \big((b^2-1)\overline R+(1+b)(a-1)\overline R\big).
 \tag{E.C2}
\]
There is a nonsplit exact sequence
\[
 0\longrightarrow\mathbb Z\langle t\rangle
  \longrightarrow L\longrightarrow
  \mathbb Z_{-}\otimes_{\mathbb ZC_4}\overline R
  \longrightarrow0,
 \tag{E.C3}
\]
where $b$ acts as $-1$ on $\mathbb Z_{-}$ and the first term
has trivial $Q$ action. Its extension class is the nonzero
element of an order-two extension group. In the connected
lift of the peripheral $T^5$ to the universal cover of $E$,
the degree-two classes $uw,ud,wd$ map respectively to
\[
                           t,\quad0,\quad0.
 \tag{E.C4}
\]
Thus the central torsion summand of the module is annihilated
by $h-1$ itself. Formula (E.C1) is a full integral module
calculation, including the central value $h=1$.  The complete
splice extension is determined separately in E.2.17, and the
smooth type by full framed handles in E.2.21; neither conclusion
is inferred from this module calculation alone.

There is also a specified chain-level peripheral comparison.
For the splitting constructed below, let $e$ generate the
free $R$ summand and let $y,z$ generate the two copies of
$\overline R$. The projection of the peripheral degree-two
chain map to (E.C1) is
\[
 \begin{array}{c|rrrrrr}
 \text{cell}&gu&gw&gd&uw&ud&wd\\ \hline
 \text{value}&0&0&6t&t&0&0
 \end{array}
 \tag{E.C19a}
\]
and, writing $\Theta$ for the outer base circle with
group image $\mu=ab$, use the cone convention: the
notation $v\Theta$ denotes the suspended cell
$\Theta\times v$, with the interval factor first.
In this convention,
\[
 \begin{array}{c|l}
  g\Theta&eh-2x-6yb+z(5-2b)\\
  u\Theta&x-yb+zb\\
  w\Theta&-y(1+b)+z\\
  d\Theta&x+2t+z(b-1).
 \end{array}
 \tag{E.C19b}
\]
These are values on chains, not assertions that all
the displayed cells are cycles in the cover. In
particular the free coefficient $h$ on $g\Theta$
is nonzero even though the free summand receives
zero from the lifted peripheral homology in (E.C4).
The original cone and its full peripheral map, not
only the map on homology, are part of the certificate.

\begin{proof}
We give the local chain comparison before the global
calculation, because the ordinary homology maps alone would
not determine (E.C1).

\emph{The cusp boundary and its derived attachment.}
Write $A=\mathbb Z[h,h^{-1}]$ and $q=h-1$. Let
$\widehat W$ be the gamma cover used in E.2.12, and let
$\widehat M$ be the corresponding boundary cover. Ordinary
chains with the deck action give finite free $A$ chain
models. The homology of $\widehat W$ is
\[
 H_0=H_1=A/(q),\quad H_2=A\oplus(A/(q))^2,
 \quad H_3=0,\quad H_4=A.
 \tag{E.C5}
\]
For the boundary, start with the torus chain complex
$C= A\otimes\bigwedge\langle g,u,w,d\rangle$,
whose differential contracts $g$ and multiplies by $q$.
The cellular monodromy has degree-one matrix
\[
 F(M_0)g=g-hd,\quad F(M_0)u=u+w,
 \quad F(M_0)w=w,\quad F(M_0)d=d.
 \tag{E.C6}
\]
Its higher matrices are exterior powers. The mapping-torus
complex is the cone of $I-F(M_0)$, with differential
\[
 d(c,e)=\big(dc+(I-F(M_0))e,-de\big).
 \tag{E.C7}
\]
For example $d(g\theta)=hd-q\theta$ and
$d(u\theta)=-w$. Cancelling these two unit pairs leaves
the degree-one classes $u$ and $\theta$, with resolutions
$q$ and a unit times $q^2$. In degree two the four surviving
classes are $uw,ud,w\theta,d\theta$, each with resolution
$q$. In particular
\[
 H_0(\widehat M)=A/(q),\quad
 H_1(\widehat M)=A/(q^2)\oplus A/(q),\quad
 H_2(\widehat M)=(A/(q))^4.
 \tag{E.C8}
\]
Geometrically this is the Heisenberg three-manifold times
a circle: the deck transformation shears the meridian
by the $d$ circle. The homology map to (E.C5) is the
identity in degree zero, projection to $u$ in degree one,
and the primitive projection
$(uw,ud,w\theta,d\theta)\mapsto(U,D,0,0)$ in degree two.
The last assertion uses the actual specialization
(B.5.16) and the literal collapsed meridian tori
(B.5.17), rather than ranks alone.

The missing extension is fixed by a relative intersection
calculation. Let $S_i$ denote the normalized hexagon
component with gamma index $i$. In its standard Picard
basis $H,E_1,E_2,E_3$, the six boundary curves are
\[
 E_1,\ H-E_1-E_2,\ E_2,\ H-E_2-E_3,
                       \ E_3,\ H-E_1-E_3.
 \tag{E.C9}
\]
The first opposite pair is self-glued in the surviving
$u$ direction; the other two pairs join $S_i$ to its
two gamma neighbours. Put $C_0=-H+E_1$, the primitive
generator of the free kernel of the normalization matrix
$B(h)$ in E.2.12. The divisor $S_i$ restricts on its
own normalization to $2C_0$: the total central divisor
is principal, and the two self-gluing preimages cancel
their terms in the normal restriction. It restricts
to $-C_0$ on each neighbouring component and to zero
on every other component. Thus the free part of the
equivariant intersection map is
\[
 H_4(\widehat W)=A\longrightarrow
 H_4(\widehat W,\widehat M)\longrightarrow A,
 \qquad 2-h-h^{-1}=-h^{-1}q^2.
 \tag{E.C10}
\]
Poincare--Lefschetz duality over $A$, with its usual
involution, gives the same coefficient, up to a unit,
for the free part of $H_2(\widehat W)\to
H_2(\widehat W,\widehat M)$. The latter relative group
is $A$: the universal-coefficient sequence in degree
four has $\operatorname{Hom}_A(H_4,A)=A$, zero $H_3$,
and no higher extension terms. Indeed every module
in (E.C5) has a projective resolution of length at
most one. The relative sequence therefore includes
\[
 0\longrightarrow A\xrightarrow{\,\text{unit}\ q^2\,}A
  \longrightarrow A/(q^2)\longrightarrow0.
 \tag{E.C11}
\]
Consequently the boundary map has a unit component in
$\operatorname{Ext}^1_A(A/(q^2),A)=A/(q^2)$.
Every unit of $A/(q^2)=\mathbb Z[\varepsilon]/(\varepsilon^2)$
lifts to a unit $\pm h^n$ of $A$.

Here is why this determines the required low chain map
without discarding a peripheral invariant. All the
modules in (E.C5) and (E.C8) have projective dimension
at most one. Their complexes split in the derived
category into their two-term projective resolutions:
the obstructions to splitting consecutive truncations
are extension groups in degree at least two and vanish.
In degrees relevant to $H_2$ of a mapping cone, a map
between these split complexes consists of the displayed
homology maps and extension components from $H_i$ to
$H_{i+1}$. Normalize the component in (E.C11) by a unit
on the free $H_2$ summand of the target. Its components
to the two torsion $H_2$ summands can then be removed
by target maps $A\to(A/(q))^2$: reduction of the unit
modulo $q$ is $\pm1$. The remaining components from
the $u$ summand are removed by target degree-one
extension maps from $H_1(\widehat W)$ to
$H_2(\widehat W)$. Finally the degree-zero component
is removed using the isomorphism on $H_0$ and a target
extension map from $H_0$ to $H_1$.

Each change just described is a unit on the diagonal
of the split complex and is triangular with extension
entries off the diagonal. It is therefore invertible.
Replacing the split resolutions by the original finite
free complexes represents it by a chain homotopy
equivalence. No realization by a diffeomorphism or a
map of the cusp space is claimed or needed. The change
is applied to the cusp target and its entire incoming
boundary map simultaneously. The other vertices and
all maps into the regular core are unchanged. It thus
gives a chain homotopy equivalence of the global
mapping-cone complex that is the identity on the
external peripheral complex. In particular the
calculation of (E.C4) retains the actual boundary
map. Tensoring with $R$ preserves these comparisons,
since $R$ is free over $A$.

One resulting cusp model in degrees zero through three
has ranks $(1,2,4,2)$, generators
\[
 (1);\quad(g,u);\quad(gu,E,U,D);\quad(R_U,R_D),
\]
and differentials $dg=q$, $d(gu)=qu$,
$dR_U=qU$, $dR_D=qD$. The map (E.C7) to this model
is, through degree two,
\[
 \begin{array}{c|ccccc}
  &g&u&w&d&\theta\\ \hline
  i_1&g&u&0&0&0
 \end{array},\qquad
 \begin{array}{c|rrrrrr}
  &gu&gw&gd&uw&ud&wd\\ \hline
  i_2&gu&-hD&E&U&D&0
 \end{array},
 \qquad i_2(v\theta)=0.
 \tag{E.C12}
\]
The signs and the factor $h$ satisfy the chain equation
directly. At $h=1$, choose
$E=-\overline C_1-\overline C_2$,
$U=\overline C_2+\overline C_3$, and
$D=\overline C_2$; this is precisely (B.5.16).

\emph{Finite caps and the global complex.}
For any marked monodromy
$A=\left(\begin{smallmatrix}1&0\\t&D\end{smallmatrix}\right)$
use the chain lift
$F(A)=\left(\begin{smallmatrix}1&0\\ht&D\end{smallmatrix}\right)$
and its exterior powers. These lifts multiply exactly
and commute with contraction by $qg^*$.

This exterior lift represents the actual equivariant
torus map. The gamma cover is
$\mathbb R_g\times T^3$, and the marked linear map is
$(g,v)\mapsto(g,Dv+gt)$. At every integer value of $g$
the translation is integral, so the vertex torus map
is exactly $D$. In a single unit strip, concentrate
the loop of translations $v\mapsto v+st$ at its
terminal endpoint by an equivariant cellular
approximation relative to the two vertex tori.
The translation prism on the minimal cellular complex
of $T^3$ is exterior multiplication by $t$: its value
on each product torus is the homology class of the
literal translation sweep, and all differentials
in that minimal torus complex are zero. The terminal
vertex has deck coefficient $h$. Consequently
\[
 v\longmapsto\bigwedge D(v),\qquad
 g\wedge v\longmapsto
 g\wedge\bigwedge D(v)+ht\wedge\bigwedge D(v),
 \tag{E.C20}
\]
which is exactly the exterior lift above. There is
only this one strip parameter. Composing two such
actual maps adds the translation loops by
$t_A+D_At_B$, the same strict multiplication law
as $F(A)F(B)=F(AB)$. Thus no uncomputed higher
coherence is being substituted for the torus map.
A regular
boundary with deck image $s$ is the cone of $sI-F(A)$.
This convention gives the group relation
$s\lambda s^{-1}=A\lambda$.

At finite cap $j$, take its adapted first fibre vector
$n_jv_j$, where
\[
 (m_1,n_1,v_1)=(3,1,(1,2,-4,0)),\qquad
 (m_2,n_2,v_2)=(4,-1,(-1,-3,3,0)).
 \tag{E.C13}
\]
It has gamma coordinate one and is fixed by $A_j$.
The other three coordinates form $L_0=\langle u,w,d\rangle$,
with monodromy $D_j=A_j|_{L_0}$. The cap is the
mapping torus of this three-torus, and its complex has
differential $sI-\bigwedge^rD_j$ on the swept cells.
The fibre-to-cap map is the identity on $\bigwedge L_0$;
on the adapted long-circle factor it is the integral
norm matrix $N_{m_jn_j}(s,D_j)$ determined by
\[
 (sI-D_j)N_{m_jn_j}(s,D_j)=(h-1)I,
 \qquad s=a\text{ or }b.
 \tag{E.C14}
\]
For positive $n_j$ the norm is
$\sum_{r=0}^{m_j-1}s^rD_j^{m_j-1-r}$, and for
$n_j=-1$ it is minus $s^{-m_j}$ times that sum.
In phase degree $r$, use
$N_{m_jn_j}(s,\bigwedge^rD_j)$, rather than the
exterior power of the norm matrix.
The boundary base-circle generator maps to the cap circle, while
the cell containing both that meridian and the adapted
long circle maps to zero. These are actual winding
maps in the mapping-torus coordinates. The linear
three-torus return fixes its origin; hence there is
no additional translation prism in this comparison.

Form the core over a bouquet of three circles with
deck images $a,b,1$ and monodromies $A_1,A_2,M_0$.
Attach the two finite complexes by (E.C14) and the
cusp by (E.C12), using their double mapping cylinders.
Through degree three the resulting free $R$ complex
has ranks $(4,20,49,62)$. Its entries and every unit
cancellation are given exactly in
\texttt{code/maps\_complement\_module.py} and its
JSON output. The program uses the integral normal
form $h^r$ times an alternating word in
$a,a^2$ and $b,b^2,b^3$; no finite group quotient or
floating-point approximation is used.

For completeness the final reduction can be checked
without reading a long matrix. Cancelling degree-zero
and degree-one unit pairs leaves ranks $(1,2,34,62)$,
with $d_1=(a-1,b-1)$. The only nonzero columns of
$d_2$ are $F$ and $F(b-1)$, where
\[
 F=\begin{pmatrix}1+a+a^2\\h(1+b+b^2+b^3)\end{pmatrix}.
 \tag{E.C15}
\]
The first component of $F$ is injective on $R$:
multiplication by $a-1$ converts it to $h-1$, and
$R$ is free over the Laurent ring $A$. Thus removing
the $F$ column after subtracting its multiple gives
a free basis of $\ker d_2$ of rank33. The corresponding
row of $d_3$ is zero by the chain equation. Unit
reductions of the remaining presentation leave five
generators $(x,t,y,z,e)$. One generator $e$ occurs
in no relation. Three pure $t$ relations are
\[
 tq=0,\qquad t(1-a^{-1})=0,
                  \qquad t(b-a^{-1})=0.
 \tag{E.C16}
\]
They generate exactly the augmentation relations on
$t$: multiply the second by $a$ to get $t(a-1)=0$,
and subtract it from the third to get $t(b-1)=0$.
Every other pure $t$ coefficient has augmentation
zero and is consequently redundant. Reducing the
remaining $t$ coefficients by augmentation, three
relations give successively $yq=0,zq=0,xq=0$.
The final three relations are
\[
 x(1+b)+3t=0,\quad x(b+b^2)+3t=0,
                   \quad x(1+b+b^2+b^3)+6t=0.
 \tag{E.C17}
\]
The last two are respectively the first multiplied
by $b$, and the first multiplied by $1+b^2$.
These identities prove (E.C1). The program verifies
the original chain equations, all Schur complements,
and the coefficient list used in (E.C16)--(E.C17).

The invariant subgroup generated by $t$ injects in
$L$. One way to see this is to start with the exact
sequence of right $\mathbb ZC_4$ modules
\[
 0\longrightarrow\mathbb Z\longrightarrow
 \mathbb Z[C_4/C_2]\longrightarrow\mathbb Z_-\longrightarrow0,
\]
where the first map is the sum of the two cosets.
Induce to $Q$ and push out its first term by the
augmentation to the trivial module $\mathbb Z$.
The resulting presentation is (E.C1), after the
change $x\mapsto x+t$ which replaces the coefficient
$3$ by $1$. Pushout preserves the displayed injection.
Restriction to $C_4$ shows nonsplitting: an invariant
correction of the lift $x$ changes this odd coefficient
by an even number. Equivalently the standard cyclic
resolution computes
$\operatorname{Ext}^1_{\mathbb ZC_4}(\mathbb Z_-,\mathbb Z)
=\mathbb Z/2$. Induction--restriction gives the same
extension group for (E.C3).

Finally insert the actual fibre cycles $uw,ud,wd$ in
the core before any cancellation and carry their
coordinate columns through the same basis changes.
The five-generator presentation gives respectively
\[
 t,\qquad t(a^{-1}-a),\qquad t(a^{-1}-1).
 \tag{E.C18}
\]
Relations (E.C16) reduce these to (E.C4). The external
boundary cover has homotopy type $T^3$ and these are
its three degree-two generators. This also checks
the ordinary coefficient independently: after taking
$K$ coinvariants, $2x+3t=0$, so $t$ is twice a
primitive class, exactly the coefficient2 in E.2.6.

To retain the remaining peripheral data, use the
boundary word $abc$, where the cusp generator $c$
has deck image1. On swept fibre chains its map to
the three core edges has blocks
\[
                 bI,\qquad F(A_1),\qquad F(A_1)F(A_2).
 \tag{E.C21}
\]
The identity
\[
 (aI-F(A_1))b+F(A_1)(bI-F(A_2))
       +F(A_1)F(A_2)(I-F(M_0))=(ab-1)I
\]
checks its chain equation. This is the usual prism
subdivision of the actual peripheral word, with
its given ribbon marking. Tracking all ten
degree-two columns through the same reductions,
and using only $t(a-1)=t(b-1)=0$ and
$xq=yq=zq=0$, gives (E.C19a)--(E.C19b).
The unreduced columns are saved in the JSON
certificate. In particular no assertion that a
homologically invisible peripheral component is
zero enters this calculation.
\end{proof}

\paragraph{Theorem E.2.17 (Proved: the full shifted homotopy module).}
Use the right-module convention of E.2.16. Write
\[
 \begin{gathered}
 R=\mathbb ZG,\qquad q_i=h_i-1,\qquad S=\mathbb Z\Gamma,\\
 \Gamma=G/\langle h_1,h_2\rangle
  =\langle a,b,c,d\mid a^3=b^4=c^3=d^4=abcd=1\rangle.
 \end{gathered}
 \tag{E.P1}
\]
where $a,b$ belong to the first complement, $c,d$ to
the second, and $\mu=ab=(cd)^{-1}$. The full integral
second-homotopy module is
\[
 \pi_2(P)\simeq R/q_2R\ \oplus\ R/q_1R
                    \ \oplus\ T_L\oplus T_R.
 \tag{E.P2}
\]
The last two modules are inflated from the following
explicit six-generator presentations over $S$.
Put $N_a=1+a+a^2$, $N_b=1+b+b^2+b^3$, and define
$N_c,N_d$ similarly.

For $T_L$ use generators $p_i,r_i$ for $0\leq i\leq2$,
put $v_i=p_i-r_i$, and define
\[
 \begin{split}
 \Omega_L={}&v_0(1+d-d^2-d^3)
       +v_1(d^3-d)+v_2(d^2+d^3),\\
 p_iN_a={}&0,\qquad r_iN_b=0\quad(0\leq i\leq2),\\
 \Omega_L(a-1)={}&\Omega_L(b-1)
       =\Omega_L(c-1)=\Omega_L(d-1)=0.
 \end{split}
 \tag{E.P3}
\]
For $T_R$ use generators $B_i,C_i$, put $w_i=B_id+C_i$,
and define
\[
 \begin{split}
 \Omega_R={}&w_0a(1-b-b^2+b^3)
              +w_1a(b^2-1)-w_2a(1+b^3),\\
 B_iN_c={}&0,\qquad C_iN_d=0\quad(0\leq i\leq2),\\
 \Omega_R(a-1)={}&\Omega_R(b-1)
       =\Omega_R(c-1)=\Omega_R(d-1)=0.
 \end{split}
 \tag{E.P4}
\]
These are complete presentations: six generators and
ten relations in each block. In particular
\[
 q_1q_2\pi_2(P)=0,\qquad
 \ker(q_1:\pi_2(P)\to\pi_2(P))
  \cap\ker(q_2:\pi_2(P)\to\pi_2(P))=T_L\oplus T_R.
 \tag{E.P5}
\]
This improves the earlier squared annihilator and
specifies its entire central-intersection summand.
Its coinvariants are $(\mathbb Z/12)^6$, so (E.P2)
recovers $\pi_2(P)_G=\mathbb Z^2\oplus(\mathbb Z/12)^6$.

Let $\Pi=\pi_1(Y)$ and let $\Sigma\simeq\pi_1(\Sigma_6)$
be its quotient by $z=h_1h_2$, as in E.2.13. Then
\[
 \pi_2(Y)\simeq(\mathbb Z\Pi)^{24}
          \oplus\operatorname{Res}^{\Gamma}_{\Sigma}(T_L\oplus T_R),
 \qquad
 \pi_2(Y)_\Pi\simeq\mathbb Z^{110}\oplus(\mathbb Z/6)^2.
 \tag{E.P6}
\]
The surface-group summand is inflated to $\Pi$, so
$z$ acts trivially on it. Formulas (E.P2)--(E.P6)
do not identify the smooth type of $P$ or $Y$.

\begin{proof}
\emph{A small chain presentation with the actual peripheral map.}
The calculation in E.2.16 permits the following model
through degree three for the first complement. Its
bases have ranks $(1,2,6,6)$, with degree-one basis
$(a,b)$ and degree-two basis $(r,x,t,y,z,e)$.
The first two differentials are
\[
 d_1=(a-1,b-1),\qquad
 d_2(r)=\begin{pmatrix}N_a\\h_1N_b\end{pmatrix},
 \quad d_2(x)=d_2(t)=d_2(y)=d_2(z)=d_2(e)=0.
 \tag{E.P7}
\]
The six degree-three columns are the six relations
\[
 t(a-1),\quad t(b-1),\quad xq_1,\quad yq_1,\quad zq_1,
                         \quad x(1+b)+3t.
 \tag{E.P8}
\]
The degree-two norm $N_a$ in (E.P7) is written with
the same expression as in (E.P3), but here $a^3=h_1$,
not $1$. The second complement is the identical
model with $(a,b,h_1)$ replaced by $(c,d,h_2)$.

The ordinary one-relator complex $C_K$ is the subcomplex
on degrees zero and one and the cell $r$. It is an
aspherical model: it is the double mapping cylinder
of the degree-three and degree-four maps of circles,
and its universal cover is a tree of contractible
lines and strips. Equivalently it is the standard
Seifert spine of the torus-knot exterior. Projection
of (E.P7)--(E.P8) onto $C_K$ retains the actual
peripheral map. In degree one this map is
\[
 g\longmapsto(N_a,0),\qquad
 u,w,d\longmapsto0,\qquad
 \Theta\longmapsto(b,1).
 \tag{E.P9}
\]
Its degree-two $r$ coefficient is $b-1$ on the swept
cell $g\Theta$ and zero on the other nine cells.
Its remaining components are exactly (E.C19).
These additional columns were tracked in the
E.2.16 certificate, rather than chosen from their
ordinary homology values.

Here and below the small complex is only a
presentation through the degree needed for $H_2$.
Replacing redundant degree-three columns by the six
generators of their image in (E.P8) changes no
degree-two cycle or boundary. Reducing a peripheral
degree-two column modulo that image changes it by
a chain homotopy with values in degree three.
Products with a torus and the mapping cones below
therefore retain the same degree-two homology.
For the comparison with an aspherical target one
may first keep the original, full complex: its
projection to $C_K$ extends by zero in degrees at
least three, since the Fox column in (E.P7) is
injective and hence the $r$ row of $d_3$ is zero.
Only after taking the kernel is its image in degree
two shortened as in (E.P8). Thus no false
acyclicity of a truncated target is used.

\emph{The exact kernel complex for the splice.}
Let $T(s)$ be the four-torus chain complex whose
first circle has deck image $s$ and whose other
three circles have trivial image. Its differential
contracts that first circle and multiplies by $s-1$.
The two vertex complexes for the actual splice are
\[
 V_1=C_{E_1}\otimes T(h_2),\qquad
 V_2=C_{E_2}\otimes T(h_1).
 \tag{E.P10}
\]
We have reversed the two product factors in the
second vertex only for notation; the Koszul signs
below account for that reversal. The boundary
complex $B$ is the nine-torus complex with ordered
coordinates
\[
 (g_1,u_1,w_1,d_1,g_2,u_2,w_2,d_2,\Theta),
 \quad\text{deck images }(h_1,1,1,1,h_2,1,1,1,\mu).
 \tag{E.P11}
\]
Its differential is the usual exterior contraction
by those deck images minus one.

For clarity all source signs can be specified
without an implicit boundary convention. Write
a boundary product cell as $\alpha\wedge\beta$
with an optional final $\Theta$, where $\alpha$
uses the first four coordinates and $\beta$ the
second four. Put $\epsilon=1$ when $\Theta$ occurs
and $0$ otherwise. Its first vertex map applies
the E.2.16 peripheral map to $(\alpha,\Theta)$,
keeps $\beta$, and has sign
$(-1)^{\epsilon(|\alpha|+|\beta|)}$.
The second map applies that peripheral map to
$(\beta,\Theta_2)$, keeps $\alpha$, and has sign
$(-1)^{|\alpha|(|\beta|+\epsilon)+\epsilon|\beta|}$.
If $\epsilon=1$, multiply its coefficient on the
right by $-\mu$. This is the exact cellular map
of $\Theta_1\mapsto\Theta_2^{-1}$, since
$\mu_2=cd=\mu^{-1}$. The suspended cells in (E.C19)
have the interval first; these signs convert from
the ordered exterior cells in (E.P11).
Denote the two maps by $j_1,j_2$.

The chain complex of the splice is the cone of
$(j_1,-j_2):B\to V_1\oplus V_2$.
Project each $C_E$ to $C_K$ and each $T(s)$ to its
first circle, and project $B$ to its three coordinates
$g_1,g_2,\Theta$. The result is the complex of the
aspherical graph of spaces
\[
 (K(K_1,1)\times S^1)\ \cup_{T^3}\
                         (K(K_2,1)\times S^1).
 \tag{E.P12}
\]
It is the model $N$ of E.2.9. Both peripheral maps
on groups are injective by E.2.10; their common
group is $\langle h_1,h_2,\mu\rangle=\mathbb Z^3$.
Consequently its universal cover is a Bass--Serre
tree of contractible vertex and edge spaces and
is contractible. Its chain complex has ranks
$(2,7,9,5,1)$, including its degree-four cell.

The indicated projection is a degreewise split
surjection of the full chain complexes. The
peripheral Fox coefficient in (E.P9) makes the
boundary square commute; higher degrees commute
as well by the product construction and the
zero higher groups of $C_K$. Let $D$ be its kernel.
The target has zero positive-degree homology, so
the long exact sequence gives
$H_2(D)\simeq H_2(\widetilde P)=\pi_2(P)$.
Using the small models for the relevant images,
the ranks of $D$ through degree three are
\[
                         (0,6,40,123).
 \tag{E.P13}
\]
Its six degree-one generators are the unused
regular-factor circles at the two vertices.
The six shifted boundary circles $u_i,w_i,d_i$
map to them by unit entries, because each circle
maps to zero in its own complement and identically
to the opposite regular torus. Cancel these six
unit pairs. This leaves a presentation of $H_2(D)$
with34 generators and123 relations, with no
uncomputed kernel.

Here is also an explicit block definition of that
matrix. In each $V_i$ keep the cells killed by
projection to $C_{K_i}\otimes C(S^1)$; their ranks
in degrees two and three are17 and45. In $B_2$
keep the33 cells containing at least one of
$u_1,w_1,d_1,u_2,w_2,d_2$. In those ordered bases
the matrix is
\[
 \mathcal M=
 \begin{pmatrix}
 (d_{V_1})_3&0&(j_1)_2\\
 0&(d_{V_2})_3&-(j_2)_2
 \end{pmatrix}:
 R^{45}\oplus R^{45}\oplus R^{33}
             \longrightarrow R^{17}\oplus R^{17}.
 \tag{E.P14}
\]
Every entry is specified by (E.P7)--(E.P11) and
(E.C19); (E.P14) itself is a complete finite
presentation of the desired module. The formula
includes the actual derived peripheral comparison.
In particular a splitting of the two cusp terms
has not been assumed from their localized homology.

\emph{Reduction to the compact presentation.}
The exact certificate
\path{code/maps_splice_module.py} checks both
peripheral chain equations and $d_2d_3=0$, cancels
the indicated unit pairs, and reduces (E.P14) to
16 generators and52 nonzero columns. It computes
in the actual amalgam
\[
 G=(K_1\times\mathbb Z)\mathop{*}_{\mathbb Z^3}
                                      (\mathbb Z\times K_2),
 \tag{E.P15}
\]
not in a finite quotient. To specify its normal
form, the torus-knot factor uses $h^r$ times an
alternating word in $a,a^2$ and $b,b^2,b^3$.
Modulo the peripheral subgroup, strip initial
$ab$ and $b^3a^2=(ab)^{-1}$, allowing their last
letter to be a partial power; choose $b^3$ in
the sole length-one tie with $a$. Inverting these
left coset representatives gives right coset
representatives. The usual alternating amalgam
normal form, followed by a final element of
$\mathbb Z^3$, is unique. These elementary word
operations are implemented exactly, including
both central exponents.

In the resulting16-generator presentation two
generators occur only in their respective single
relations $q_2e_1=0$ and $q_1e_2=0$. They give the
first two direct summands of (E.P2). For the other
14 generators, two generators $t_0,t_1$ have
trivial action of $a,b,c,d$. There are three
generators initially killed by $q_2$ whose pure
relations include $h_1N_b$; the identity
$h_1N_b(b-1)=1-h_1$ kills them by $q_1$ as well.
There are three initially killed by $q_1$ with
pure $N_c$ relations, and three more with pure
$h_2N_d$ relations. They are killed by $q_2$.
The remaining generators already have both
central relations. Thus the whole14-generator
summand is a module over $S$, without imposing
an additional quotient. The pure norm columns
are exact differences of the recorded columns
and are checked before setting either central
element equal to1.

For a readable final reduction, the14 generators
separate into two seven-generator blocks. The
first has generators $U_0,U_1,U_2,X,Y,Z,t_1$,
with $t_1$ invariant, $U_iN_b=0$, and
\[
 \begin{split}
 (U_0+X-Yd+Zd)N_a&=0,\\
 (U_1-Y(1+d)+Z)N_a&=0,\\
 (U_2+X+2t_1+Z(d-1))N_a&=0,\\
 X(1+d)+3t_1&=0.
 \end{split}
 \tag{E.P16}
\]
The second has generators $B_i,C_i,t_0$, with
$B_iN_c=C_iN_d=0$, $t_0$ invariant, and the single
additional relation $t_0=\Omega_R$. These formulas
follow by collecting the displayed16-generator
columns; the certificate retains every column
and its unit reductions.

In the first block put
\[
 v_0=X-Yd+Zd,\quad v_1=-Y(1+d)+Z,
 \quad v_2=X+2t_1+Z(d-1),\qquad
 r_i=U_i,\quad p_i=U_i+v_i.
 \tag{E.P17}
\]
The linear part of this change on $(X,Y,Z)$ is
\[
 \begin{pmatrix}1&0&1\\-d&-1-d&0\\d&1&d-1\end{pmatrix},
 \quad\det=1,
\]
over $\mathbb Z[d,d^{-1}]$. It is therefore an
integral group-ring basis change. Its inverse gives
\[
 X=v_0(1-d^2)+v_1(d^2-d)+(v_2-2t_1)d^2.
\]
Substitute in the last relation of (E.P16), using
the invariance of $t_1$. The result is
$t_1=\Omega_L$. Eliminate $t_1$ with this unit
coefficient. Its four invariance relations become
the four relations in (E.P3). Eliminating $t_0$
in the second block gives (E.P4). This proves the
compact presentation in (E.P2)--(E.P4), including
the integral central-intersection information.

For the two cyclic axis summands, the other
central element is a non-zero-divisor: their
group rings are free over its Laurent polynomial
ring. Thus neither contributes to the common
kernel of $q_1,q_2$. This proves (E.P5). On
coinvariants, the four invariance relations in
each block vanish, while the norm relations
become three copies of multiplication by3 and
three copies of multiplication by4. Each block
therefore has coinvariants $(\mathbb Z/12)^3$.
This is an independent recovery of E.2.11.

\emph{Restriction to the actual circle fibre.}
The universal covers of $P$ and $Y\times\mathbb R$
are $\Pi$-equivariantly homotopy equivalent by
E.2.13. The map $\Pi\to G/\langle h_i\rangle$
is injective, since $\chi(h_i)=\pm12$, and has
index12. Hence each of the first two group-ring
summands in (E.P2) is free of rank12 over
$\mathbb Z\Pi$. The last two factor through the
index12 surface subgroup $\Sigma<\Gamma$ and
have trivial $z$ action. This proves the module
formula in (E.P6).

For its exact coinvariants use the regular
12-sheet permutation representation with
exponents $(4,-3,-4,3)$ on $(a,b,c,d)$.
Evaluating the compact12-by20 presentation gives
an integral144-by240 matrix. Its Smith form has
56 unit entries and two entries6; its free
cokernel rank is86. The certificate computes
the same result from the independently retained
14-by26 presentation. Adding the24 free
coinvariants proves the last formula in (E.P6).
All these are integral module and homotopy
statements. No smooth surgery or smooth
classification is inferred from them.
\end{proof}


\paragraph{Proposition E.2.18 (Proved: the cusp quadratic attachment).}
Keep the right-module presentation of E.2.16 and write
\[
 R=\mathbb ZK,\qquad M=\pi_2(E)=eR\oplus T,
 \qquad T=\overline R^{\,2}\oplus L.
\]
For an abelian group $A$, let $\Gamma(A)$ be Whitehead's universal
quadratic group, with quadratic generator $q(a)$ and symmetric
bilinear cross effect
$[a,b]=q(a+b)-q(a)-q(b)$.  In particular $[a,a]=2q(a)$.
There is a lift of the normalized cusp component
$S=\mathbb {CP}^2\#3\overline{\mathbb {CP}}^{\,2}$ to the
universal cover $\widetilde E$ whose fundamental class has image
\[
 \pm\bigl([e(h-1),t]-12q(t)\bigr)
       \quad\hbox{under}\quad
 H_4(\widetilde E;\mathbb Z)\longrightarrow\Gamma(M).
 \tag{E.Q1}
\]
The sign is solely the orientation chosen on the normalized
component.  Its mixed component in
\[
 \Gamma(M)=\Gamma(eR)\oplus(eR\otimes_{\mathbb Z}T)
                                      \oplus\Gamma(T)
 \tag{E.Q2}
\]
is nonzero.  Consequently $E$ is not homotopy equivalent to
$E_0\#(S^2\times S^4)$ for any connected compact six-manifold
$E_0$, even without requiring the equivalence to preserve the
boundary.  Thus the free $R$ summand of its second homotopy group
does not come from a sphere-product connected summand.
Likewise the actual nine-dimensional fibre $Y$ is not
homotopy equivalent to $Y_0\#(S^2\times S^7)$ for any
connected closed nine-manifold $Y_0$, despite the twenty-four
free $\mathbb Z\Pi$ summands in (E.P6).

\begin{proof}
We use the natural, deck-equivariant part of Whitehead's exact
sequence
\[
 H_4(\widetilde E)\xrightarrow{\omega_E}\Gamma(\pi_2(E))
       \longrightarrow\pi_3(E)\longrightarrow
 H_3(\widetilde E)\longrightarrow0.
 \tag{E.Q3}
\]
The exactness and identification of its quadratic term are
\cite[Theorems 1 and 5]{mapsWhitehead50}; the direct-sum rule
(E.Q2) is proved there in Section~2.  These statements apply
to the simply connected universal cover, with no finiteness
requirement.  Naturality, rather than a choice of a word for
one attaching cell, will specify the element in (E.Q1).

For the normalization calculation take the negatives of
the complex-oriented Picard homology generators and denote
them by $(H,E_1,E_2,E_3)$.  Their intersection matrix is
still $\operatorname{diag}(1,-1,-1,-1)$.  This common reversal
matches the phase-cell orientations of (B.5.16): the toric
line with class $-(uw+6\gamma\delta)$ has complex curve
degrees $(-6,0,1)$, whereas its values on the corresponding
phase classes are $(6,0,-1)$; the support function $rs$
similarly fixes the middle sign.  The exact wall-degree
calculation in D.1.9 gives this comparison.  With these
orientations the six conductor classes are
\[
 E_1,\ H-E_1-E_2,\ E_2,\ H-E_2-E_3,
                                      E_3,\ H-E_1-E_3.
\]
The first opposite pair is self-glued in the surviving $u$
direction.  In the gamma cover its relation is therefore
$E_1=H-E_2-E_3$.  The other two pairs connect adjacent gamma
translates.  Their relations, after eliminating $H$, are
$(h-1)E_2=(h-1)E_3=0$.  Thus the image of normalization
homology is, with these actual curve generators,
\[
 A E_1\oplus(A/(h-1))E_2\oplus(A/(h-1))E_3,
 \qquad A=\mathbb Z[h,h^{-1}],\qquad H=E_1+E_2+E_3.
 \tag{E.Q4}
\]
This also follows directly from the transpose of the matrix
$B(h)$ of E.2.12, with the intersection signs in the Picard
basis retained.  It is the whole $H_2$ of the gamma cover:
the degree-zero normalization row is the cellular complex
of an infinite cylinder, whose second homology is zero.

We fix the comparison with the cusp basis of (E.C12), rather
than identify its free summand only abstractly.  In the ordinary
central fibre the first conductor class $\overline C_1$ is the
self-$u$ pair; the other two are the images of $E_2,E_3$,
in either order.  The chosen cusp basis is
\[
 E=-\overline C_1-\overline C_2,\qquad
 U=\overline C_2+\overline C_3,\qquad D=\overline C_2.
\]
The two torsion coordinates in (E.Q4) are determined by
this ordinary comparison, because each is annihilated by
$h-1$.  The normalization map is onto the whole group in
(E.Q4), and $E_2,E_3$ generate its entire $A$-torsion
submodule.  Consequently the coefficient of $E_1$ in the
free quotient is a generator of $A$, not merely a nonzero
element: it is a unit $\pm h^n$.  The ordinary comparison
fixes its sign.  Translating the entire chosen normalization
component by $h^{-n}$ makes this coefficient $-1$ and leaves
both torsion coordinates fixed.  The common Picard sign
convention above changes every generator simultaneously;
no arbitrary shear of the free generator is introduced.
The resulting
normalization map in the basis used for the full peripheral
comparison is
\[
 E_1\longmapsto-E-D,\qquad
 E_2\longmapsto D,\qquad E_3\longmapsto U-D,
 \qquad H\longmapsto-E+U-D.
 \tag{E.Q5}
\]
Exchanging $E_2,E_3$ has no effect below.  This argument
retains the full normalization map on $H_2$; it does not
suppose that the derived chain changes in E.2.16 are induced
by diffeomorphisms.

The intersection form of $S$ is
$\operatorname{diag}(1,-1,-1,-1)$.  Its usual four-cell
description, or the Hopf invariant of its two-dimensional
classes, therefore gives
\[
 \omega_S[S]=q(H)-q(E_1)-q(E_2)-q(E_3).
\]
Substitute $H=E_1+E_2+E_3$ and then (E.Q5).  Bilinearity
of the cross effect gives
\[
 \omega_S[S]\longmapsto
 [E_1,E_2]+[E_1,E_3]+[E_2,E_3]
                      =-[E,U]-2q(D).
 \tag{E.Q6}
\]
The normalization map $S\to E$ lifts to $\widetilde E$
because $S$ is simply connected.  Naturality of (E.Q3)
shows that the last expression is the image of its actual
four-dimensional fundamental class.  In particular no
unrecorded three-cell can invalidate the assertion that
this element belongs to $\operatorname{im}\omega_E$.

The exact global cancellation certificate now retains
the three cusp two-cycles in addition to the peripheral
chains.  It gives, in the generators $(x,t,y,z,e)$ of
E.2.16,
\[
 E\longmapsto(1-h)e+6t,\qquad
 U\longmapsto t,\qquad D\longmapsto0.
 \tag{E.Q7}
\]
For clarity, before the last relation reductions its
$t$ coefficient for $E$ is
$-3h+4ha+4ha^2+h^2$, whose augmentation is six;
its free coefficient is $1-h$.  Its coefficient for $D$
is $2h^{-1}a^2-1-a$, which has augmentation zero and hence
annihilates the invariant generator $t$.  The other
coordinates vanish.  The program
\path{code/maps_cusp_quadratic.py} checks these identities
using the identical integral group-ring cancellations as
E.2.16.  Applying $\Gamma$ to (E.Q7) in (E.Q6) gives
\[
 -[(1-h)e+6t,t]
                =[e(h-1),t]-12q(t),
\]
as asserted.  In ordinary coinvariants $t=2v$ and
$E=12v$ for a primitive class $v$.  The quadratic image
is then $-48q(v)$, consistent with the degree-two cup
square evaluated on the normalization; this is a check,
not the proof of (E.Q1).

We finish with the connected-sum obstruction, allowing
every possible free summand.  Multiplication by $h-1$
is injective on $R$: a finitely supported group-ring
element constant along an infinite $h$ orbit is zero.
It annihilates $T$, so
\[
 \operatorname{Ann}_M(h-1)=T,
                       \qquad\operatorname{Hom}_R(T,R)=0.
 \tag{E.Q8}
\]
Suppose that a copy $fR$ is a direct summand of $M$.
Write $f=eu+t_0$ and let a retraction send $e$ to $v$.
It vanishes on $T$, and hence $vu=1$.

The ring $\mathbb ZK$ is directly finite.  One elementary
verification uses the index-twelve subgroup of $K$:
its quotient by the central $\mathbb Z$ is the free group
of the punctured genus-three surface, so it is
$F_6\times\mathbb Z$.  This subgroup, and hence $K$, is
residually finite.  In every finite-group quotient the
regular representation is finite dimensional over
$\mathbb Q$, so $vu=1$ implies $uv=1$.  Residual
finiteness separates the finite support of $uv-1$ and
proves the same equality in $R$.  Thus $u$ is a unit,
the kernel of the retraction is exactly $T$, and every
free splitting is obtained by replacing $e$ by
$eu+t_0$.

Under such a change the mixed component of (E.Q1) is
$f u^{-1}(h-1)\otimes t$.  It is nonzero: $t$ has
infinite order by (E.C3), and $u^{-1}(h-1)$ is nonzero.
The central generator may also be reversed by a group
automorphism, which multiplies $h-1$ by a unit and
does not affect the argument.

If $E\simeq E_0\#(S^2\times S^4)$, its universal cover
has, through dimension four, the homotopy type of
$\widetilde E_0$ wedged with one copy of $S^2\vee S^4$
for each deck element.  This follows by lifting the
connected-sum balls; deleting or regluing their boundary
five-spheres does not alter the four-skeleton.  The new
two-spheres generate a free $R$ summand.  The four-spheres
have trivial attaching maps, and every remaining
four-dimensional Whitehead image lies in
$\Gamma(\pi_2(E_0))$.  There can therefore be no mixed
component between this free summand and its complement.
This contradicts (E.Q1)--(E.Q8) for every choice of the
summand, and proves the stated obstruction.

For the assertion about $Y$, first retain the actual maps
from the two complements to the splice.  The certificate
\path{code/maps_splice_quadratic.py} tracks the four cycles
$(t_0,e_0,t_1,e_1)$ through every unit cancellation used in
E.2.17.  In its sixteen-generator presentation they become
the four unit vectors in rows $(4,3,12,15)$, respectively
(rows are numbered from zero).  The subsequent seven-generator
block changes therefore give
\[
 t_0=\Omega_R,\qquad t_1=\Omega_L,
 \qquad e_0\in R/q_2R,\quad e_1\in R/q_1R
 \tag{E.Q9}
\]
with exactly the axes in (E.P2).  These inclusion maps have
not been selected from an abstract direct-sum decomposition.
Both $t_i$ are invariant under $G$: the four displayed
relations for each $\Omega$ in (E.P3)--(E.P4) annihilate
it by $a-1,b-1,c-1,d-1$.  The certificate also checks these
as the eight pure columns numbered
$(7,8,11,13)$ for $t_0$ and $(1,3,17,18)$ for $t_1$
in the retained fourteen-generator central presentation,
before eliminating the invariant generators.
On $\Gamma$-coinvariants
their respective images are
\[
 -2(B_2+C_2),\qquad 2(p_2-r_2),
 \tag{E.Q10}
\]
and each has order exactly six: the relevant two coordinates
have independent orders three and four in (E.P3)--(E.P4).
In particular an integer annihilating either $t_i$ in the
full module must be divisible by six.

Put $A_\Pi=\mathbb Z\Pi$ and $z=h_1h_2$.  Choose the twelve
coset representatives giving the free restriction of each
axis in (E.P6), and denote the resulting twenty-four generators
by $f_j$.  The normalized cusp surfaces and all their chosen
deck translates are actual maps to $\widetilde P$.  Transfer
them to $\widetilde Y$ using the $\Pi$-equivariant homotopy
equivalence $\widetilde P\simeq\widetilde Y\times\mathbb R$.
On the first axis $h_2$ is the identity, so multiplication
by $h_1-1$ equals multiplication by $z-1$; the analogous
statement holds on the second axis.  Naturality of (E.Q3)
therefore supplies, for every $j$, an element in the image
of the Whitehead map of $\widetilde Y$ of the form
\[
 \pm\bigl([f_j(z-1),t_{i(j)}]-12q(t_{i(j)})\bigr),
 \tag{E.Q11}
\]
where $i(j)=0$ on the first twelve axes and $1$ on the
remaining twelve.  Invariance of $t_i$ under $G$ removes
any extra coset action in this formula.

Write $\pi_2(Y)=F\oplus T_Y$, with $F=A_\Pi^{24}$ and
$(z-1)T_Y=0$ as in (E.P6).  Suppose a free copy $fA_\Pi$
is a direct summand, with a retraction $r:\pi_2(Y)\to A_\Pi$
satisfying $r(f)=1$.  Multiplication by $z-1$ is injective
on $A_\Pi$, so $r(T_Y)=0$.  Write $r(f_j)=v_j$ and
$f=\sum_j f_j u_j+t$; then
\[
                         \sum_j v_j u_j=1.
 \tag{E.Q12}
\]
In the splitting $fA_\Pi\oplus\ker r$, the mixed component
of (E.Q11) is $f v_j(z-1)\otimes t_{i(j)}$.  If all these
mixed components were zero, every integer coefficient of
$v_j(z-1)$ would be divisible by six, by (E.Q10).  The
group-ring basis makes this implication immediate even if
$t_i$ has larger or infinite order.  Multiplication by
$z-1$ is also injective on $(\mathbb Z/6)[\Pi]$: along an
infinite $z$ orbit a finitely supported coefficient function
with zero difference is zero.  Hence every $v_j$ belongs
to $6A_\Pi$, contradicting (E.Q12).

Finally, a connected sum with $S^2\times S^7$ has, through
dimension four in its universal cover, the original space
wedged with a free deck orbit of two-spheres.  All images
of its four-dimensional Whitehead map consequently have
zero mixed component with that free summand.  The preceding
argument rules this out for every free summand of the actual
module.  This proves the assertion for $Y$.
\end{proof}

This computes a specified part of the actual quadratic
attaching data.  It does not identify every higher handle
attachment or the smooth type of the shifted fivefold.


\paragraph{Proposition E.2.19 (Proved: the relative smooth cusp model).}
For the original marked cusp, let
\[
 Z(q)=\frac{\log q}{2\pi i}B_0+C(q),\qquad
 B_0=\begin{pmatrix}0&1\\-1&0\end{pmatrix},
\]
where $C(q)$ is holomorphic and bounded on a disc about zero.
For a sufficiently small closed disc the actual six-dimensional
cusp tube is orientation-preservingly diffeomorphic, relative to
its full marked torus-bundle boundary, to the unitless toric
quotient defined by $C=0$.  The latter has the following finite
normal-crossings plumbing description.

Start with the canonical complex line bundle of
$S=\mathbb {CP}^2\#3\overline{\mathbb {CP}}^{\,2}$ and its six
toric curves in cyclic order
\[
 \begin{aligned}
 C_0&=E_1,& C_1&=H-E_1-E_2,& C_2&=E_2,\\
 C_3&=H-E_2-E_3,& C_4&=E_3,& C_5&=H-E_1-E_3.
 \end{aligned}
 \tag{E.T1}
\]
Pair $C_i$ with $C_{i+3}$ for $i=0,1,2$.  Both the surface-normal
line and the canonical line restricted to every $C_i$ have degree
$-1$.  In the resulting $\mathcal O(-1)\oplus\mathcal O(-1)$
coordinates $(z,a,b)$, with $z=0$ at the endpoint meeting
$C_{i+1}$ and $b$ the surface-normal coordinate, the pairing is
\[
             (z,a,b)\longmapsto(z^{-1},zb,za).
 \tag{E.T2}
\]
The function giving the cusp parameter is $q=zab$ in this
coordinate chart.  Formula (E.T2), not merely the induced map
of zero sections, specifies the normal framing of each plumbing.
The two triple points are the endpoint orbits
$\{0,2,4\}$ and $\{1,3,5\}$, where endpoint $i$ means
$C_i\cap C_{i+1}$.  Ordered local complex triples undergo a
cyclic permutation around each orbit, with threefold composition
the identity.  The plumbing is formed with collars and rounded
corners; its distinguished outer boundary is $|q|=\eta$.

This statement supplies the relative smooth cusp input for the
ordinary framed handle construction of E.2.21.  It does not
promote the homological cancellations of E.2.16 to smooth
cancellations.

\begin{proof}
We first construct the comparison by a smooth family, rather
than use the logarithmic homeomorphism in (B.5.7).  Set
$C_s(q)=(1-s)C(q)$ for $0\leq s\leq1$.  The affine toric cover
is fixed.  On its dense torus the deck translation for
$\lambda\in\mathbb Z^2$ becomes
\[
 (x,q,s)\longmapsto
 \bigl(q^{B_0\lambda}
       \exp(2\pi i C_s(q)\lambda)x,\ q,\ s\bigr).
 \tag{E.T3}
\]
The exponential is componentwise.  The powers of $q$ permute
the toric charts by translation of the $A_2$ triangulation,
while the exponential factors are nowhere-zero analytic units
on every chart.  Thus (E.T3) defines a smooth family of
holomorphic deck actions on a fixed smooth toric covering
space.  On a central stratum the lattice translation permutes
its bounded fan cells without fixed points.  Off the central
fibre, freeness and properness are exactly the nondegeneracy
of the real period lattice.

These assertions hold uniformly after shrinking $\eta$.
Indeed put $T=-\log|q|/(2\pi)$, and let $K$ bound
$\|\operatorname{Im}C_s(q)\|$ on the compact parameter interval
and closed disc.  For $T>K+1$ the matrix
$TB_0+\operatorname{Im}C_s(q)$ is invertible, since $B_0$ is
orthogonal.  A fundamental parallelogram has logarithmic
position
$B_0a+T^{-1}\operatorname{Im}C_s(q)a$ with
$a\in[0,1]^2$; it meets a uniformly bounded number of toric
charts.  More explicitly choose a triangle containing the
unperturbed point $B_0a$ and let $P$ be its unimodular vertex
matrix.  Its barycentric coordinates $t_j$ are nonnegative,
and the corresponding chart coordinates satisfy
$\log|z_j|=t_j\log|q|+
O(2\pi\|P^{-1}\|K\sqrt2)$.
There are only finitely many such triangles, so with
$c=2\pi\sqrt2\max\|P^{-1}\|$ every chosen chart coordinate
has modulus at most $\exp(cK)$, uniformly in $q,s$ and
all compact phases.  Their closed polydiscs have compact
closure in the affine charts.  These finitely many compact
chart images prove
compactness of the quotient over $[0,1]$ and $|q|\leq\eta$.
The central charts and their incidence relations are fixed.
Projection of this quotient to $s$ is consequently a proper
smooth submersion, including on every normal-crossings stratum.
Its only boundary is $|q|=\eta$ together with the parameter
end faces, and its boundary restriction is also a submersion.

On $|q|=\eta$ prescribe the marked real coordinates
$x=\exp(2\pi i(Z_s(q)a+b))$, $(a,b)\in T^4$.
They give a smooth product trivialization in $s$ and precisely
the prescribed mapping-torus marking.  Extend its horizontal
vector field over a collar, and then over the compact family.
Local toric product charts provide extensions tangent to
every central stratum; a partition of unity preserves that
property and the lift of $\partial_s$.  Integrating this
field proves the asserted relative diffeomorphism.  Its
orientation is the complex orientation on each fibre.
This is relative Ehresmann transport in an actual smooth
family; no smoothness is inferred from a logarithmic formula
at $q=0$.

It remains to identify the endpoint $C=0$ with the framed
plumbing data.  In the lattice $\mathbb Z^2$ use cyclic rays
\[
 (1,0),\ (0,1),\ (-1,1),\ (-1,0),\ (0,-1),\ (1,-1).
 \tag{E.T4}
\]
The star of the vertex zero in the $A_2$ triangulation consists
of the six triangles spanned by zero and consecutive rays.
The corresponding three-dimensional cones have primitive rays
$(v,1)$ and determinant one.  Their common divisor at the
central ray is the toric surface $S$.  A principal central
fibre has normal line
\[
             \mathcal O_S(-C_0-\cdots-C_5)=K_S.
 \tag{E.T5}
\]
The cyclic surface fan in (E.T4) also proves this identification
directly: it is the fan of the three torus-fixed blow-ups of
$\mathbb {CP}^2$, and each neighboring-ray relation gives
$C_i^2=-1$.  In the Picard basis (E.T1),
$K_S=-3H+E_1+E_2+E_3$, so $K_S\cdot C_i=-1$.

Near a double curve the two adjacent unimodular charts give
the total space of $\mathcal O(-1)\oplus\mathcal O(-1)$.
Its two base charts have coordinates
\[
 (z,a,b),\qquad
 (w,A,B)=(z^{-1},za,zb),\qquad q=zab=wAB.
 \tag{E.T6}
\]
The period translation taking a neighboring component back
to the chosen component exchanges the two normal directions.
It reverses the chosen endpoint coordinate on the opposite
toric curve.  In the first chart this is exactly (E.T2);
in the target infinity chart it is simply $(z,b,a)$, so it
extends smoothly over the endpoint.  The map is involutive,
preserves $q$, and its complex Jacobian determinant is $1$.
These equalities specify the normal-disk attachment, including
its sign and angular twisting.

The endpoint $C_i\cap C_{i+1}$ is sent to
$C_{i+2}\cap C_{i+3}$.  At the first endpoint write the
ordered ambient coordinates as the surface-normal coordinate,
the $C_i$-normal coordinate and the $C_{i+1}$-normal coordinate.
The next endpoint has their order cyclically permuted.
After three crossings the same coordinates return, not just
the same unframed planes.  Thus the two endpoint orbits give
ordinary triple charts $q=z_0z_1z_2$, with no additional
clutching loop at either triple point.

Equivalently, for a triangle with vertex matrix $P$ whose
columns are $(v_i,1)$, translation by $n\in\mathbb Z^2$
replaces it by $L_nP$, where
\[
 L_n=\begin{pmatrix}1&0&n_1\\0&1&n_2\\0&0&1\end{pmatrix}.
 \tag{E.T7}
\]
It acts as the identity in the corresponding translated
unimodular coordinates, since
$(L_nP)^{-1}L_nP=I$.  This verifies the preceding plumbing
and triple identifications in the actual toric charts.
More explicitly, order the star matrix as
$P_i=((0,1),(r_i,1),(r_{i+1},1))$ and reverse its
columns when using the coordinate order $(z,a,b)$.
For every $i$ the matrix of translation by $-r_i$ from
this chart to the opposite following-endpoint chart is
\[
 \begin{pmatrix}-1&0&0\\1&0&1\\1&1&0\end{pmatrix}.
 \tag{E.T9}
\]
Its monomials are exactly (E.T2).  In normal-first order,
$P_{i+2}^{-1}L_{-r_i}P_i$ is the cyclic permutation
matrix.  Thus the framing and triple-cycle assertions
are checked for all six actual fan identifications,
not imposed on an abstract plumbing diagram.
Taking collars of the component core, three double-curve
annuli and two triple charts, and rounding their outer
corners, gives the asserted finite plumbing.  These are
subcharts of the smooth toric quotient, so the construction
is Hausdorff and the collar identifications are the indicated
smooth ones; it is not a quotient formed by identifying
arbitrary closed disk-bundle neighborhoods.

Finally the period translations for $g,u$ are $-e_2,e_1$.
On $q=\eta e^{-2\pi i\theta}$ the marked expression is
$x=\exp(2\pi i((-\theta+iT)B_0a+b))$.
At $\theta=1$ it agrees with $\theta=0$ and
$b$ replaced by $b-B_0a$.  This is exactly the inverse
endpoint convention for
\[
 M_0=\begin{pmatrix}
 1&0&0&0\\0&1&0&0\\0&1&1&0\\-1&0&0&1
 \end{pmatrix}.
 \tag{E.T8}
\]
The comparison therefore preserves the full outer marking,
not merely the isomorphism class of the boundary manifold.
The six unimodular charts, all normal degrees, the plumbing
Jacobian and involution, both triple cycles and (E.T8) are
checked exactly in \path{code/maps_cusp_smooth_plumbing.py}.
\end{proof}


\paragraph{Proposition E.2.20 (Proved: the explicit smooth plumbing presentation).}
The following finite smooth plumbing presentation uses only standard
pieces and specified attaching maps. Its conversion into ordinary
framed handles and the relative smooth identification are given in E.2.21.

Let $F=\mathbb R^4/\mathbb Z^4$, with basis $(g,u,w,d)$, and use the
two matrices $A_j$ and fixed vectors $v_j$ of A.1.1.  Set
\[
 (m_1,m_2)=(3,4),\quad(\kappa_1,\kappa_2)=(1,-1),\quad
 w_j=\kappa_jv_j,\qquad L=\ker\gamma,
\]
so that $\Lambda=\mathbb Zw_j\oplus L$ and $A_jw_j=w_j$.
On $L$ the two matrices are
\[
 B_1=\begin{pmatrix}0&1&0\\-1&-1&0\\1&0&1\end{pmatrix},
 \qquad
 B_2=\begin{pmatrix}0&-1&0\\1&0&0\\0&1&1\end{pmatrix}.
 \tag{E.U1}
\]
Define the standard smooth four-manifolds
\[
 \mathcal S_j=(T^3\times\mathbb R)/
       ((l,\sigma+1)\sim(B_j^{-1}l,\sigma)),
 \qquad \mathcal N_j=\mathcal S_j\times D^2_r.
 \tag{E.U2}
\]
In particular the cap's oriented normal disk is part of the
definition, not an unspecified disk bundle.

The cap boundary is identified with $M_{A_j}$ by the following
explicit map.  Write $x=\phi w_j+l$ and use the interval marking
$(x,1)\sim(A_j^{-1}x,0)$ on $M_{A_j}$.  Then
\[
 [x,\theta]\longmapsto
 \left([l,\sigma=m_j\kappa_j\phi+\theta],
                         \ r=e^{2\pi i\kappa_j\phi}\right).
 \tag{E.U3}
\]
All angles on the right are taken in their stated mapping-torus
quotients.  In the ordered angular variables $(\phi,\theta)$
and $(\sigma,\alpha)$, where $r=e^{2\pi i\alpha}$, the map is
\[
 \begin{pmatrix}\sigma\\\alpha\end{pmatrix}
 =
 \begin{pmatrix}m_j\kappa_j&1\\\kappa_j&0\end{pmatrix}
 \begin{pmatrix}\phi\\\theta\end{pmatrix},
 \qquad
 \begin{pmatrix}\phi\\\theta\end{pmatrix}
 =
 \begin{pmatrix}0&\kappa_j\\1&-m_j\end{pmatrix}
 \begin{pmatrix}\sigma\\\alpha\end{pmatrix}.
 \tag{E.U4}
\]
These are integral unimodular matrices.

To define a compact six-manifold $\mathcal E_{\rm pl}$, take the
flat $F$-bundle over an oriented disc with three open interior
discs removed.  Its three marked inner loop monodromies are
$A_1,A_2,M_0$, with $A_1A_2M_0=I$.  Equivalently cut the planar
surface along three disjoint paths to its outer boundary, take
the product with $F$, and reglue the two banks by these three
integral linear diffeomorphisms.  This specifies smooth band
attachments, including their product normal coordinates.
Fill the first two boundary components by (E.U2)--(E.U4), and
the third by the finite canonical-bundle plumbing of E.2.19,
using its complete marked $M_0$ boundary map.  The remaining
boundary is marked as $F\times S^1_\mu$, with the outer meridian
orientation of (E.S1).  The core uses its zero section, so no
affine translation remains unspecified on that boundary.

For the actual original regular-fibre complement there is an
orientation-preserving diffeomorphism
\[
                  E\ \cong\ \mathcal E_{\rm pl}
                 \quad\hbox{relative to }F\times S^1_\mu .
 \tag{E.U5}
\]
Consequently the shifted fivefold has the explicit smooth
presentation
\[
 P\ \cong\
 (\mathcal E_{\rm pl}\times T^4)
 \ \cup_{\ (x_1,x_2,\theta)\mapsto(x_1,x_2,-\theta)\ }\
 (T^4\times\mathcal E_{\rm pl}).
 \tag{E.U6}
\]
Every piece in this presentation is a standard torus product,
a linear three-torus mapping torus times an oriented disk, or
the canonical disk bundle of the explicitly specified toric
blow-up with the normal-disk plumbings (E.T2).  All attachment
maps have been specified on their normal coordinates.

\begin{proof}
We verify first that (E.U2)--(E.U4) are the actual finite caps,
including the normal framing.  The original cap is
\[
 (F\times D^2_s)/
 \bigl((x,s)\sim(A_jx+v_j/m_j,e^{-2\pi i/m_j}s)\bigr).
 \tag{E.U7}
\]
The good-reduction real flat marking is smooth on the entire
root disc, so these are smooth coordinates even at $s=0$.
Because $\gamma A_j=\gamma$ and $\gamma(v_j)=\kappa_j=\pm1$,
\[
                  r=s\,e^{2\pi i\kappa_j\gamma(x)}
 \tag{E.U8}
\]
is invariant under the displayed action.  It is single-valued
on $F$, has the same absolute value as $s$, and is complex
linear on every normal disk.  The induced map
$[x,s]\mapsto([x],r)$ identifies (E.U7) smoothly with
$S_j\times D^2$.  An inverse is obtained by choosing any
representative $x$ and setting
$s=r e^{-2\pi i\kappa_j\gamma(x)}$; invariance makes this
independent of the representative in the quotient.

In the integral basis $(w_j,u,w,d)$ the affine generator is
\[
                  (\phi,l)\longmapsto
                     (\phi+\kappa_j/m_j,B_jl).
 \tag{E.U9}
\]
Thus $\sigma=m_j\kappa_j\phi$ gives exactly the mapping torus
in (E.U2).  This also supplies a direct inverse to the product
description: choose a representative $(l,\sigma)$, put
$x=(\kappa_j\sigma/m_j)w_j+l$ and
$s=r e^{-2\pi i\sigma/m_j}$.  Replacing that representative
by the defining mapping-torus relation changes $(x,s)$ by
the finite affine generator and hence not its class.

The actual boundary parametrization from (B.5.3) is
\[
 [x,\theta]\longmapsto
       [x+\theta v_j/m_j,e^{-2\pi i\theta/m_j}].
 \tag{E.U10}
\]
Substitute (E.U10) into (E.U8)--(E.U9).  The coefficient of
$\theta$ in the normal angle is
$-1/m_j+\kappa_j^2/m_j=0$, while the base angle is
$m_j\kappa_j\phi+\theta$.  This proves (E.U3) and (E.U4).
At $\theta=1$ that base angle advances by one, precisely
replacing $l$ by $B_j^{-1}l$; this checks the mapping-torus
seam.  A lattice change of $\phi$ changes the base angle by
$m_j\kappa_j$, and $B_j^{m_j}=I$, so it introduces no further
ambiguity.

The function of the cap that gives its original base parameter
is, in these smooth coordinates,
\[
                         q=r^{m_j}e^{-2\pi i\sigma}.
 \tag{E.U11}
\]
On the marked boundary it has angle
$m_j\alpha-\sigma=-\theta$, the original clockwise base
convention.  A positive normal meridian at fixed
$(l,\sigma)$ has $(\Delta\phi,\Delta\theta)=(\kappa_j,-m_j)$.
Since $\kappa_jw_j=v_j$, it is exactly $t_{v_j}\rho_j^{-m_j}$,
recovering the actual filling relation
$\rho_j^{m_j}=t_{v_j}$.  The normal circle is therefore not
silently exchanged for a once-around base loop.  The exact
integer checks of all these seam, framing and meridian
identities are in \path{code/maps_finite_cap_framings.py}.

Over the punctured planar core, the actual family has the
smooth real flat marking with linear monodromies $A_j$ and
its zero section.  The logarithmic finite gluing has exactly
the drift in (E.U10); after subtracting it, the remaining
translation has a single-valued real lift and is removed by
the collar isotopy already proved in B.5.1.  No nonzero
lattice winding is discarded.  At the cusp use the actual
relative smooth diffeomorphism of E.2.19.  Each comparison
has precisely the prescribed boundary value and can be
made product on a smaller collar.  They therefore glue to
the relative diffeomorphism (E.U5).  In particular the
normal framing of the unfilled ordinary fibre is unchanged.

Finally use the two hemispheres in the actual shifted
fibre-product decomposition of E.2.6.  On each hemisphere
the other factor is a regular marked torus bundle over a
disc, hence a product preserving its zero section and flat
marking.  Their common boundary has unchanged torus
coordinates and opposite normal angles.  Taking the product
of (E.U5) with those regular factors and gluing gives
(E.U6), with its induced complex orientation.  This is a
comparison of all smooth pieces and all boundary maps,
not an identification inferred from group or homology.
\end{proof}

E.2.21 explicitly converts the canonical-bundle self-plumbing and
all the exterior attachments into ordinary framed handles.
In the absolute planar-core description used here, the band part
contains three products
$T^4\times D^1\times D^1$ attached by the specified linear maps;
its product handles include index five.  Those high handles
cannot all disappear: if a compact six-manifold with this boundary
had only handles of index at most four, its dual relative
decomposition would make $\pi_1(\partial E)\to\pi_1(E)$ surjective,
contradicting (E.S1).  The displayed plumbing presentation is
therefore not being confused with a four-dimensional thickening
or a chain cancellation.


\paragraph{Theorem E.2.21 (Proved: ordinary framed handles and the smooth type).}
The following finite ordinary framed-handle prescription defines a
compact oriented six-manifold $\mathcal H_6$ with marked boundary
$T^4\times S^1$, and a closed oriented ten-manifold $\mathcal H_{10}$.
There are orientation-preserving diffeomorphisms
\[
 E\cong\mathcal H_6\quad\hbox{relative to the full marked boundary},
 \qquad P\cong\mathcal H_{10}.
 \tag{E.H0}
\]
The respective numbers of ordinary handles, in increasing index, are
\[
 (2,5,14,22,21,10,2),\qquad
 (4,31,119,290,485,574,485,290,119,31,4).
\]
The attaching embeddings and their normal coordinates are specified
below; these numbers alone are not the specification.  No handle
cancellation, smooth recognition from a chain complex, or uncomputed
homotopy-sphere summand is used.
The balls and complex normal charts carry their complex orientations;
orient the product blocks by their displayed coordinate embeddings.
For example the dense block below has orientation
$-d\arg X\wedge d\arg Y\wedge d\mu_1\wedge d\mu_2
\wedge d\operatorname{Re}q\wedge d\operatorname{Im}q$.
This fixes the orientation convention even when the chosen order of
angle and moment coordinates differs from the complex coordinate order.

\emph{1. The two triple-point balls.}
Use two complex coordinate balls, denoted $V_+$ and $V_-$, with
ordered coordinates $Z=(Z_0,Z_1,Z_2)$ and $W=(W_0,W_1,W_2)$.
The boundary faces are rounded coordinate faces of
\[
 |Z_i|\leq\epsilon_v,\quad |Z_0Z_1Z_2|\leq\eta,
 \qquad
 |W_i|\leq\epsilon_v,\quad |W_0W_1W_2|\leq\eta.
 \tag{E.H1}
\]
They are smooth six-balls.  Throughout the prescription, rounding is
in the nonnegative radii, preserves the angular coordinates, and is
fixed on the attaching collars.  Positive changes of the radii only
give the corresponding collar isotopy.  The complex products in
(E.H1) are the same function $q$.
The coordinate formulas may first be read on manifolds with corners;
apply the same radial collar reparametrization on the two sides of
each attaching face before smoothing their union.  This convention
also specifies the maps at a corner meeting $|q|=\eta$.

\emph{2. The three annular blocks, with their endpoint framings.}
For $i=0,1,2$ use a block $S^1\times D^1\times D^4$ with coordinates
\[
 z=\exp(s+2\pi i\theta),\quad -R\leq s\leq R,
 \qquad (a,b)\in D^4,\qquad q=zab.
 \tag{E.H2}
\]
The normal four-ball here is a rounded radial domain
$|a|,|b|\leq\epsilon_e$, $|ab|\leq\eta/|z|$;
its positive radial trivialization identifies it with $D^4$.
The ends attach to framed tubes about the following coordinate
circles.  A table entry $(Z_j;Z_k,Z_l)$ means that the curve
coordinate is $Z_j$ and the two ordered complex normal coordinates
are $Z_k,Z_l$; the same convention applies to $W$.
\[
\begin{array}{c|c|c}
 i&z=0\hbox{ end: }(z;a,b)&z=\infty\hbox{ end: }(w;A,B)\\ \hline
0&(Z_2;Z_1,Z_0)&(W_0;W_1,W_2)\\
1&(W_2;W_1,W_0)&(Z_1;Z_2,Z_0)\\
2&(Z_0;Z_2,Z_1)&(W_1;W_2,W_0)
\end{array}
\tag{E.H3}
\]
The full-dimensional endpoint map between the two end charts is
\[
                 (w,A,B)=(z^{-1},za,zb).
 \tag{E.H4}
\]
In particular the second attaching circle has angle $-\theta$,
and its two complex normal vectors each rotate by $+2\pi\theta$.
The factors $|z|$ are the specified positive radial rescalings.
They do not erase either angular rotation.  Every coordinate circle
of the two balls occurs exactly once in (E.H3).  The maps are
embeddings of $S^1\times\partial D^1\times D^4$, not just attaching
loops with an unspecified choice of normal framing.

\emph{3. The dense torus block and its full sector maps.}
Attach $T^2\times D^2\times D^2_q$ along
$T^2\times\partial D^2\times D^2_q$.  Here is a completely explicit
choice of its base disk and the attaching map.  Use the rays
$r_i$ of (E.T4) and the hexagon
\[
 \mathcal P=\{v\in\mathbb R^2:\langle v,r_i\rangle\geq-1\},
 \quad
 \operatorname{Vert}(\mathcal P)=
 \{(-1,-1),(0,-1),(1,0),(1,1),(0,1),(-1,0)\}.
 \tag{E.H5}
\]
On the positive real part of $(\mathbb C^*)^2$, use the coordinates
given by the toric moment map
\[
 \mu(X,Y)=
 \frac{\displaystyle\sum_{v\in\operatorname{Vert}(\mathcal P)\cup\{0\}}
                 |X^{v_1}Y^{v_2}|^2v}
      {\displaystyle\sum_{v\in\operatorname{Vert}(\mathcal P)\cup\{0\}}
                 |X^{v_1}Y^{v_2}|^2}.
 \tag{E.H6}
\]
Choose a small inward parallel hexagon, with its corners rounded
symmetrically, $\mathcal P_\epsilon\subset\operatorname{int}\mathcal P$.
The inverse image is $T^2\times\mathcal P_\epsilon$; its two circle
coordinates are $\arg X,\arg Y$.  On the boundary disk circle put
$(X,Y)$ equal to the positive inverse of $\mu$, multiplied by these
two phases.  On the six sectors of that boundary, the attaching
map into the preceding blocks is the following monomial map, with
the whole complex normal coordinate $q$ included:
\[
\begin{array}{c|ccc}
 i& n_i&a_i&a_{i+1}\\ \hline
0&q/(XY)&X&Y\\
1&q/Y&XY&X^{-1}\\
2&qX&Y&(XY)^{-1}\\
3&qXY&X^{-1}&Y^{-1}\\
4&qY&(XY)^{-1}&X\\
5&q/X&Y^{-1}&XY
\end{array}
\qquad q=n_i a_i a_{i+1}.
\tag{E.H7}
\]
Use the cyclic identifications
$(n_2,a_2,a_3)=(Z_1,Z_2,Z_0)$,
$(n_4,a_4,a_5)=(Z_2,Z_0,Z_1)$ and the analogous two cyclic
permutations of $W$ for vertices $3,5$.  Vertices $0,1$ are $Z,W$.
The chart at an edge is $(z,a,b)=(a_{i+1},a_i,n_i)$, with (E.H4)
at its other end.  These rules specify (E.H7) also where a sector
meets an annular block.  The transition on a common sector is the
monomial map $P_{i+1}^{-1}P_i$, where
$P_i=((0,1),(r_i,1),(r_{i+1},1))$; the sum of its exponent rows
is $(1,1,1)$, so it preserves $n_i a_i a_{i+1}$.  Thus the sector maps specify the
normal framing as well as the core map.  They have no unassigned
phase, clutching map, or choice of normal line.

\emph{4. Expansion into ordinary handles.}
We use the following fixed elementary rule on every parameterized
block in this prescription:
\[
 T^r\times D^k\times D^{n-r-k}
 \quad\hbox{attached on}\quad
 T^r\times\partial D^k\times D^{n-r-k}.
 \tag{E.H8}
\]
Cut each ordered parameter circle into its two closed semicircles,
with small endpoint collars.  The first semicircle is its zero
handle and the second is its one handle.  For a bit vector
$\varepsilon\in\{0,1\}^r$, the resulting ordinary handle is of
index $k+|\varepsilon|$.  Attach by increasing $|\varepsilon|$,
then lexicographically.  On the original attaching face use
the displayed full-dimensional attaching embedding of the block;
on a newly exposed parameter face use the identity endpoint map
in that circle.  The normal coordinates are the remaining ordered
semicircle coordinates followed by the original normal coordinates.
Corner rounding and the permutation putting negative coordinates
first complete an ordinary framed attaching embedding
$S^{k+|\varepsilon|-1}\times D^{n-k-|\varepsilon|}$.
This is a finite rule for the full maps, not an instruction to choose
some attaching map in a stated homotopy class.  For one parameter
it is the literal two-piece decomposition in
\cite[Section~4.1, equation~(16), p.~103]{mapsVogel09}.
Its iteration requires no classification theorem.

Apply (E.H8) first to (E.H2), with $(r,k,n)=(1,1,6)$, and then
to the core, with $(r,k,n)=(2,2,6)$.  Together with the two balls
this gives cusp handle counts $(2,3,4,2,1)$.

\emph{5. The torus bands and the finite caps.}
Take the marked flat torus bundle of E.2.20 over the disc with
three holes, regarding the cusp boundary as the incoming boundary.
Cut along the two marked disjoint arcs from the two finite boundary
circles to the outer circle.  The cut surface is an annulus;
its torus bundle is the collar of the already constructed cusp
boundary.  Restoring the two strips attaches two blocks
\[
                 T^4\times D^1\times D^1.
 \tag{E.H9}
\]
To fix their maps without a planar drawing, use the oriented ribbon
polygon whose successive based inner boundary words are
$A_1,A_2,M_0$ and whose outer word is their product.  The negative
bank of the $j$th cut has coordinate $(x,t)$, the positive bank
has $(A_jx,t)$, and the strip uses the product normal coordinate.
Reversing a bank orientation replaces that matrix by its inverse.
This is the cut-polygon definition of the flat bundle in E.2.20;
the identity $A_1A_2M_0=I$ specifies the unfilled outer boundary.
There is no affine term.  Thus (E.H9) has its full attaching
embedding, including the transverse interval, fixed by these
linear maps.  Expand each by (E.H8), with $(r,k,n)=(4,1,6)$.
The unused boundary circles are precisely the two finite boundaries
and the marked ordinary boundary.
For clarity, the initial bank coordinates on the already attached
cusp boundary are the explicit coordinates of E.2.19:
\[
 \begin{gathered}
 q=\eta e^{-2\pi i\theta},\qquad T=-\log\eta/(2\pi),\\
 (X,Y)=\exp\bigl(2\pi i((-\theta+iT)B_0a+b)\bigr),
 \qquad a=(g,u),\ b=(w,d).
 \end{gathered}
 \tag{E.H9a}
\]
The exponential is componentwise; at an interval seam use
$b\mapsto b-B_0a$.  Equations (E.H7) and (E.H9a) therefore
give the actual bank embeddings into the existing handle
boundary, including their derivatives.  They do not merely
identify that boundary abstractly as a mapping torus.

For the $j$th finite cap use its already explicit
$\mathcal S_j\times D^2_r$ and the entire boundary embedding
(E.U3).  The mapping torus $\mathcal S_j$ has the following product
block construction: a block $T^3\times D^1$ over the first
semicircle of $\sigma$, followed by a block $T^3\times D^1$ over
the second semicircle, with endpoint maps $I$ and $B_j^{-1}$.
The first is a three-parameter zero handle, the second a
three-parameter one handle.  As a filling relative to
$\mathcal S_j\times S^1_r$, the normal disk supplies two negative
directions.  Consequently the relative cap consists of
\[
 \begin{array}{c|c|c}
 \hbox{base-circle block}&(r,k,n)&\hbox{full attaching data}\\ \hline
 0&(3,2,6)&\hbox{(E.U3), the first semicircle, and the $r$-disk}\\
 1&(3,3,6)&\hbox{(E.U3), endpoint maps $I,B_j^{-1}$, and the $r$-disk}.
 \end{array}
 \tag{E.H10}
\]
More explicitly an attaching face in the normal direction is sent
to (E.U3) at $r=e^{2\pi i\alpha}$, and an attaching face in a
base direction is the identical face of the preceding base block,
crossed with $D^2_r$.  At the base seam use $B_j^{-1}$ on $l$ and
the identity on $r$.  The inverse angular formula (E.U4) supplies
the full map back to the boundary being filled.  In particular
the normal meridian is $v_j-m_j\rho_j$, not the base meridian.
Apply (E.H8) to the two blocks in (E.H10).  Each cap adds the
ordinary relative counts $(1,4,6,4,1)$ in indices $2,\ldots,6$.
This completes the definition of $\mathcal H_6$.

\emph{6. The closed ten-dimensional handlebody.}
For each handle of $\mathcal H_6$, take its product with the four
ordered torus circles and apply (E.H8).  This gives
$\mathcal H_6\times T^4$ with its full boundary framing.  On the
other side take the dual relative handle decomposition of
$T^4\times\mathcal H_6$: reverse the handle order and exchange
the negative and positive disks of every handle.  A handle of
index $j$ becomes one of index $10-j$, and its attaching map is
the actual belt embedding of the original handle.  Glue its
initial boundary by the explicit map
\[
                 (x_1,x_2,\theta)\longmapsto
                                    (x_1,x_2,-\theta).
 \tag{E.H11}
\]
All subsequent dual attaching maps are obtained by this fixed
exchange of the two coordinate disks, in reverse order, on the
previously specified embeddings.  Use the factor order $E,T^4$
on the first half and $T^4,E$ on the second; their interchange
has sign $(-1)^{24}=1$.  The boundary-angle reversal in (E.H11)
is the orientation-reversing boundary gluing.  This defines
$\mathcal H_{10}$ entirely by ordinary handles.  If
\[
 c(t)=2+5t+14t^2+22t^3+21t^4+10t^5+2t^6,
 \qquad
       c_{10}(t)=(c(t)+t^6c(t^{-1}))(1+t)^4,
 \tag{E.H12}
\]
the coefficients are the counts in the statement.  The expanded
lists of all $76$ and $2432$ handles, including their parameter
bits, index, coindex, side, duality and attaching-map macro, are
\path{code/maps_ordinary_handles.out.json}.  Equations
(E.H1)--(E.H11), (E.U3)--(E.U4) and the displayed integer matrices
are the definition of those macros.

\begin{proof}
The main geometric point is that the first three steps give
\emph{actual smooth blocks of the cusp tube}.  We give the collar
argument explicitly.  Work in the smooth toric quotient of E.2.19.
Choose disjoint small coordinate neighborhoods of the two triple
points.  The three incident coordinate axes in each neighborhood
are exactly the six ends in (E.H3).  Remove smaller endpoint disks
from each normalized double curve.  Its remaining annulus has
the two complex normal lines $\mathcal O(-1)|_{\rm annulus}$,
which are trivialized by $(a,b)$ in (E.T6).  The resulting normal
domain in the actual tube is exactly the domain in (E.H2).
Every ray in $\mathbb C^2$ meets its rounded boundary once and
transversely.  Indeed the three defining radius functions have
positive radial derivative wherever the corresponding boundary
face is active.  Positive radial rescaling therefore gives a
smooth family of four-balls as $z$ ranges over this compact
annulus.  Choose the radial change to be constant linear scaling
on a common sufficiently small ball about the origin and interpolate
monotonically only away from that ball; it is then smooth at the
origin as well as at the boundary and preserves the stated normal
frames.  The same argument for the coordinate and product
radius functions in (E.H1) gives six-balls at the triple points.
Rounding can be chosen constant on the end attaching collars.

There is no unverified compatibility of these end frames.  At
vertex $i$ the normal-first coordinates are $(n_i,a_i,a_{i+1})$.
At vertex $i+2$ the deck identification is their cyclic
permutation, and the threefold composition is the identity by
E.2.19.  At the two ends of $C_i$ the coordinates are
$(z,a,b)=(a_{i+1},a_i,n_i)$ and
$(w,A,B)=(a_{i-1},a_i,n_{i-1})$.  The ordinary
$\mathcal O(-1)^2$ transition is (E.H4).  Substituting the two
cyclic permutations gives exactly (E.H3).  This checks all six
ends, not just one representative.  Triviality of a normal bundle
over the annulus is not being used to replace its relative
endpoint framing by a constant framing.

We next choose the stratum collars so that their remaining
interface is the entire attaching face of the dense block.
On the normalized toric surface take a small collar of its
hexagonal toric boundary, choosing the corner neighborhoods
first and the edge strips second.  Choose the strips narrower
than the corner collars and invariant under the opposite-edge
inversion and the two cyclic triple identifications.  The
uncovered moment region is the rounded interior disk
$\mathcal P_\epsilon$.  This can be done directly with six
inward parallel affine inequalities in (E.H5), smoothing each
corner inside its already chosen corner neighborhood.  Thus the
common frontier in the smooth part of the normalization is
exactly $T^2\times\partial\mathcal P_\epsilon$, not an additional
collection of annuli or disks.

The map (E.H6) gives the claimed coordinates: in logarithmic
positive coordinates it is the gradient of a log-sum-exponential
function.  Its Hessian is twice the covariance matrix of the
seven affinely spanning exponent vectors, hence positive
definite.  Strict convexity and the growth in each supporting
direction give a diffeomorphism onto
$\operatorname{int}\mathcal P$.  The inverse on the compact
disk and its frontier is smooth.  The canonical line on this
free torus locus is trivial; more specifically the toric
function $q$ is a nonzero toric factor times its normal
coordinate.  Hence \emph{$q$ itself}, including its phase,
trivializes that normal line.  The portion of the truncated
tube over this disk is therefore exactly
\[
 (T^2\times\mathcal P_\epsilon)\times D^2_q,
 \quad\hbox{with interface}\quad
 (T^2\times\partial\mathcal P_\epsilon)\times D^2_q.
 \tag{E.H13}
\]
In extending the preceding moment cuts off the zero fibre,
hold their equations constant in this $q$ coordinate on an
open neighborhood of the compact frontier.  On their ends
use the product cuts of the annular and triple charts.  A
collar interpolation between these choices is an interpolation
of positive radial cut functions, fixed on each already chosen
end collar.  Their radial derivatives remain positive, so its
flow is a collar isotopy.  The cuts can consequently be chosen
to agree with (E.H13) exactly.  In particular no twisting of
the $q$ disk is introduced on the six sector seams.

These are blocks in a sufficiently small actual truncated
neighborhood of the compact central fibre.  Their central
parts cover that fibre: two triple points, the three open
double annuli and the remaining normalized surface core.
Compactness from E.2.19 permits one common sufficiently small
$\eta$: the union of their open chart neighborhoods contains
the compact zero fibre, so properness implies that it contains
the entire inverse image of a sufficiently small closed
$q$-disc.  The nested cuts partition this entire inverse
image, rather than just a chosen sub-neighborhood of it.
If a larger originally chosen tube is used, the
difference is an exterior collar: on the compact annular
region without critical points choose a vector field lifting
the radial direction of $q$, and integrate it.  On the
prescribed outer collar choose the lift in the marked real
torus coordinates of E.2.19.  It is then the identity on all
those marked coordinates during radial collar transport.
This proves that no additional component or handle, and no
unrecorded outer-boundary diffeomorphism, is being omitted.

The actual embedding in (E.H13) is (E.H7).  To check this,
the exponent matrix for a unimodular chart is $P_i^{-1}$,
applied to $(X,Y,q)$.  This gives each of the six displayed
rows.  Consecutive changes are $P_{i+1}^{-1}P_i$ and telescope
to the identity.  The translations identifying the two
triple orbits are exactly those already used in (E.H3).
All six normal coordinate multipliers in (E.H7) are thereby
fixed, and the sector seams are equal maps on overlaps,
not just homotopic maps.  Positive rescaling of the chosen
radial domains leaves these monomials unchanged before
the collar reparametrization.  This completes the smooth,
fully framed identification of the cusp with steps 1--3.

For completeness the handle expansion rule is itself a
geometric equality.  For one circle, write $S^1=I_0\cup I_1$.
First attach $I_0\times D^k\times D^{n-k-1}$; after rounding it
is a $k$-handle.  The remaining block is
$(I_1\times D^k)\times D^{n-k-1}$, a $(k+1)$-handle.
Its attaching face consists of the original negative face
and its two endpoint faces in the first handle.  The stated
maps on these three faces agree on their corners because
they are restrictions of the same product coordinates.
Iterate this argument on the ordered circle factors.
It proves (E.H8), including all newly created attaching
faces and their normal framings.

The planar core relative to its incoming circle has two
ordinary surface one-handles.  Its flat bundle is a product
over each strip, so restoring the two cut strips is exactly
(E.H9).  The zero section and the full lattice marking fix
the normal coordinates on these attachments.  Each cap in
(E.H10) is the actual smooth product of E.2.20.  Its relative
handle claim can also be checked by the function
$-|r|^2+\varepsilon f(l,\sigma)$, with a cutoff near the
incoming boundary: choose $f$ by a two-critical-point
function of $\sigma$, perturbed on its two critical tori
by the product Morse function on $T^3$.  On the remaining
arcs its base derivative dominates the perturbation, so
there are no further critical points.  The indices are
two plus the base indices, as in (E.H10).  The product
description above fixes its handles without needing this
alternative Morse verification.  Summing the actual blocks
gives
\[
 (2,3,4,2,1,0,0)+2(0,1,4,6,4,1,0)
                     +2(0,0,1,4,6,4,1)
       =(2,5,14,22,21,10,2).
\]
Every comparison has its prescribed boundary value, so the
relative diffeomorphism of E.2.20 gives the first assertion
of (E.H0).  Taking the products, the dual relative handles
and the actual boundary flip of E.2.20 gives the second.
This is an identification by explicit full-dimensional
embeddings of ordinary handles.  It neither asserts an
unproved connected-sum normal form nor leaves a cap called
``the original complement'' inside the prescription.

The exact checks of all exponent matrices, the six endpoint
placements, both normal angular rotations, all sector seams,
the total linear monodromy, and the expanded handle lists
are in \path{code/maps_ordinary_handles.py}.  Its saved output
has Euler numbers $2$ and $0$, respectively, as a check;
the smooth identification follows from the preceding block
and collar proof, not from these numerical equalities.

We finish by stating precisely why these data determine the smooth
type, including its top handles.  The relevant relative uniqueness
results are Milnor's Theorem~3.13, which identifies an elementary
cobordism from its full characteristic embedding, Theorem~1.4,
which makes smooth composition unique with the marked boundary
fixed, and the isotopy-extension Theorem~5.8
\cite{mapsMilnor65}.  In the present notation the characteristic
embedding is the whole map
$S^{j-1}\times D^{n-j}\longrightarrow\partial W$, not just
the attaching sphere or its homology class.  Its normal coordinates
are part of (E.H3)--(E.H11).  The elementary consequence needed
here can also be seen directly.  If $\phi_0,\phi_1$ are joined
by the specified isotopy of full attaching embeddings, extend
the velocity field on their images to a slightly larger boundary
collar and multiply it by a cutoff which is zero on the protected
marked boundary.  Its flow extends the attaching isotopy to an
ambient boundary isotopy.  Extending this flow across a product
collar of $W$, and using the identity on the new coordinate handle,
identifies $W\cup_{\phi_0}(D^j\times D^{n-j})$ with
$W\cup_{\phi_1}(D^j\times D^{n-j})$.  The same cutoff argument
applies to the small radial corner changes above.  This derives
the relative, framed version from the stated isotopy-extension
theorem, without assuming an extension of an arbitrary sphere
diffeomorphism over a disk.

Apply this comparison successively in the displayed block and
sector order.  The initial ball comparisons are their coordinate
identifications.  At each next stage the actual attaching map
and the prescribed attaching map agree in the computed normal
coordinates, or differ by the explicitly described positive
collar isotopy.  The comparison of that stage is consequently
the identity on the new coordinate handle and the product-collar
extension on the previous stages.  Theorems~3.13 and~1.4 give an
inductive relative diffeomorphism of the whole six-manifold.
The final ordinary boundary coordinates are unchanged.  The
same argument, with the actual belt embeddings and the specified
boundary flip, applies to every handle on the dual ten-dimensional
side.  In particular a top handle retains its entire attaching
map $S^{n-1}\to\partial W$; none is replaced by an arbitrary
capping disk after forgetting that map.  There is therefore no
remaining disk-gluing or exotic-sphere ambiguity.  The computed
ordered full framed attaching data, under this relative uniqueness
theorem, are the classifying data used for (E.H0); the preceding
algebraic invariants are not being used as a smooth classification.
\end{proof}




\section{F. The rank-three search and a parity obstruction}

\paragraph{Theorem F.1.1 (Proved: Euler number of a Kuhn cusp).}
Let $g:M^8\to S^2$ be a proper triangulable map, a $T^6$-bundle
off finitely many values.  Suppose its unique non-quotient special
fibre is the rank-three Kuhn toroidal fibre with period-lattice
index $d$, and all its other special reduced fibres are free
finite quotients of $T^6$.  Then $\chi(M)=6d$.

\begin{proof}
A fundamental cube has six unimodular tetrahedra.  Taking the
quotient by an index-$d$ sublattice multiplies every cell count
by $d$, so the quotient cell counts are $d(1,7,12,6)$.
A cell of dimension $r$ corresponds to a torus stratum
$(\mathbb C^*)^{3-r}$.  Only the $6d$ point strata contribute
to Euler characteristic.  Euler integration over the regular
base gives zero since $\chi(T^6)=0$, and a finite free torus
quotient has Euler characteristic $0/|G|=0$.  Additivity now
gives the formula.  This checks the rank-three instance of the
localization mechanism in K15.  It requires the specified
toroidal filling, not just its unipotent monodromy matrix.
\end{proof}

\paragraph{Theorem F.1.2 (Known, with derivation: characteristic numbers).}
For a closed almost complex eight-manifold with $c_1=c_3=0$
rationally,
\[
 6\mid\chi,\qquad
 3c_2^2=45\sigma-14\chi,\qquad
 \operatorname{Td}=(3\sigma-\chi)/48.
 \tag{F.1}
\]
If it is three-connected and $\chi=6$, then
\[
 H_4(M;\mathbb Z)=\mathbb Z^4,\quad
 \sigma=2,\quad Q_M=3\langle1\rangle\oplus\langle-1\rangle,
 \quad c_2^2=2,\quad\operatorname{Td}=0.
 \tag{F.2}
\]

\begin{proof}
Write $\chi^1$ for the index of the almost complex Dolbeault
symbol twisted by the complex cotangent bundle.  The index
formulas, expanded in \cite[(1.2)--(1.4)]{ffEsoleKang}, give
\[
 \operatorname{Td}=\frac{3c_2^2-c_4}{720},\qquad
 \chi^1=\frac{3c_2^2-31c_4}{180},\qquad
 \sigma=\frac{3c_2^2+14c_4}{45}.
\]
They remain index identities for an almost complex structure;
integrability is not needed for integrality of the symbol index.
Subtracting four times the first from the second gives
$\chi^1-4\operatorname{Td}=-\chi/6$, proving the divisibility.
Eliminating $c_2^2$ gives (F.1).

Three-connectivity and Poincare duality leave cohomology only
in degrees zero, four and eight, and force $H_4$ to be free.
Thus $b_4=\chi-2=4$.  The Todd congruence gives
$\sigma\equiv2\pmod{16}$, and $|\sigma|\leq4$ forces
$\sigma=2$.  The middle pairing is unimodular.  In rank four
the even unimodular alternatives have signature zero, while
the odd alternatives are diagonal, by
\cite[Theorem 2.2.1]{ffPetersSterk}.  Its positive and negative
indices are three and one, giving (F.2).  The remaining numbers
follow by substitution in (F.1).
\end{proof}

\paragraph{Theorem F.1.3 (Known, with derivation: Mandal's example).}
The oriented manifold
$M_{\ell,m}=\ell\mathbb{HP}^2\#m\overline{\mathbb{HP}^2}$
admits an orientation-compatible almost complex structure exactly
when $\ell=2m+1$ and $m$ is odd.  In particular $M_{3,1}$ is
almost complex and has the invariants (F.2).

\begin{proof}
This is the exact criterion of Mandal's Corollary 1.4
\cite{ffMandal}; we read both that statement and its proof in
Section 5.  Here is its application in the normalization needed
here.  The tangent invariants in Mandal's normalized basis are
$p_1=2\sum x_i$, $p_2=7(\ell-m)u$ and
$e=(\ell+m+2)u$.  The stable-complex congruence of his
Theorem 1.1 becomes
$7(\ell-m)\equiv\ell-m\pmod4$, or $\ell-m$ even.
The rank-eight criterion of his Theorem 1.3(iii) becomes
$6(\ell-m)=2(\ell+m+2)$, or $\ell=2m+1$.
Together they say precisely that $m$ is odd.  At $(3,1)$ we
obtain $\chi=6$, $\sigma=2$, $p_1^2=8$ and $c_2^2=2$.
An integrable complex structure is not asserted by that theorem.
The later preprint arXiv:2610.02147 appeared in the literature
search, but its full text could not be opened in this run and
is not used as an authority.
\end{proof}

\paragraph{Proposition F.1.4 (Proved: the diagonal obstruction).}
If $A\in\operatorname{GL}_n(\mathbb Z)$ has no eigenvalue one,
every affine map $x\mapsto Ax+b$ of $\mathbb R^n/\mathbb Z^n$
has a fixed point.  Thus the diagonal elliptic finite monodromies
in Lead 7 admit no free affine lifts.

\begin{proof}
The real map $I-A$ is invertible.  A real solution of
$(I-A)x=b$ maps to a fixed point in the torus.  The diagonal
copies of the order-four and order-six elliptic rotations have
no eigenvalue one, so the conclusion applies to every translation
part, not only to integral or invariant choices.
\end{proof}

\paragraph{Proposition F.2.1 (Proved: the cusp normal form).}
Under the integral index-one condition in F.2, the cusp matrix is
conjugate in $\operatorname{GL}_6(\mathbb Z)$ to
\[
 U=\begin{pmatrix}I_3&I_3\\0&I_3\end{pmatrix}.
\]
Its integral centralizer consists exactly of
$\left(\begin{smallmatrix}P&Q\\0&P\end{smallmatrix}\right)$
with $P\in\operatorname{GL}_3(\mathbb Z)$ and
$Q\in\operatorname{Mat}_{3\times3}(\mathbb Z)$.

\begin{proof}
Let $N=M_0-I$.  The index-one condition says that the image
is the saturated kernel lattice and that
$\Lambda/\ker N\to\operatorname{im}N$ is an isomorphism.
Choose a basis $a_1,a_2,a_3$ of this image and lifts $b_i$ with
$Nb_i=a_i$.  Saturation shows that the six vectors form an
integral basis.  It gives the displayed matrix.  Commuting a
block matrix with $N$ forces its lower-left block to vanish
and its two diagonal blocks to agree.  Invertibility over
$\mathbb Z$ is then equivalent to unimodularity of that block.
\end{proof}

\paragraph{Proposition F.2.2 (Proved: a complete finite-order filter).}
The linear order of a finite matrix in $\operatorname{GL}_6(\mathbb Z)$
which admits a free affine lift belongs to
\[
 \{1,2,3,4,5,6,8,10,12\}.
 \tag{F.3}
\]
There are 38 possible degree-six cyclotomic characteristic
polynomials with an eigenvalue one, and 19 of them commute
with some real complex structure.  These are necessary types,
not 38 or 19 conjugacy classes of admissible tuples.

\begin{proof}
Proposition F.1.4 requires a nonzero one-eigenspace.  The other
cyclotomic factors have total degree at most five.  The elementary
identity
\[
 \frac{\phi(n)^2}{n}=
 \prod_{p^a\parallel n}p^{a-2}(p-1)^2\geq\frac12
\]
implies $n\leq50$ when $\phi(n)\leq5$.
The complete enumeration through 50 gives
$n=1,2,3,4,5,6,8,10,12$.  Enumerating their multiplicities
with total degree six and at least one factor $T-1$ gives 38
polynomials.  A commuting complex structure exists on a real
finite-order module exactly when the multiplicities of the
real eigenvalues $1$ and $-1$ are even; nonreal pairs already
carry such a structure.  Imposing those two parity conditions
leaves 19.  Taking least common multiples of their root orders
gives exactly (F.3), also for the 38-polynomial list.
The code prints every polynomial and every exterior invariant
rank.  This is a bounded exhaustive enumeration for the
stated filter.  It does not bound affine multiplicities when
the affine action has larger order than its linear part.
\end{proof}

\paragraph{Proposition F.2.3 (Refuted: determination by peripheral data alone).}
Monodromy matrices and affine twist vectors alone do not
determine the homology of a compact glued manifold.  A cusp
filling, or an equivalent relative topological specification,
is additional data.

\begin{proof}
Use the constructed fourfold $Z$ and blow up a smooth point
of the smooth locus of its cusp fibre.  The operation takes
place entirely inside that filling.  It leaves every regular
fibre, every finite quotient model, every peripheral monodromy
and every affine twist unchanged.  The total space remains
smooth and compact, and its map to the same base remains
holomorphic.  Blowing up a point in complex dimension four
replaces a ball by the disc bundle of $\mathcal O_{\mathbb P^3}(-1)$.
Mayer--Vietoris along their common $S^7$ boundary adds a copy
of $\mathbb Z$ in degrees two, four and six.  Thus the Euler
number increases by three and the homology changes.  The
peripheral index-one unipotent matrix is unaffected.
The stronger specification of a particular height-one Kuhn
filling excludes this example.  With that extra specification
the question remains a substantive classification problem;
we do not claim to refute it.
\end{proof}

\paragraph{Theorem F.3.1 (Proved: the self-dual parity obstruction).}
Suppose a compact smooth eight-manifold has the torus fibration
and the specified Kuhn cusp of F.1.1 with index one, and its
other reduced fibres are free quotients of complex three-tori
by finite holomorphic affine actions.  If its rational
$H^1$ monodromy representation is self-dual, then it cannot
satisfy $b_1=b_2=0$.  In particular no symplectic-monodromy
datum of F.2 can give the three-connected fourfold required
in F.4, for either $k=2$ or $k=3$ (or any number $k$ of finite
points).

\begin{proof}
Let $\mathcal V$ be the rank-six cohomology local system on
the punctured sphere and let $j$ be its inclusion.  On a finite
quotient fibre, averaging gives $R^1g_*\mathbb Q=(j_*\mathcal V)$
at that point.  At the Kuhn cusp, its normalization components
are simply connected and the degree-one branch complex is
the cochain complex of its dual $T^3$.  Specialization takes
its three deck classes to the covectors annihilating the
angular lattice, exactly the three-dimensional monodromy
invariant subspace.  Hence the same equality holds at the cusp.
It holds everywhere.

Put $a=\dim H^0(B,j_*\mathcal V)$ and
$b=\dim H^2(B,j_*\mathcal V)$.  The only differential that
can kill a class in $H^0(R^1)$ or in $H^2(R^0)$ is
$d_2:H^0(R^1)\to H^2(R^0)=\mathbb Q$.  If $b_1=b_2=0$,
it must be an isomorphism, so $a=1$.  Moreover
$H^1(B,j_*\mathcal V)$ must vanish because its position
$(1,1)$ has no possible incoming or outgoing differential.
Self-duality and compactly supported duality give $b=a=1$.

At the $i$th finite point, let $r_i$ be the dimension of the
fixed subspace in $H^1$.  The derivative is complex linear,
so its one-eigenspace has even real dimension.  Each $r_i$
is therefore even.  The cusp fixed dimension is three.
Euler additivity gives
\[
 \chi(B,j_*\mathcal V)=6(1-k)+3+\sum_{i=1}^k r_i,
\]
an odd integer.  Consequently
\[
 \dim H^1(B,j_*\mathcal V)=a+b-\chi(B,j_*\mathcal V)
\]
is odd and cannot vanish.  This is the required contradiction.
A preserved nondegenerate alternating form identifies a
representation with its dual, so symplectic monodromy is
included.  The explicit form
$(6\gamma\wedge\delta+u\wedge w)\oplus\omega_E$
checks this hypothesis for the example of Part E.
\end{proof}

\paragraph{Proposition F.3.2 (Proved: the remaining rank patterns).}
If a possibly non-self-dual holomorphic datum with the same
local models had $b_1=b_2=0$, put
$b=\dim H^0(B,\mathcal V^*)$.  For $k=2$, the only possible
unordered pairs $(r_1,r_2;b)$ are
\[
 (2,2;0),\qquad(2,4;2).
\]
For $k=3$, they are
\[
 (2,2,6;0),\ (2,4,4;0),\ (2,4,6;2),\ (4,4,4;2).
 \tag{F.4}
\]

\begin{proof}
The preceding Leray argument still gives $a=1$ and vanishing
$H^1(B,j_*\mathcal V)$.  Euler additivity therefore gives
$b=\sum r_i+3-1-6(k-1)$.  Each $r_i$ belongs to
$\{2,4,6\}$ because a free affine lift requires an eigenvalue
one and complex linearity makes its multiplicity even.
Also $0\leq b\leq\min(3,r_1,\ldots,r_k)$, by inclusion in
each local invariant subspace.  Enumerating these three-value
lists gives exactly (F.4), as independently printed by the
complete bounded filter.  No existence of a tuple or of a
period family with these ranks is asserted.
\end{proof}

\paragraph{Theorem F.2.5 (Proved: a non-self-dual topological candidate).}
There is a closed smooth oriented eight-manifold with the prescribed
index-one Kuhn cusp, two free affine fillings of multiplicities $5,4$,
trivial fundamental group, and rational Betti numbers
\[
 (1,0,0,0,4,0,0,0,1). \tag{F.5}
\]
It admits an almost complex structure.  Its signature is two and
its middle intersection lattice modulo torsion is
$3\langle1\rangle\oplus\langle-1\rangle$.
Integral $H_2,H_3$ and a holomorphic period realization are not
asserted here.

\begin{proof}
Define the two cohomological matrices
\[
A=\begin{pmatrix}
1&0&0&1&-1&1\\0&1&-1&-1&1&0\\0&1&0&-1&0&0\\
0&2&-1&-1&0&1\\0&2&0&-1&0&0\\0&1&1&-1&0&0
\end{pmatrix},\quad
B=\begin{pmatrix}
1&0&0&0&0&0\\0&0&-1&0&1&0\\0&0&-1&0&0&1\\
0&0&-1&0&0&0\\0&1&-1&0&0&0\\0&0&-1&1&0&0
\end{pmatrix}.
\]
Use homology monodromies $M_1=A^t$, $M_2=B^t$ and
$M_0=(M_1M_2)^{-1}$, with twists
\[
 v_1=(1,3,-2,-1,2,0)^t,\qquad v_2=(-1,0,0,0,0,0)^t .
\]
Exact multiplication gives orders five and four.  For $N=M_0-I$,
$N^2=0$, $\operatorname{rank}N=3$, and the nonzero Smith factors
are $1,1,1$.  The finite characteristic polynomials are
$(T-1)^2\Phi_5(T)$ and $(T-1)^2(T+1)^2(T^2+1)$.
Thus each finite matrix commutes with an orientation-compatible
complex structure.  The primitive invariant covector
$\ell(x)=x_0$ has $\ell(v_1)=1$ and $\ell(v_2)=-1$.
For $\tau_i(x)=M_ix+v_i/m_i$, every nonidentity power is free,
because its image under $\ell:T^6\to\mathbb R/\mathbb Z$ is
translation by a nonintegral multiple of $\pm1/m_i$.

The finite fillings
$(T^6\times\Delta)/\langle(\tau_i,\zeta_{m_i})\rangle$ are smooth.
Take the stipulated Kuhn filling at the third point, marked by
Proposition F.2.1.  The torus bundle over the intervening pair of
pants has monodromies $M_1,M_2,M_0$.  Translations are isotopic to
the identity, so its boundary bundles identify with those of the
three fillings.  Glue using the displayed twists and zero cusp
clutching.  This proves existence and compact smoothness.

For van Kampen, the normal closure of the cusp angular lattice is
already generated by the columns of
$[\,N\mid M_1N\mid M_2N\,]$.
Its Hermite basis is exactly $e_1,\ldots,e_5$.  The residual fibre
group is therefore the central group $\mathbb Z\langle h\rangle$,
and the remaining meridians satisfy
$5a=h$, $4b=-h$, $a+b=0$.
Their three-by-three relation matrix is unimodular.  Hence
$\pi_1=1$, with no uncomputed finite quotient.

The exact rational Leray input is
\[
\begin{array}{c|rrrrrrr}
q&0&1&2&3&4&5&6\\ \hline
h^0(j_*\bigwedge^q\mathcal V)&1&1&0&0&0&0&1\\
h^1(j_*\bigwedge^q\mathcal V)&0&0&0&3&0&0&0\\
h^2(j_*\bigwedge^q\mathcal V)&1&0&0&0&0&1&1
\end{array}
\]
where the cohomology monodromies are $A^{-1},B^{-1}$, whose
simultaneous invariant/coinvariant calculations agree with $A,B$.
The finite local invariant ranks are $(1,2,3,4,3,2,1)$ at both
points.  The cusp row is $(1,3,9,9,9,3,1)$, and Theorem E.1.3
adds one degree-four cusp kernel.  The degree-one transgression
is an isomorphism since $\pi_1=1$.  Poincare duality forces the
degree-six transgression to be an isomorphism.  No other rational
differential is possible, giving (F.5).

For the almost complex assertion, prescribe the existing complex
structures on the three fillings.  A constant linear complex
structure along each regular fibre is a section of the associated
bundle with fibre
$\operatorname{GL}_6^+(\mathbb R)/\operatorname{GL}_3(\mathbb C)$.
This fibre deformation retracts to
$\operatorname{SO}(6)/\operatorname{U}(3)\simeq\mathbb{CP}^3$.
It is connected and simply connected.  Relative obstruction
theory over the two-dimensional pair of pants therefore extends
the prescribed boundary sections.  A smooth section and compatible
horizontal distribution give an almost complex structure agreeing
with the fillings.  Rational $H^2,H^6$ vanish, so $c_1,c_3$
vanish rationally.  Theorem F.1.2 and the unimodular lattice
classification give the signature and free form.

The exact program \texttt{fourfolds\_nonselfdual.py} verifies every
matrix, Smith form, order, and exterior rank above.  It verifies
this one tuple, not a bounded simultaneous-conjugacy classification.
\end{proof}

\paragraph{Theorem F.2.6 (Proved: integral three-connectivity of the candidate).}
The manifold of Theorem F.2.5 has free integral homology with
Betti numbers (F.5).  In particular $H_2=H_3=0$, and it is
three-connected.

\begin{proof}
We identify the integral torsion differentials explicitly.
For the finite fillings, Smith forms give
$H_1(B_5)=\mathbb Z^2$ and
$H_1(B_4)=\mathbb Z^2\oplus\mathbb Z/2$.
The covering images on integral degree-two and degree-three
cohomology have index one.  In degree two use
$B_m\to S^1$, given by $m\ell$, with fibre $T^5$.
A primitive invariant vertical loop $\lambda$ supplies the
remaining suspension two-cycle.  An invariant integral
two-form $\omega$ descends exactly when
$\omega(v,\lambda)/m$ is integral; vertical periods are
already integral.  These three values on integral invariant
bases are $(-1,0,-1)$ at five and $(0,0,1)$ at four.
In degree three the determinants of the invariant middle
skew pairings are $5^4$ and $4^4$.  Pullback multiplies the
unimodular downstairs pairing by $m$, so the covering indices
are one.  The degree-two/degree-four invariant pairing has
determinant of absolute value $5^3$, respectively $4^3$.
Since the degree-two covering index is one, the degree-four
index is one as well.  Theorem E.1.3 gives integral
surjectivity at the cusp.

For coefficient exterior degrees two, three and four, the
global coboundary matrix obtained by stacking
$\bigwedge^qA-I$ and $\bigwedge^qB-I$ has one Smith factor
two and all its other nonzero Smith factors one.
Its torsion cocycle restricts nontrivially at the order-four
point.  Thus the integral parabolic groups in these degrees
are torsion free, of ranks $0,3,0$.  In degree one the global
coboundary image is primitive; the finite covering cokernels
are $\mathbb Z/5,\mathbb Z/4$, and the invariant covector
$\ell$ maps to a generator of each.  The Chinese remainder
theorem therefore gives $H^1(B,R^1g_*\mathbb Z)=0$.
The global coinvariants in coefficient degrees one through
four are each $\mathbb Z/2$.

For the crucial torsion differentials, write
$M_2=\operatorname{diag}(1,D)$.  The unimodular matrix
\[
P=\begin{pmatrix}
-1&0&-1&0&1\\-1&2&-1&1&3\\-1&1&-1&2&3\\
0&-1&0&-1&-1\\1&-1&2&-1&-3
\end{pmatrix}
\]
conjugates $D$ to a four-cycle permutation and a negative
one-dimensional summand.  Hence $B_4$ is a mapping torus of
a cellular torus automorphism.  Its torsion class in degree
$q$ is represented by
$w\in\operatorname{Tor}\operatorname{coker}
(\bigwedge^{q-1}D^t-I)$.
The boundary is $B_4\times S^1$ with original base coordinate
$\theta=a+4k$.  Write $h=k$ on its finite-cover fibre; then
$a=-4h$.  The fibre has the vertical $T^5$ and the loop
$(-4,1)$ in the $(a,k)$ coordinates.
The restriction of the suspension class to that fibre is
the norm $(I+D^t+(D^t)^2+(D^t)^3)w$, which is zero because
a torsion coinvariant representative lies rationally in
$\operatorname{im}(D^t-I)$.
The remaining component in the boundary Wang sequence is
exactly the original vertical coinvariant $w$.
This description follows on cellular chains by taking the
four consecutive fundamental prisms; the monomial marking
makes their boundary the displayed norm.

The base Mayer--Vietoris connecting map therefore identifies
the point-supported torsion Leray differential with
\[
d_2(w)=[w]\in
\frac{\bigwedge^{q-1}\mathbb Z^6}
 {(\bigwedge^{q-1}A-I)\bigwedge^{q-1}\mathbb Z^6+
  (\bigwedge^{q-1}B-I)\bigwedge^{q-1}\mathbb Z^6}.
\tag{F.6}
\]
Indeed fibre restriction vanishes, so its first nonzero
boundary restriction is the just-computed class in
$H^1(S^1,\bigwedge^{q-1}\mathcal V)$.
The connecting map quotients by the two global monodromy
relations, which is the two-step differential of the
Cech--fibre double complex.
Changing the sheaf at finitely many points by its finite
covering cokernels or point-supported kernels does not
change this degree-two base cohomology group.
Inverse monodromy conventions change the representatives
by units and do not change the nonzero classes below.

The exact Smith transformations in the code give the
following vertical torsion representatives, in lexicographic
exterior bases of $\mathbb Z^5$:
\[
\begin{split}
q=2:\ &(0,0,0,-1,-2),\\
q=3:\ &(0,0,0,0,0,0,0,1,0,-1),\\
q=4:\ &(0,0,0,0,0,0,1,0,0,0).
\end{split}
\]
Each has residue one in its global $\mathbb Z/2$ quotient.
Thus all three differentials in (F.6) are isomorphisms.
The relevant integral Leray groups are
\[
\begin{array}{c|ccccc}
q&0&1&2&3&4\\ \hline
H^0(R^q)&\mathbb Z&\mathbb Z&\mathbb Z/2&
 \mathbb Z/2&\mathbb Z\oplus\mathbb Z/2\\
H^1(R^q)&0&0&0&\mathbb Z^3&0\\
H^2(R^q)&\mathbb Z&\mathbb Z/2&\mathbb Z/2&
 \mathbb Z/2&\mathbb Z/2 .
\end{array}
\]
The free degree-one transgression is onto because $\pi_1=1$
and rational $H^2=0$: the universal coefficient theorem makes
$H^2$ torsion free, so the finite cokernel of this map
$\mathbb Z\to\mathbb Z$ must vanish.  Together with the three torsion
isomorphisms this gives $H^2=H^3=0$ and $H^4=\mathbb Z^4$;
the two free filtration pieces of $H^4$ have ranks three
and one.  The universal coefficient theorem now gives
$H_2=H_3=0$.  Hurewicz gives three-connectivity, and Poincare
duality determines all remaining integral groups.
\end{proof}

\paragraph{Theorem F.3.4 (Proved: exclusion of the abelian-surface source).}
The Engel construction proposed in F.3, starting with a polarized
algebraic family of abelian surfaces with a toric-rank-two cusp,
taking a fibrewise degree-zero line bundle, and quotienting its
punctured bundle by a lifted translation, cannot have $b_1=b_2=0$
with the stipulated rank-three square-zero Kuhn cusp.
The two sections need not be proportional.

\begin{proof}
Let $W$ be the rank-four rational homology local system of the
abelian surfaces.  It is semisimple by
\cite[Definition 4.2.4, Proposition 4.2.5 and Theorem 4.2.6]{ffDeligne}.
Its toric-rank-two cusp has
$N_W^2=0$ and $\ker N_W=\operatorname{im}N_W$ of dimension two.
Semisimplicity splits $W=W^\Gamma\oplus W'$.
The image of $N_W$ lies in $W'$ whereas $W^\Gamma$ lies in its
kernel.  The equality of image and kernel forces $W^\Gamma=0$.
The coinvariants vanish as well.

The rank-six rational homology system $T$ after the two Engel
operations has a natural filtration
\[
 0\subset L\subset E\subset T,\qquad
 L\simeq\mathbb Q,\quad E/L\simeq W,\quad T/E\simeq\mathbb Q.
\]
The lower line is the circle in the punctured line bundle and
the upper quotient is the lifted-translation generator.
No polarization on $T$ is assumed.
The triviality of $T/E$ gives $\operatorname{im}N_T\subset E$.
Maximal square-zero cusp monodromy gives
$\ker N_T=\operatorname{im}N_T$; thus every globally fixed vector
lies in $E$.  Since $W^\Gamma=0$, the fixed vectors of $E$ are
exactly $L$.  Hence $\dim T^\Gamma=1$.
The dual filtration, whose middle term is $W^*$, gives
$\dim(T^*)^\Gamma=1$ by the identical argument.

Consequently the two ranks $a,b$ in Theorem F.3.1 are both one.
Its Euler-parity calculation applies without requiring self-duality
of $T$, and contradicts $b_1=b_2=0$.
This excludes the proposed algebraic abelian-surface source and,
more generally, every such filtration with semisimple middle term.
It does not exclude an arbitrary unpolarized torus period family.
\end{proof}

\paragraph{Theorem F.2.7 (Proved: the orbifold quotient).}
For the specified torus fillings in F.2, with affine
multiplicities $m_1,\ldots,m_k$, the fundamental group of the
total space has a canonical surjection onto
\[
 \langle x_1,\ldots,x_k\mid
       x_i^{m_i}=1,\ x_1\cdots x_k=1\rangle.
 \tag{F.27}
\]
Consequently a simply connected example with $k=2$ requires
$\gcd(m_1,m_2)=1$.  There is no simply connected example with
$k=3$ and $m_1,m_2,m_3\geq2$.  An affine multiplicity-one
point is an ordinary smooth point and reduces to the problem
with one fewer multiple fibre.

\begin{proof}
Over the complement of small discs at the special values,
the torus bundle projects surjectively on fundamental groups
to the punctured base.  Quotient its fundamental group by the
normal closure of all regular-fibre loops.  At a finite
filling its root cover is a smooth torus bundle over a disc,
and its free deck group is cyclic of order $m_i$.
The covering exact sequence shows that the quotient of the
filling group by its fibre subgroup is precisely this cyclic
group.  The boundary base meridian maps to its generator.
Thus adjoining that filling imposes exactly
$x_i^{m_i}=1$ after the fibre loops have been killed.
At the index-one Kuhn cusp the fundamental group is generated
by the three deck loops of its dual torus; these are regular
fibre loops.  Its base meridian contracts.  Hence its filling
imposes the ordinary meridian relation and leaves no extra
generator in this quotient.  The base relation gives the
product relation in (F.27).  Van Kampen proves that (F.27) is
the resulting quotient, in particular the asserted surjection.

For two points, substitute $x_2=x_1^{-1}$; the quotient is
$\mathbb Z/\gcd(m_1,m_2)$.
For three points of orders at least two it is a nontrivial
triangle group.  One can see nontriviality without assuming
a classification of the total space.  Take a geodesic triangle
with angles $\pi/m_1,\pi/m_2,\pi/m_3$ in the sphere, Euclidean
plane or hyperbolic plane according as their sum is greater
than, equal to or less than $\pi$.
In the spherical case the possible orders are, after sorting,
$(2,2,n),(2,3,3),(2,3,4),(2,3,5)$, and the corresponding
spherical triangles exist.  The Euclidean and hyperbolic
angle conditions give the other triangles directly.
If $r_1,r_2,r_3$ are reflections in the opposite sides,
the rotations $r_2r_3,r_3r_1,r_1r_2$ have the stated
orders and product one.  They therefore give a nontrivial
representation of (F.27) into the orientation-preserving
isometry group.  Its quotient cannot be trivial.
Finally a free cyclic action of order one is the identity
construction on a root disc of order one, hence an ordinary
torus bundle filling and not a multiple fibre.
\end{proof}

\paragraph{Proposition F.2.15 (Proved: automatic local complex compatibility).}
Every simply connected datum of F.2 with two finite affine
fillings automatically has finite linear actions commuting
with real complex structures.  This assertion concerns the
individual finite actions and does not produce simultaneous
global periods.

\begin{proof}
Let $m_i$ be the order of the free cyclic affine action
defining the $i$th filling.  Its linear order divides $m_i$.
By F.2.7 the integers $m_1,m_2$ are coprime, so at least one
of them is odd.  The corresponding linear determinant is
one, because a finite real matrix of odd order cannot have
determinant minus one.  The relation
$M_1M_2M_0=I$ and the unipotence of $M_0$ imply
$\det M_1\det M_2=1$, so both determinants are one.

A finite-order real matrix is semisimple.  Its nonreal
eigenvalues occur in conjugate pairs, and its only possible
real eigenvalues are $1$ and $-1$.  Determinant one says
that the $-1$ eigenspace has even real dimension.  Since the
total dimension is six, the $1$ eigenspace also has even
real dimension.  Choose a complex structure on each of
these real eigenspaces by pairing a real basis.  On each
nonreal isotypical real summand, multiplication by $i$ in
its complex eigenspace supplies a commuting complex
structure.  Their direct sum commutes with $M_i$.
The fixed eigenspace is nonzero because the affine action
is free, as in F.1.4, so its even dimension is at least two.

Each chosen structure makes $\mathbb R^6/\mathbb Z^6$ a
complex torus and the given affine map holomorphic; thus it
supplies precisely the local complex-compatibility
hypothesis used in the subsequent reductions.  For three genuine
finite fillings F.2.7 has already excluded simple
connectivity without this argument.
\end{proof}

\paragraph{Proposition F.2.8 (Proved: exclusion of an involutory finite part).}
A simply connected datum with $b_2=0$, $k=2$ and the stated
finite affine fillings has neither finite linear
monodromy of order one or two.  Its unordered pair of finite
linear orders must therefore belong to
\[
 \{(3,4),(3,5),(3,8),(3,10),(4,5),
       (5,6),(5,8),(5,12)\}.
 \tag{F.28}
\]
The affine multiplicities may still be larger than these
linear orders.

\begin{proof}
An identity finite generator would make the product of the
two finite matrices finite, contrary to its nontrivial
unipotent cusp.  Suppose the cohomology generators are
$A,B$ with $A^2=I$, where the other finite linear order is
odd.  It must be odd because the linear orders divide the
coprime affine multiplicities of Theorem F.2.7.
Write $U=AB$; the cusp condition gives $(U-I)^2=0$.
Multiplication of this identity on the left by $A$ and on
the right by $B^{-1}$ gives
\[
 BA+AB^{-1}=2I.
 \tag{F.29}
\]
For $v\in\ker(B-I)$ we obtain $(B+I)Av=2v$.
As $B$ has odd order, $B+I$ is invertible, so $Av=v$.
Every $B$-fixed vector is therefore globally fixed.
The finite free affine action has a nonzero fixed
linear subspace, and F.2.15 makes its real dimension even; it is
at least two.  Consequently the global cohomology invariant
rank is at least two.  The Leray argument of Theorem F.3.1
shows that $b_1=b_2=0$ requires that rank to be exactly one,
a contradiction.  Inverse or transpose monodromy conventions
give the same identity and fixed-space conclusion.

Finally take the complete finite-order list (F.3), discard
one and two, and keep its coprime unordered pairs.
The eight pairs in (F.28) are exactly the result of that
finite enumeration.  This is only a necessary filter, not
a classification of relative positions or affine twists.
\end{proof}

\paragraph{Proposition F.2.9 (Proved: the only first-cohomology pattern for two points).}
For a simply connected two-multiple-fibre datum with $b_2=0$,
the rational
cohomology system $V$ has one invariant, no coinvariants,
and fixed rank two at each finite generator.  Its quotient
$W=V/\mathbb Qe$ by the invariant line has neither invariants
nor coinvariants.  If one finite characteristic polynomial
on $W$ is $(T-1)\Phi_d(T)$ with $\deg\Phi_d=4$, then $W$
is irreducible over $\mathbb Q$.

\begin{proof}
Write $a=\dim V^\Gamma$, $b=\dim(V^*)^\Gamma$ and $f_i$
for the finite fixed ranks.  Theorem F.3.1 gives $a=1$,
vanishing first parabolic cohomology, and
$f_1+f_2=4+b$.  By F.2.15 each $f_i$ is positive and even, so $b$ is
even.  The maximal square-zero cusp has
$\ker N=\operatorname{im}N$.  Therefore every globally
fixed vector belongs to $\operatorname{im}N$, whereas every
globally fixed covector kills $\operatorname{im}N$.
The global invariant and coinvariant spaces pair to zero.

Restriction of a globally invariant covector to the fixed
space of a finite generator is injective: averaging that
generator is the projection onto its fixed space, and the
covector is unchanged by averaging.  Its restriction kills
the $a$ global invariants in that space.  Hence
$a+b\leq f_i$ for each $i$.
If $b\geq2$ is even, the even ranks satisfy
$f_i\geq b+2$, so $4+b=f_1+f_2\geq2b+4$, a contradiction.
Thus $b=0$ and $f_1=f_2=2$.

The quotient $W$ has no invariant covectors, since its dual
embeds into $V^*$.  If it had a fixed line, its preimage in
$V$ would be an extension of the trivial line by the trivial
line.  Every finite generator on this two-dimensional
space would be unipotent of finite order, hence identity.
This would give two invariants in $V$, a contradiction.
Thus $W$ also has no invariant vectors.

If one finite characteristic polynomial is
$(T-1)\Phi_d(T)$ with the second factor irreducible of degree
four, every proper nonzero rational submodule has dimension
one or four.  In dimension one that generator acts as the
identity.  Since the product of the two finite generators
is unipotent, their induced one-dimensional characters are
inverse; hence the other generator also acts trivially.
This would be an invariant line.  In dimension four the
one-dimensional quotient is trivial by the same argument,
giving a coinvariant.  Both have just been excluded.
\end{proof}

\paragraph{Theorem F.2.10 (Proved: exclusion of three further order pairs).}
No simply connected datum with $k=2$, $b_2=0$, and finite linear orders
$(3,5)$, $(3,8)$ or $(3,10)$ has the required maximal
square-zero cusp.  Among such data, the
remaining necessary unordered linear order pairs are
\[
 \{(3,4),(4,5),(5,6),(5,8),(5,12)\}.
 \tag{F.30}
\]
When one order is five, the other finite characteristic
polynomial cannot be $(T-1)^2\Phi_4(T)^2$ or
$(T-1)^2\Phi_6(T)^2$.

\begin{proof}
Let the second order be $d\in\{5,8,10\}$.
Proposition F.2.9 gives an irreducible rational rank-five
quotient $W$, with finite characteristic polynomials
\[
 (T-1)(T^2+T+1)^2,
 \qquad (T-1)\Phi_d(T).
\]
Since the invariant line of $V$ lies in $\operatorname{im}N_V$,
the quotient cusp has Jordan type $(2,2,1)$.
All three local conjugacy classes are self-dual.  Their
centralizer dimensions are respectively $9,5,13$.
For the last value, the conjugate partition is $(3,2)$,
so the dimension is $3^2+2^2$; the finite values are the
sums of squares of their eigenvalue multiplicities.
Consequently the invariant-sheaf Euler formula applied to
$\operatorname{Hom}(W,W^*)$ gives
\[
 \chi\bigl(\mathbb P^1,j_*\operatorname{Hom}(W,W^*)\bigr)
      =-25+9+5+13=2.
 \tag{F.31}
\]
If $W$ were not self-dual, irreducibility would make both
$\operatorname{Hom}_\Gamma(W,W^*)$ and
$\operatorname{Hom}_\Gamma(W^*,W)$ zero.  These are the
degree-zero group and the dual of the degree-two group in
(F.31).  The Euler characteristic would then be nonpositive,
a contradiction.  Thus $W$ has an invariant nondegenerate
rational bilinear form.  Its symmetric part is nonzero,
since an alternating form cannot be nondegenerate in odd
dimension; irreducibility then makes that symmetric part
nondegenerate.  Both generators have determinant one.

Over $\mathbb C$ use
$\operatorname{Spin}_5\simeq\operatorname{Sp}_4$ and its
explicit degree-two representation on the primitive exterior
square.  Choose finite semisimple lifts $R,S$ of the two
orthogonal generators, choosing one of their signs so that
$RS$ is the unipotent lift of their unipotent product.
The orthogonal Jordan type $(2,2,1)$ lifts to a rank-one
transvection.  This can be verified directly, as follows.
For a rank-two square-zero skew-adjoint operator on a
five-dimensional quadratic space, its image is a maximal
isotropic plane, and its nonzero skew form on the dual plane
determines one orbit under the complex orthogonal group.
A rank-one symplectic transvection in dimension four maps,
under the primitive exterior square, to precisely that orbit:
its logarithm has two length-two blocks and one trivial block.
The exact exterior matrices checking these ranks and square
zero are included in \texttt{fourfolds\_order\_pairs.py}.

The first lift has, up to a common sign, eigenvalues
$1,1,\zeta_3,\zeta_3^{-1}$.  Indeed a symplectic torus with
eigenvalues $x,x^{-1},y,y^{-1}$ has primitive exterior
eigenvalues $1,xy,x/y,y/x,(xy)^{-1}$; substitution of the
two equal primitive-third pairs gives the asserted list.
In particular $R$ has a repeated eigenvalue.
The lift $S$ has order dividing $2d$, so it has no eigenvalue
of order three or six.  Nor can it have eigenvalue $1$ or
$-1$: a symplectic pair $x=x^{-1}=\pm1$ would make two of
the nontrivial primitive exterior eigenvalues repeat, whereas
$\Phi_d$ has four distinct roots.
Thus $R^{-1}$ and $S$ have coprime characteristic polynomials.

We finish with an elementary consequence of the rank-one
companion argument.  Two matrices with coprime characteristic
polynomials whose difference has rank one are both cyclic.
To prove this, let $u$ span the difference image.
A common invariant subspace either meets the image nontrivially,
in which case the two induced quotient matrices agree, or is
killed by the difference, in which case their restrictions
agree.  Coprimality forces that subspace to be respectively
the whole space or zero.  The span of the iterates of $u$
under either matrix is stable under both and is nonzero;
it is therefore the whole space, proving cyclicity.
Apply this to $R^{-1},S$, since
$R^{-1}-S=R^{-1}(I-RS)$ has rank one.
But a finite semisimple matrix with a repeated eigenvalue
is not cyclic.  This contradicts the displayed spectrum of
$R$, and excludes all three order pairs.
The same proof gives the final assertion: replace the two
primitive-third pairs on $W$ by two primitive-fourth or
primitive-sixth pairs.  The local centralizer dimension
is still nine, so the Hom Euler characteristic is again two.
The spin lift has, up to sign, eigenvalues
$1,1,\zeta_j,\zeta_j^{-1}$ for $j=4$ or $6$.
The other lift has order dividing ten and has no eigenvalue
$1$ or $-1$ because its primitive exterior spectrum is
simple.  Their spectra are disjoint, and cyclicity gives
the same contradiction.
\end{proof}

\paragraph{Proposition F.2.11 (Proved: rational uniqueness of the extension).}
Let $W$ be a rational rank-five representation of a
twice-punctured affine line, with two finite generators
having fixed rank one and product with unipotent Jordan type
$(2,2,1)$.  Suppose $W$ has neither invariants nor
coinvariants.  Up to rescaling the invariant line, there is
exactly one rational extension
$0\to\mathbb Qe\to V\to W\to0$ in which the two finite
generators retain their finite orders and the cusp has
square-zero logarithm of rank three.  Thus, for the
two-point candidates of F.2.9, the remaining rational
representation problem is exactly the problem of the
rank-five quotients.  This statement does not identify
integral lattices in those representations.

\begin{proof}
Write the two quotient generators as $T_1,T_2$ and the
extension generators in block form
$\left(\begin{smallmatrix}1&c_i\\0&T_i\end{smallmatrix}\right)$.
If $T_i$ has order $m_i$, finite order of the extension is
equivalent to $c_i(I+T_i+\cdots+T_i^{m_i-1})=0$.
Over $\mathbb Q$ this says $c_i\in\operatorname{row}(T_i-I)$,
a space of dimension four.  The two rows give eight
parameters.  Changes of splitting by a row $z\in W^*$
give $c_i\mapsto c_i+z(T_i-I)$; their space has dimension
five because $W^*$ has no invariants.  Hence extensions
split at the two finite cyclic groups form a vector space
of dimension three.

Restriction to the cusp takes this space to
$W^*/\operatorname{row}(T_0-I)$, also of dimension three.
Its kernel is the first parabolic cohomology of $W^*$.
That kernel vanishes: the invariant-sheaf Euler characteristic
is $-5+1+1+3=0$, and its degree-zero and degree-two groups
vanish by the assumed absence of invariants and coinvariants.
Thus cusp restriction is an isomorphism over $\mathbb Q$.

Put $N=T_0-I$ on $W$, so $N^2=0$ and $\operatorname{rank}N=2$.
The cusp extension has nilpotent part
$\left(\begin{smallmatrix}0&c\\0&N\end{smallmatrix}\right)$.
Its square is zero exactly when $cN=0$.
The corresponding subspace of cusp extension classes is
\[
 \frac{\operatorname{ann}(\operatorname{im}N)}
      {\operatorname{row}N},
 \tag{F.32}
\]
which has dimension $3-2=1$.
A class in (F.32) is zero exactly when $c$ vanishes on
$\ker N$; those are precisely the extensions whose cusp
nilpotent part still has rank two.
Every nonzero class raises that rank to three and has square
zero.  Rescaling $e$ acts transitively on the nonzero rational
classes in this one-dimensional space.  This proves both
existence and uniqueness in the stated rational category.
\end{proof}


\paragraph{Theorem F.2.12 (Proved: an integral obstruction for a rational $(5,12)$ family).}
Let a rational rank-six cohomology representation have generators
\[
 A=\begin{pmatrix}1&0\\0&A_5\end{pmatrix},\qquad
 B=\begin{pmatrix}1&c\\0&B_5\end{pmatrix},\qquad
 c=(0,0,0,1,0),
\]
where
\[
 A_5=\begin{pmatrix}
 1&0&1&0&0\\1&0&0&0&0\\-1&0&0&1&1\\
 1&-1&0&-1&0\\0&1&0&0&0
 \end{pmatrix},\qquad
 B_5=\begin{pmatrix}
 0&1&0&0&0\\0&0&0&0&1\\1&-2&0&0&0\\
 0&1&0&-1&-1\\0&0&1&1&2
 \end{pmatrix}.
 \tag{F.33}
\]
For every invariant integral lattice with an index-one maximal
square-zero cusp, and every free affine filling of multiplicities
$m_1,m_2$ divisible respectively by five and twelve, a simply
connected total space has
\[
 \operatorname{Tor}_3 H_2
       \simeq\mathbb Z/3^{v_3(m_2)}.
 \tag{F.34}
\]
In particular none of these data is three-connected.
The standard lattice nevertheless has a primitive cusp,
free affine fillings of orders five and twelve, and a simply
connected total space with rational Betti vector
$(1,0,0,0,4,0,0,0,1)$.

\begin{proof}
All matrix identities and integral lattices used below are
printed and checked exactly in
\texttt{fourfolds\_512\_probe.py}.
The generators have orders five and twelve, characteristic
polynomials $(T-1)^2\Phi_5$ and $(T-1)^2\Phi_{12}$,
and $(AB-I)^2=0$ with three nonzero Smith factors equal to one.
Their common invariant line is the first coordinate line.
For homology matrices $A^t,B^t$ put
\[
 u_1=(1,0,0,0,0,0)^t,\qquad
 u_2=(1,-1,1,-1,0,-1)^t,
 \qquad \ell=(1,0,0,0,0,0).
\]
They are invariant vectors with $\ell(u_i)=1$.
The cusp vanishing lattice and its two finite translates
generate the saturated lattice $\ker\ell$.
Choose twists $v_1=-2u_1$, $v_2=5u_2$.
Their $\ell$-values are units modulo five and twelve, proving
every nonidentity affine power free.
After killing the vanishing lattice the group relations are
$5a=-2h$, $12b=5h$, $a+b=0$; their determinant has absolute
value $|-24+25|=1$.  Thus the total space is simply connected.
The exact rational exterior calculations give
\[
 h^0=(1,1,0,0,0,0,1),\qquad
 h^1=(0,0,0,3,0,0,0),\qquad
 h^2=(1,0,0,0,0,1,1)
\]
for the invariant sheaves, and the same rational Leray
argument as in F.2.5 gives the asserted Betti vector.

We first prove that changing the lattice cannot remove the
obstruction at the prime three.  Write $R=\mathbb Z_3$ and
$K=\mathbb Q_3$.  The standard quotient lattice $W_R$ defined
by $A_5,B_5$ has irreducible reduction modulo three.
Indeed $\Phi_5$ is irreducible over $\mathbb F_3$, since
three has order four modulo five.  Every proper nonzero
$A_5$-stable subspace therefore has dimension one or four.
On a one-dimensional subrepresentation or quotient its
$A_5$ character is one; the unipotent product forces the
$B_5$ character to be one too.  The stacked and horizontal
matrices for $A_5-I,B_5-I$ both have rank five modulo three,
as certified in the code, ruling out those two possibilities.

This irreducible reduction makes $W_R$ the unique stable
$R$-lattice up to homothety.  To see this, rescale another
lattice into $W_R$ but not into $3W_R$.  Its reduction is
a nonzero invariant subspace, hence all of $W_R/3W_R$.
Nakayama's lemma, or the elementary Smith form of the
inclusion, then makes it equal to $W_R$.

Normalize the quotient of any stable lattice in $V_K$ to
be $W_R$.  Its intersection with the invariant line is
$3^rRe$ for some integer $r$.
The representation of the free group on the two generators
has extension class
$\alpha\in H^1(\Gamma,W_R^*)$.
This cohomology group is torsion free: the coboundary map
from $W_R^*$ into two copies has full column rank modulo
three, since $W_R^*$ has no invariants modulo three.
The class $\alpha$ is primitive in this free $R$-module.
Otherwise the extension would split modulo three and $V$
would have a trivial quotient there.  The horizontal
matrix $[A-I\ B-I]$ instead has Smith factors
$(1,1,1,1,1,2)$, so it has full rank modulo three.

In a basis starting with $3^re$, the extension class is
$3^{-r}\alpha$.  Integrality is therefore possible exactly
when $r\leq0$.  Changes of splitting introduce no further
classes: the coboundary image is primitive, so an integral
coboundary over $K$ is already an integral coboundary over
$R$.  Thus these lattices are represented by multiplying
the extension rows by $3^{-r}$.
On the quotient cusp $N_W$ has primitive rank-two image.
The extra cusp map is the nonzero map
$\ker N_W/\operatorname{im}N_W\to Re$.
It is a unit for the displayed lattice, because its full
cusp has Smith factors $(1,1,1)$; for the lattice indexed
by $r$ it is multiplied by $3^{-r}$.
Hence the full cusp is primitive exactly when $r=0$.
We have proved uniqueness of the primitive-cusp lattice
over $R$, without relying on a bounded lattice search.

It remains to compute the obstruction on that lattice.
We first justify the arbitrary-twist comparison integrally.
For any global invariant lattice take the primitive invariant
homology covector $\ell$ and quotient the van Kampen group
by every fibre loop in $\ker\ell$.
The finite presentations become
$[a,h]=[b,h]=1$, $a^{m_1}=h^{r_1}$,
$b^{m_2}=h^{r_2}$, where $r_i=\ell(v_i)$.
The cusp meridian gives $ab=h^{-\kappa}$ for an integer
$\kappa$ determined by its translation clutching
($\kappa=0$ for the stated standard clutching).
The cusp angular lattice lies in $\ker\ell$, so it adds
no further relation to this quotient.  In particular this
is an abelian quotient with determinant
$m_1m_2\kappa+m_2r_1+m_1r_2$.
Simple connectivity makes that determinant $\pm1$.
Reduction modulo each $m_i$ proves both
$\gcd(m_1,m_2)=1$ and $\gcd(r_i,m_i)=1$.

The two finite fillings have free $R$-cohomology in the
degrees needed here.  In the invariant marking whose first
homology vector over $R$ has $\ell$-value one, the map
$m_i\ell$ from the reduced finite fibre to $\mathbb R/\mathbb Z$
has connected fibre $T^5$: the affine generator permutes
its $m_i$ preimage components by addition of $r_i$ modulo
$m_i$, transitively by the proved coprimality.
Its monodromy on vertical homology is $D_i^{s_i}$, where
$s_ir_i\equiv1\pmod {m_i}$, possibly followed by a translation.
Thus it is a mapping torus of a power prime to the linear
order.  The translation is isotopic to the identity and
does not change its Wang maps.  At the second point, over
$R$ we may take the invariant first vector to be $v_2/r_2$;
its denominator is prime to three and does not change the
$R$-homology lattice.  At the first point the covering
degree $m_1$ is prime to three, so averaging and the transfer
identify its $R$-cohomology with the invariant fibre
cohomology directly.
In degrees one and two the matrices
$D_i-I$ and $\bigwedge^2D_i-I$ have only unit Smith factors.
This proves the stated freeness by the Wang sequence.

The finite degree-two covering cokernel is read from
integral periods.  Vertical cycles have integral periods.
The remaining primitive two-cycle is the commuting torus
formed by a base suspension and a primitive fixed vertical
loop $\lambda_i$; its covering period is
$\omega(u_i,\lambda_i)/m_i$ for an invariant integral
two-form $\omega$.
More explicitly the homological Wang sequence is
\[
 0\longrightarrow\operatorname{coker}(D_i^{s_i}-I\mid H_2(T^5;R))
 \longrightarrow H_2(B_i;R)
 \longrightarrow\ker(D_i^{s_i}-I\mid H_1(T^5;R))
 \longrightarrow0.
\]
The last group has rank one, generated by $\lambda_i$.
The commuting suspension torus maps to that generator;
at the second point the invariant vector $v_2/r_2$ may be used after taking
$R$-linear combinations of integral cycles.
Therefore vertical cycles and this one torus generate
$H_2(B_i;R)$, proving sufficiency of the period test.
For integral bases of the invariant two-forms, the numerator
rows at the two points are
\[
 (0,0,5),\qquad (3,1,1).
 \tag{F.35}
\]
The vertical Wang quotient has no torsion, so those cycles
generate the full free second homology and the period test
is sufficient as well as necessary.

Since $3\mid m_2$, the already-proved coprimality gives
$3\nmid m_1$ and makes $r_2$ a unit in $R$.
Changing that unit twist or adding a fixed vertical vector
does not change the three-primary cokernel in (F.35): the
former multiplies its generating period by a unit, and the
latter has zero wedge with the fixed vertical loop.
The first finite point has zero covering cokernel over $R$;
the second has $R/m_2R$.

The global invariant degree-two group is zero, and the
degree-two global coboundary image is primitive, with
vanishing rational parabolic first cohomology.  Therefore
$H^1(\mathbb P^1,R^2g_*R)=R/m_2R$.
The global degree-one coinvariants vanish over $R$, and
there are no degree-two or degree-three point-supported
kernels: the finite cap groups just proved free, and the
Kuhn cusp has its sole kernel in degree four.
The degree-three global invariant group also vanishes.
Consequently the only nonzero total-degree-three Leray
term over $R$ is $E_2^{1,2}=R/m_2R$; it has no possible
differential.  Thus $H^3(g^{-1}(\mathbb P^1);R)=R/m_2R$.
The universal coefficient theorem gives precisely (F.34).
\end{proof}

\paragraph{Theorem F.2.13 (Proved: two further rational families excluded integrally).}
Let $C_f,C_g$ denote the degree-four companion matrices.
Take $f=\Phi_5$ and either
\[
 \text{case A: }g=\Phi_{12},\qquad
 \text{case B: }g=\Phi_3\Phi_6.
 \tag{F.36}
\]
Let $W$ be the quotient of
$\bigwedge^2(C_f^{-1},C_g)$ by its invariant line, and
let $V$ be its unique rational extension of F.2.11.
For every invariant integral lattice in $V$ with primitive
maximal square-zero cusp, every simply connected free affine
filling, with multiplicities $m_1,m_2$ divisible by five and
six respectively, fails three-connectivity.
More precisely, put $r=v_3(m_2)\geq1$.
In case A, $\operatorname{Tor}_3 H_3$ has order
$3^{2r-1}$.  In case B,
$\operatorname{Tor}_3 H_2\simeq\mathbb Z/3^r$.
For $m_2=6$ the respective obstructing group is $\mathbb Z/3$.

\begin{proof}
The complete integer representatives are supplied by the
companion and Smith-basis construction in
\texttt{fourfolds\_56\_probe.py} and
\texttt{fourfolds\_56b\_probe.py}.
In each case the source cusp is a rank-one transvection;
the primitive exterior quotient has cusp type $(2,2,1)$,
finite fixed ranks one, and no invariants or coinvariants.
The second finite characteristic polynomial on $V$ is
$(T-1)^2(T+1)^2\Phi_3$ in case A and
$(T-1)^2(T+1)^2\Phi_6$ in case B.
The exact extension equations give an integral primitive
rank-three cusp in both cases.

For completeness these are not merely rational candidates
with impossible twist lattices.  In the displayed integral
bases of the probe, change lattice by the rational basis
matrix $Q$ whose diagonal entries are two and whose only
off-diagonal entries are $Q_{1,4}=Q_{1,6}=1$.
The resulting generators $Q^{-1}AQ,Q^{-1}BQ$ remain
integral and have primitive cusps.
The invariant homology vectors of $\ell$-value one can be
chosen as
\[
\begin{array}{c|cc}
 &u_1&u_2\\\hline
 \mathrm A&(1,3,-3,0,1,0)&(1,-1,1,2,-1,1)\\
 \mathrm B&(1,0,0,1,0,1)&(1,0,0,1,0,0).
\end{array}
\]
Use twists $u_1,-u_2$ and multiplicities five and six.
All nonidentity affine powers are free, and the cusp normal
closure is the saturated $\ker\ell$.
The residual determinant $6-5=1$ proves simple connectivity.
The rational Leray ranks are the same as in F.2.5.

The complete three-adic lattice argument in F.2.12 applies
to each case, not only to these representatives.
The reason is checked by the same exact integer matrices:
the rank-five quotient has no invariants or coinvariants
modulo three, its order-five polynomial is
$(T-1)\Phi_5$ with irreducible $\Phi_5$ modulo three, the
global extension class is primitive, and the rank-six
cusp has unit Smith factors.
Thus each representation has exactly one primitive-cusp
$\mathbb Z_3$-lattice up to homothety.
Both the quotient argument forcing $r_i$ prime to $m_i$
and the finite mapping-torus comparison of F.2.12 apply.

Set $R=\mathbb Z_3$.  All finite cap cohomology groups
through degree four are free over $R$: the nonunit Smith
factors of their vertical exterior differences, when
present, are two.  The global coboundary and coinvariant
Smith factors in degrees one through four likewise have
no factor three.  The rational invariant and parabolic
ranks are those printed in F.2.12.  In particular,
the global invariant groups in degrees two, three and
four vanish, their relevant coinvariants vanish over $R$,
and the degree-two parabolic first cohomology is zero.

For the transformed integral lattices just specified, the
finite degree-two suspension periods at order six are
\[
 \mathrm A:\ (-2,-1,-1),\qquad
 \mathrm B:\ (1/3,0,1/3).
 \tag{F.37}
\]
The order-five periods are integral and its covering degree
is a unit in $R$ for any simply connected multiplicities.
In case B the numerator row at six is $(2,0,2)$; replacing
six by $m_2$ gives the cyclic cokernel $R/3^rR$.
The total-degree-three Leray argument of F.2.12 gives
$H^3(P;R)=R/3^rR$, proving the asserted $H_2$ obstruction.

In case A the obstruction instead occurs in degree four.
The invariant rank-four middle cohomology lattice of the
sixfold torus cover has skew-pairing determinant $6^2$.
This determinant is computed exactly, with the integral
invariant basis, in \texttt{fourfolds\_56a\_newlattice.py}.
The free middle pairing downstairs is unimodular over $R$.
If $I$ is the index of the pullback lattice in the invariant
lattice, pullback degree $m_2$ gives
\[
 I^2\,6^2=m_2^4
 \quad\text{up to a unit of }R.
 \tag{F.38}
\]
Thus the degree-three finite covering cokernel has
three-primary order $3^{2r-1}$.
There are no global invariant degree-three forms to
quotient this cokernel out.  The integral parabolic first
cohomology in degree three is free over $R$, so the torsion
of $H^1(B,R^3g_*R)$ is precisely that finite cokernel.
There are no differentials to or from this middle base
column.  In total degree four the other surviving term
is the free rank-one point-supported cusp class;
the degree-two base term is zero over $R$.
The resulting extension by a free module splits.
Therefore $\operatorname{Tor}_3H^4(P;R)$ has order
$3^{2r-1}$, and the universal coefficient theorem identifies
it with $\operatorname{Tor}_3H_3(P;\mathbb Z)$.
This proves the claimed obstruction in case A.
\end{proof}

\paragraph{Theorem F.2.14 (Proved: four order pairs and eight local spectral pairs).}
Every simply connected two-finite-point datum with $b_2=0$ has
one finite linear order equal to five and the other equal
to four, six, eight or twelve.
After suppressing the common factor $(T-1)^2$, its unordered
pair of finite characteristic polynomials is exactly one
of the following necessary possibilities:
\[
\begin{array}{c|c|l}
 \text{orders}&\text{order-five factor}&\text{other factor}\\\hline
 (5,4)&\Phi_5&\Phi_2^2\Phi_4\\
 (5,6)&\Phi_5&\Phi_2^2\Phi_3,\ \Phi_2^2\Phi_6,\ \Phi_3\Phi_6\\
 (5,8)&\Phi_5&\Phi_8\\
 (5,12)&\Phi_5&\Phi_{12},\ \Phi_3\Phi_4,\ \Phi_4\Phi_6.
\end{array}
 \tag{F.39}
\]
Its rank-five quotient is rationally irreducible.
In particular all order-three pairs, including $(3,4)$,
are impossible.  The table is a necessary spectral filter;
it does not identify the relative rational or integral
positions of the matrices having those polynomials.

\begin{proof}
By F.2.15 the finite actions are locally complex-compatible.
By F.2.9, the rank-six cohomology system $V$ has one
invariant, no coinvariants, and finite fixed ranks two.
Its quotient $W$ has neither invariants nor coinvariants.
Every rational one-dimensional subrepresentation or
quotient character of $W$ is trivial: its two finite
orders divide the coprime affine multiplicities of F.2.7,
so at least one is odd and that character value is one;
the unipotent product forces the other value to be one.
Thus $W$ and $W^*$ have no invariant rational line of any
character.

The space of invariant alternating forms on $W$ has
dimension at most one.  Indeed a nonzero such form cannot
have rank four, because its radical would be a stable
rational line in $W$.  It therefore has rank two and is
a decomposable element of $\bigwedge^2W^*$.
If two independent invariant rank-two forms existed,
their sum would also have rank at most two; hence their
wedge product would be zero.  Their two supporting
planes in $W^*$ would be distinct stable planes meeting
in a line.  That stable rational line contradicts the
preceding paragraph.  This proves the bound without an
irreducibility assumption.

The dual exterior sequence
\[
 0\longrightarrow\bigwedge^2W^*
 \longrightarrow\bigwedge^2V^*
 \longrightarrow W^*\longrightarrow0
\]
and $(W^*)^\Gamma=0$ show that the coinvariant dimension
$b_2'$ of $\bigwedge^2V$ is at most one.
On the other hand its invariant dimension $a_2'$ is zero.
To see this, the corresponding $H^0(R^2g_*\mathbb Q)$
term has no point-supported rational kernel at a finite
fibre or in cusp degree two, and its only possible outgoing
differential has target $H^2(R^1g_*\mathbb Q)=0$.
It therefore survives in total degree two, which vanishes.

Let $d_i$ be the dimension fixed by the $i$th finite
generator on $\bigwedge^2V$.  The cusp fixed dimension is
nine, so the invariant-sheaf Euler identity gives
\[
 -15+d_1+d_2+9
       =a_2'-h^1(j_*\bigwedge^2V)+b_2'\leq1.
 \tag{F.40}
\]
Each finite generator has fixed rank two on $V$.
Its complementary four-dimensional real representation
commutes with a complex structure and has no fixed vectors.
It either has two distinct nonreal inverse pairs, one
nonreal pair with multiplicity two, two negative real
eigenvalues and a nonreal pair, or four negative eigenvalues.
The respective invariant dimensions in degree two are
$3,5,3,7$.  In particular $d_i$ is odd and at least three.
The left side of (F.40) is therefore nonnegative and even,
so it must be zero.  Thus $d_1=d_2=3$.

The complete cyclotomic filter F.2.2 now leaves precisely
the finite polynomials with two fixed eigenvalues and
the indicated degree-two fixed rank three.
Keeping coprime unordered pairs gives the eight rows
expanded in (F.39).  This finite enumeration is reproduced
from the complete saved 38-type list by
\texttt{fourfolds\_order\_pairs.py}.
An order-three generator would instead have characteristic
polynomial $(T-1)^2\Phi_3^2$, hence $d_i=5$, explaining
its exclusion directly.
Every remaining pair contains $(T-1)^2\Phi_5$, so F.2.9
proves rational irreducibility of $W$.
\end{proof}

Theorem F.2.10 supplies an independent spin-representation
proof for three of the excluded order-three pairs and for
the repeated fourth- and sixth-root cases.
Within the third row of the order-six spectral possibilities,
$\Phi_3\Phi_6$, no exclusion follows from F.2.13: that theorem
treats particular rational representations with the first
two order-six spectra.  Likewise F.2.12 excludes one
rational representation of order pair $(5,12)$, not its
entire spectral locus.


\paragraph{Theorem F.2.16 (Proved: an unbounded integral family).}
There are infinitely many simultaneous rational conjugacy classes
of integral rank-six tuples of the linear type $(5,4)$ in F.2,
with primitive maximal square-zero cusp.  An explicit infinite
subfamily has free affine lifts of orders five and four.
Apart from the parameter zero, the rank-five irreducible quotient
in the family is not self-dual.  Thus no finite bound on rational
conjugacy classes is valid for the unrestricted input tuples.
This does not assert infinitely many three-connected survivors.

\begin{proof}
All the following expressions are explicit.  Let $v\in\mathbb Z$,
$q=v^2-v+1$, and let $A$ be the permutation matrix
$Ae_j=e_{j+1}$ on $\mathbb Z^5$, with subscripts taken modulo five.
Put
\[
 X=(e_0,\ e_2+ve_4),\qquad
 Y=\begin{pmatrix}
 0&-q&-v^2&-1&v\\
 0&1-vq&-vq&v^2-v&q
 \end{pmatrix},\qquad N=XY,
 \tag{F.U1}
\]
and set $C=(I+N)A$, $B=C^{-1}=A^{-1}(I-N)$.
Direct multiplication gives
\[
 YX=0,\quad Y(A+I)^{-1}X=I_2,\quad
 A^5=C^4=I_5,\quad
 \chi_C(x)=(x-1)(x+1)^2(x^2+1).
 \tag{F.U2}
\]
The image of $X$ is primitive, and $Y$ is onto $\mathbb Z^2$.
For the latter assertion, $2(A+I)^{-1}X$ is integral, so the
cokernel of $Y$ is killed by two; reduction of (F.U1) in the
two cases $v\equiv0,1\pmod2$ has rank two.  Thus $N$ has
rank two, $N^2=0$, and its two nonzero Smith factors are one.

The common invariant space of $A,B$ is zero: the only
$A$-fixed line is generated by $\mathbf1=(1,1,1,1,1)^t$,
and the first entry of $Y\mathbf1$ is $-2q\ne0$.
There are no common invariant covectors either, since
$\mathbf1^tX=(1,v+1)$ and $Y$ has rank two.
The characteristic polynomial of $A$ is $(x-1)\Phi_5(x)$.
A proper rational invariant subspace must therefore have
dimension one or four.  A one-dimensional subrepresentation
or quotient has $A$ acting as one; the unipotent product then
forces $B$ also to act as one.  The two vanishing assertions
exclude both possibilities.  The rank-five pair is irreducible.

Here is an explicit primitive rank-six extension, including
the integral lattice.  Let $P$ have columns
$\mathbf1/q,e_1,e_2,e_3,e_4$, and put
$A'=P^{-1}AP$, $B'=P^{-1}BP$.  These matrices are integral:
$A\mathbf1=\mathbf1$ and
\[
 Y\mathbf1=2q\binom{-1}{1-v}
\]
show that the index-$q$ overlattice is stable.  Define
\[
 s=v^2+v-1,\qquad r=-\frac{(v+1)(v^2-2)}2,
 \qquad
 \beta=(s,\ s-vr,\ v^2s+vr,\ r,\ -vs-r).
 \tag{F.U3}
\]
The number $r$ is integral for every integer $v$.
Set
\[
 \widehat A=\begin{pmatrix}1&0\\0&A'\end{pmatrix},
 \qquad
 \widehat B=\begin{pmatrix}1&\beta\\0&B'\end{pmatrix}.
 \tag{F.U4}
\]
They have orders five and four, and
$(\widehat A\widehat B-I)^2=0$.
To check primitivity without a parameter bound, write
$X'=P^{-1}X$, $Y'=YP$.  Two rows of $X'$ have determinant
$-1$.  The quotient of $\mathbb Z^5$ by its image has
coordinates indexed by $0,3,4$, and the induced matrix $Y'$
there is
\[
 \begin{pmatrix}-2&-1&v\\2(1-v)&v^2-v&q\end{pmatrix}.
\]
Its kernel has primitive generator
$(-v^3+v-1,\,2,\,2(1-v^2))^t$: its first coordinate is
odd and the second is two.  Lift this to
\[
 z=(-v^3+v-1,0,0,2,2(1-v^2))^t.
\]
The same three entries are its signed two-by-two minors,
whose gcd is one; hence $Y'$ is onto $\mathbb Z^2$.
Then $Y'z=0$ and
\[
 \beta z=s(v^3-v-1)+2v^2r=1.
 \tag{F.U5}
\]
Thus the rank-two primitive cusp in the quotient acquires a
primitive third image direction in the extension.  All three
nonzero Smith factors of $\widehat A\widehat B-I$ are one.

For later homology comparisons use the following constant
change of lattice:
\[
 L=\frac14\begin{pmatrix}
 -2&3&-5&-1&5&-3\\
 0&2&-2&2&-2&2\\
 0&-4&0&0&4&-4\\
 0&0&0&-4&4&-4\\
 0&0&4&-4&4&-4\\
 0&0&0&-4&4&0
 \end{pmatrix},\quad
 A_v=L^{-1}\widehat A L,\quad B_v=L^{-1}\widehat B L.
 \tag{F.U6}
\]
Here $\det L=-1/4$, $L^{-1}$ is integral, and all entries of
$A_v,B_v$ are polynomials in $v$ with integer coefficients.
The primitive-cusp assertion is unchanged at odd primes;
at two, direct reduction for $v=0,1$ gives rank three.
At $v=0$, (F.U6) is exactly the pair in Theorem F.2.5.

For the infinite subfamily $v\in20\mathbb Z$, take homology
monodromies $M_1=A_v^t$, $M_2=B_v^t$ and
$M_0=(M_1M_2)^{-1}$.  The reversed product has the same
primitive square-zero cusp type.  Write
$a=(1,3,-2,-1,2,0)^t$ and $e=(1,0,0,0,0,0)^t$, and set
\[
 w_1=\frac15\sum_{j=0}^4(A_v^t)^j a,
 \qquad
 w_2=-\frac14\sum_{j=0}^3(B_v^t)^j e.
 \tag{F.U7}
\]
These vectors are integral.  Indeed, the polynomial matrices
reduce to their parameter-zero values modulo five and four,
respectively, and $a,e$ are fixed by the corresponding
parameter-zero matrices.  They are fixed by $M_1,M_2$, with
first coordinates $1,-1$.  Both monodromies preserve the
first-coordinate covector $\ell$.  Thus the affine actions
$x\mapsto M_i x+w_i/m_i$, with $(m_1,m_2)=(5,4)$, are free:
every nontrivial power translates the $\ell$-circle nontrivially.

Finally, let $H=NA|_{\operatorname{im}N}$ and
$H'=NA^{-1}|_{\operatorname{im}N}$.  The construction gives
\[
 \det H=1-v^2q,\quad \det H'=1+v^2q,
 \qquad\operatorname{tr}(A^2C^2)=v^5-v^4+3.
 \tag{F.U8}
\]
These are simultaneous-conjugacy invariants; the last is
unbounded on the positive multiples of twenty.  If the
irreducible rank-five representation were self-dual, its
invariant bilinear form would have a nonzero symmetric part,
which is nondegenerate by irreducibility; an alternating
nondegenerate form is impossible in odd dimension.  Relative
to that symmetric form, $A,C$ are orthogonal and
$N^*=-N$, so $I+tN$ is orthogonal for every $t$.
The odd-dimensional characteristic polynomial of $(I+tN)A$
is then anti-reciprocal.  Its quadratic coefficient in $t$ is
\[
 x^2\bigl((\det H)x-\det H'\bigr),
\]
by the determinant lemma, so anti-reciprocity forces
$\det H=\det H'$.  Formula (F.U8) contradicts this whenever
$v\ne0$.

All the matrix identities, coefficient integrality, both
mod-two checks, and the affine formula are checked symbolically
by \texttt{code/fourfolds\_unbounded\_family.py}.  Its separate
Smith diagnostics over $-6\le v\le10$ are explicitly bounded
and are not the basis for the universal assertions above.
\end{proof}

\paragraph{Proposition F.2.17 (Proved: an odd torsion obstruction in this lattice).}
For the integral tuple $(A_v,B_v)$ in (F.U6), any simply
connected filling with free finite affine caps and the fixed
primitive Kuhn cusp has a subgroup
\[
 \mathbb Z/(v^2-v+1)\ \subset\operatorname{Tor}H_2(M;\mathbb Z).
 \tag{F.U9}
\]
In particular, in this specified family of lattices only
$v=0$ and $v=1$ can be three-connected, regardless of the
affine multiplicities and twists.  The stronger exclusion for every stable lattice in each
rational class is proved in Theorem F.2.19 below.

\begin{proof}
First compute the global coinvariants in the hatted lattice.
In the rank-five quotient, the $A'$ relations identify
$f_1=f_2=f_3=f_4=t$ and $qf_0=5t$.
The number $q=v^2-v+1$ is relatively prime to five, as its
five residue values show.  Thus these coinvariants are
infinite cyclic, with $f_0=5h$ and $t=qh$.
The $B'$ relations become $q\kappa_j h$, where
\[
 (\kappa_0,\ldots,\kappa_4)=
 (2v^2,\ v(v^3+v-1),\ v(v^3+v+1),
       -v^3+v+1,\ -v^3-v-1).
 \tag{F.U10}
\]
Their gcd is one: it divides $2v$ by subtracting the second
and third entries, and it divides $v^3-v-1$ by the fourth;
the latter is odd and is $-1$ modulo $v$.
Hence the rank-five global coinvariants are $\mathbb Z/q$.
The extension line is killed in rank-six coinvariants by
(F.U5), since
$(\widehat A\widehat B-I)(0,z)^t=e$.
Consequently the hatted rank-six coinvariants are also
$\mathbb Z/q$.  The lattice change (F.U6) is an isomorphism
over every $\mathbb Z_{(p)}$ with $p$ odd, so the odd part of
the actual rank-six global coinvariants is precisely
$\mathbb Z/q$.

Let $V$ be the degree-one fibre cohomology local system.
The only difference between $R^1g_*\mathbb Z$ and $j_*V$
is a finite cokernel at finite caps; the cusp degree-one
specialization is onto.  Thus
\[
 H^2(\mathbb P^1,R^1g_*\mathbb Z)
 =H^2(\mathbb P^1,j_*V)=V_\Gamma,
\]
including its $q$-primary part just computed.  The space of
global degree-two invariants is zero: the irreducible odd
rank-five quotient has no invariant alternating tensor,
and the extension filtration
$0\to W\to\bigwedge^2V\to\bigwedge^2W\to0$
gives $(\bigwedge^2V)^\Gamma=0$.

It remains to check every possible source of the differential
$d_2:H^0(\mathbb P^1,R^2g_*\mathbb Z)
\to H^2(\mathbb P^1,R^1g_*\mathbb Z)$.
The cusp has no kernel in degree two.  At a finite cap the
kernel is torsion in $H^2$ of that cap, namely torsion in its
first homology.  Simple connectivity of the total filling
forces the distinguished affine numerator to be prime to the
actual affine multiplicity, by the determinant argument of
F.2.12.  The cap is therefore a mapping torus of a torus
automorphism whose order divides the linear order, five or
four.  Its first-homology torsion is supported only at the
primes five or two, respectively: averaging over the finite
order kills the torsion of the vertical coinvariants.
Hence the source of this $d_2$ has no torsion at any prime
dividing $q$, since $q$ is odd and prime to five.

The subgroup $\mathbb Z/q$ in bidegree $(2,1)$ survives.
There are no longer differentials on a two-dimensional base,
and this terminal filtration term is a subgroup of
$H^3(M;\mathbb Z)$.  Universal coefficients identify its
torsion with a subgroup of $\operatorname{Tor}H_2(M;\mathbb Z)$.
This proves (F.U9).  Finally, $v^2-v+1=1$ for an integer $v$
exactly when $v=0$ or $v=1$.
\end{proof}

\paragraph{Theorem F.2.18 (Proved: a survivor with a non-self-dual quotient).}
The parameter $v=1$ in (F.U6), with affine orders five and four
and twists
\[
 w_1=(1,3,-2,-1,2,0)^t,\qquad
 w_2=(-1,0,0,0,0,0)^t,
\]
gives a three-connected smooth eight-manifold with
$H_4=\mathbb Z^4$ and all other reduced homology zero except
$H_8=\mathbb Z$.  It is almost complex and is oriented
diffeomorphic to $3\mathbb{HP}^2\#(-\mathbb{HP}^2)$.
Its irreducible rank-five quotient has no invariant bilinear
form, so integral three-connectivity does not force that
quotient to be orthogonal.  It has bounded-cusp holomorphic
periods by Theorem F.3.11; simultaneous finite holomorphic
extensions are not asserted here.

\begin{proof}
The first cohomological matrix is the same $A$ as in
Theorem F.2.5.  The second is
\[
 B_1=\begin{pmatrix}
 1&0&0&0&0&0\\
 0&0&-1&0&1&-1\\
 0&0&-1&0&0&0\\
 0&0&-2&0&1&-2\\
 0&1&-2&0&1&-3\\
 0&0&-1&1&0&-1
 \end{pmatrix}.
 \tag{F.U11}
\]
Take $M_1=A^t$, $M_2=B_1^t$ and
$M_0=(M_1M_2)^{-1}$.  The displayed twists are fixed integral
vectors with $\ell(w_1)=1$, $\ell(w_2)=-1$.
They therefore give free affine actions.  The cusp is
primitive by F.2.16.  Exact Hermite reduction gives
\[
 \operatorname{col}_{\mathbb Z}
 [\,M_0-I\mid M_1(M_0-I)\mid M_2(M_0-I)\,]
       =\langle e_1,\ldots,e_5\rangle.
\]
Van Kampen leaves only $5a=h$, $4b=-h$, $a+b=0$,
whose relation matrix is unimodular.  This proves
$\pi_1=1$ for the specified gluing, including zero cusp
clutching.

For completeness we give the full new integral certificates
needed to reuse the already proved differential calculation
in F.2.6.  The rational exterior invariant and parabolic ranks
are exactly the table there.  The finite cap first homology
groups are $\mathbb Z^2$ and $\mathbb Z^2\oplus\mathbb Z/2$.
On integral invariant two-form bases their suspension periods
divided by the covering orders are, respectively,
$(-1,0,-1)$ and $(0,0,1)$.  The determinants of the invariant
degree-three skew pairings are $5^4$ and $4^4$, while the
degree-two/degree-four pairings have absolute determinants
$5^3$ and $4^3$.  Thus all the finite covering indices used
in degrees two, three and four are one.

For exterior coefficient degrees two, three and four, the
global coboundary Smith form has one factor two and all other
nonzero factors one.  Its torsion class has peripheral
triviality flags $(\mathrm{true},\mathrm{false},\mathrm{true})$,
so the parabolic groups are torsion free, of ranks $0,3,0$.
The degree-one coboundary is primitive and the degree-one
finite covering cokernels are killed by the common invariant
covector and the Chinese remainder theorem, exactly as in
F.2.6.  The global coinvariants in degrees one through four
are each $\mathbb Z/2$.

Write $M_2=\operatorname{diag}(1,D)$.  This time use the
unimodular marking
\[
 P_1=\begin{pmatrix}
 -1&1&-1&0&-1\\
 -1&2&-1&0&-1\\
 -1&0&0&2&-1\\
 1&-1&0&-1&1\\
 0&0&2&-1&1
 \end{pmatrix},\qquad \det P_1=-1.
\]
It satisfies $DP_1=P_1\operatorname{diag}(\text{four-cycle},-1)$.
The cap therefore has the same cellular mapping-torus model
used to prove (F.6).  With the unchanged twist $w_2=-e_0$,
the boundary coordinates remain $\theta=a+4k$, $h=k$;
the suspension representative restricts to $w$ in the base
direction and to $-\operatorname{Norm}_4(w)=0$ in the
auxiliary fibre direction.  Consequently the coefficient in
the torsion differential is again one.

The new vertical representatives in lexicographic exterior
bases are
\[
\begin{split}
 q=2:\ &(0,-1,-1,-1,0),\\
 q=3:\ &(0,0,1,0,0,0,0,0,0,1),\\
 q=4:\ &(0,0,0,0,0,0,0,-1,-1,0).
\end{split}
\]
Each has residue one in its global coinvariant
$\mathbb Z/2$.  Thus all three torsion differentials in (F.6)
are isomorphisms.  The integral Leray table and its free
degree-one differential are therefore exactly those of
F.2.6.  They give $H^2=H^3=0$ and $H^4=\mathbb Z^4$.
Universal coefficients, Hurewicz and Poincare duality give
the stated integral homology and three-connectivity.

The relative almost-complex extension over the pair of pants
in F.2.5 applies: the finite matrices have the same even
real eigenspace dimensions and the cusp is the same marked
Kuhn type.  The purely smooth characteristic-vector and
normal-invariant argument in F.4.1 then gives the stated
oriented diffeomorphism type.  No integrability at the finite
punctures is used for this smooth classification.

Finally, (F.U8) has $(\det H,\det H')=(0,2)$ at $v=1$,
so its irreducible rank-five quotient is not self-dual.
The exact program
\texttt{code/fourfolds\_v1\_homology.py} runs every matrix,
Smith, Hermite, suspension-period, pairing and differential
certificate just displayed.  Its monomial marking was found
inside the stated finite cube and then verified directly;
the proof uses that explicit marking, not search completeness.
\end{proof}

\paragraph{Theorem F.2.19 (Proved: a residual obstruction for every lattice).}
Let a rational rank-six datum have finite linear orders five
and four, a primitive maximal cusp on a stable integral lattice,
and invariant line $L$, with irreducible rank-five quotient $W$
and no rational global coinvariants.  Suppose that at a prime
$p\ne2,5$ the semisimplified reduction of a stable lattice in
$W$ is the sum of the trivial module and an absolutely
irreducible module of dimension four.  Then no stable integral
lattice in that rational datum, with primitive maximal cusp,
and no free finite affine twists can give
$\pi_1=1$ and $H_2=0$ simultaneously.

For every $v\notin\{0,1\}$, the rational family in F.2.16
satisfies this hypothesis at each prime dividing
$v^2-v+1$.  Consequently the exclusion of those parameters
holds for \emph{all} stable lattices and all twists, not just
the lattice in (F.U6).  The two remaining parameters are the
verified survivors of F.2.6 and F.2.18.

\begin{proof}
We first give the lattice-independent mechanism.  For any
stable lattice $V_\mathbb Z$, use its primitive invariant
line $L_\mathbb Z$ and quotient lattice $W_\mathbb Z$.
The semisimplified reduction of a rational representation is
independent of its stable lattice.  One elementary way to see
this is to compare two commensurable lattices by their finite
quotient $T$.  The exact sequence obtained by reduction modulo
$p$ has end terms $T[p]$ and $T/pT$, which have the same
composition factors by the filtration of $T$ by powers of $p$.
Hence the two lattice reductions have equal classes in the
Grothendieck group of finite-length modules.  Chaining such
comparisons proves the assertion for any two lattices.

Thus $W_\mathbb Z/pW_\mathbb Z$ has length two, with one
trivial constituent and one irreducible four-dimensional
constituent.  It has either a trivial submodule or a trivial
quotient (and possibly both).

Suppose first that it has a trivial submodule.  Its inverse
image in $V_\mathbb Z/pV_\mathbb Z$ is an extension of a
trivial module by a trivial module.  Each finite generator
acts on this two-dimensional space by a unipotent matrix.
Since its order is five or four, prime to $p$, it acts as
the identity.  There are therefore at least two global
invariant vectors in degree-one fibre cohomology modulo $p$.
Write their dimension as $a_p\ge2$; this is also the
dimension of the global coinvariants of fibre homology over
$\mathbb F_p$, by the perfect pairing between fibre homology
and cohomology.

The abelianized van Kampen presentation now shows
\[
 b_1(M;\mathbb F_p)\ge a_p-1.
 \tag{F.U12}
\]
Indeed, after the monodromy relations the fibre contributes
$a_p$ dimensions, and the two finite meridians contribute
two more.  The two finite power relations and the single
base relation impose at most three linear relations.
The cusp introduces no further fibre relation in this
abelianization.  On the dual homology lattice its nilpotent
operator is still primitive and maximal, so its angular lattice is
$\ker(M_0-I)=\operatorname{im}(M_0-I)$ by primitivity, and
the latter image already vanishes among the global
monodromy relations.  Its base meridian is the base relation
just counted.  Changing affine twists changes coefficients
of these three relations, not their number.  Thus (F.U12)
contradicts simple connectivity.

Suppose instead that the reduction of $W$ has a trivial
quotient.  Its integral global coinvariants have nonzero
$p$-primary torsion, since they have rational rank zero.
The natural map from the global coinvariants of $V$ to
those of $W$ is an isomorphism.  To check its kernel, a
primitive generator $e$ of $L_\mathbb Z$ lies in
$\ker(M_0-I)=\operatorname{im}(M_0-I)$, and hence is already
zero among global coinvariants.  Therefore $H^2(R^1)$ has
nonzero $p$-primary torsion.  If the total space were simply
connected, the finite-cap argument of F.2.17 would make
every possible torsion source in $H^0(R^2)$ supported only
at two and five.  The invariant part of $\bigwedge^2V$ is
zero by irreducibility of the odd-dimensional quotient.
Consequently the nonzero $p$-primary subgroup survives to
$H^3(M;\mathbb Z)$ and then to
$\operatorname{Tor}H_2(M;\mathbb Z)$, again a contradiction.
This proves the general criterion.

For its application to F.2.16, let $p\mid q=v^2-v+1$.
We have $p\ne2,5$.  In the original cycle lattice of (F.U1),
$Y\mathbf1=2q(-1,1-v)^t$, so $\mathbf1$ gives a trivial
submodule after reduction.  On the four-dimensional quotient
the first finite generator has characteristic polynomial
$\Phi_5$, and $C=B^{-1}$ has polynomial
$(x+1)^2(x^2+1)$.  This quotient is absolutely irreducible.
Here is a characteristic-independent proof over
$\overline{\mathbb F}_p$.

A one-dimensional composition factor would make the
eigenvalues of $A$ and $C$ equal, because $CA^{-1}$ is
unipotent.  An eigenvalue of order dividing both five and
four must be one, which is absent from $\Phi_5$.
If the quotient were reducible, its only remaining possible
composition-factor dimensions would be two and two.
On each factor, the determinants of $A$ and $C$ must agree
and have orders dividing five and four, so both determinants
are one.  The two $C$ eigenvalue pairs must therefore be
$(-1,-1)$ and $(i,-i)$ in some order.  On the former
factor $C=-I$, since its order is prime to $p$.
Then the unipotence of $CA^{-1}=-A^{-1}$ would force the
eigenvalues of $A$ to be $-1$, impossible for $\Phi_5$.
This proves absolute irreducibility and establishes the
required residual composition factors at every $p\mid q$.
When $v\notin\{0,1\}$, the integer $q$ is greater than one
and has such a prime divisor.  The criterion applies.
\end{proof}

\paragraph{Theorem F.2.20 (Proved: the complete rational $(5,4)$ normal form).}
Consider ordered rational rank-five pairs $(A,C)$ with
\[
 \chi_A=(x-1)\Phi_5,
 \qquad \chi_C=(x-1)(x+1)^2(x^2+1).
\]
Put
\[
 N=CA^{-1}-I,
 \qquad N^2=0,\quad\operatorname{rank}N=2,
\]
where $A^5=C^4=I$ and the pair is irreducible over $\mathbb Q$.
Their simultaneous rational conjugacy classes are parametrized
bijectively by
\[
 \bigl\{[r:s:t]\in\mathbb P^2(\mathbb Q):D(r,s,t)\ne0\bigr\}
 \ \sqcup\ \mathbb A^1(\mathbb Q),
 \tag{F.R1}
\]
where
\[
\begin{split}
 D={}&r^3-4r^2s-r^2t+7rs^2+rst-3rt^2\\
     &{}-7s^3+11s^2t-11st^2+7t^3.
\end{split}
 \tag{F.R2}
\]
The second summand lies over the single exceptional point
$[4:1:3]$ of $D=0$.  The construction below is explicit and
uses only a ten-variable linear system.  By F.2.11 it also
parametrizes the rational rank-six maximal-cusp extensions
of all these pairs, with $B=C^{-1}$.
This is a rational classification; the stable integral
lattices, affine twists and surviving homology still require
the arithmetic conditions of F.2.

\begin{proof}
Write $K=\mathbb Q(\zeta)$, $\zeta^4+\zeta^3+\zeta^2+
\zeta+1=0$, and identify the rational space with
$\mathbb Q\oplus K$.  The matrix of $A$ is multiplication by
$(1,\zeta)$:
\[
 A=\operatorname{diag}(1,J),\qquad
 J=\begin{pmatrix}
 0&0&0&-1\\1&0&0&-1\\0&1&0&-1\\0&0&1&-1
 \end{pmatrix},\qquad R=(A+I)^{-1}.
 \tag{F.R3}
\]
Its rational centralizer is $\mathbb Q^*\times K^*$.
Let $I_N=\operatorname{im}N$.  If $X$ is any basis matrix
for this plane and $N=XY$, then $YX=0$.
The identity
\[
 C+I=(A+I)+XYA
\]
shows that its two-dimensional kernel has the same dimension
as the kernel of
$I_2+YARX=I_2-YRX$.  Hence
\[
 YRX=I_2,
 \qquad I_N\cap RI_N=0.
 \tag{F.R4}
\]
Conversely the two matrix equations in (F.R4), together with
$YX=0$, give precisely that two-dimensional $-1$ eigenspace.

The plane $I_N$ projects nontrivially to $\mathbb Q$:
otherwise the first-coordinate covector is invariant under
both $A$ and $C$, contrary to irreducibility.
It contains no nonzero vector of $\mathbb Q\oplus0$,
because $R$ acts as $1/2$ on that line and (F.R4) would fail.
Write a basis of the plane as $(0,u),(1,v)$, where
$u\in K^*$.  Multiplication on $K$ by $u^{-1}$ puts it into
\[
 I_N=\langle(0,1),(1,\beta)\rangle_{\mathbb Q}.
\]
Changing the second basis vector by a rational multiple of
the first changes $\beta$ by a rational constant.
Two such planes are in the same centralizer orbit exactly
when $\beta'=a\beta+b$, with $a\in\mathbb Q^*$ and
$b\in\mathbb Q$.  Indeed a centralizer element carrying
one to the other must preserve the kernel of the projection
to $\mathbb Q$, so its $K^*$ component is itself rational;
the displayed affine change follows, and every such change
is realized.  We may therefore take
\[
 \beta=r\zeta+s\zeta^2+t\zeta^3,
 \qquad [r:s:t]\in\mathbb P^2(\mathbb Q),
 \qquad
 X=\begin{pmatrix}0&1\\1&0\\0&r\\0&s\\0&t\end{pmatrix}.
 \tag{F.R5}
\]
The triple is nonzero; a rational $\beta$ would make the
plane contain the forbidden first-coordinate line.
The stabilizer of such a plane in the centralizer consists
only of common scalar matrices.  Thus it introduces no
further identifications of the resulting nilpotent operators.

Write $Y$ as an unknown two-by-five matrix.  It is determined
by the following ten linear equations:
\[
 YX=0,\qquad YRX=I_2,\qquad
 \operatorname{tr}(YAX)=-1,\qquad
 \operatorname{tr}(YA^{-1}X)=1.
 \tag{F.R6}
\]
The first eight are (F.R4).  The last two are exactly
$\operatorname{tr}C=\operatorname{tr}C^{-1}=-1$, since
$\operatorname{tr}A=\operatorname{tr}A^{-1}=0$ and
$C^{-1}=A^{-1}(I-XY)$.
They suffice to obtain the required characteristic
polynomial: the two $-1$ roots and determinant one leave a
monic cubic of constant term $-1$, and the two traces fix
that cubic to $(x-1)(x^2+1)$.
These remaining three roots are distinct, and the $-1$
eigenspace already has dimension two, so $C$ is semisimple
and $C^4=I$.  Also $YRX=I_2$ ensures that $N$ has rank two.
Thus every solution of (F.R6) gives all the asserted local
linear conditions.

In row-major order on the ten entries of $Y$, the determinant
of the coefficient matrix in (F.R6) is $D/4$.
For $D\ne0$ this proves uniqueness and gives the explicit
formula $Y=\operatorname{reshape}(E^{-1}b)$, where $E$ is
the coefficient matrix of the displayed linear equations
and $b$ their displayed right side.
The exact expanded matrix and all ten rational functions
are printed by \texttt{fourfolds\_relative\_normalform.py}.
There is no search or degree cutoff in this calculation.

We check the entire exceptional locus, rather than discard it.
Set
\[
 Q_1=r^2-rs-rt+s^2+t^2,
 \qquad Q_2=r^2-2rs-3rt+4s^2-4st+4t^2.
\]
The first column of the unique generic solution is
\[
 y_{10}=\frac{2Q_1Q_2}{D},\qquad
 y_{20}=\frac{2(r-s-t)Q_1}{D}.
 \tag{F.R7}
\]
Here $y_{10},y_{20}$ are the first column of $Y$.
The two equalities after multiplication by $D$ are polynomial
linear combinations of (F.R6), by its adjugate matrix.
The program prints the corresponding two polynomial rows
and verifies those identities, so they remain necessary
also when $D=0$.
Moreover
\[
 2Q_1=(r-s)^2+(r-t)^2+s^2+t^2>0
\]
for every nonzero rational triple.  Consequently consistency
on $D=0$ forces $r=s+t$ and $Q_2=0$.
Substitution gives
\[
 D(s+t,s,t)=(t-3s)(2t-s)^2,
 \qquad Q_2(s+t,s,t)=(3s-t)(s-2t).
 \tag{F.R8}
\]
There are precisely two resulting projective points,
$[3:2:1]$ and $[4:1:3]$.
At the former, $[X\mid RX]$ has rank less than four, so
(F.R4) is impossible.  At the latter the coefficient and
augmented matrices of (F.R6) both have rank nine, and every
solution is
\[
 Y_z=\begin{pmatrix}
 4z+16/5&0&-2z-1&4/5-2z&2z\\
 2z-2/5&0&-z&2/5-z&z
 \end{pmatrix},\qquad z\in\mathbb Q.
 \tag{F.R9}
\]
Direct substitution proves both necessity and sufficiency.
This establishes the exceptional affine line in (F.R1).

Every pair just constructed is irreducible over $\mathbb Q$.
An $A$-stable proper rational subspace has dimension one or
four.  A one-dimensional subrepresentation or quotient must
have trivial character, since its finite eigenvalues of
orders dividing five and four must agree under the
unipotent product.  The only $A$-invariant covector is the
first-coordinate covector.  Its product with $N$ is the
second row of $Y$, which is nonzero by $YRX=I_2$.
The only $A$-invariant vector is the first-coordinate vector;
its image under $N$ vanishes exactly if both entries in
(F.R7) vanish.  For $D\ne0$, positivity of $Q_1$ and
(F.R8) rule this out.  On the exceptional line the two
entries $4z+16/5$ and $2z-2/5$ never vanish simultaneously.
There are therefore neither invariant vectors nor covectors,
which excludes both possible dimensions.

Finally, rescaling $(r,s,t)$ is exactly the centralizer
change already described; it gives a conjugate pair.
Conversely any conjugacy of two normalized pairs centralizes
$A$ and hence identifies their projective parameters as proved
above.  At the exceptional point its stabilizer is scalar,
so distinct $z$ give distinct pairs.  This proves the stated
bijection, including completeness and uniqueness.
The saved program checks all determinant, adjugate,
characteristic-polynomial and finite-order identities over
the rational function field, and checks the exceptional
line symbolically for every $z$.
\end{proof}

\paragraph{Corollary F.2.21 (Proved: all good-prime lattices are forced).}
For a $(5,4)$ datum with $\pi_1=1$ and $H_2=0$, the
rank-five quotient has absolutely irreducible reduction
at every prime $p\ne2,5$.  At each such prime its stable
$\mathbb Z_p$-lattice is unique up to homothety.
The same is true of the rank-six extension among lattices
with primitive maximal cusp.  Thus, once a rational pair
in F.2.20 admits one survivor, all remaining integral-lattice
ambiguity is supported at the primes two and five.
This assertion supplies no finite bound on those bad-prime
lattices or on the rational parameters.

\begin{proof}
Work first over an algebraic closure of $\mathbb F_p$.
Both finite generators are semisimple because $p\ne2,5$.
A one-dimensional composition factor must be trivial:
its $A$ and $C$ eigenvalues agree by the unipotent product,
and their orders divide five and four.
There can be at most one such factor, since the $1$
eigenvalue of $A$ occurs once.
If it occurs, the remaining four-dimensional module is
absolutely irreducible by the spectral argument in F.2.19.
The only other possible reducible composition pattern is
$2+3$.
On each factor, determinants of $A$ and $C$ agree and
have orders dividing five and four, hence are one.
The two-dimensional $C$ eigenvalue pair is therefore
$(-1,-1)$ or $(i,-i)$.
The first would give $C=-I$ on that factor, contradicting
the unipotence of $CA^{-1}$ and the absence of $-1$
among the eigenvalues of $A$.
In the second case the three-dimensional factor has
$C$ eigenvalues $1,-1,-1$, so $C^2=I$ there; it has
an $A$-fixed vector because the other factor has determinant
one and cannot contain the simple $1$ eigenvalue of $A$.

Here is the required elementary involutory identity in the
present characteristic.  Put $B=A^{-1}$ on this factor.
From $(CB-I)^2=0$ and $C^2=I$ one obtains
\[
 BC+CB^{-1}=2I.
\]
For $Bv=v$ this says $(B+I)Cv=2v$.
The operator $B+I$ is invertible, since $B^5=I$ and
$p\ne2$, so $Cv=v$.  The resulting common fixed line
contradicts simplicity of the three-dimensional factor.
We have proved that any reducible residual representation
has exactly the constituents $1$ and an absolutely
irreducible four-dimensional module.  Theorem F.2.19
excludes that possibility for a survivor, proving absolute
irreducibility at every indicated prime.

Let $R=\mathbb Z_p$ and $K_p=\mathbb Q_p$.
For two stable quotient lattices, multiply one by a power
of $p$ until it is contained in the other but not in its
$p$-multiple.  Its nonzero residual image must then be the
whole quotient by absolute irreducibility.  Nakayama's
lemma, or the Smith form of the inclusion, makes the
lattices equal.  This proves quotient-lattice uniqueness.

For completeness the extension-lattice comparison is also
integral.  Fix a primitive-cusp lattice $V_R$, with invariant
line $Re$ and quotient $W_R$.
The coboundary map
\[
 W_R^*\longrightarrow W_R^*\oplus W_R^*
\]
is primitive, because its reduction has trivial kernel.
Thus its cokernel, the degree-one cohomology of the free
group on the two generators, is torsion free, and an
integral coboundary over $K_p$ is already integral over $R$.
The extension class $\alpha$ is primitive in this cokernel.
Indeed, divisibility by $p$ would make the extra cusp map
\[
 \ker N_W/\operatorname{im}N_W\longrightarrow Re
\]
divisible by $p$, contrary to the three unit Smith factors
of the full cusp.  Here $N_W$ has primitive image:
it is the quotient of the primitive image of $N_V$ by
$Re\subset\operatorname{im}N_V$.

Normalize any other stable lattice by the already proved
quotient homothety so that its quotient is $W_R$.
Its invariant line is $p^rRe$ for some integer $r$.
The extension class in that basis is $p^{-r}\alpha$;
primitivity makes it integral only when $r\leq0$.
Primitivity of the coboundary map eliminates any additional
rational change of splitting.  The extra cusp map is then
multiplied by $p^{-r}$, so its Smith factor is a unit exactly
when $r=0$.  The lattices coincide after the common homothety.
This proves the full assertion and retains the bad-prime
and rational-parameter problems explicitly.
\end{proof}

\paragraph{Theorem F.2.22 (Proved: a terminating lattice algorithm for each rational pair).}
Let $(A,C)$ be any irreducible rational rank-five pair in
one of the eight spectral cases of F.2.14, with
$A^5=I$, $B=C^{-1}$, and $(AB-I)^2=0$ of rank two.
There is an explicit terminating algorithm which decides
whether the pair preserves an integral lattice, and, if so,
lists all such lattices up to simultaneous integral conjugacy.
It then lists all rank-six extensions with the same finite
orders and primitive maximal square-zero cusp, up to
simultaneous integral conjugacy.  The bounds below depend
on the input rational pair and are proved to lose no lattice.
They do not bound the rational parameter space or classify
all affine twists and resulting homology.

\begin{proof}
We give the algorithm and its termination proof.
The fixed line of $A$ has dimension one.  Choose a nonzero
vector $e$ on it and a nonzero fixed covector $\varphi$.
Irreducibility implies that the vectors $ge$, with $g$ a
word in $A,C$, span the whole rational space; the analogous
statement holds for the covectors $\varphi h$.
In fact words of length at most four suffice.  The spans of
words of lengths at most $j$ increase from dimension one;
if one step does not increase the span, it is invariant
under both generators and hence already the whole space.
There are only four possible increases.

Choose words $g_1,\ldots,g_5$ and $h_1,\ldots,h_5$ giving
bases in this way, and put
\[
 P=I+A+\cdots+A^4,\qquad T_{ij}=g_iPh_j.
 \tag{F.L1}
\]
The rank-one operator $P$ is a nonzero scalar multiple of
$e\varphi$.  Hence the 25 matrices $T_{ij}$ are a rational
basis of $\operatorname{Mat}_5(\mathbb Q)$.
Each belongs to $\mathcal R=\mathbb Z[A,C]$.
Let $\mathcal R_0$ be their integral span and set
\[
 \mathcal R_0^\#=
 \{X:\operatorname{tr}(XT)\in\mathbb Z
                     \text{ for all }T\in\mathcal R_0\}.
 \tag{F.L2}
\]
The trace pairing on the full matrix algebra is
nondegenerate, as is seen from the matrix units, so this
is another full rational lattice, computed by inverting
the 25-by-25 Gram matrix.

If any stable lattice exists, every element of $\mathcal R$
acts integrally on it.  Consequently
$\operatorname{tr}(XY)\in\mathbb Z$ for all
$X,Y\in\mathcal R$, and necessarily
\[
 \mathcal R_0\subset\mathcal R\subset\mathcal R_0^\#.
 \tag{F.L3}
\]
A nonintegral Gram entry therefore proves nonexistence.
Otherwise let
$d=[\mathcal R_0^\#:\mathcal R_0]$, the absolute Gram
determinant.  Start with the lattice generated by
$\mathcal R_0$ and $I$, and repeatedly replace a current
lattice $L$ by $L+AL+CL$, using exact Hermite reduction.
If it leaves $\mathcal R_0^\#$, (F.L3) proves that there
is no stable lattice.
While it remains inside, every strict enlargement reduces
its index by a factor at least two.
There can be at most $\lfloor\log_2d\rfloor$ such
enlargements starting from $\mathcal R_0$, so a bound of
$1+\lfloor\log_2d\rfloor$ tests after adjoining $I$ is
sufficient.  Stabilization gives exactly $\mathcal R$:
the stable lattice contains $I$ and all words in $A,C$,
and every element introduced belongs to their ring.
It is then an order, and
$\Lambda_0=\mathcal R\mathbb Z^5$ is an explicitly
computed stable lattice.  This proves the decision step,
including termination in the nonexistence case.

Use a basis of $\Lambda_0$ from now on, so
$\mathcal R\subset\operatorname{Mat}_5(\mathbb Z)$.
Smith reduction gives an integer $c\ge1$ with
\[
 c\operatorname{Mat}_5(\mathbb Z)\subset\mathcal R;
 \tag{F.L4}
\]
one may take the largest Smith factor of the inclusion
of the order in the full matrix lattice.
Normalize any stable lattice by a rational homothety so
that it is contained in $\mathbb Z^5$ and the gcd of all
coordinates of a lattice basis is one.
The operators $cE_{ij}$ preserve it.  The coordinate values
of its vectors generate $\mathbb Z$, so Bezout's identity,
applied to $cE_{ij}$, shows that it contains every $ce_i$.
Therefore every normalized lattice lies in the finite interval
\[
 c\mathbb Z^5\subset\Lambda\subset\mathbb Z^5.
 \tag{F.L5}
\]
Conversely every lattice in this interval stable under
$A,C$ is retained by the algorithm.
The only rational operators commuting with the pair are
scalars: they act by a scalar on the one-dimensional fixed
line of $A$, and then by the same scalar on the spanning
vectors $g_ie$.  Thus rational equivalence of two stable
lattices is precisely homothety.  The coordinate-gcd
normalization is unique, so (F.L5) gives the integral
conjugacy classification without further identifications.

Here is a finite enumeration bound which is also implemented.
For each $p^a\Vert c$, enumerate all subspaces of
$\mathbb F_p^5$ in reduced echelon form.
Beginning with $\mathbb Z^5$, pass through $a$ levels of
stable lattices between $pH$ and $H$, retaining only those
containing $p^a\mathbb Z^5$ and identifying equal Hermite
forms.  Every desired lattice occurs: for any target
$\Lambda$, the chain
$\Lambda+p^j\mathbb Z^5$, $0\le j\le a$, is such a
chain and ends at $\Lambda$.
If $s_p$ denotes the number of all subspaces of
$\mathbb F_p^5$, there are at most $s_p^a$ branches before
deduplication.  Intersect the resulting lattices for the
different primes and discard those of coordinate gcd
greater than one.  The Chinese remainder theorem proves
that this enumerates (F.L5) completely.

For each quotient lattice $W_{\mathbb Z}$ so obtained,
put $N_W=AB-I$.
Its image must be primitive: for a primitive full cusp,
$\operatorname{im}N_V=\ker N_V$ contains the primitive
invariant line, and quotienting that image by the line
gives a primitive image in $W_{\mathbb Z}$.
Reject a quotient whose two nonzero Smith factors are
not one.  For the remaining quotients use extension rows
$c_1,c_2\in\mathbb Z^{1\times5}$ and impose
\[
\begin{gathered}
 c_1(I+A+\cdots+A^4)=0,
 \qquad c_2(I+B+\cdots+B^{m-1})=0,\\
 (c_2+c_1B)N_W=0,
 \qquad m=\operatorname{ord}(B).
\end{gathered}
 \tag{F.L6}
\]
These are integral linear equations.  Their integral kernel
modulo the coboundaries
\[
 z\longmapsto\bigl(z(A-I),z(B-I)\bigr),
 \qquad z\in\mathbb Z^{1\times5},
 \tag{F.L7}
\]
is a finitely generated group $G$ of rank one by F.2.11.
Compute its free and torsion factors by Smith reduction.
Choose an integral generator $w$ of the free rank-one
quotient $\ker N_W/\operatorname{im}N_W$.
The homomorphism
\[
 \kappa:G\longrightarrow\mathbb Z,
 \qquad [(c_1,c_2)]\longmapsto(c_2+c_1B)w
 \tag{F.L8}
\]
is well defined and nonzero: (F.L6) kills the image, and
(F.L7) kills $w$; rational nonvanishing is F.2.11.
The full cusp has unit Smith factors exactly when
$\kappa=\pm1$.
Indeed the quotient rank-two image is already primitive,
and the remaining invariant-line image is generated by
that one integer.
If $\operatorname{im}\kappa\ne\mathbb Z$, reject the
lattice.  Otherwise the inverse image of $1$ is a finite
coset of $\operatorname{Tor}G$, and lists all extensions.
Changing the sign of the invariant-line generator identifies
$\kappa=-1$ with $\kappa=1$.
There are no further identifications: an integral isomorphism
on the quotient is $\pm I$, on the invariant line is
$\pm1$, and, after $\kappa=1$ is fixed, their ratio is one.
The remaining row is exactly an integral change of splitting
(F.L7).
This proves completeness and termination of the extension
step as well.

The order, trace-dual bound, conductor and finite-interval
steps are implemented in
\texttt{fourfolds\_lattice\_algorithm.py};
the extension-kernel and Smith steps are implemented in
\texttt{fourfolds\_two\_classes\_lattices.py} with the
finite order as an argument.  Their applications and exact
exhaustive bounds are recorded next.
\end{proof}

\paragraph{Theorem F.2.23 (Proved: complete lattice classes for the two family survivors).}
For each of the two rational representations $v=0$ and $v=1$
of F.2.16, there are exactly two homothety classes of stable
rank-five quotient lattices, and exactly three integral
conjugacy classes of rank-six extensions with primitive
maximal cusp.  Exactly one of those three can have simply
connected finite affine fillings.  It is the integral class
already realized in F.2.6, respectively F.2.18.
Thus arbitrary stable-lattice changes produce no further
linear survivor classes anywhere in the family F.2.16.
This assertion retains the classification of all affine
twists and multiplicities as a separate problem.

\begin{proof}
Apply the preceding algorithm in the original five-cycle
lattice (F.U1).  For both parameters the order is obtained
after four strict enlargements.  Its Smith factors in the
full matrix lattice are 21 factors one and four factors two;
thus the conductor $c=2$ in (F.L4) is exact.
All 374 subspaces of $\mathbb F_2^5$, including zero and
the whole space, are enumerated, with no height restriction
or neighbor-depth assumption.
After imposing stability and coordinate gcd one, the two
lattice representatives are
\[
 \mathbb Z^5,\qquad
 \Lambda_* =2\mathbb Z^5+\mathbb Z(1,1,1,1,1)^t.
 \tag{F.L9}
\]
The latter has index $16$.
The generic prime-by-prime implementation of (F.L5) and
an independent bit-vector enumeration give exactly the
same two Hermite forms.

For $\mathbb Z^5$, the extension group in (F.L6)--(F.L7)
is $\mathbb Z$ and $\kappa$ is a unit on its generator.
For $\Lambda_*$, that group is
$\mathbb Z\oplus\mathbb Z/2$, again with primitive
$\kappa$.  Hence there are exactly one and two primitive
extensions, respectively.
For each extension the program computes the image of the
primitive global invariant homology covector $\ell$ on
the two finite fixed lattices.  The resulting positive
divisibilities are
\[
\begin{array}{c|c|c}
 \text{quotient lattice}&\text{extension class}&
       (\operatorname{div}_1,\operatorname{div}_2)\\\hline
 \mathbb Z^5&\text{unique}&(1,2)\\
 \Lambda_*&\text{first torsion coset}&(1,1)\\
 \Lambda_*&\text{second torsion coset}&(1,2).
\end{array}
 \tag{F.L10}
\]
The table is identical for $v=0$ and $v=1$.
All six full matrices, the integral kernel bases, the
Smith factors and the finite divisibilities are computed
exactly by \texttt{fourfolds\_two\_classes\_lattices.py}
and saved in its output and JSON certificate.

For a simply connected filling the determinant quotient
of F.2.12 requires $\gcd(\ell(v_i),m_i)=1$ at each
finite cap.  The second affine multiplicity is divisible
by four.  A divisibility of two in (F.L10) therefore
excludes that class for every choice of twists and larger
multiplicities.
The sole remaining class has an explicit integral
intertwiner of determinant $\pm1$ with the corresponding
pair in (F.U6), as printed and checked by the program.
Those pairs already have the verified simply connected,
three-connected fillings.  F.2.19 excludes every other
rational parameter for all lattices.
Together these facts prove all the claims, including the
absence of any additional linear survivor class in the
unbounded family.
\end{proof}

\paragraph{Theorem F.2.24 (Proved: every surviving twist in the unbounded family).}
Use the fixed Kuhn cusp and the affine translation clutching
convention of F.2.12.  In the entire rational family F.2.16,
allow every stable lattice, every free affine finite lift,
and every pair of actual finite multiplicities.
The data with $\pi_1=1$ and $H_2=H_3=0$ are precisely the
following, up to the integral linear classification already
proved and the translation identifications described below:
\[
\begin{gathered}
 v\in\{0,1\},\qquad
 \text{the unique eligible integral class of F.2.23},\\
 (m_1,m_2)=(5,4),\qquad
 20\kappa+4r_1+5r_2=\pm1.
\end{gathered}
 \tag{F.T1}
\]
Here $r_i=\ell(v_i)$ is the invariant affine numerator, and
$\kappa\in\mathbb Z$ is the invariant coordinate of the
cusp translation clutching, with signs fixed by the displayed
van Kampen relation below.  The remaining fixed-direction
components of the twists impose no condition.
Every manifold in (F.T1) is oriented diffeomorphic to its
corresponding known survivor and hence to
$3\mathbb{HP}^2\#(-\mathbb{HP}^2)$.
This completes the survivor classification for this entire
unbounded rational family; it does not classify the other
rational points of F.2.20 or the other spectral loci.

\begin{proof}
F.2.19 excludes all other rational parameters, and F.2.23
excludes all other integral linear classes.
We therefore work with either known integral pair, in the
marking of (F.U6).  Its primitive invariant homology
covector is $\ell(x)=x_0$.
Each finite fixed lattice has an integral basis
$u_i,\lambda_i$, where $\ell(u_i)=1$,
$\ell(\lambda_i)=0$, and $\lambda_i$ is primitive in
$\ker\ell$.
We may write the affine numerator as
$v_i=r_i u_i+s_i\lambda_i$.
This covers arbitrary finite affine lifts: translating the
origin makes their translation vector invariant by averaging,
and its $m_i$th iterate is then the integral vector $v_i$.
The resulting boundary translation change is included in the
specified affine clutching.

For both linear pairs, the normal closure of the cusp angular
lattice in the fibre lattice is exactly $\ker\ell$, by
F.2.5 and F.2.18.  Van Kampen therefore leaves the abelian
presentation
\[
 m_1a=r_1h,\qquad m_2b=r_2h,\qquad
 a+b+\kappa h=0.
 \tag{F.T2}
\]
This also defines the sign of $\kappa$.
Its determinant is, up to sign,
\[
 p=m_1m_2\kappa+m_2r_1+m_1r_2.
 \tag{F.T3}
\]
Thus simple connectivity is exactly $p=\pm1$.
It forces $\gcd(m_1,m_2)=1$ and
$\gcd(r_i,m_i)=1$ by reduction modulo each multiplicity.

We first prove that a surviving datum has the minimal
multiplicities.  Set $n_1=5$, $n_2=4$ and $m_i=n_i k_i$.
In the splitting $\Lambda=\mathbb Zu_i\oplus\ker\ell$,
write the finite linear map as $1\oplus D_i$.
The reduced finite fibre is a mapping torus of $T^5$ over
the circle $m_i\ell$; its fibre is connected because
$r_i$ is prime to $m_i$.
Its vertical monodromy is a power of $D_i$ prime to $n_i$,
possibly followed by a translation, which changes no
integral Wang group.

The primitive fixed vertical loop $\lambda_i$ gives the
only additional free two-cycle beyond vertical cycles.
Consequently an invariant integral two-form $\omega$ on
the covering torus descends integrally exactly when
\[
 \frac{\omega(v_i,\lambda_i)}{m_i}
       =\frac{r_i\omega(u_i,\lambda_i)}{m_i}
       \quad\text{is an integer}.
 \tag{F.T4}
\]
For precision, the torus made from the affine generator
and $\lambda_i$ projects with degree $r_i$ to the mapping-
torus base, whereas the ordinary $u_i$ suspension projects
with degree $m_i$.  A Bezout combination is a primitive
suspension.  Its integrality condition is equivalent to
(F.T4), because the ordinary covering periods are already
integral and $\gcd(r_i,m_i)=1$.
The Wang sequence shows that these vertical cycles and one
primitive suspension generate the free second homology;
there is no omitted period condition.

The integral invariant two-form computations in F.2.6 and
F.2.18 give
\[
 \gcd_{\omega}\omega(u_i,\lambda_i)=n_i.
 \tag{F.T5}
\]
They are also independently printed for both pairs by
\texttt{fourfolds\_family\_twists.py}.
Hence the covering-image cokernel in degree two at the
$i$th finite point is exactly $\mathbb Z/k_i$.
The global invariant group of $\bigwedge^2\mathcal V$
and its integral first parabolic group both vanish.
The skyscraper cokernel sequence therefore gives
\[
 H^1(B,R^2g_*\mathbb Z)
       =\mathbb Z/k_1\oplus\mathbb Z/k_2.
 \tag{F.T6}
\]
Point-supported torsion kernels do not change this degree-one
base cohomology, and the cusp specialization is onto.
The term (F.T6), in bidegree $(1,2)$, survives every Leray
differential on the two-dimensional base.
All rational terms of total degree three vanish, independently
of the multiplicities and twists, by the same exterior
local-system table as F.2.6.  Thus $H^3$ is finite.
If either $k_i>1$, its surviving nonzero filtration quotient
makes $H^3$ nonzero; universal coefficients then give
nonzero torsion in $H_2$.
This proves $m_1=5$, $m_2=4$ for every survivor.

We now identify all twists with these multiplicities.
The following integral fact permits removal of each $s_i$:
\[
 \operatorname{Norm}_{n_i}:
 \ker\ell\longrightarrow\mathbb Z\lambda_i
 \quad\text{is onto}.
 \tag{F.T7}
\]
Here the norm is $I+D_i+\cdots+D_i^{n_i-1}$.
One can read (F.T7) directly from (F.T5).
Let $\varphi_i$ generate the primitive invariant covector
lattice on $\ker\ell$.
The splitting by $u_i$ shows that the gcd in (F.T5) is
$|\varphi_i(\lambda_i)|$: its forms are
$\ell\wedge\varphi_i$ plus invariant vertical forms.
Thus that value is $n_i$.
The norm has the form
$x\mapsto d_i\varphi_i(x)\lambda_i$ for an integer
$d_i$; evaluating $\varphi_i$ gives
$|d_i|n_i=n_i$, hence $|d_i|=1$.
This proves surjectivity.  The saved program additionally
gives integral preimages $q_i\in\ker\ell$ with
$\operatorname{Norm}_{n_i}q_i=\lambda_i$.

Put
\[
 w_i=-\frac{s_i}{n_i}\sum_{j=0}^{n_i-1}jM_i^j q_i.
\]
Then
\[
 (I-M_i)w_i=-s_i\lambda_i/n_i+s_iq_i.
 \tag{F.T8}
\]
Translation by $w_i$ conjugates the finite affine action
with numerator $r_i u_i+s_i\lambda_i$ to that with
numerator $r_i u_i$: the remaining $s_iq_i$ is integral.
It gives a diffeomorphism of the entire finite cap.
In the standard linear boundary marking, its translation
section is
$w_i+\theta s_i\lambda_i/n_i$.
Its endpoint mismatch with the linear deck map is exactly
$s_iq_i$ by (F.T8), so its integral boundary residue has
$\ell$-coordinate zero.

We explain why these local changes extend simultaneously,
including the prescribed cusp marking.
Let $P$ be the complementary pair of pants and
$\mathcal T\to P$ its flat torus group bundle.
A section on the boundary extends over the relative
one-skeleton because the torus is connected.
The only remaining obstruction is the two-cell obstruction
with values in the fibre fundamental-group local system:
\[
 H^2(P,\partial P;\Lambda)
        \simeq H_0(P;\Lambda)=\Lambda_\Gamma.
 \tag{F.T9}
\]
The equality follows from the oriented relative cellular
complex, equivalently integral Poincare--Lefschetz duality
with local coefficients.  Its boundary residues map to
the indicated global coinvariants.
For both linear pairs the horizontal monodromy relation
matrix has five unit Smith factors, so
$\Lambda_\Gamma=\mathbb Z$, detected exactly by $\ell$.
This is independently checked in the saved program.
Thus the residues $s_iq_i$, with the identity section at the
cusp, have zero obstruction and extend over $P$.
A section may be smoothed relative to its boundary collars;
translation by it is a fibre-preserving diffeomorphism.
Together with the finite-cap maps it removes all $s_i$
without changing the cusp.
The same argument removes all components of translation
clutching whose $\ell$-coordinate is zero.

It remains to handle the integer numerators and the remaining
cusp coordinate.  If $p=-1$, the congruences from (F.T3)
give
\[
 r_1=1+5a,\qquad r_2=-1+4b,
 \qquad a+b+\kappa=0.
 \tag{F.T10}
\]
Adding $5a u_1$ or $4b u_2$ changes no finite affine action;
in its linear boundary marking it changes the translation
residue by $a u_1$ or $b u_2$.
The difference of the cusp marking has residue of
$\ell$-coordinate $\kappa$, by (F.T2).
Their total obstruction in (F.T9) is therefore
$a+b+\kappa=0$.
The same extension argument identifies the glued manifold
with the original twists $(u_1,-u_2)$ and zero cusp
clutching.
If $p=1$, the identical calculation starts from
$(-u_1,u_2)$.

Finally simultaneous fibre inversion identifies these two
sign choices by an orientation-preserving diffeomorphism.
It conjugates every affine translation to its negative.
It extends over the Kuhn cusp: the lattice involution
$(x,h)\mapsto(-x,h)$ preserves its toric fan, since a
simplex given by a chain in a unit cube is taken to the
chain with reversed coordinate order in the opposite cube.
It normalizes the cusp lattice action and therefore descends
to its quotient, while fixing the base parameter.  The resulting toric map is holomorphic,
and in the regular real six-torus inversion has determinant
$(-1)^6=1$.  All translation maps also preserve orientation.
Thus every datum in (F.T1) gives the stated oriented smooth
type and has $H_2=H_3=0$, proving sufficiency as well as
necessity.
\end{proof}

\paragraph{Theorem F.2.25 (Proved: the universal degree-two suspension obstruction).}
Every datum of F.2 with $\pi_1=1$ and $H_2=H_3=0$ has
exactly two genuine finite fillings and has minimal affine
multiplicities.  After ordering the order-five point first,
its finite characteristic polynomials are one of precisely
the following two remaining necessary possibilities:
\[
\begin{array}{c|c|c}
 (m_1,m_2)&\chi_{M_1}&\chi_{M_2}\\\hline
 (5,4)&(x-1)^2\Phi_5&(x-1)^2\Phi_2^2\Phi_4\\
 (5,6)&(x-1)^2\Phi_5&(x-1)^2\Phi_2^2\Phi_3.
\end{array}
 \tag{F.S1}
\]
In particular all order-eight and order-twelve spectral
loci, and the other two order-six spectral loci of F.2.14,
are excluded for every rational relative position, every
stable lattice and every free affine twist.
There is no multiplicity search bound in this statement:
the only possible actual affine multiplicities are the
two displayed pairs.

\begin{proof}
Theorem F.2.7 excludes three genuine finite fillings.
For two fillings, F.2.9 and F.2.14 give a primitive global
invariant homology covector $\ell$, finite fixed rank two,
and an irreducible rank-five cohomology quotient.
In particular
\[
 H^0(B,j_*\bigwedge^2\mathcal V)=0.
 \tag{F.S2}
\]
This is also immediate from the exterior filtration and
irreducibility in odd dimension, as in F.2.17.
Simple connectivity gives the determinant quotient of
F.2.12, so each invariant affine numerator $r_i=\ell(v_i)$
is prime to the actual multiplicity $m_i$, and the two
multiplicities are coprime.
As before one may conjugate any finite affine lift by a
translation so that its numerator $v_i$ is integral and
fixed by its linear part.

Consider one finite point.  Let $M$ have linear order $n$,
so $n\mid m$, and let its complementary degree-four factor
be $f(x)$, with
$\chi_M=(x-1)^2f(x)$.
The image of $\ell$ on the finite fixed lattice is $g\mathbb Z$
for some $g\mid n$: averaging $I+M+\cdots+M^{n-1}$ on
any vector of $\ell$-value one proves this divisibility.
Since $g\mid r$ and $\gcd(r,m)=1$, we have $g=1$.
Choose a fixed integral vector $u$ with $\ell(u)=1$.
It gives an integral invariant splitting
\[
 \Lambda=\mathbb Zu\oplus H,
 \qquad H=\ker\ell,\qquad M=1\oplus D,
 \qquad\chi_D=(x-1)f(x).
 \tag{F.S3}
\]
The order of $D$ is $n$.
Let $\lambda$ generate its primitive fixed lattice and
let $\varphi$ generate its primitive invariant covector
lattice.  Choose their signs so that
\[
 h=\varphi(\lambda)>0.
 \tag{F.S4}
\]
Their rational pairing is nonzero, since the finite map
is semisimple and its averaging projection has rank one.

Both $\operatorname{Norm}_n(D)$ and $f(D)$ are integral
operators, have image on the line $\mathbb Q\lambda$,
and have invariant defining covector proportional to
$\varphi$.  Primitivity of $\lambda$ and $\varphi$
therefore implies
\[
 \operatorname{Norm}_n(D)=\frac n h\lambda\varphi,
 \qquad
 f(D)=\frac{f(1)}h\lambda\varphi,
 \qquad h\mid n,\quad h\mid f(1).
 \tag{F.S5}
\]
For example, the integer coefficient in the second formula
is seen by applying the integral operator to a vector on
which $\varphi$ is one; evaluating on $\lambda$ determines
that coefficient.  No lattice splitting over a localization
or semisimplification modulo a prime is used here.

Write $v=ru+s\lambda$.  The quotient finite fibre is a
connected $T^5$ mapping torus over the circle $m\ell$,
because $r$ is prime to $m$.
Its vertical cycles already have integral periods for an
invariant integral two-form on the covering torus.
The Wang sequence leaves one additional primitive free
suspension cycle, from $\lambda$.
Exactly as in (F.T4), descent of an invariant form $\omega$
is equivalent to
$m\mid r\omega(u,\lambda)$, and hence to
$m\mid\omega(u,\lambda)$.
In the integral splitting (F.S3), every invariant two-form
is an integral multiple of $\ell\wedge\varphi$ plus an
invariant vertical two-form.  Its possible values on
$(u,\lambda)$ are therefore exactly $h\mathbb Z$.
It follows that the degree-two specialization image at
this point has cokernel
\[
 C^2_p\simeq\mathbb Z/(m/h).
 \tag{F.S6}
\]
The integer $m/h$ is defined because $h\mid n\mid m$.
Torsion in the finite fibre cohomology is a specialization
kernel, not an additional restriction on these free periods.
The primitive suspension and vertical cycles generate the
free second homology, so (F.S6) is the full cokernel.

Let $I^2$ be the image sheaf of
$R^2g_*\mathbb Z\to j_*\bigwedge^2\mathcal V$.
The cusp specialization is onto by E.1.3, and all kernels
and cokernels at finite points are skyscraper sheaves.
Thus $H^1(R^2g_*\mathbb Z)=H^1(I^2)$.
Equation (F.S2) and the cokernel exact sequence give an
injection
\[
 \bigoplus_{p\text{ finite}}C^2_p
           \lhook\joinrel\longrightarrow H^1(B,R^2g_*\mathbb Z).
 \tag{F.S7}
\]
The Leray term of bidegree $(1,2)$ has no incoming or
outgoing differential on a real two-dimensional base.
If $H_2=H_3=0$, the universal coefficient theorem gives
$H^3=0$, so that term must vanish.
Every cyclic cokernel in (F.S7) is therefore zero.
Combining (F.S5)--(F.S6) yields the universal necessary
condition
\[
 m=h=n,\qquad n\mid f(1).
 \tag{F.S8}
\]
This both fixes the multiplicity and tests the spectrum.

For completeness, apply (F.S8) to the entire eight-case
list of F.2.14.  The order-five factor gives $\Phi_5(1)=5$.
At the other point the exact values are
\[
\begin{array}{c|c|r|r}
 n&f&f(1)&n/\gcd(n,f(1))\\\hline
 4&\Phi_2^2\Phi_4&8&1\\
 6&\Phi_2^2\Phi_3&12&1\\
 6&\Phi_2^2\Phi_6&4&3\\
 6&\Phi_3\Phi_6&3&2\\
 8&\Phi_8&2&4\\
 12&\Phi_{12}&1&12\\
 12&\Phi_3\Phi_4&6&2\\
 12&\Phi_4\Phi_6&2&6.
\end{array}
 \tag{F.S9}
\]
The last column is a lower bound for the nonzero local
covering-cokernel order already at the smallest allowed
multiplicity.  Only the first two rows pass.
The complete table is checked by
\texttt{fourfolds\_suspension\_filter.py}; its input is
exactly these eight spectra, with no matrix-entry or
relative-position cutoff.  This proves (F.S1) and all
claimed universal exclusions.
\end{proof}

\paragraph{Theorem F.2.26 (Proved: the complete rational remaining $(5,6)$ locus).}
The irreducible rational rank-five pairs in the second row
of (F.S1), with square-zero rank-two cusp, are classified
up to simultaneous rational conjugacy by
\[
 \{[r:s:t]\in\mathbb P^2(\mathbb Q):
             D(r,s,t)\ne0,\ Q_-(r,s,t)\ne0\}
       \ \sqcup\ \mathbb A^1(\mathbb Q),
 \qquad Q_-=r^2-rs+st,
 \tag{F.R10}
\]
where $D$ is the same cubic as in (F.R2).
The exceptional line again lies over $[4:1:3]$.
Thus the two rational spectral loci left by F.2.25 both
have complete explicit normal forms, including their
entire rational exceptional loci.

\begin{proof}
Use precisely the space $\mathbb Q\oplus\mathbb Q(\zeta_5)$,
matrices $A,R,X$ and projective plane normalization in
(F.R3)--(F.R5).  The centralizer and image-plane arguments
in F.2.20 require only irreducibility and the two-dimensional
$-1$ eigenspace of $C$; both hold here.
Consequently every pair has a unique projective image-plane
parameter, with only scalar stabilizer.
The unknown $Y$ now satisfies
\[
 YX=0,\qquad YRX=I_2,\qquad
 \operatorname{tr}(YAX)=-2,\qquad
 \operatorname{tr}(YA^{-1}X)=2.
 \tag{F.R11}
\]
Its coefficient matrix is unchanged from (F.R6), so has
determinant $D/4$.
The two traces and the $-1$ eigenspace imply
\[
 \chi_C=(x-1)(x+1)^2(x^2+x+1).
\]
Its other three roots are distinct, and the $-1$ eigenspace
has its full multiplicity, so $C^6=I$.
The cusp $XY$ still has square zero and rank two.
For $D\ne0$, the ten linear equations give the unique
explicit solution.

The two first-column adjugate identities become
\[
 D y_{10}=2Q_-Q_2,
 \qquad D y_{20}=2(r-s-t)Q_-,
 \tag{F.R12}
\]
with the same $Q_2$ as in (F.R7).
For $D\ne0$, these entries vanish simultaneously exactly
when $Q_-=0$.
Indeed, if $Q_-\ne0$ they would instead force
$r=s+t$ and $Q_2=0$, and (F.R8) would give $D=0$.
The invariant-covector argument of F.2.20 is unchanged.
Thus the pair is rationally irreducible exactly off the
conic $Q_-=0$ in the generic locus.

No nonzero rational triple satisfies $D=Q_-=0$.
If $s=0$, the conic forces $r=0$, and then $D=7t^3$.
If $s\ne0$, rescale to $s=1$ and put
$r=u$, $t=u-u^2$.  Direct substitution gives
\[
 D(u,1,u-u^2)=
 -(7u^2-11u+7)(u^4-u^3+u^2-u+1).
 \tag{F.R13}
\]
The first factor has discriminant $-75$ and positive
leading coefficient.  The second is strictly positive,
as the identity
\[
 u^4-u^3+u^2-u+1
 =\left(u^2-\frac u2\right)^2
  +\frac34\left(u-\frac23\right)^2+\frac23
\]
shows.  In particular neither has a rational zero.
On the exceptional cubic the necessary identities (F.R12)
therefore force the same two points $[3:2:1]$ and $[4:1:3]$
as before.  The first is again impossible because
$[X\mid RX]$ has rank less than four.
At the second, the complete solution is
\[
 Y_z=\begin{pmatrix}
 4z+26/5&0&-2z-2&-2z-1/5&2z+1\\
 2z-2/5&0&-z&2/5-z&z
 \end{pmatrix},\qquad z\in\mathbb Q.
 \tag{F.R14}
\]
Its first-column entries cannot vanish simultaneously,
so each pair on this line is irreducible.
There are no further identifications, by the scalar
stabilizer argument already proved in F.2.20.
This establishes the claimed bijection.

The exact program is
\texttt{fourfolds\_relative\_six\_probe.py}.
It checks the modified characteristic polynomials, finite-order
identities, first-column adjugate expressions, conic
substitution and entire exceptional line symbolically.
Its auxiliary trace-parameter computations for the other
repeated-negative spectrum are recorded as well, although
that spectrum is now excluded integrally by F.2.25.
No finite sample is used to infer a rational classification.
The finite lattice and extension algorithm of F.2.22 applies
to every point of (F.R10); it does not by itself decide
which rational parameters survive the remaining integral
homology tests.
\end{proof}

\paragraph{Theorem F.2.27 (Proved: the entire remaining order-six locus is excluded).}
Every simply connected datum in the second row of (F.S1),
at every rational relative position and every stable integral
lattice with primitive maximal cusp, has
\[
                 \operatorname{Tor}_3 H_3(M;\mathbb Z)\ne0
\]
if its degree-two suspension cokernels vanish.
Consequently no such datum is three-connected.  Together with
F.2.25 this leaves only the minimal order pair $(5,4)$ as a
possible three-connected datum in the full problem F.2.
\begin{proof}
We use cohomology with coefficients in the discrete valuation
ring $R=\mathbb Z_3$.  Suppose the degree-two suspension
cokernels vanish; otherwise F.2.25 already gives nonzero
$H_2$ or $H_3$.  That theorem forces $m=n=h=6$ at the
order-six point, in the notation (F.S3)--(F.S5).
We first determine its local integral representation over $R$.

The splitting by the fixed vector $u$ is integral:
$\Lambda_R=Ru\oplus H_R$.  On $H_R$, the finite map $D$
has characteristic polynomial
$(x-1)(x+1)^2\Phi_3(x)$.  Since two is a unit in $R$,
the idempotents $(I\pm D^3)/2$ give an integral direct sum
\[
 H_R=L_+\oplus L_-,\qquad
 \operatorname{rank}L_+=3,\quad
 \operatorname{rank}L_-=2.
\]
On $L_-$ the map is exactly $-I$, by its rational spectrum.
On $L_+$ it has order three.  The primitive fixed vector
$\lambda$ lies in $L_+$, and the primitive invariant
covector $\varphi$ vanishes on $L_-$, so their restrictions
remain primitive.  Equation (F.S5) with $h=6$ gives
\[
 2(I+D+D^2)|_{L_+}=\lambda\varphi|_{L_+}.
 \tag{F.X1}
\]
This operator is nonzero modulo three.  But modulo three
$I+D+D^2=(D-I)^2$ and $(D-I)^3=0$.  The reduction of $D$
on the rank-three space therefore has one Jordan block
of length three.  Choose $q\in L_+$ whose image has
$(D-I)^2q\ne0$ modulo three.  Then $q,Dq,D^2q$ are a
basis modulo three and hence an $R$-basis.  In this basis
$D$ is the three-cycle permutation matrix, since $D^3=I$.
We have proved, without a global lattice classification,
that the full finite representation over $R$ is
\[
                   1\oplus\operatorname{Perm}_3\oplus(-I_2).
 \tag{F.X2}
\]
Its dual has the same form, by reversing the three-cycle.

Let $u,a,b,c,d,e$ be a cohomology basis for (F.X2), with
$a,b,c$ cycled and $d,e$ negated.  A full integral basis
of the invariant degree-three lattice is
\[
 u(ab+bc+ca),\qquad abc,\qquad ude,\qquad(a+b+c)de.
 \tag{F.X3}
\]
Here juxtaposition denotes exterior multiplication.
To check completeness, terms with an odd number of
$d,e$ are negated by the cube of the finite map and
have no invariants.  The terms with zero or two such
factors give exactly the four orbit sums in (F.X3),
each primitive.  With volume $uabcde$, their skew
pairing has one block with coefficient three and one
block with coefficient a unit.  Thus its determinant
has three-adic valuation two.  The exact certificate
\texttt{fourfolds\_universal\_56\_obstruction.py} checks
the full invariant kernel, its primitivity, and the
pairing independently by exterior multiplication and
six-column determinants; its Smith factors are
$(1,1,3,3)$.

Let $B_6$ be the actual reduced finite fibre.  Its torus
cover has degree six.  Its free middle cohomology has
rank four and unimodular skew pairing, by integral
Poincare duality.  If $I$ denotes the index of its
pullback lattice in the invariant lattice, multiplication
of pairings by the covering degree gives
\[
 2v_3(I)+2=4v_3(6)=4,
 \qquad v_3(I)=1.
 \tag{F.X4}
\]
Consequently the three-primary degree-three covering
cokernel is $R/3R$.  Translations of the affine action
do not enter this calculation: the cover degree,
linear invariant lattice and oriented pairings are
those of the actual free quotient for every allowed
affine numerator.

We explain why this local group must appear in total
homology even when the global integral lattice varies.
Let $\mathcal V$ be the rational rank-six cohomology
local system.  It is an extension of the irreducible
rank-five system $W$ by a trivial line.  The determinant
of $W$ is trivial.  Irreducibility in odd dimension
implies that neither $\bigwedge^2 W$ nor its dual
has a nonzero invariant: an invariant alternating
form would have an invariant radical and would have
to be nondegenerate, impossible in odd dimension.
Also $\bigwedge^3W\simeq\bigwedge^2W^*$ by its
trivial determinant.  The exterior filtration therefore
gives
\[
 (\bigwedge^3\mathcal V)^\Gamma=0,\qquad
 ((\bigwedge^2\mathcal V)^*)^\Gamma=0.
 \tag{F.X5}
\]
These rational vanishing statements remain true over
$\mathbb Q_3$, since they are kernels of rational
matrices.

As in (F.S7), take the image of
$R^3g_*R\to j_*\bigwedge^3\mathcal V_R$.
Its kernel and cokernel are point-supported sheaves.
The first vanishing in (F.X5) makes the local group
$R/3R$ inject into $H^1(B,R^3g_*R)$.  Bidegree $(1,3)$
has no incoming or outgoing Leray differential on
the two-dimensional base, so this injection survives
to $E_\infty^{1,3}$.

There is one extension issue, and it cannot remove
this torsion.  The second vanishing in (F.X5) says
that $H^2(B,R^2g_*R)$ has rational rank zero.
Point-supported changes do not alter this degree-two
base cohomology; it is the coinvariant module, as
in (E.9).  It is a finitely generated torsion $R$-module
and hence a finite group.  Therefore
$E_\infty^{2,2}$ is finite as well.  In the filtration
of $H^4(M;R)$, take the inverse image of the injected
$R/3R\subset E_\infty^{1,3}$.  It is an extension
of $R/3R$ by the finite group $E_\infty^{2,2}$.
It is thus a nonzero finite three-primary subgroup
of $H^4(M;R)$; in particular it cannot be absorbed
as a finite quotient of a free filtration kernel.
Universal coefficients identify this torsion with
the three-primary torsion of $H_3(M;\mathbb Z)$.
This proves the claim at every rational position and
every integral lattice, and completes the exclusion.
\end{proof}

\paragraph{Theorem F.2.28 (Proved: universal fundamental group and twist normalization).}
In the remaining $(5,4)$ locus, fix the marked standard Kuhn
cap and allow affine translation clutching.
Let $\Lambda=\mathbb Z^6$ be fibre homology and let
$\Lambda_\Gamma$ be its global monodromy coinvariants.
Write $\Lambda_\Gamma=\mathbb Zh\oplus T$, with $T$ finite,
and let $\ell$ be the primitive invariant covector detecting
the free factor.  The following are complete reductions for
any datum that could satisfy $\pi_1=1$ and $H_2=H_3=0$.

The actual multiplicities must be $(5,4)$, and the local
vertical fixed pairings must be $h_1=5,h_2=4$, as in F.2.25.
For any affine twists and clutching the fundamental group
is abelian.  If its free presentation determinant is nonzero,
its order is
\[
 |\pi_1|=|20\kappa+4r_1+5r_2|\,|T|,
 \qquad r_i=\ell(v_i).
 \tag{F.T11}
\]
If that determinant is zero the group is infinite.
In particular simple connectivity is equivalent to
\[
 T=0,\qquad 20\kappa+4r_1+5r_2=\pm1.
 \tag{F.T12}
\]
Whenever these conditions and $h_1=5,h_2=4$ hold, every
such twist choice is oriented diffeomorphic to the normalized
filling with twists $(u_1,-u_2)$ and zero cusp clutching,
where $u_i$ is any integral fixed vector with $\ell(u_i)=1$.
Thus each eligible integral linear datum has exactly one
normalized smooth filling to test for the remaining
homology conditions; no affine-twist enumeration remains.

\begin{proof}
The image of the fibre fundamental group is a normal
abelian subgroup in the van Kampen presentation.
After quotienting by it, the two finite meridians satisfy
$a^5=b^4=ab=1$.
That quotient is trivial, because five and four are coprime.
The entire fundamental group is therefore equal to the
image of the fibre group and is abelian.
This argument does not assume that the normal closure of
cusp vanishing cycles already exhausts $\ker\ell$.

In abelianization the monodromy relations give
$\Lambda_\Gamma$.
The primitive cusp angular lattice is the image of its
nilpotent operator, and this image is already zero among
those monodromy relations.
The remaining presentation is
\[
 5a=v_1,\qquad4b=v_2,\qquad a+b+v_0=0
 \quad\text{in }\Lambda_\Gamma\oplus\mathbb Za\oplus\mathbb Zb,
 \tag{F.T13}
\]
where $\ell(v_0)=\kappa$ fixes the sign convention.
Use any Smith presentation of $T$ and split the free part
$\mathbb Zh$.
The full relation matrix is block triangular, with the
presentation of $T$ as one diagonal block and the
three-by-three free matrix of determinant
$\pm(20\kappa+4r_1+5r_2)$ as the other.
Its determinant gives (F.T11); when the latter free
matrix has rank less than three the group has positive
rank and is infinite.
For determinant $\pm1$, elimination of the three free
variables leaves precisely $T$, even if the twists have
nonzero torsion components.
This proves the exact simple-connectivity criterion.

We give the global twist comparison to make its scope
precise.  Under (F.T12), the finite fixed covector image
is one at each point, so choose integral fixed $u_i$ with
$\ell(u_i)=1$.  The primitive vertical fixed vector
$\lambda_i$ and invariant covector $\varphi_i$ satisfy
$\varphi_i(\lambda_i)=m_i$, by the assumed maximal local
pairings.  Formula (F.S5) makes the finite norm onto
$\mathbb Z\lambda_i$.
Consequently every vertical numerator component
$s_i\lambda_i$ is removed by the explicit cap translation
(F.T8), whose boundary residue has $\ell$-coordinate zero.

For the complementary pair of pants the full relative
obstruction group for extending boundary translation
sections is
$H^2(P,\partial P;\Lambda)=\Lambda_\Gamma$.
Now $T=0$, so it is exactly $\mathbb Z$ detected by
$\ell$, not just a rational or free quotient.
Prescribing the identity translation at the cusp therefore
allows all the vertical finite changes to extend over $P$.
The same argument removes the zero-$\ell$ component of
any cusp translation clutching.
For determinant $-1$, write
$r_1=1+5a$, $r_2=-1+4b$; then (F.T12) says
$a+b+\kappa=0$.
The two integer changes of finite boundary marking and
the remaining cusp change thus also have zero total
relative obstruction and extend.
This gives the normalized twists $(u_1,-u_2)$.
For determinant $1$ the same argument gives their negatives.
The orientation-preserving fibre inversion extending over
the Kuhn fan, proved in F.2.24, identifies the two choices.
All maps used in the comparison are cap translations or
translations by smooth torus-bundle sections over $P$;
the prescribed cusp is retained throughout.
This proves the claimed complete twist normalization.
\end{proof}

\paragraph{Theorem F.2.29 (Proved: a complete finite integral survivor test).}
Take any rational pair in F.2.20 and any of its integral,
primitive-cusp rank-six extensions supplied by F.2.22.
Write $M_1,M_2$ for geometric transport on fibre homology
$\Lambda=\mathbb Z^6$, in the peripheral convention
$M_0M_1M_2=I$.
First require the preliminary conditions
\[
 \Lambda_\Gamma=\mathbb Z,\qquad
 \ell(\Lambda^{M_i})=\mathbb Z,\qquad h_1=5,\quad h_2=4 .
 \tag{F.HC1}
\]
Here $\ell$ is the primitive global invariant covector, and
$h_i$ is the integral fixed-vector/fixed-covector pairing
on $\ker\ell$ from F.2.25.
Failure of any condition excludes the datum.
If they hold, use the unique normalized affine filling of
F.2.28, with twists $(u_1,-u_2)$ and zero cusp clutching.
The following six finite integer-matrix tests are necessary
and sufficient for this filling to have $H_2=H_3=0$.

Let
\[
 T_i=M_i^{-t},\quad T_0=(T_1T_2)^{-1},\qquad
 G_q=\operatorname{coker}
 [\,\bigwedge^qT_1-I\mid\bigwedge^qT_2-I\,]
 \quad(1\le q\le4).
 \tag{F.HC2}
\]
All $G_q$ are finite.  Let $P_q$ be the integral parabolic
cohomology group of the exterior local system
$\bigwedge^q\mathcal V$, for $q=2,4$.
Let $t_q\in G_{q-1}$, $2\le q\le5$, be the natural global
coinvariant image of the unique nonzero torsion suspension
class at the order-four cap, as specified below.
Let $c_4\in G_3$ be the natural image of the unique nonzero
torsion class in
$\operatorname{coker}(\bigwedge^3T_0-I)$.
Then the six tests are
\[
\begin{array}{c|ccc}
 &\text{first test}&\text{second test}&\text{third test}\\ \hline
 H^3=0&G_1=\langle t_2\rangle&P_2=0&t_3\ne0\\
 H^5=0&G_3=\langle t_4,c_4\rangle&P_4=0&t_5\ne0 .
\end{array}
\tag{F.HC3}
\]
Every group, class, and equality here is computed by
finitely many integer kernels, Smith forms, and exterior
minors.  This is a decision criterion for every rational
parameter and every integral lattice in the complete
normal form; it makes no bounded-height assertion about
the rational parameter set.

\begin{proof}
We first establish the uniform local calculations needed
for the criterion, so the test does not assume that a
general finite matrix is globally conjugate to the example
in F.2.6.
The fixed vectors $u_i$ split the homology lattice as
$\mathbb Zu_i\oplus H_i$, where $H_i=\ker\ell$.
Write $D_i=M_i|_{H_i}$.
The maximal pairing $h_i=m_i$ gives
\[
 \operatorname{Norm}_{m_i}(D_i)=\lambda_i\varphi_i ,
 \tag{F.HC4}
\]
with both factors primitive.

At five, work over $\mathbb Z_5$.
Modulo five the norm is $(D_1-I)^4$ and is nonzero.
Choose $q$ with $\operatorname{Norm}_5q$ primitive.
The five vectors $q,D_1q,\ldots,D_1^4q$ are independent
modulo five: otherwise the nilpotent Jordan block generated
by $q$ would have length less than five.
They are therefore a $\mathbb Z_5$ basis.
Thus $D_1$ is the regular five-cycle permutation locally.

At four, work over $\mathbb Z_2$.
The norm is $(D_2-I)^3$ modulo two and is nonzero.
The vectors $q,D_2q,D_2^2q,D_2^3q$ similarly generate
a primitive regular four-cycle submodule $Q$ of rank four.
The quotient is rank one with $D_2=-1$, by the
characteristic polynomial.
For a lift $z$ of its generator, put $a=(D_2+I)z\in Q$.
Since $D_2^4=I$,
\[
 (D_2^3-D_2^2+D_2-I)a=0.
\]
On the regular four-cycle module the kernel of the
displayed alternating norm is exactly $(D_2+I)Q$,
integrally.  One can see this by writing its four
coefficients; equivalently $D_2+I$ has three nonzero
Smith factors, all one, and the indicated rational image.
Subtracting a suitable element of $Q$ from $z$ makes
$D_2z=-z$.
Consequently
\[
 M_1\sim 1\oplus\operatorname{perm}_5\quad\text{over }\mathbb Z_5,
 \qquad
 M_2\sim1\oplus\operatorname{perm}_4\oplus(-1)
 \quad\text{over }\mathbb Z_2 .
 \tag{F.HC5}
\]
These local comparisons apply to cap cohomology, not just
to abstract representation lattices.
Indeed the integral lattice of intertwiners is the kernel
of an integer matrix and its tensor with $\mathbb Z_p$
is the lattice of local intertwiners.
Approximate a local invertible intertwiner by an integral
one whose determinant is a $p$-adic unit.
Keeping the $u_i$ summand fixed produces an equivariant
torus isogeny of degree prime to $p$, preserving the
normalized translation $\pm u_i/m_i$.
The resulting cap and boundary coverings induce
isomorphisms on $\mathbb Z_p$ cohomology: the vertical
torus maps are cohomology isomorphisms over $\mathbb Z_p$,
so the mapping-torus Wang sequences give the assertion.
At primes not dividing $m_i$, transfer already makes
finite-cover pullback an isomorphism onto invariants.

Here are the full required integral results for these
two monomial models.
The ranks of their invariant exterior lattices in
degrees zero through six are
$(1,2,3,4,3,2,1)$.
The order-five cap has no cohomology torsion in degrees
two through five.
The order-four cap has exactly $\mathbb Z/2$ in each
of those four degrees.
This follows from the Wang sequence for the mapping
torus of the vertical $D_i$:
its torsion in degree $q$ is
\[
 \operatorname{Tor}\operatorname{coker}
       (\bigwedge^{q-1}D_i^t-I).
 \tag{F.HC6}
\]
In particular it introduces no additional torsion primes.

Integral pullback from either finite cap is onto the
invariant lattice in degrees two through five.
For degree two this is precisely the $m_i/h_i=1$
suspension calculation of F.2.25.
For completeness, the determinants of the pairings between
the invariant lattices in complementary degrees
$(1,5),(2,4),(3,3)$ are respectively
$m_i,m_i^3,m_i^4$ in absolute value.
The degree-one cover image has index $m_i$.
The free downstairs complementary pairings are unimodular,
and pullback multiplies each pairing by the covering
degree $m_i$.
The middle pairing first gives degree-three index one.
The degree-two index one then gives degree-four index one;
the degree-one index $m_i$ gives degree-five index one.
These determinant identities and all the Wang Smith
forms are verified by the fixed, exhaustive computation
in \texttt{fourfolds\_local\_caps.py}, for both orders and
all exterior degrees zero through six.
Thus no unexamined finite matrix is being inferred from
an example.

We next identify the boundary maps.
Choose any unimodular matrix $P_2=[u_2,H_2]$ realizing
the splitting at four, and put
$J=P_2^{-t}|_{\mathbb Z^5}$ for the inclusion of vertical
cohomology.
Choose the unique nonzero torsion representative
\[
 w_q\in\operatorname{coker}
            (\bigwedge^{q-1}D_2^t-I),\qquad
 t_q=[(\bigwedge^{q-1}J)w_q]\in G_{q-1}.
 \tag{F.HC7}
\]
A Smith transformation supplies an integral $w_q$.
Its norm is exactly zero, since it lies rationally in
the image of $\bigwedge^{q-1}D_2^t-I$.
The finite boundary coordinates of F.2.6 satisfy
$\theta=a+4k$, $h=k$, hence $a=\theta-4h$.
The pullback of the suspension cocycle has $\theta$
component $w_q$ and $h$ component
$-\operatorname{Norm}_4(w_q)=0$.
This holds in every exterior degree, including $q=5$.
It is an integral cellular assertion for the monomial
model, where the torus cellular differentials vanish;
the prime-to-two isogeny comparison transports it to
every eligible cap.
The base connecting map therefore gives exactly the
class $t_q$ in (F.HC7).
Different representative choices change it by a global
monodromy relation.  Any sign change is immaterial for
an order-two class.

There is one further boundary map, from the free
degree-four cusp kernel.
The primitive cusp is integrally conjugate to
$\left(\begin{smallmatrix}I&I\\0&I\end{smallmatrix}\right)$.
Its degree-three exterior difference has ten Smith
factors one and one Smith factor two; its cokernel is
$\mathbb Z^9\oplus\mathbb Z/2$.
The other exterior differences in degrees one, two,
four and five have torsion-free cokernels.
The exact matrices and a torsion representative are
saved by \texttt{fourfolds\_cusp\_boundary.py}.
Let $C$ be the smooth cusp cap and $Y=\partial C$.
The boundary Wang sequence gives
$\operatorname{Tor}H^4(Y)=\mathbb Z/2$.
In the Poincare--Lefschetz exact sequence its image in
\[
 H^5(C,Y)\cong H_3(C)=\mathbb Z^9
\]
is zero, so this torsion class lifts to $H^4(C)$.
Its restriction to an ordinary fibre is zero.
By E.1.3 the kernel of integral specialization is
a primitive copy $K=\mathbb Z$ in degree four.
The lift is in $K$, and the image of this cyclic group
in $H^4(Y)$ contains the nonzero torsion class.
A cyclic image containing nonzero torsion cannot
also have a nonzero free part.
Its primitive generator therefore restricts to the
unique nonzero torsion class of $H^4(Y)$.
Consequently the cusp contribution to $d_2$ is
exactly $c_4$ as defined above.  In particular this
column must not be discarded merely because it
vanishes in the first examples.

It remains to prove that these are all the tests.
For the irreducible odd-dimensional quotient $W$
of $\mathcal V$, invariant alternating forms and
bivectors vanish: a nonzero one would give an
isomorphism $W\cong W^*$ with an odd-dimensional
nondegenerate alternating pairing.
The exact exterior sequences for
$0\to\mathbb Q\to\mathcal V_{\mathbb Q}\to W\to0$,
together with $\det W=1$, give
\[
 H^0(B,j_*\bigwedge^q\mathcal V)=0
       \quad(2\le q\le5),\qquad
 G_q\otimes\mathbb Q=0\quad(1\le q\le4).
 \tag{F.HC8}
\]
For the endpoint assertions, the extension cannot split:
on $\mathbb Q\oplus W$ a square-zero nilpotent has rank at
most two, whereas the stipulated rank is three.
An invariant covector of $\mathcal V_{\mathbb Q}$ would
split this extension, since $W$ has no trivial quotient.
Thus $(\mathcal V_{\mathbb Q}^*)^\Gamma=0$, which gives
both $H^0(\bigwedge^5\mathcal V_{\mathbb Q})=0$ and
$G_1\otimes\mathbb Q=0$.
Here and below the first assertion is integral as well,
since the coefficient lattices are torsion free.
The local invariant ranks in exterior degrees two
and four are $3,3,9$ at the three punctures.
The Euler formula for the rank-fifteen local system
therefore gives $P_2\otimes\mathbb Q=P_4\otimes\mathbb Q=0$.
In degree one it similarly gives rational parabolic
cohomology zero.
The only global degree-one section is the invariant
line.  Simple connectivity and the universal
coefficient theorem make its free transgression
onto $H^2(B,\mathbb Z)$; equivalently the rational
second Betti number is zero, as in F.2.6.
Thus the normalized filling has rational $b_2=b_3=0$.

All specializations in exterior degrees two through
five are onto, by the local calculations and E.1.3.
Their only point-supported kernels are
\[
\begin{array}{c|cccc}
q&2&3&4&5\\ \hline
\text{kernel}&\mathbb Z/2&\mathbb Z/2&
                 \mathbb Z\oplus\mathbb Z/2&\mathbb Z/2 .
\end{array}
\tag{F.HC9}
\]
The free summand in degree four is at the cusp; the
others are at four.
Consequently $H^1(B,R^qg_*\mathbb Z)=P_q$ in the
degrees used here, while
$H^2(B,R^qg_*\mathbb Z)=G_q$.
The latter equality also holds in degree one:
changing a sheaf at isolated points changes no
degree-two cohomology.  The equality with coinvariants
follows from compactly supported top cohomology on
the punctured sphere and the sequence $j_!\to j_*$.

The base has real dimension two.
Thus the only possible nonzero differential in the
integral Leray spectral sequence is
$d_2:E_2^{0,q}\to E_2^{2,q-1}$.
The complete associated graded groups in total
degrees three and five are therefore
\[
\begin{split}
 \operatorname{gr}H^3&:
   G_1/\langle t_2\rangle,\quad P_2,\quad
   \ker(\mathbb Z/2\xrightarrow{\,t_3\,}G_2),\\
 \operatorname{gr}H^5&:
   G_3/\langle t_4,c_4\rangle,\quad P_4,\quad
   \ker(\mathbb Z/2\xrightarrow{\,t_5\,}G_4).
\end{split}
\tag{F.HC10}
\]
These are all finite.  A finite filtered group
vanishes exactly when every associated graded group
vanishes, regardless of extension classes.
Since $\pi_1=1$ and rational $b_2=b_3=0$, the universal
coefficient theorem gives
$H^3=\operatorname{Ext}(H_2,\mathbb Z)$, so
$H_2=0$ exactly when $H^3=0$.
Integral Poincare duality gives $H^5\cong H_3$.
This proves (F.HC3).

For reproducibility, no sheaf calculation is left
as a black box in the integer test.
Given $A=\bigwedge^qT_1$, $B=\bigwedge^qT_2$ on a
free lattice $L$, form the saturated integer kernel
\[
 K_q=\ker[\,A-I\mid A(B-I)\mid -(AB-I)\,]\subset L^3.
 \tag{F.HC11}
\]
Its image under
$(x,y,z)\mapsto((A-I)x,(B-I)y)$ is precisely the
integral parabolic cocycle lattice.
Quotient by the columns of
$\left(\begin{smallmatrix}A-I\\B-I\end{smallmatrix}\right)$
to obtain $P_q$.
This uses the actual integral images of peripheral
differences, not their saturations.
Smith transformations also supply every $w_q$, the
cusp order-two representative, and the tests for
their images and generated subgroups in $G_q$.

The program \texttt{fourfolds\_survivor\_test.py}
implements exactly these operations.
Its executed test set is precisely $v=0,1,2$ from
F.2.16.  For $v=0,1$ all six tests pass;
$G_1,\ldots,G_4$ are each $\mathbb Z/2$,
$P_2=P_4=0$, all four finite-cap images are nonzero,
and the cusp image is zero.
For $v=2$ they are each $\mathbb Z/6$,
$P_2=P_4=\mathbb Z/3$, and the first, second,
fourth and fifth tests fail.  The computation thus
checks both accepted and rejected data.

We finish by making the convention dictionary explicit.
The imported family matrices called $A,B$ in that
program are geometric cohomology pullbacks, so the
geometric homology inputs are $(A^t,B^t)$.
Their genuine left cohomology transport in (F.HC2)
is $(A^{-1},B^{-1})$, which is what the code uses.
For comparison with computations written using
$(A,B)$, simultaneous reversal of the two peripheral
generators sends
\[
 g_1\mapsto g_1^{-1},\qquad
 g_2\mapsto g_2^{-1},\qquad
 g_0=(g_1g_2)^{-1}\mapsto
     g_2g_1=g_1^{-1}(g_1g_2)g_1 .
 \tag{F.HC12}
\]
This is an orientation-reversing marking of the
pair of pants fixing its three ends: each peripheral
conjugacy class is inverted.
Thus the new cusp is conjugate to the inverse old
cusp, not an independently chosen third matrix.
Inversion preserves invariant and coinvariant
lattices, and the indicated conjugating generator
acts trivially on global coinvariants.
All boundary classes in (F.HC3) have order two, so
reversal signs do not change the tests.
The implementation uses (F.HC2) directly and hence
does not depend on this comparison for its validity.
\end{proof}

\paragraph{Corollary F.2.30 (Proved: constructive classification of all survivors).}
For the marked Kuhn cap and affine translation clutchings
used throughout F.2, the following is a complete constructive
classification of the data with
$\pi_1=1$ and $H_2=H_3=0$, up to integral conjugacy of their
linear tuples and oriented diffeomorphism of their fillings.

There are no data with three genuine finite fillings.
With two fillings their actual affine orders are five and
four.  Order the order-five point first.
The complete, unbounded parameter set is the rational
projective normal form of F.2.20 together with its
exceptional rational affine line.
For each parameter, run the terminating order,
conductor, stable-lattice, and primitive-extension algorithm
of F.2.22.  Discard precisely the integral classes failing
(F.HC1) or one of (F.HC3).
The classes retained by this procedure are exactly all
surviving integral linear tuples.
For each retained class all simply connected affine
twists are exactly the choices satisfying (F.T12);
their fillings are all oriented diffeomorphic to its
normalized filling.
If ordered labels with the order-four point first are
desired, interchange the finite labels and recompute
$M_0=(M_1M_2)^{-1}$.

Every retained datum has
\[
 \pi_1=1,\qquad
 H_j=\begin{cases}
  \mathbb Z,&j=0,8,\\
  \mathbb Z^4,&j=4,\\
  0,&\text{otherwise},
 \end{cases}
 \tag{F.HC13}
\]
and its oriented smooth manifold is
$3\mathbb{HP}^2\#(-\mathbb{HP}^2)$.
It has bounded-cusp holomorphic periods in the literal
sense of F.3, by F.3.11.
The classification here concerns the survivor subproblem.
It does not classify the nonsurviving tuples or compute
their full integral homology, as also requested in F.2.

\begin{proof}
The exclusions F.2.7, F.2.14, F.2.25, and F.2.27
leave exactly the stated minimal orders and rational
quotient problem.
The rational normal form is injective on conjugacy
classes, including the exceptional line.
The finite interval and extension-coset enumeration
in F.2.22 gives every integral class over each rational
point, with all bounds proved for that input.
No distinct rational points can become integrally
conjugate, since an integral conjugacy induces a
rational quotient conjugacy.
Theorem F.2.28 and the necessary-and-sufficient
criterion F.2.29 prove both directions of the
retention rule and account for every affine twist.
This is also an exhaustive enumeration if desired:
enumerate primitive projective integer triples with
first nonzero coordinate positive by increasing height,
and enumerate reduced rational exceptional parameters
by increasing numerator and denominator height.
Every finite stage has finitely many inputs, each
terminates by F.2.22, and every rational input occurs.
No finite cutoff is imposed or needed.

Hurewicz makes the retained simply connected manifolds
three-connected.  The only singular contribution to
Euler number is the Kuhn cusp, whose value is six
by E.1.3; torus bundles and the free finite caps have
Euler number zero.
Integral Poincare duality and the universal coefficient
theorem then give (F.HC13), including torsion freeness.

For the smooth classification, F.2.15 supplies a
compatible complex structure at each finite cap, and
the cusp already has its complex model.
Choose their common fibre orientation to be the one
induced by the cusp.
Commuting complex structures at the finite caps can
be chosen in that component: their two-dimensional
fixed real summands allow the orientation to be
reversed when necessary.
The fibre of the bundle of orientation-compatible
complex structures retracts to $\mathbb{CP}^3$.
Its connectedness and simple connectivity extend
the prescribed boundary structures across the
two-dimensional pair of pants, exactly as in F.2.5.
A compatible horizontal distribution gives an
almost complex structure on the total manifold.
Theorem F.1.2 therefore gives signature two,
intersection lattice $I_{3,1}$, and $c_2^2=2$.
Since $H^2=0$, $p_1/2=-c_2$ is primitive,
characteristic, and has square two.
The complete Wall--Nikulin--Nagy argument in F.4.1
depends only on these almost complex characteristic
numbers and three-connectivity.
It applies without alteration and gives the claimed
exact oriented diffeomorphism.
Finally F.3.11 applies to the resulting determinant-one
monodromy and primitive cusp.
It asserts the requested punctured-base periods with
bounded correction, not an additional unproved
simultaneous good extension at all finite points.
\end{proof}

\paragraph{Theorem F.2.31 (Proved: complete fundamental-group data for every input).}
Let $k=2$ or $3$, and let an integral rank-six tuple satisfy the
finite-order and primitive maximal-cusp conditions of F.2.  No
irreducibility, determinant-one, or prospective-survivor hypothesis is
imposed.  Specify each actual finite affine action as
\[
 \tau_i(x)=M_ix+v_i/m_i,\qquad v_i\in\Lambda=\mathbb Z^6,
 \quad \tau_i^{m_i}=1\text{ on }\Lambda_{\mathbb R}/\Lambda,
 \tag{F.G1}
\]
and require its cyclic action to be free.  The numerator need not be
fixed by $M_i$.  Define its actual integral power vector and the cusp
lattice by
\[
 w_i=\frac1{m_i}\sum_{j=0}^{m_i-1}M_i^jv_i\in\Lambda^{M_i},
 \qquad U=\operatorname{im}\bigl((M_1\cdots M_k)^{-1}-I\bigr).
 \tag{F.G2}
\]
Use left conjugation transport for the peripheral generators $s_i$,
and record the integral translation clutching by $\kappa\in\Lambda$
in $s_1\cdots s_kx^\kappa=1$.  Then the full fundamental group is
\[
 \left\langle x^\lambda,s_1,\ldots,s_k\ \left|\
 \begin{array}{l}
 x^\lambda x^\mu=x^{\lambda+\mu},\quad x^u=1\ (u\in U),\\
 s_ix^\lambda s_i^{-1}=x^{M_i\lambda},\quad
 s_i^{m_i}=x^{w_i},\quad s_1\cdots s_kx^\kappa=1
 \end{array}\right.\right\rangle .                         \tag{F.G3}
\]
Only a basis of $\Lambda$ and of $U$ is needed in this finite
presentation.  The program below also determines exactly which
quotient of the apparent fibre subgroup injects.

Put
\[
 L=\sum_{g\in\langle M_1,\ldots,M_k\rangle}gU,
 \qquad A=\Lambda/L.                                      \tag{F.G4}
\]
This is an integral sum, without saturation.  For two points, put
$d=\gcd(m_1,m_2)$, $M=M_1$, and
$S_j(M)=I+M+\cdots+M^{j-1}$.  Choose integers $a,b$ with
$am_1+bm_2=d$, and set, in $A$,
\[
 \begin{split}
 w'_2&=-w_2-S_{m_2}(M)\kappa,\\
 e&=(m_2/d)w_1-(m_1/d)w'_2,\qquad
 z=aw_1+bw'_2,\qquad \overline A=A/\langle e\rangle .
 \end{split}                                             \tag{F.G5}
\]
There is an exact sequence and a complete marked extension description
\[
 1\longrightarrow\overline A\longrightarrow\pi_1
   \longrightarrow C_d\longrightarrow1,
 \qquad txt^{-1}=M(x),\quad t^d=z.                         \tag{F.G6}
\]
Thus its order is $d|\overline A|$ if $\overline A$ is finite;
otherwise it is infinite.  In particular $d=1$ gives precisely
the abelian group $\overline A$.

For three genuine points, so $m_i\ge2$, let
\[
 \Gamma=\langle q_1,q_2,q_3\mid q_i^{m_i}=1,
                         \ q_1q_2q_3=1\rangle,
 \qquad \epsilon=\sum_i1/m_i-1.                           \tag{F.G7}
\]
The action $q_i\mapsto M_i$ is well defined on $A$.
If $\epsilon\le0$, the exact sequence is
\[
                    1\longrightarrow A\longrightarrow\pi_1
                         \longrightarrow\Gamma\longrightarrow1. 
 \tag{F.G8}
\]
If $\epsilon>0$, $\Gamma$ is the finite spherical triangle group.
Using its action on $A$, form the exact transfer vector
\[
 e=\sum_{g\in\Gamma}g\kappa+
       \sum_{i=1}^3\ \sum_{g\in\Gamma/\langle q_i\rangle}g w_i.
 \tag{F.G9}
\]
The second sums use left coset representatives and are independent
of their choice.  In this case the exact sequence has kernel
$A/\langle e\rangle$, with the same lifted relations (F.G3).
Consequently (F.G3), (F.G6), (F.G8), and (F.G9) give every requested
fundamental group by finite operations, including the precise
injected abelian kernel and all extension data.  These statements
do not claim a conjugacy classification of the original tuples.

For every input, the full first integral homology is the Smith
cokernel of the integer matrix
\[
 \begin{pmatrix}
 M_1-I&\cdots&M_k-I&-w_1&\cdots&-w_k&\kappa\\
 0&\cdots&0&&\operatorname{diag}(m_1,\ldots,m_k)&&\boldsymbol1
 \end{pmatrix}.                                           \tag{F.G10}
\]
Here the bottom left zero block has $k$ rows.  The code records
the full matrix, not only its invariant factors.

\begin{proof}
\emph{Actual finite and cusp caps.}
The affine lift on $\mathbb R^6$ has $m_i$th power translation by
$w_i$, as the displayed sum shows.  It commutes with that power,
so $M_iw_i=w_i$.  The group generated by lattice translations
and this lift is the fundamental group of the reduced flat
six-manifold.  Its presentation consists of the lattice, its
$s_i$ conjugation relations, and $s_i^{m_i}=w_i$.
The finite cap deformation retracts onto that reduced manifold
by the equivariant radial contraction of its root disc.
Thus these are the actual cap relations, even for a numerator
which is not an invariant lattice vector.

For the primitive Kuhn cap, the compact toric strata in its
normalization and all their intersection components are simply
connected.  Applying van Kampen successively to their tubular
neighborhoods identifies the fundamental group of the central
fibre with that of its dual three-torus.  The nearby fibre map
kills the angular lattice $U$ and sends the long-period lattice
$\Lambda/U$ isomorphically onto this dual-torus group.  This is
also the explicit toric collapse used in E.1.3.  The primitive
index-one condition gives $U=\ker(M_0-I)$ and makes $\Lambda/U$
free of rank three.  A section through a smooth central stratum
bounds the base meridian by its parameter disc.  The smooth cap
retracts onto the central fibre, so it adds exactly these
relations, with no extra meridian relation left unspecified.

Over the punctured sphere the bundle has a section up to homotopy,
because the base retracts to a graph.  Its group is the appropriate
semidirect product.  Choose the cap sections and retain their
relative translation clutching as $\kappa$.  Van Kampen with the
actual boundary maps just described gives exactly (F.G3).
Changing the transport convention causes no group ambiguity:
reverse the ordered list and replace each affine map by its
inverse, and replace $\kappa$ by $-\kappa$.  The new linear
matrices are $M_k^{-1},\ldots,M_1^{-1}$, the power vectors are
$-w_k,\ldots,-w_1$, and the cusp image lattice is unchanged.
The substitutions $s_i\mapsto s_{k+1-i}^{-1}$ give the explicit
isomorphism between these two presentations.  This supplies the
dictionary when geometric transport is recorded with opposite
path-composition order.

\emph{Integral closure and two-point elimination.}
In a normal abelian fibre subgroup, the normal closure of $U$
is precisely the integer sum (F.G4).  Start with its column
Hermite basis and repeatedly append its images under every
$M_i$ and $M_i^{-1}$, then take Hermite normal form again.
If the rational rank increases, it can do so at most three
times.  At the first nonincrease the rational span is invariant
under every generator and its inverse, hence never increases
again.  If the current lattice has saturation index $D$ in
that span, every subsequent strict enlargement divides this
index by an integer at least two.  There can therefore be at
most $\lfloor\log_2D\rfloor$ such enlargements.  This proves
termination and gives the computable bound used by the code;
it never replaces a finite-index lattice by its saturation.

On $A$ one has $M_2=M^{-1}$ in the two-point case.  The product
relation eliminates $s_2=t^{-1}x^{-\kappa}$, where $t=s_1$.
Multiplying this word $m_2$ times gives
\[
        (t^{-1}x^{-\kappa})^{m_2}
            =t^{-m_2}x^{-S_{m_2}(M)\kappa}.
\]
Thus the two power relations become $t^{m_1}=w_1$ and
$t^{m_2}=w'_2$.  Both right sides are fixed by $M$ on $A$.
Also $M^{m_1}=M^{m_2}=1$, hence $M^d=1$ there.
Bezout elimination gives exactly $e=0$ and $t^d=z$ in
(F.G5).  Replacing the Bezout pair changes $z$ by a multiple
of $e$ and therefore changes no extension datum.

To prove the asserted injection, construct
$\overline A\rtimes_M\mathbb Z$ and quotient by its central
element $t^dz^{-1}$.  This cyclic subgroup has nonzero
projection $d\mathbb Z$ to the infinite cyclic factor, so it
intersects $\overline A$ trivially.  The quotient has precisely
the presentation obtained by elimination.  This proves (F.G6),
including the absence of any further hidden fibre relation.

\emph{The three-point kernel and the finite transfer.}
For completeness this step computes the fibre injection; it
does not infer it from the appearance of a normal subgroup in
a presentation.  Let $\mathcal B$ be the orbifold classifying
space obtained from the three-holed sphere by attaching, at
its three boundary circles, the mapping cylinders of
$S^1\to BC_{m_i}$.  Over the punctured sphere take the bundle
with fibre $BA=K(A,1)$ and the specified action.  For each
boundary, the cyclic extension
\[
 E_i=\langle A,s_i\mid s_i a s_i^{-1}=M_i(a),\ s_i^{m_i}=w_i\rangle
\]
contains $A$ injectively: the same semidirect-product construction
as above applies because $M_i^{m_i}=1$ and $M_iw_i=w_i$.
The associated fibration $BA\to BE_i\to BC_{m_i}$ restricts
on $S^1$ to the required mapping-torus fibration.  Replace the
maps by their mapping-path fibrations before homotopy gluing,
so that the fibre is preserved on every attaching cylinder.
Glue these three fibrations, using the specified relative clutching
$\kappa$ on the punctured sphere.  Van Kampen gives (F.G3)
for the resulting total space.  Thus the homotopy sequence
of this fibration over $\mathcal B$ calculates precisely
the image of $A$ in the group under discussion.

When $\epsilon\le0$, the universal orbifold cover is the plane
or hyperbolic plane.  Its contractibility identifies
$\mathcal B$ with $B\Gamma$: indeed its Borel construction
has the just-described boundary-cylinder model and contractible
universal fibre.  The Euclidean and hyperbolic alternatives
are the oriented two-orbifold classification in
\cite[Theorems 2.3--2.4]{rootScott83}.
Since $\pi_2(B\Gamma)=0$, the homotopy sequence proves (F.G8).

When $\epsilon>0$, the universal orbifold cover is $S^2$.
Pulling the constructed $BA$ fibration to this finite cover
trivializes its local system.  Its transgression
$\pi_2(S^2)\to A$ is its integral clutching vector.
There is one lifted ordinary base cell for each $g\in\Gamma$,
contributing $g\kappa$, and one lifted cone cap for each left
coset of $\langle q_i\rangle$, contributing $gw_i$.
To check the sign, the base relation has fibre term $+\kappa$
and each cone relation has term $-w_i$; the oriented sphere
chain subtracts the cone cells from the base cell.  The result
is therefore the plus sign in both sums of (F.G9).
Changing a coset representative by $q_i^j$ does not change
its contribution because $w_i$ is fixed.  The result is
$\Gamma$ invariant, as also checked by the code.  The long
homotopy sequence now says that the only additional kernel
is exactly $\langle e\rangle$.  This proves the spherical
assertion, with the orientation fixed by this cell calculation.

There is no implicit infinite search for a spherical group.
After sorting its orders the possibilities are
$(2,2,n),(2,3,3),(2,3,4),(2,3,5)$, with groups of orders
$2n,12,24,60$, respectively.  These follow directly from
the positive spherical-triangle angle condition.
The code constructs the dihedral, tetrahedral, octahedral,
or icosahedral permutation group and finds the marked
peripheral generators in that finite group.  Breadth-first
words supply their actions on $A$; every generator edge is
checked modulo the actual lattice $L$.  It then computes
the two finite sums in (F.G9) and their integer Smith form.

An order-one site is not a genuine cone.  Its matrix is the
identity and its relation is $s_i=x^{w_i}$.  Eliminate it
from (F.G3); moving this vector to the right of the remaining
peripheral product replaces the clutching by
\[
       \kappa+(M_{i+1}\cdots M_k)^{-1}w_i.                 \tag{F.G11}
\]
This reduces the three-point case with such a site to the
two-point computation.  A primitive rank-three cusp cannot
remain with only one finite linear monodromy: a finite-order
unipotent matrix is the identity.  Thus no smaller case has
been omitted within the hypotheses of this theorem.

\emph{Freeness, abelianization, and executed certificates.}
Freeness of (F.G1) is itself decidable exactly.  A nontrivial
stabilizer in a cyclic group contains a subgroup of prime
order, so it is enough to check $\tau_i^{m_i/p}$ for every
prime $p\mid m_i$.  For an affine map $Rx+b$ of a torus,
a fixed point exists precisely when every integral covector
in $\ker(R^t-I)$ takes an integral value on $b$.
Smith form proves this criterion: the image of $R-I$ on
the torus is the connected subtorus annihilated by those
primitive covectors.  The code computes the powers in the
rational affine matrix group and performs these finite tests.

Abelianizing (F.G3) yields (F.G10).  No cusp column was lost:
\[
 \operatorname{im}(M_0-I)
       \subset\sum_i\operatorname{im}(M_i-I)
\]
follows by telescoping the product on the right; inversion
preserves the image of its difference from the identity.
The Smith form therefore computes every free and torsion
summand of $H_1$ exactly.

The program \texttt{fourfolds\_vankampen.py} saves all input
matrices, power vectors, integral closure traces, group
extension data, and homology presentation matrices.
It checks the nine explicitly listed three-point parameters
of its original test set, four two-point order/clutching
choices, a genuinely nonfixed affine numerator, and a
finite-index closure step.  It also checks a nonabelian
rank-six two-point tuple, a spherical $(2,4,2)$ example
with injected kernel $\mathbb Z^2\oplus\mathbb Z/4$, and
the infinite perfect hyperbolic example with orders
$(5,4,19)$.  The nonabelian test deliberately belongs to
the full $\operatorname{GL}_6$ scope and has orientation-
reversing finite matrices; it is not claimed to admit a
complex compactification.  All bounds above are derived
from each input, rather than imposed as a search cutoff.
\end{proof}

\paragraph{Theorem F.2.32 (Proved: a complete three-point logarithmic family).}
Start with the compact complex fourfold of F.3.7 and its
normalized order-five and order-four affine fillings.
Let $\ell$ be the primitive invariant covector on fibre
homology, with the sign chosen in F.2.28.
At any additional ordinary fibre choose
$v\in\mathbb Z^6$ and an integer $m\ge2$ such that
\[
 r=\ell(v),\qquad \gcd(r,m)=1.
\tag{F.LT1}
\]
Perform the marked logarithmic transformation with
translation $v/m$ described below.
This gives a smooth compact complex fourfold $X_{m,v}$
with Euler number six and three genuine finite affine
fillings, whose linear orders are five, four and one
and whose actual affine orders are five, four and $m$.
Its fundamental group is
\[
 \pi_1(X_{m,v})=
 \langle a,b,c,h\mid h\text{ central},\
   a^5=h,\ b^4=h^{-1},\ c^m=h^r,\ abc=1\rangle .
\tag{F.LT2}
\]
This is the group of the oriented Seifert manifold
with invariants
$(0;(5,-1),(4,1),(m,-r))$, in the convention
$x_i^{\alpha_i}h^{\beta_i}=1$ and product $x_i=1$.
Its base orbifold is hyperbolic, and its central fibre
$h$ has infinite order.

Put $C_{m,r}=\mathbb Z/(20r-m)\mathbb Z$, with
$\mathbb Z/0\mathbb Z=\mathbb Z$, and let
$\varepsilon=1$ when $20r=m$ and zero otherwise.
All the integral homology groups are
\[
\begin{array}{c|c}
j&H_j(X_{m,v};\mathbb Z)\\ \hline
0,8&\mathbb Z\\
1,6&C_{m,r}\\
2&\mathbb Z^\varepsilon\oplus(\mathbb Z/m)^5\\
3&(\mathbb Z/m)^{10}\\
4&\mathbb Z^4\oplus(\mathbb Z/m)^{10}\\
5&(\mathbb Z/m)^5\\
7&\mathbb Z^\varepsilon .
\end{array}
\tag{F.LT3}
\]
Under (F.LT1) the exceptional free-rank case is exactly
$(m,r)=(20,1)$.
In particular $(m,r)=(19,1)$ and $(21,1)$ have
$H_1=0$ but nontrivial, infinite fundamental groups.
More generally this occurs whenever $|20r-m|=1$.
The result treats every vertical component of $v$
in $\ker\ell$, not only a chosen fixed representative.

\begin{proof}
We first specify the complex construction and its marking.
Over a small disc at the ordinary value use a holomorphic
period marking
$T_t=\mathbb C^3/\Lambda_t$.
Let $\omega_v(t)$ be the holomorphic period vector
corresponding to $v$.
On the root disc $t=s^m$ quotient the pulled-back family by
\[
 (s,Z)\longmapsto
   (\zeta_m s,\ Z+\omega_v(s^m)/m).
\tag{F.LT4}
\]
Condition (F.LT1) makes $v$ have exact order $m$
in $\mathbb Z^6/m\mathbb Z^6$.
Every nontrivial power therefore acts freely, including
over $s=0$.
The quotient is smooth, has a proper map to the disc,
and has central fibre of multiplicity $m$.
On the punctured disc its holomorphic identification with
the original family is
\[
 (s,Z)\longmapsto
 \left(s^m,\ Z-\frac{\log s}{2\pi i}\omega_v(s^m)\right).
\tag{F.LT5}
\]
Changing the branch of $\log s$ changes the second
coordinate by a lattice period.
Changing $s$ by $\zeta_m s$ changes its logarithm by
$2\pi i/m$ and cancels the translation in (F.LT4).
Thus (F.LT5) is a well-defined holomorphic isomorphism
of the quotient with the original torus family off zero.
Gluing on a fixed annulus supplies the asserted compact
complex manifold: its charts are the original charts
and the free finite quotient charts (F.LT4), and it is
the union of compact pieces over closed base regions.
The gluing is an isomorphism on collars, so the total
space is Hausdorff and smooth.
The Euler number is unchanged, since the removed
ordinary torus cap, the new cap, and their boundary
all have Euler number zero.
This marking fixes the new relation $c^m=v$ and
zero extra cusp clutching.

In the original datum the invariant normal closure
of the cusp angular lattice is exactly $\ker\ell$,
as proved in F.2.5.
Adding the identity linear generator does not change
that lattice.
The fibre image in van Kampen therefore reduces to
the central generator $h$, while the original two
finite relations remain $a^5=h$, $b^4=h^{-1}$.
The new relation is $c^m=h^{\ell(v)}$ and the base
relation is $abc=1$.
This proves (F.LT2), including the absence of any
dependence on the vertical part of $v$.

One may verify the Seifert identification directly:
take a product circle bundle on a three-holed sphere
and attach its three solid tori with primitive meridian
slopes $(5,-1),(4,1),(m,-r)$.
They are primitive by (F.LT1), and van Kampen gives
exactly (F.LT2).
Its orbifold base has
$1/5+1/4+1/m<1$ for every $m\ge2$.
Pulling back the circle fibration to its contractible
hyperbolic orbifold cover gives a circle bundle on
the plane, hence the trivial circle bundle.
The subsequent universal cover is the plane times
$\mathbb R$, and the fibre subgroup is infinite cyclic.
This also proves directly that none of these groups
is trivial.
The rational Euler number in these conventions is
$r/m-1/20$.

In abelianization the relations $5a=h$, $4b=-h$
give $a=4x$, $b=-5x$, $h=20x$ because five and four
are coprime.
The product relation gives $c=x$.
The remaining relation is $(m-20r)x=0$.
Thus $H_1=C_{m,r}$ exactly, including its integral
generator and the exceptional infinite-cyclic case.

We now compute the higher integral groups, including
all extensions.
Write $\mathcal R^q$ for the higher direct-image
sheaf of the original fourfold and
$\mathcal R_m^q$ for that of $X_{m,v}$, after identifying
the bases and their ordinary local systems.
At the new point the reduced finite quotient is a
torus with lattice
$\mathbb Z^6+\mathbb Z(v/m)$.
Its covering by the original torus is cyclic of
degree $m$.
The degree-one pullback has Smith factors
$1,1,1,1,1,m$, so its degree-$q$ cokernel is
\[
 C^q=(\mathbb Z/m)^{\binom5{q-1}},
 \qquad 1\le q\le6 .
\tag{F.LT6}
\]
There is no cohomology torsion in this torus cap.
The sheaf comparison, which does not assert a map
between the total manifolds, is therefore
\[
 0\longrightarrow\mathcal R_m^q
  \longrightarrow\mathcal R^q
  \longrightarrow C^q_{\{p\}}\longrightarrow0 .
\tag{F.LT7}
\]
It is the identity away from the new point and the
covering image at that point.

For $2\le q\le5$, all original global sections of
$\mathcal R^q$ are supported at the original finite
point or the cusp, by (F.HC8)--(F.HC9).
Their map to $C^q_{\{p\}}$ is zero.
Consequently their $H^0$ groups are unchanged,
their $H^2$ groups are unchanged, and
\[
 0\longrightarrow C^q
  \longrightarrow H^1(\mathcal R_m^q)
  \longrightarrow H^1(\mathcal R^q)
  \longrightarrow0 .
\tag{F.LT8}
\]
The original degree-one base cohomology of these
sheaves is zero for $q=2,4,5$ and is $\mathbb Z^3$
for $q=3$.
The assertion for $q=5$ also follows immediately
from the original $H^6=0$: its bidegree $(1,5)$
has no incoming or outgoing differential.
Thus (F.LT8) gives $C^q$ in degrees two, four and
five, and $\mathbb Z^3\oplus C^3$ in degree three.

The original four order-two torsion transgressions
remain isomorphisms.
Indeed their boundary suspension representatives and
their images in global coinvariants are the ones in
F.2.29; adding the identity monodromy changes none
of those coinvariant quotients.
These point-supported boundary maps do not depend
on the free degree-one transgression, which changes
under the new logarithmic gluing.
The free degree-four cusp column likewise remains
zero, since its boundary torsion representative has
zero image in that same unchanged global quotient.
It follows that the full filtrations in degrees
three and five have just one nonzero term:
\[
 H^3=C^2=(\mathbb Z/m)^5,\qquad
 H^5=C^4=(\mathbb Z/m)^{10}.
\tag{F.LT9}
\]
In degree four there is an exact sequence
\[
 0\longrightarrow
       \mathbb Z^3\oplus(\mathbb Z/m)^{10}
  \longrightarrow H^4
  \longrightarrow\mathbb Z\longrightarrow0.
\]
It splits as an abelian-group extension because
its quotient is free, giving the stated middle group.
Thus no torsion extension has been inferred from
rational ranks.

For the remaining free ranks,
$H^1(X_{m,v})=\operatorname{Hom}(C_{m,r},\mathbb Z)$
has rank $\varepsilon$.
The degree-one global section group has rank one,
and the only rational differential involving it is
$d_2:H^0(\mathcal R_m^1)\to H^2(\mathcal R_m^0)$.
Its kernel has rank $\varepsilon$, hence its
cokernel also has rank $\varepsilon$.
The degree-one parabolic group is rationally zero,
and there are no rational degree-two global sections.
Therefore $b_2=\varepsilon$.
For an integral check of this degree-one row, its
finite sheaf cokernel is
$\mathbb Z/\gcd(20,m)$: the original invariant
section restricts as $20\ell$, and its value at
the new covering cokernel is $20r$ modulo $m$.
This finite group changes no rational rank.

Poincare duality gives
$H_3=H^5$, $H_4=H^4$, and $H_5=H^3$.
The universal coefficient theorem now identifies
the torsion of $H_2$ with $H^3$ and its free rank
with $\varepsilon$.
It also gives $H^2=C_{m,r}$: when $20r\ne m$ this
is $\operatorname{Ext}(H_1,\mathbb Z)$; when
$20r=m$ it is the free rank-one group.
Duality gives $H_6=C_{m,r}$ and
$H_7=\mathbb Z^\varepsilon$.
This proves the complete table (F.LT3).

The program \texttt{fourfolds\_three\_point\_family.py}
checks the original four torsion images and zero
cusp image once, then computes the actual cyclic
covering Smith forms in all exterior degrees one
through six and the exact van Kampen abelianization
for the thirty specified pairs
$m\in\{2,3,4,5,6,8,10,19,20,21\}$,
$r\in\{-1,1,m+1\}$.
It includes a nonzero vertical component in every
test vector.
The unbounded assertion follows from (F.LT6)--(F.LT9),
not from extrapolation of those finite tests.
\end{proof}

\subsection*{Effective integral chains for the local filling pairs}

\paragraph{Theorem F.2.33 (Proved: arbitrary finite affine caps).}
Let $g(x)=Mx+v/m$ be a free cyclic affine action of actual order
$m\geq2$ on $T^r$, with $M\in\mathrm{GL}(r,\mathbb Z)$ and
$v\in\mathbb Z^r$.  For every specified degree bound $d$, there is
an explicit terminating integer algorithm producing chains and
boundary maps through degree $d+1$ for the pair
\[
 B=(D^2\times T^r)/\langle(s,x)\mapsto
 (e^{2\pi i/m}s,g(x))\rangle,
 \qquad Y=\partial B.                                      \tag{F.CA1}
\]
The cap complex has at most $2^r$ generators in each degree and
the boundary complex at most $2^{r+1}$.  These are complexes of
finite rank in each degree; only their indicated finite prefixes
are required to compute $H_j(B)$, $H_j(Y)$ and $H_j(B,Y)$ for
$j\leq d$.  Their maps retain the affine translations and all
integral higher coherence terms.  In particular, the construction
does not replace an equivariant torus chain complex by its
homology without recording the correction terms.

\emph{Proof.}
We give the complete recursive algorithm.  All chains below are
finite integer dictionaries, so the displayed operations are
effective over $\mathbb Z$.

First consider the normalized bar resolution of $\mathbb Z^r$
over $R=\mathbb Z[t_1^{\pm1},\ldots,t_r^{\pm1}]$ and its Koszul
resolution $K=R\otimes\bigwedge^*\mathbb Z^r$.
The latter differential is contraction with $(t_i-1)_i$.
Write $t^a=t_1^{a_1}\cdots t_r^{a_r}$ and put
$S_n(t)=(t^n-1)/(t-1)$, interpreted as its finite integral
Laurent sum for positive, zero or negative $n$.
The following $\mathbb Z$-linear contraction of the augmented
Koszul complex fixes a completely explicit choice:
\[
 s_K(t^a e_I)=
 \sum_{j<\min I}
 t_{j+1}^{a_{j+1}}\cdots t_r^{a_r}S_{a_j}(t_j)
 e_j\wedge e_I,                                           \tag{F.CA2}
\]
where $\min\varnothing=r+1$ and indices in $I$ are increasing.
The bar contraction is
$s_B(t^a[g_1|\cdots|g_q])=[a|g_1|\cdots|g_q]$; terms with a
zero bar entry are zero.  Set $p_0=1$ and define the
$R$-linear projection and homotopy on coefficient-one bar
generators recursively by
\[
 p_q=s_Kp_{q-1}d_B,
 \qquad
 h_q=s_B(1-ip_q-h_{q-1}d_B).                               \tag{F.CA3}
\]
Here $i(e_I)$ is the alternating sum of the bars of the ordered
unit vectors in $I$.  The recursion decreases degree.  The
identities for the two contractions give, by induction,
$pi=1$, $dh+hd=1-ip$, $ph=hi=h^2=0$.
For example $s_K^2=s_B^2=0$, $ps_B=s_Kp$, and the alternating
unit-vector formula gives $pi=1$; substituting these into
(F.CA3) proves the three side identities in the same induction.
Augmenting $R\to\mathbb Z$ gives a strong chain contraction
of the ordinary torus bar complex to $\bigwedge^*\mathbb Z^r$.

Translations must enter before this contraction.  Let
$b_j=\sum_{a=0}^{j-1}M^av/m$ and let $p_j$ be its coordinatewise
fractional part.  The $m$ points $p_j$ are distinct because the
action is free.  Use the full fundamental groupoid of $T^r$
on these points: an arrow from object $a$ to object $b$ is an
integer vector $\lambda$, represented by the affine path from
$p_a$ to $p_b+\lambda$.  Composition adds the integer vectors.
The generator acts strictly by shifting objects and, for its
$s$th power, by
\[
 \lambda\longmapsto M^s\lambda+z_s(b)-z_s(a),
 \qquad z_s(a)=M^sp_a+b_s-p_{a+s}\in\mathbb Z^r.             \tag{F.CA4}
\]
The subscripts on objects are taken modulo $m$.
The nerve of this groupoid is equivariantly homotopy equivalent
to the torus: the simplex with successive arrow labels
$\lambda_i$ maps affinely to the torus with lifted vertices
$p_{a_j}+\sum_{i\leq j}\lambda_i$.
Forgetting the objects contracts its chains to the ordinary
torus bar complex.  An explicit homotopy is the negative
alternating prism which, at each possible position, inserts
the zero-labelled arrow to object zero and replaces all
subsequent objects by zero.  Combining this with (F.CA3)
gives a contraction $(i,p,h)$ of the groupoid chains to
$\bigwedge^*\mathbb Z^r$ with the same side identities.

Use the normalized cyclic bar construction on this strict
groupoid action.  Its differential is $d_v+\delta$, where
$\delta$ is the ordinary group-bar boundary, multiplied by
$(-1)^q$ on vertical degree $q$; its first face applies (F.CA4).
The vertical contraction transfers this to cyclic bar chains
with exterior coefficients by the finite formulas
\[
 I=(1+h\delta)^{-1}i,\qquad
 P=p(1+\delta h)^{-1},\qquad D=p\delta I.                  \tag{F.CA5}
\]
Inverses mean their alternating geometric series.  On a chain
of cyclic bar degree $p$ each series has at most $p+1$ terms,
since $\delta$ lowers that degree.  Multiplication and the
contraction identities directly give $(d_v+\delta)I=ID$ and
$P(d_v+\delta)=DP$,
and a chain homotopy inverse.  Equivalently these are the
finite homological perturbation identities, here derived by
expanding the finite series.  The first term of $D$ is the
cyclic bar differential for the exterior action of $M$.
Every remaining term lowers cyclic bar degree by at least two.

For completeness, the second contraction used in the algorithm
is also explicit.  The cyclic periodic resolution has one
$\mathbb Z[C_m]$ generator in each degree, with differential
$g-1$ in odd degree and $1+g+\cdots+g^{m-1}$ in even degree.
On $g^a$, $0\leq a<m$, its augmented contraction is
$\sum_{j=0}^{a-1}g^j$ in even degree and is $1$ precisely when
$a=m-1$ in odd degree, zero otherwise.
Apply (F.CA3) with this contraction to compare it with the
cyclic bar resolution.  Tensor its coefficients with the
exterior action of $M$.  Transferring the higher terms of $D$
once more by (F.CA5) terminates: each higher term followed by
the horizontal homotopy decreases cyclic degree by at least
one.  The resulting basis in total degree $n$ is
$\{(p,I):p+|I|=n,\ I\subseteq\{1,\ldots,r\}\}$.
This proves the claimed $2^r$ bound and termination in every
specified degree.

The free action makes the Borel construction equivalent to
$T^r/C_m$, and radial contraction identifies this with $B$.
For $Y$ use exactly the same algorithm on $T^{r+1}$ with
matrix $\operatorname{diag}(1,M)$ and numerator $(1,v)$.
The projection forgetting its first circle is equivariant and
represents $Y\to B$ under the radial contraction.
More generally an equivariant integral matrix $L$ acts on
groupoid arrows as
$\lambda\mapsto L\lambda+k_b-k_a$, where
$k_a=Lp_a-p'_a\in\mathbb Z^{r'}$.
Transferring this strict bar map by $P'LI$ gives the required
integer chain map.  Its mapping cone computes the pair.
Every series is finite with the proved degree bound, and
integer kernels and Smith normal forms finish the calculation.
\hfill$\square$

\emph{Executed certificates.}
The programs \path{code/fourfolds_bar_koszul.py} and
\path{code/fourfolds_affine_cap_chains.py} implement these
formulas, including the prime-subgroup freeness check of F.2.31.
The first checks all 281 abelian and 498 cyclic bar simplices
in its explicitly stated alphabets and degree ranges.
The second checks translated circles of orders two and three,
two affine presentations of the Klein bottle (one with a
nonfixed numerator), and the three-dimensional half-turn
manifold.  The separate rank-six program checks every degree
zero through seven for the order-five and order-four caps,
against the independent integral Wang sequence of their
five-torus mapping-torus descriptions.
The boundary-map program checks both translated circles and
the Klein bottle, including the full chain-map identities and
the integral Thom-shift homology of each cone.  These are
finite verification sets, not an asserted exhaustive test over
all input matrices; the terminating formulas prove generality.

\paragraph{Theorem F.2.34 (Proved: the marked Kuhn cusp pair).}
The fixed primitive rank-three Kuhn cusp admits an explicit
finite integer chain complex for its cap $W$, its boundary
$Y$, and the inclusion $Y\to W$.  It is obtained from a
category with 26 objects, 170 nonidentity morphisms, and
respectively $26,170,288,144$ strict flags of lengths
$0,1,2,3$.  The construction records the full integral
collapse map.  In particular it gives
\[
\begin{array}{c|rrrrrrrrr}
j&0&1&2&3&4&5&6&7&8\\\hline
\operatorname{rank}H_j(W)&1&3&9&9&10&3&1&0&0\\
\operatorname{rank}H_j(Y)&1&4&12&18&18&12&4&1&0\\
\operatorname{rank}H_j(W,Y)&0&0&1&3&10&9&9&3&1
\end{array}                                                \tag{F.CK1}
\]
The only torsion in these three rows is
$\operatorname{Tor}H_3(Y)=\mathbb Z/2$.

\emph{Proof and exact algorithm.}
Let $\mathcal K$ be the periodic Kuhn triangulation of
$\mathbb R^3$.  Normalize every simplex by translating its
lexicographically first vertex to zero.  The objects are its
simplex orbits; a morphism $\tau\to\sigma$ is every strict
translated face inclusion $\tau+w\subset\sigma$.
Keep distinct translated inclusions as distinct morphisms.
Composition adds the translation vectors.  Enumerating all
nonempty faces of the six chains in the unit cube gives
$(1,7,12,6)$ simplex orbits in dimensions zero through three,
and the displayed morphism and flag counts.  The category
nerve is the barycentric cell model of the dual decomposition
of $\mathbb R^3/\mathbb Z^3$; the direction of the face order
is reversed when interpreted as dual cells.

For an object $\tau$ let $L_\tau\subset\mathbb Z^3$ be the
primitive span of its vertex differences.  Its primitivity
follows from unimodularity of every Kuhn simplex.  Choose by
Smith reduction a quotient map
$Q_\tau:\mathbb Z^3\to\mathbb Z^{3-\dim\tau}$ with kernel
$L_\tau$, and an integral section $R_\tau$.
Define two strict diagrams of tori on this category:
\[
\begin{array}{c|c|c}
 &\text{object at }\tau&\text{map for }\tau+w\subset\sigma\\\hline
\mathcal W&T^{3-\dim\tau}&Q_\sigma R_\tau\\
\mathcal Y&T^4&S_w=\begin{pmatrix}1&0\\ w&I_3\end{pmatrix}.
\end{array}                                                \tag{F.CK2}
\]
The first maps compose because $Q_\sigma$ vanishes on
$L_\tau$, and the second compose because translations add.
The maps $[\,0\ Q_\tau\,]$ form a natural transformation
$\mathcal Y\to\mathcal W$: the only additional term is
$Q_\sigma w=0$, since $w$ is a vertex of $\sigma$.

Here these diagrams describe the actual local topology.
The toric moment quotient on the dual cell to $\tau$ has
angular torus $T^3/(L_\tau\otimes\mathbb R/L_\tau)$;
passing to a face collapses exactly the next angular
subtorus.  Barycentrically subdividing the dual cells gives
the homotopy colimit of $\mathcal W$ for the central fibre.
For its real oriented logarithmic boundary, retain the
argument $\theta_0$ of the base parameter and the three
angular coordinates $\theta$ of the toric chart.  Translating
the chart by $w$ changes them to
$(\theta_0,\theta+w\theta_0)$, which is precisely $S_w$.
Thus the logarithmic boundary is the homotopy colimit of
$\mathcal Y$, and forgetting the logarithmic phases on each
stratum is exactly the natural transformation just given.
This can also be read directly in the unimodular charts
$x_0\cdots x_k=t$: the moment coordinates form the dual
simplex, the four logarithmic phases remain, and the ordinary
stratum retains their quotient by the vertex differences.
The finite face identifications give (F.CK2), so there are
no further gluing constants in this fixed marked cap.

The Kuhn map with its divisor log structures is proper,
log smooth and vertical after restricting to a small disc.
Achinger--Ogus, Theorem~4.1.1, identifies its topological germ
with the mapping cylinder of this logarithmic collapse
\cite{ffAchingerOgus20}.  Consequently the preceding
homotopy-colimit diagrams and natural transformation model
the pair $(W,Y)$ itself.  Fix the marking by letting $b_i$
be the category loop with translation label $e_i$.  A
homotopy-colimit edge with map $S_{e_i}$ has the relation
$b_i^{-1}\theta_0b_i=\theta_0+a_i$; equivalently
$\theta_0b_i\theta_0^{-1}=b_i+a_i$.  This specifies the
required primitive cusp
$U=\begin{pmatrix}I_3&I_3\\0&I_3\end{pmatrix}$.

We specify finite integer chains, rather than only the
homology spectral sequence of these diagrams.  Replace every
torus by its ordinary normalized bar complex.  The affine
simplex map from the bar construction to the torus is natural
for all the integer homomorphisms in (F.CK2), so this replacement
does not change either homotopy colimit or its collapse map.
Take the simplicial replacement over all strict flags of the
finite category.  Its vertical differential is the torus bar
boundary; its horizontal differential composes adjacent
arrows, with the first face applying the diagram map, and
has the sign $(-1)^q$ in vertical degree $q$.
Apply the first contraction and the finite series (F.CA5)
object by object.  There are no flags longer than three,
so every series has at most four terms.  The resulting
finite basis consists of a strict flag beginning at $\tau$
and an exterior basis vector of its torus.  Transfer the
natural transformation by $P_{\mathcal W}FI_{\mathcal Y}$.
This proves both the chain-map identity and completeness of
the collapse map, including all integral coherence corrections.

The raw cap chain ranks in degrees zero through six are
\[
 (26,199,688,1231,1226,648,144),
\]
and the raw boundary ranks in degrees zero through seven are
\[
 (26,274,1124,2420,3010,2186,864,144).
\]
The program enumerates every object, every translated
incidence and every flag, verifies composition, checks
$d^2=0$ on every generator, and checks the collapse-chain
identity on every generator.  Integral cancellation of a
unit pivot is a chain homotopy equivalence: clear its row by
integer column operations, then remove the resulting
contractible two-generator summand.  The adjacent differential
updates are forced by $d^2=0$.  After respectively 2063,
4988 and 7087 such cancellations for the cap, boundary and
collapse cone, Smith reduction gives (F.CK1).  All matrices
and the remaining differentials are saved by
\texttt{fourfolds\_kuhn\_cap\_chains.py}; the reduced cap and
relative differentials vanish, and the boundary has one
nonzero Smith factor, equal to two in degree four to three.
This computes the integral groups without a filtration
extension problem.\hfill$\square$

\emph{Scope of the chain result.}
The local finite pairs and the full cusp pair now have
effective marked chain models.  To turn them into a universal
integral homology calculation for arbitrary nonsurviving F.2
data, one must also implement the maps from these boundary
models to the common torus-bundle complement, with its
specified peripheral word and clutching translation.  The
local results alone do not assert that this remaining global
comparison has been completed.

\paragraph{Theorem F.2.35 (Proved: complete integral homology of a marked gluing).}
For every fully marked F.2 datum there is a terminating exact
integer algorithm computing every integral homology group of
the resulting eight-manifold, including every torsion invariant
factor.  The input consists of the finite homology matrices
$M_1,\ldots,M_k$, free affine numerators $v_i/m_i$ of their
actual orders, the integral clutch vector $\kappa$, and the
chosen primitive cusp marking
\[
 M_0=(M_1\cdots M_k)^{-1},\qquad
 M_0P=P\begin{pmatrix}I_3&I_3\\0&I_3\end{pmatrix},
 \qquad P\in\mathrm{GL}(6,\mathbb Z).                       \tag{F.GC1}
\]
Here $k=2$ or $3$ and the genuine orders are at least two;
order-one sites are first eliminated as in F.2.31.
If a marking is not supplied, the algorithm makes and records
one by Smith reduction.  It does not identify that choice
with an unspecified prescribed marking.
The complete chain prefix through degree nine has at most
4003 generators in any degree.  No spectral-sequence
extension, saturation convention, orientation assumption or
unbounded search enters the homology calculation.

\emph{Proof.}
We supply both global boundary maps which were left outstanding
after F.2.34.  Let $F_k$ be the free group with ordered
generators $s_1,\ldots,s_k$, and let $M(w)$ be the product
specified by a reduced word $w$.  The torus-bundle complement
over the punctured sphere retracts to the corresponding bundle
over a graph.  Its fundamental group and multiplication are
\[
 \Gamma=\mathbb Z^6\rtimes F_k,
 \qquad (x,w)(y,u)=(x+M(w)y,wu).                            \tag{F.GC2}
\]
The fibre and base are aspherical and the base has no
two-cells.  Hence a chosen section identifies the complement
up to homotopy with $B\Gamma$.  These section and peripheral
choices are fixed by the input clutch convention below.

\emph{A finite complement complex with its full bar projection.}
Use reduced words for $F_k$.  The tree resolution of the free
group has $K_0=\mathbb Z[F_k]$ and
$K_1=\bigoplus_i\mathbb Z[F_k]e_i$, with
$de_i=s_i-1$, and no higher terms.  Its explicit augmented
contraction sends a word to its signed path in the Cayley
tree: for each positive letter $s_i$ add its preceding prefix
times $e_i$, and for a negative letter $s_i^{-1}$ subtract
the prefix including that letter times $e_i$.
This is the Fox derivative, and its boundary is $w-1$.
Together with the bar insertion contraction it gives, by
the recursion (F.CA3), a strong contraction of the free-group
bar resolution to this tree resolution.  Every recursion
decreases bar degree and every word is finite.

For a bar simplex $[(x_1,w_1)|\cdots|(x_n,w_n)]$ of $\Gamma$,
put $y_j=M(w_1\cdots w_{j-1})x_j$.
The exact Alexander--Whitney chain map to the torus/free-group
double bar complex is
\[
 \sum_{p+q=n}(-1)^{pq}
 [w_1|\cdots|w_p]\otimes[y_{p+1}|\cdots|y_n].              \tag{F.GC3}
\]
Terms with an identity bar entry are zero.  The horizontal
first face acts on torus vectors by
$\rho(w)=M(w)^{-1}$, and the horizontal boundary has sign
$(-1)^q$.  Notice that $\rho$ is a right action:
$\rho(wu)=\rho(u)\rho(w)$.  This order is essential when
the matrices do not commute.  Formula (F.GC3) follows from
the diagonal bisimplicial bar construction; alternatively,
its chain-map identity follows by pairing adjacent interior
faces, with the first face using
$M(w_1)^{-1}y_j=M(w_2\cdots w_{j-1})x_j$.
The usual shuffle inverse is finite in each degree, so the
map is a chain equivalence.

Apply the torus contraction and (F.CA5), then the free-group
tree contraction and the same finite perturbation formula.
Higher terms decrease group bar degree at least two, whereas
the tree complex only has degrees zero and one.  Thus the
resulting complement complex is
\[
 C_n(E)=\bigwedge^n\mathbb Z^6\ \oplus\
 \bigoplus_{i=1}^k\bigwedge^{n-1}\mathbb Z^6,
 \qquad
 d(e_i\otimes\xi)=(-1)^{n-1}
   (\bigwedge^{n-1}M_i^{-1}-1)\xi.                         \tag{F.GC4}
\]
Its differential alone has this short form; its projection
from group bars retains the finite higher perturbation terms.
Those projection terms are used for every boundary map.

\emph{Finite affine boundary maps.}
For the $i$th finite cap use the orbit groupoid of F.2.33.
Its one-edge homotopy colimit for the affine generator is
the boundary mapping torus.  Vertically contracting it gives
a bounded complex with
$C_n(Y_i)=\bigwedge^n\mathbb Z^6\oplus
\bigwedge^{n-1}\mathbb Z^6$ and differential
$(-1)^{n-1}(\bigwedge^{n-1}M_i-1)$ on the second summand.
Its inclusion into the cyclic bar construction, followed by
the cap projection in (F.CA5), is the full map $Y_i\to B_i$.
In particular this one-edge model avoids an unnecessary
periodic resolution of the seven-dimensional boundary.

Here is an explicit group-bar map for its other inclusion.
Send a torus-groupoid arrow labelled $\lambda$ to the
translation $x^\lambda\in\Gamma$.  With the integral offsets
$z_1(a)$ in (F.CA4), put
$\eta_a=s_i^{-1}x^{z_1(a)}$.
For an arrow $f:a\to b$, the exact identity
\[
 \eta_a\,\psi(gf)=\psi(f)\,\eta_b                        \tag{F.GC5}
\]
is the natural-transformation identity: it is just
$s_i^{-1}x^{M_i\lambda+z_1(b)}
 =x^\lambda s_i^{-1}x^{z_1(b)}$.
For a $q$-simplex, the edge map is $(-1)^q$ times the
alternating prism which inserts $\eta$ at every position,
uses the original arrows before the insertion and the
transformed arrows afterwards.  Its boundary is precisely
the signed difference of the two end maps, by (F.GC5).
Compose this prism with the finite vertical lift of the
one-edge complex, then with (F.GC3) and the complement
projection.  This proves the full chain-map identity.
The offsets record the actual integer affine lift: replacing
$v_i$ by $v_i+m_i\lambda$ changes the offsets and the map,
even if the affine map on the torus has not changed.
An optional unimodular local basis simply applies its matrix
to every translation in (F.GC5).

We make the distinction between a bar cell and the geometric
positive root explicit.  Put
$a_i=z_1(0)=\lfloor v_i/m_i\rfloor$, coordinatewise, and
let $\tau_i$ be the positive one-edge bar cell closed using
the zero-labelled return arrow to object zero.  With the
first-face convention above its geometric loop is
\[
 \tau_i=g_i^{-1}x^{a_i},\qquad
 F_E(\tau_i)=-e_i+M_i^{-1}a_i\in C_1(E),\qquad
 g_i=x^{a_i}\tau_i^{-1}.                                  \tag{F.GC5a}
\]
Thus the bar cell has the opposite peripheral orientation,
with the displayed integer correction when the numerator
has a nonzero integer part.  Formula (F.GC5) sends its
group element to $s_i^{-1}x^{a_i}$; consequently the
geometric positive root in the last expression of (F.GC5a)
maps to $s_i$ exactly.  This is the stated conversion between
the boundary and complement chain conventions.  In a local basis
$Q$, replace $a_i$ in the group and complement formulas by
$Q\lfloor Q^{-1}v_i/m_i\rfloor$.
Writing $\mathcal N_i=\sum_{j=0}^{m_i-1}M_i^j$, the cyclic
degree-two cell has boundary
$m_i\tau_i+w_i-\mathcal N_i a_i$.
Indeed $\tau_i^{m_i}=x^{\mathcal N_i a_i-w_i}$ follows
directly from $g_i^{m_i}=x^{w_i}$.
This independently checks the affine power relation and
the degree-one prism sign, including negative numerators.

There is an optional exact geometric shortcut.  If a verified
local integral basis gives $M_i=\operatorname{diag}(1,D)$ and
$v_i=\epsilon e_0$, $\epsilon\in\{1,-1\}$, eliminate the
first fibre generator from the cap group using
$g^{m_i}=u^\epsilon$.  The cap is the mapping torus with
group $\mathbb Z^5\rtimes_D\mathbb Z$.
The fibre vector $(u,a)$ maps to $x^a g^{\epsilon m_i u}$,
and the root generator maps to $g$.  Substituting this
homomorphism for $\psi$ in the same prism gives the boundary
map to this bounded geometric cap model.  The shortcut is
used only after checking the displayed integral normal form;
the general periodic construction remains available for
every other input.

To check the peripheral convention directly, the universal
cover of $Y_i$ has the root generator $g_i$ with linear part
$M_i$ and affine displacement $v_i/m_i$.  Its integral
translations and $g_i$ map to the same fibre translations
and $s_i$ in (F.GC2).  The full base-circle translation obeys
$g_i^{m_i}=x^{w_i}t$, where $w_i$ is the power vector of
F.2.31.  Filling the disc kills $t$, giving exactly
$s_i^{m_i}=x^{w_i}$.  The straight paths between the chosen
orbit points give (F.GC5), so the prism is this marked map,
including numerators with nonzero integer parts.

\emph{The marked cusp boundary map.}
Write $P=[A\ B]$ in (F.GC1), and set
\[
 w_0=(s_1\cdots s_k)^{-1},\qquad
 c=w_0x^{-\kappa}\in\Gamma.
                                                               \tag{F.GC6}
\]
Send the three angular generators to the columns of $A$,
the three tropical generators to the columns of $B$, and
the base-phase generator to $c$.  Relation (F.GC1) gives
$c b_i c^{-1}=b_i+a_i$ and all the other nilpotent boundary
relations.  Filling kills the angular generators and $c$,
so (F.GC6) gives exactly
$s_1\cdots s_kx^\kappa=1$, in agreement with F.2.31.

This homomorphism has a direct chain implementation on the
full Kuhn diagram.  For an incidence with translation label
$w$, assign the group element $x^{Bw}$.  A vertical phase
$(n,a)\in\mathbb Z\oplus\mathbb Z^3$ maps to $c^nx^{Aa}$.
The incidence shear is $S_w(n,a)=(n,a+n w)$, and
\[
 (c^nx^{Aa})x^{Bw}=x^{Bw}(c^nx^{A(a+n w)}).                \tag{F.GC7}
\]
For a flag of $p$ incidence arrows and a vertical bar
simplex of length $q$, sum over all $(p,q)$ shuffles.
Each incidence step contributes its assigned group element;
each vertical step contributes its phase transformed by
all preceding incidence shears.  Give the shuffle the sign
$(-1)^{\#\{\text{incidence steps before vertical steps}\}}$.
Interior face cancellations use exactly (F.GC7), and adjacent
incidence faces use addition of their translation labels.
This is therefore a chain map from the double bar complex
to $C_*(B\Gamma)$.  Precompose with the finite Kuhn lift of
F.2.34 and postcompose with (F.GC3) and the complement
projection.  The other cusp map is the full logarithmic
collapse already computed in F.2.34.
Both the cusp boundary and the complement are aspherical,
so the displayed marked homomorphism also verifies the
homotopy class of this geometric inclusion.

\emph{The final finite integer calculation.}
Take the mapping cone of the just constructed maps
\[
 \bigoplus_{i=0}^k C_*(Y_i)\longrightarrow
 C_*(E)\oplus\bigoplus_{i=0}^k C_*(B_i),
 \qquad y\longmapsto(f_E(y),-f_B(y)),                     \tag{F.GC8}
\]
where $B_0=W$ is the Kuhn cap.  Gluing collars makes this
the homotopy pushout model of the actual closed manifold.
Consequently its cone computes integral homology directly;
there is no remaining filtration extension to choose.

For clarity the exact finite size bound can be read off
before running any matrix calculation.  Put
$t_n=\sum_{q=0}^{\min(n,6)}\binom6q$, with binomial
coefficients outside their range interpreted as zero.
Let $w_n$ and $y_n$ be the raw cap and boundary ranks of
F.2.34, zero outside their listed ranges.  Then the cone rank
in degree $n$ is
\[
 \binom6n+k\binom6{n-1}+k t_n+w_n+
 k\binom7{n-1}+y_{n-1}.                                 \tag{F.GC9}
\]
For $0\leq n\leq9$ its maximum is 3890 when $k=2$ and
4003 when $k=3$.  Every bar comparison and perturbation
series used to produce these matrices terminates by its
degree filtration.  Word reductions and affine orbit
enumerations are finite, with no imposed size cutoff.
Computing integer kernels followed by Smith normal forms
now gives $H_0,\ldots,H_8$; higher groups vanish by dimension.
Unit cancellations are optional exact integral chain
equivalences which substantially reduce the matrices before
Smith reduction.  This proves the algorithm and the bound.
\hfill$\square$

\emph{Verification and implementation.}
The implementation is in \path{code/fourfolds_global_chains.py}
and \path{code/fourfolds_total_homology.py}.
The generic JSON driver \path{code/fourfolds_marked_gluing.py}
retains the complete supplied marking and affine data, computes
the full group structure of F.2.31, and independently compares
its abelianization with the chain calculation of $H_1$.
The Fox contraction is checked on all 156 explicitly listed
free-group bar simplices through degree three; the exact
Alexander--Whitney and finite-projection identities are
checked on 156 specified semidirect-product simplices with
two noncommuting rank-two matrices.
Four finite boundary comparisons using the independent
seven-dimensional quotient model have acyclic integral
comparison cones.  The more efficient one-edge model is
checked in six stated low-rank cases, including negative,
nonfixed, and nontrivial integer-part numerators; both its
complement comparison and its full Thom cone pass.
The program \path{code/fourfolds_peripheral_signs.py} also
checks (F.GC5a) and the full cyclic degree-two relation in
seven cases, including both actual rank-six caps.  It
reconstructs the positive geometric root as a group element
and checks that its complement bar projection is $+e_i$.
For the Kuhn boundary, every one of its 10048 raw generators
passes the chain-map identity in all degrees zero through
seven.  Its comparison to the corresponding one-generator
complement has an integrally acyclic cone, eliminated by
5088 unit cancellations.  Thus the comparison also checks
the nonzero boundary order-two torsion.
The full eight-dimensional integration for the original
order-$(5,4)$ survivor has passed.  Every marked boundary
column and the complete Mayer--Vietoris differential are
checked.  After 7372 integral unit cancellations, its Smith
homology is $\mathbb Z$ in degrees zero and eight,
$\mathbb Z^4$ in degree four, and zero otherwise, exactly
as in the independent calculation F.2.6.
The corresponding three-point run with logarithmic order
$m=2$ also passes, after 7470 integral unit cancellations.
It gives $H_1=H_6=\mathbb Z/18$,
$H_2=H_5=(\mathbb Z/2)^5$, $H_3=(\mathbb Z/2)^{10}$,
$H_4=\mathbb Z^4\oplus(\mathbb Z/2)^{10}$,
$H_0=H_8=\mathbb Z$ and $H_7=0$, exactly as in the
independent theorem F.2.32.  It reuses only the identical
first-two finite boundary maps: their group bars involve only
the fibre and their own peripheral generator, so adjoining
a third generator changes neither their columns nor their
target row indices.  The third cap, the new cusp word,
the full cone and every Smith group are recomputed.

A third full integration uses
$A=I_3$, $B=\operatorname{diag}(0,1,0)$ in (F.IV2),
actual affine orders $4,4$, numerators $a_1,a_2$, and zero
clutching.  These two linear actions reverse orientation.
Its 7345 integral unit cancellations give free ranks
\[
                  (1,0,4,11,10,1,6,3,0).
\]
The nonzero torsion groups are
\[
\begin{aligned}
 \operatorname{Tor}H_1=\operatorname{Tor}H_5
       &=(\mathbb Z/2)^3\oplus\mathbb Z/4,&
 \operatorname{Tor}H_2&=(\mathbb Z/2)^2,\\
 \operatorname{Tor}H_3=\operatorname{Tor}H_7
       &=\mathbb Z/2,&
 \operatorname{Tor}H_4&=(\mathbb Z/2)^4\oplus(\mathbb Z/4)^5.
\end{aligned}
\]
The independent van Kampen $H_1$ agrees.  Its Euler number
is six, its mod-two Betti vector is
$(1,4,10,14,20,14,10,4,1)$, and its integral top groups
have the required nonorientable form.  Thus this integration
also tests the orientation-reversing part of the stated scope.

The \texttt{--fresh} option bypasses all saved map checkpoints.
The free-bar and semidirect tests, all four periodic-boundary
comparisons, all six one-edge comparison/Thom tests, all seven
peripheral-sign checks, the complete 10048-generator cusp
comparison, and the cached generic survivor computation have
also been independently rerun successfully.  The
reduced-chain, van Kampen $H_1$, Euler and mod-two duality
checks have additionally been independently replayed on
all three complete integration outputs.  The general
algorithm is proved by the explicit chain equivalences and
finite operations above; these finite execution sets are
reported with their exact scope.

\paragraph{Theorem F.2.36 (Proved: simultaneous integral conjugacy for every input).}
There is an unconditional exact algorithm which decides simultaneous
$\mathrm{GL}_6(\mathbb Z)$ conjugacy of any two marked matrix tuples of
F.2 and returns a conjugating matrix when one exists. The algorithm
does not impose a bound on their entries, on the affine multiplicities,
or on the integral conjugator. More generally it applies to any integral
tuple containing a matrix $1+N$ with $N^2=0$ and $\operatorname{rank}N=3$;
the primitive-image assumption is unnecessary for this statement.

\begin{proof}
We use column vectors and seek $P$ satisfying $PX_i=Y_iP$ for every
marked index, including the cusp index. Put
\[
\begin{aligned}
 H&=\{P\in\operatorname{Mat}_6(\mathbb Z):PX_i=Y_iP\ (\forall i)\},\\
 \mathcal O&=\{C\in\operatorname{Mat}_6(\mathbb Z):CX_i=X_iC\ (\forall i)\},
 \qquad E=\mathbb Q\mathcal O.
\end{aligned}                                                    \tag{F.IC1}
\]
Both lattices are saturated integer kernels of explicitly displayed
linear equations in thirty-six unknowns. Thus Hermite or Smith
reduction computes full bases, rather than sublattices of finite index.

We first decide rational isomorphism. For a basis $h_1,\ldots,h_s$ of
$H$, form the polynomial $f(c)=\det(\sum c_jh_j)$. If it is zero,
there is no rational isomorphism. Otherwise successively substitute
one of $0,1,\ldots,6$ for each variable while keeping the remaining
polynomial nonzero. At each stage one of these seven values works,
because the degree in that variable is at most six. The final matrix
$D\in H$ is invertible over $\mathbb Q$. This is a terminating exact
test, including the negative case, and not a search for a lucky
invertible intertwiner.

The full lattice $L=D^{-1}H\subset E$ is a right $\mathcal O$-module.
Let $J$ be the Jacobson radical of $E$, let $B=E/J$, and use bars
for images under the quotient map. We next verify the splitting
hypothesis needed for the principal-ideal algorithm. In a rational
$N$-adapted basis,
\[
 N=\begin{pmatrix}0&I_3\\0&0\end{pmatrix},\qquad
 E\subseteq
 \left\{\begin{pmatrix}A&T\\0&A\end{pmatrix}:A,T\in
                                      \operatorname{Mat}_3(\mathbb Q)\right\}.
                                                               \tag{F.IC2}
\]
Projection to $A$ has square-zero kernel and image an algebra
$S\subseteq\operatorname{Mat}_3(\mathbb Q)$. A nilpotent ideal is
contained in the radical, so $E/J\cong S/J(S)$.

Choose a composition series of the faithful $S$-module $\mathbb Q^3$.
The kernel of its action on the associated graded module lowers this
filtration and is therefore nilpotent. Conversely the radical kills
every simple factor. Thus this kernel is exactly $J(S)$ and every
simple factor of $S/J(S)$ acts nontrivially on a composition factor.
For a simple algebra $\operatorname{Mat}_n(D)$, where $D$ has centre
$K$ and degree $d$ over that centre, every nonzero simple module has
$\mathbb Q$-dimension $n d^2[K:\mathbb Q]$. If $D$ were
noncommutative this would be at least four, which is impossible in
$\mathbb Q^3$. Consequently
\[
 B\cong\prod_j\operatorname{Mat}_{n_j}(K_j),\qquad
 n_j[K_j:\mathbb Q]\le3.                                    \tag{F.IC3}
\]
Every noncommutative factor is therefore
$\operatorname{Mat}_2(\mathbb Q)$ or
$\operatorname{Mat}_3(\mathbb Q)$; the remaining factors are number
fields of degree at most three. This proves hypothesis (H) of
\cite[Algorithm 8.3]{rootBHJ22} for every tuple under consideration.

For clarity this decomposition is effective without a split-algebra
oracle. Rational linear algebra computes $J$ as the radical of the
regular trace pairing and computes the quotient multiplication table.
Its simple factors are obtained from its centre. A commutative factor
is a number field. In a noncommutative factor of dimension $n^2$, we
enumerate integral combinations $e$ of a rational basis until the
left ideal $Be$ has dimension $n$. Such an $e$ exists because the
factor is $\operatorname{Mat}_n(\mathbb Q)$, so its rank-one matrices
are rational points; clearing denominators expresses one among the
enumerated combinations. Left multiplication on $Be$ then supplies
an explicit isomorphism to $\operatorname{Mat}_n(\mathbb Q)$.
The program verifies multiplication on every pair of basis elements.

The images $\overline{\mathcal O}$ and $\overline L$ are a full order
and a full right lattice in $B$. If the right order of
$\overline L$ differs from $\overline{\mathcal O}$, this lattice
cannot be principal over that order. Otherwise Algorithm 8.3 of
\cite{rootBHJ22}, in its right-ideal form, decides principality and
returns a generator $\overline\alpha$ when one exists. The cited
algorithm passes through a maximal order and its central conductor,
tests local freeness and principality over the maximal order, and
solves the remaining reduced-norm congruence in a finite quotient.
Its correctness proof uses its Theorem 4.5(b) and Theorem 7.9; its
Theorem 8.4 specifies the number-field, finite-field and integer
factorization subproblems. These have unconditional exact algorithms.
In the recorded implementation the class and unit groups of every
number-field component are certified with \texttt{GRH=false} before
the principal-ideal decision is used. The negative decision thus
does not presume the generalized Riemann hypothesis.

If $\overline L$ is not principal, $L$ cannot be principal. If it is
principal, solve the integer linear equations expressing
$\overline\alpha$ as a combination of the images of the chosen basis
of $L$. This produces a lift $\alpha\in L$ and an integral
intertwiner $P=D\alpha\in H$. We claim that a simultaneous integral
conjugacy exists if and only if $|\det P|=1$.

The sufficient direction is immediate. For the necessary direction,
suppose $P_0\in H$ is unimodular. Every $h\in H$ has
$P_0^{-1}h\in\mathcal O$, so
$L=\alpha_0\mathcal O$, where $\alpha_0=D^{-1}P_0$.
Write $\alpha=\alpha_0u$ with $u\in\mathcal O$.
Because both image elements generate $\overline L$,
$\overline u$ is a unit of $\overline{\mathcal O}$.
Lift its inverse to $v\in\mathcal O$. Then $uv=1+j$ with
$j\in J\cap\mathcal O$, and the finite geometric series for
$(1+j)^{-1}$ shows that $u^{-1}=v(1+j)^{-1}\in\mathcal O$.
Thus $u$ is an integral unimodular matrix and
$P=P_0u$ is unimodular. This also proves that any chosen lift of the
image generator gives the same yes-or-no answer. It is the
right-module form of \cite[Theorem 3.4 and Propositions 8.1--8.2]{rootBHJ22},
with the needed determinant condition written out.

Every branch either proves nonexistence or returns a matrix whose
integrality, determinant and all intertwining equations are verified
exactly. This completes the decision procedure.
\end{proof}

The implementation is
\path{code/fourfolds_integral_conjugacy.jl}, with the pinned runtime
setup and runner in the same directory. It is an application of the
published lattice-isomorphism algorithm to the special centralizer
forced by the cusp; we claim no priority for integral-conjugacy
decidability. The accompanying source-reading record identifies the
primary theorem statements and the precise implementation paths.
This comparison algorithm is distinct from a completed tabulation
of all nonsurviving tuples or from an arbitrary-height truncation
of such a tabulation.


The exact execution certificate has twelve comparisons. Six algebraic
tests exercise the full $\operatorname{Mat}_3(\mathbb Q)$ quotient,
$\operatorname{Mat}_2(\mathbb Q)\times\mathbb Q$, a rational
multiplicity obstruction, a principal quotient whose lift has
determinant four, a nonprincipal ideal over $\mathbb Q(\sqrt{-5})$
with unconditionally certified class number two, and a cubic
centre with a certified rank-one unit group. Six further comparisons
use the actual admissible order-$(5,4)$ tuples with parameters
$0,1,2$: each is recovered after an explicit large unimodular
conjugation, and all three distinct pairs are rejected. Their word
traces $\operatorname{tr}(M_1M_2^3)=2,1,-10$ independently certify
the latter negative answers. All returned positive matrices satisfy
every exact intertwining equation and have determinant one.

\paragraph{Theorem F.2.37 (Proved: a complete canonical-representative algorithm).}
Every marked linear tuple in F.2, for either $k=2$ or $k=3$,
has a computable canonical integral representative and a
returned simultaneous unimodular conjugator. Two inputs have
the same representative if and only if they are simultaneously
$\mathrm{GL}_6(\mathbb Z)$ conjugate. The algorithm uses a proved
finite bound depending on the input. No fixed global entry bound
can represent every class.

This is an effective classification, not a claim that its
infinite output has been printed in a finite table. The
ordered involution sector has the separate complete finite
normal form F.2.38, and all of its $56$ classes have been
computed. The general bounded fallback is potentially very
expensive; its finiteness is a theorem, not a runtime estimate.

\begin{proof}
We first remove the existential condition about affine lifts
from the admissibility test for a linear tuple. A finite-order
matrix $M$ admits a free finite affine lift on $T^6$ if and
only if $\ker(M-1)\ne0$. If this kernel is zero, $1-M$ is a
surjective endomorphism of the torus, so every affine map with
linear part $M$ has a fixed point. Conversely let the order of
$M$ be $d>1$ and choose a primitive $u\in\ker(M-1)\cap\mathbb Z^6$.
The map $x\mapsto Mx+u/d^2$ has actual order $d^2$: its
$d$th power is translation by $u/d$, of exact order $d$.
For each prime $p\mid d$, the generator of its order-$p$
subgroup is translation by $u/p$, which has no fixed point.
Every nontrivial stabilizer contains a prime-order subgroup,
so the whole cyclic action is free. If $M=1$, translation
by a primitive half-vector is a free lift of order two.
The invariant integral kernel is primitive, so such a
primitive vector is obtained by Smith reduction.

Finite order is also effectively bounded at this linear
stage. Its minimal polynomial is a product of distinct
cyclotomic polynomials, and eigenvalue one uses at least one
of the six dimensions. Every other cyclotomic factor therefore
has degree at most five and their total degree is at most
five. For a cyclotomic index $d$, each prime divisor $p$ has
$p-1\le5$, and its prime-power factors satisfy
$p^{a-1}(p-1)\le5$. Enumerating these factors and the possible
degree sums gives exactly
\[
                 d\in\{1,2,3,4,5,6,8,10,12\}.             \tag{F.CN1}
\]
Here $d$ is the order of the matrix, the least common multiple
of its cyclotomic indices. Multiplicities of a factor can only
increase the dimension and introduce no new orders. Each
listed order occurs with a fixed vector after adding trivial
summands. Thus taking the powers in (F.CN1), checking the
fixed kernel, the product relation, and the cusp Smith
factors gives an exact admissibility test. The unbounded
actual affine multiplicities are not bounded by (F.CN1).

For an admissible full tuple $T=(M_1,\ldots,M_k,M_0)$ let
$H(T)$ be its maximum absolute entry. If a three-point linear
tuple has an identity at a finite position, temporarily delete
that identity, classify the remaining two-point tuple, and
reinsert the identity at the same marked position. There can
be only one such identity, since otherwise a finite matrix
would equal the inverse of the nontrivial unipotent cusp.
Deleting a fixed identity changes neither height nor the
lexicographic comparison of tuples in that conjugacy-invariant
sector. This is only a linear normalization: a nontrivial
affine order at the identity-linear site is retained in the
gluing datum. In the two-involution sector use the canonical
representative defined in F.2.38. This sector is invariant under simultaneous
conjugacy. In every other sector, order full integer tuples
first by $H$, then lexicographically by their row-major
entries in the displayed marked order. Define the canonical
representative to be the first admissible tuple in the
simultaneous conjugacy class.

To compute it, enumerate heights $1,\ldots,H(T)$ and the
integer entries at that height. There are at most
\[
                     (2H(T)+1)^{36(k+1)}                  \tag{F.CN2}
\]
tuples in this complete box. The product relation allows
$M_0$ to be calculated from the first $k$ matrices, so the
implementation enumerates fewer entries. It also rejects
matrices with different characteristic polynomials or orders
before calling the conjugacy decision F.2.36. These filters
cannot reject an equivalent tuple. At least one comparison
succeeds: the supplied tuple itself belongs to its own box.
The first successful comparison returns the canonical
representative and its exact conjugating matrix.

Every simultaneous orbit has a least element in this order.
Its height is at most the height of any supplied member,
which proves that (F.CN2) loses no class representative for
that input. If two inputs are equivalent they therefore
return the same least element. Conversely the returned
conjugators show that equal representatives imply equivalent
inputs. The casewise involution normal form has the same
property by F.2.38, and an involution tuple cannot be
equivalent to a tuple in the other branch. This proves the
classification assertion, including its finite stopping rule.

Finally, F.2.16 already gives an unbounded admissible integral
family. In its homology convention direct symbolic
multiplication gives
\[
             \operatorname{tr}(M_1M_2^3)
                       =-v^4+v^3-v^2+2.                 \tag{F.CN3}
\]
This is unbounded even on the previously constructed free
affine subfamily $v\in20\mathbb Z$. It is invariant under
simultaneous conjugacy. A fixed finite entry box contains
only finitely many tuples and hence only finitely many such
traces; it cannot contain a representative of every class.
The input-dependent nature of (F.CN2) is therefore necessary
for this particular exhaustive procedure.
\end{proof}

The public implementation is
\path{code/fourfolds_tuple_classification.py}.
Its finite involution branch has been run on all $56$
representatives after nontrivial simultaneous integral
conjugation, and returned every exact representative and
conjugator. F.2.38 separately records $112$ normalization
checks. The marked identity-insertion branch was also run
on all $168$ conjugated inputs from the 56 classes and the
three possible identity positions; every representative and
conjugator passed. The general height-box fallback has not been
exhaustively run on a noninvolution tuple; no such computation
is claimed. Its underlying exact comparison procedure was
run on the twelve instances listed in F.2.36, including
three actual order-$(5,4)$ classes. The symbolic trace
calculation is saved with its executable program and output.
The all-twist group algorithm is F.2.31, and the full
necessary-and-sufficient three-connected-survivor tests are
F.2.30. These uses of the classifier do not presume that
arbitrarily many classes can be tabulated in finite time.

When transporting a specified glued manifold, the returned
matrix $P$ also sends each affine numerator to $Pv_i$, the
clutch vector to $P\kappa$, and the supplied integral cusp
marking to $P P_{\mathrm{cusp}}$. The multiplicities are
unchanged. These are the actual linear changes of fibre
coordinates and therefore carry every boundary gluing map.
The linear canonical representative does not authorize
discarding the cusp marking and independently choosing a
different Smith marking for the filled manifold.

\paragraph{Theorem F.2.38 (Proved: the complete ordered involution stratum).}
There are exactly 56 simultaneous $\mathrm{GL}_6(\mathbb Z)$
conjugacy classes of ordered tuples $(M_1,M_2,M_0)$ satisfying
the primitive maximal-cusp conditions of F.2 with both finite
linear orders equal to two.  An explicit integral normalizer
returns a canonical representative and a checked unimodular
conjugating matrix for every such input, without any bound on
its entries.  Every representative has free affine lifts of
actual order four at both finite points.  None can give a
simply connected filling, for any choices of free affine
orders or twists.

\begin{proof}
First choose an integral cusp basis, so that
\[
 M_0=I_6+N,\qquad N=\begin{pmatrix}0&I_3\\0&0\end{pmatrix}.
                                                               \tag{F.IV1}
\]
This is effective by Smith reduction.  In fact choose vectors
$\beta_i$ whose images $\alpha_i=N\beta_i$ form a primitive
basis of $\operatorname{im}N=\ker N$.
The index-one hypothesis says that their classes form a basis
of the quotient, so $[\alpha_1,\alpha_2,\alpha_3,
\beta_1,\beta_2,\beta_3]$ is unimodular.

Since $M_2=M_1(I-N)$ and $M_1^2=M_2^2=I$, we have
$M_1NM_1=-N$.  Solving this block equation gives precisely
\[
 M_1=\begin{pmatrix}A&B\\0&-A\end{pmatrix},\qquad
 M_2=\begin{pmatrix}A&B-A\\0&-A\end{pmatrix},\qquad
 A^2=I_3,\quad AB=BA.                                     \tag{F.IV2}
\]
Every integral conjugator preserving (F.IV1) is
$\begin{pmatrix}C&D\\0&C\end{pmatrix}$, with
$C\in\mathrm{GL}_3(\mathbb Z)$ and $D$ arbitrary integral.
Thus one first classifies $A$ integrally; after fixing it,
$B$ is considered modulo $AD+DA$ and modulo conjugation by
the integral centralizer of $A$.

We include the integral involution classification needed here.
For the lattice $L=\mathbb Z^3$, let $L_+$ and $L_-$ be the
primitive kernels of $A-I$ and $A+I$.  They satisfy
$2L\subseteq L_+\oplus L_-\subseteq L$.
The quotient $L/(L_+\oplus L_-)$ is a vector space over
$\mathbb F_2$.  Its projections into $L_+/2L_+$ and
$L_-/2L_-$ are injective: a glue vector with only one
nonzero projection would give a new integral vector in the
corresponding eigensublattice, contradicting its definition.
It is consequently the graph of an isomorphism between two
subspaces of equal dimension $r$.
Unimodular changes of bases in $L_\pm$ put this graph in
coordinate form.  For each glued pair $e_+,e_-$ the vectors
$(e_++e_-)/2$ and $(e_+-e_-)/2$ form a swap block.
The unglued vectors give trivial and sign blocks.  Hence
there is a unique type
\[
 A\cong I_p\oplus(-I_q)\oplus
 \begin{pmatrix}0&1\\1&0\end{pmatrix}^{\oplus r},
 \qquad p+q+2r=3.                                         \tag{F.IV3}
\]
The dimensions of the two rational eigenspaces and the
index $[L:L_+\oplus L_-]=2^r$ prove uniqueness.
Here $r$ is zero or one, so this normalization only requires
integer kernels, one Smith reduction and, in the glued case,
one explicit half-sum pair whose integrality is checked.

For fixed (F.IV3), compute the quotient of the commuting
matrices by the subgroup $\{AD+DA\}$.
On the trivial and sign diagonal blocks the norm is
respectively $2D$ and $-2D$.  The commuting off-diagonal
blocks between the trivial and sign summands are zero.
Every commuting block involving a swap summand is a norm.
For example a trivial-to-swap block is a row $(a,a)$ and
is obtained as $(d,e)+(e,d)$; a sign-to-swap block has the
same calculation with a difference.  For two swap summands,
the commuting block is
$\begin{pmatrix}a&b\\b&a\end{pmatrix}$, and applying the
norm to $\begin{pmatrix}d&e\\f&g\end{pmatrix}$ gives
$\begin{pmatrix}e+f&d+g\\d+g&e+f\end{pmatrix}$.
The reverse mixed blocks have the transposed calculations.
Therefore the full integral quotient is exactly
\[
 \operatorname{End}_{\mathbb Z[A]}(L)/\{AD+DA\}
 \cong \operatorname{Mat}_p(\mathbb F_2)
       \oplus\operatorname{Mat}_q(\mathbb F_2).             \tag{F.IV4}
\]
No finite-index sublattice has been substituted for a norm
image; the displayed block calculations are surjective
integer formulas.

The centralizer action on (F.IV4) is precisely independent
similarity in the two matrix factors.  Indeed, in a product
of commuting block matrices, every contribution factoring
through a swap block is a norm.  The remaining diagonal
blocks are just the reductions of the trivial and sign
blocks of the centralizer matrix.  If that matrix is an
integral unit, its inverse shows that both reductions are
invertible.  Conversely block-diagonal unimodular changes
on the trivial and sign summands induce every element of
$\mathrm{GL}_p(\mathbb F_2)\times
\mathrm{GL}_q(\mathbb F_2)$; elementary matrices give
unimodular lifts.  There is no additional identification
between the two factors, since an intertwiner of a trivial
and a sign lattice is zero over $\mathbb Z$.

It remains a genuinely finite classification.  Enumerate
all $2^{n^2}$ matrices of size $n=0,1,2,3$ over $\mathbb F_2$
and conjugate by every invertible matrix of that size.
Their respective numbers of similarity classes are
$1,2,6,14$.  Taking the smallest bit-coded representative
in each orbit gives a fixed canonical choice.
The six integral involution types and their counts are
\[
\begin{array}{c|rrrrrr}
 (p,q,r)&(0,3,0)&(1,2,0)&(2,1,0)&(3,0,0)&(0,1,1)&(1,0,1)\\\hline
 \text{classes}&14&12&12&14&2&2
\end{array}
\]
and their sum is 56.  A representative is (F.IV2) with
$B$ equal to the chosen $0$--$1$ representatives in its
trivial and sign blocks, and zero in every other block.

This also gives an effective normalizer for an arbitrary
input, not merely a list of representatives.  Normalize its
cusp and then its involution by the preceding Smith and
glue constructions.  Use the finite-field orbit lookup to
normalize its two reductions.  In dimensions at most three
the binary lift of an invertible binary matrix is already
unimodular: its determinant is odd and has absolute value
at most two.  The latter bound follows by adjoining a row
and column of ones and applying the $4\times4$ Hadamard
bound.  Finally solve the full integral linear equation
$AD+DA=B-B_{\rm can}$ by Smith reduction.  The preceding
norm calculation proves solvability.  Conjugating by
$\begin{pmatrix}I&D\\0&I\end{pmatrix}$ subtracts this norm
from the upper-right block.  The product of these explicitly
recorded conjugators is integral unimodular and carries all
three marked matrices to their canonical representative.

Both matrices in (F.IV2) have a three-dimensional fixed
space.  Choose a primitive integral fixed vector $v_i$.
Then $g_i(x)=M_ix+v_i/4$ has actual order four:
$g_i^2$ is the nontrivial translation $v_i/2$.
This order-two subgroup is free, and every nontrivial
subgroup of a cyclic group of order four contains it, so
the full action is free.  For any other finite affine lift
the actual order is still even, since its linear part has
order two.  The two-point orbifold quotient of F.2.31 is
therefore the nontrivial cyclic group of order
$\gcd(m_1,m_2)$.  Simple connectivity is impossible for
every choice of twists and actual orders.
\end{proof}

The complete finite enumeration and arbitrary-input
normalizer are implemented in
\path{code/fourfolds_involution_classes.py}.
The program verifies the full norm-quotient Smith form in
each of the six integral types, all finite-field similarity
orbits, all 56 representatives, and a nontrivial global
integral conjugate of each representative.  It returns and
checks 112 unimodular conjugators.  It also checks the 112
chosen order-four affine lifts and computes the complete
fundamental-group structure and $H_1$ of all 56 fillings
using F.2.31.  The exact matrices, numerators, orbit tables
and group data are saved in
\path{records/fourfolds_involution_classes.json}.
These executed finite sets exhaust this involution stratum;
they are not a height truncation of the other finite-order
strata.

\paragraph{Theorem F.2.39 (Proved: the full effective classification and all marked gluings).}
The linear classification and all-twist topological computation
asked for in F.2 have the following complete effective answer.
For $k=2$ and $k=3$, every admissible ordered integral tuple
has the canonical representative of F.2.37, with an explicit
simultaneous unimodular conjugator.  Equality of these
representatives is necessary and sufficient for simultaneous
integral conjugacy.  There is no fixed global height cutoff;
the finite stopping bound for an input of height $H$ is
$(2H+1)^{36(k+1)}$ candidates before the proved rejection
filters.  The separate involution normal form has exactly
56 classes and is computed completely in F.2.38.

For every choice of actual free affine orders and integer
numerators, integral clutch vector, and specified primitive
Kuhn cusp marking, F.2.31 and F.2.35 compute the full
fundamental group and every integral homology group of the
gluing.  This is the fully marked correction to F.2.3;
peripheral matrices are not substituted for a choice of
filling.  The algorithms have no cutoff on affine orders,
twist entries or lattice indices.  The complete list of data
with $\pi_1=1$ and $H_2=H_3=0$ is the constructive
parametrization of F.2.30, with its exact lattice enumeration
and necessary-and-sufficient finite tests.  Each retained
filling has the homology and oriented smooth type stated
there.

\begin{proof}
The input is finite, so the prime-subgroup test in F.2.31
decides whether every claimed affine cyclic action is free
and has the stated actual order.  The integral cusp test and
its adapted basis are F.2.1.  The complete linear admissibility
filter and the canonical representative, including the proved
finite stopping bound, are F.2.2 and F.2.36--38.  The
characteristic-polynomial filters used by that algorithm are
conjugacy invariants, so none can remove a possible
representative.  The unbounded trace family (F.CN3) proves
that no fixed finite box could replace the input-dependent
bound.

We spell out how the linear classification is used with a
specified gluing.  If $R$ is its returned conjugator and
$P$ is the supplied cusp marking, change the data by
\[
 M_i'=RM_iR^{-1},\qquad v_i'=Rv_i,\qquad
 \kappa'=R\kappa,\qquad P'=RP.                           \tag{F.SY1}
\]
The actual orders are unchanged.  The coordinate map
$x\mapsto Rx$ on the torus intertwines the finite quotient
charts, since it sends $M_ix+v_i/m_i$ to
$M_i'x'+v_i'/m_i$.  At the cusp, its coordinates in the
marked Kuhn model are unchanged:
$(P')^{-1}x'=P^{-1}x$.  Hence it intertwines every toric
chart, translated face identification and collapse map of
that fixed filling.  It also carries the clutch path to
the path with vector $\kappa'$.  These maps agree on the
collars and give a diffeomorphism of the two fully marked
gluings.  Thus homology can be computed in the original
coordinates, while the canonical data and their group
presentation are also returned.  Independently replacing
$P'$ by another Smith marking is not this comparison.

An actual order-one site is ordinary, even if its integer
numerator has been recorded.  Delete it by the explicit
comparison of F.2.31: if it is the $i$th site, replace
$\kappa$ by
$\kappa+(M_{i+1}\cdots M_k)^{-1}v_i$ before deletion.
The cusp matrix and its supplied marking are unchanged.
There cannot be an order-one site in an admissible two-point
tuple: its other finite matrix would be the inverse of a
nontrivial unipotent matrix.  Nor can a three-point tuple
have two such sites.  Thus this reduction leaves precisely
the genuine two- or three-point inputs of F.2.35 and loses
no case from the prompt.

For each remaining input, the exact unsaturated normal
closure, cyclic or triangle-group extension and exceptional
spherical Euler relation of F.2.31 give the full fundamental
group.  The finite cap, complement and marked Kuhn chain
maps of F.2.33--35 give a finite integer Mayer--Vietoris cone
through degree nine.  Integer kernels and Smith reduction
therefore give all $H_0,\ldots,H_8$ with their complete
torsion factors.  This is a calculation for arbitrary
supplied twists, rather than a restriction to the affine
orders of one displayed example.  Its independent
abelianization check compares the resulting $H_1$ with the
van Kampen answer.

Finally F.2.7--30 prove both directions of the survivor
classification.  Genuine three-point orbifold quotients
exclude simple connectivity.  The two-point filters reduce
every survivor to the minimal order-$(5,4)$ spectrum; the
rational normal form, finite stable-lattice and extension
algorithms include every integral representation above it.
The universal affine normalization includes every twist,
and the six exact integral tests decide retention.
Conversely every retained datum satisfies all three
topological conditions by those tests.  F.2.30 computes
its common integral homology and applies the almost-complex
and smooth-classification arguments to identify its oriented
type.  There is no residual untested rational-height bound
in this parametrization.  The canonical-representative
algorithm can remove duplicates from it if desired, but
no finite global enumeration is asserted or needed for
these complete per-input decisions.
\end{proof}

The unified entry point is
\path{code/fourfolds_classify_gluing.py}.
It records the original and transported markings, the exact
linear normalizer, both group presentations and the complete
integer homology calculation.  Its general noninvolution
canonical branch has the proved finite stopping rule of
F.2.37 and the same stated practical limitation: that
astronomically large fallback is not claimed exhaustively
executed.  The executed finite catalogue, conjugacy tests,
full marked homology examples and their exact bounds are
those recorded in F.2.33--38.
The identity-insertion branch has additionally been run on
all 168 inputs obtained by placing an identity in each of
the three marked positions of each of the 56 involution
classes, after a nontrivial simultaneous conjugation.
Every exact representative and unimodular conjugator passes.
The full unified driver has been run on the orientation-reversing
involution gluing of F.2.35 and on the same gluing presented
with an additional actual order-one site.  In the latter
test its numerator is $(0,1,0,1,0,0)^T$ and the original
clutch is $(1,-1,0,1,0,0)^T$; the suffix comparison reduces
the clutch to zero.  Both executions return the complete
canonical data, transported cusp marking, exact group and
the full integral homology of that example, with the
independent abelianization check.  Their outputs and full
JSON certificates are retained.

\subsection{A geometric realization of the non-self-dual periods}

We now construct periods for the specific integral datum of
Theorems F.2.5--F.2.6.  The source below is a family of genus-two
curves, but the rank-six system is an extension of its primitive
exterior square.  It does not have the filtration excluded in
Theorem F.3.4.

\paragraph{Theorem F.3.5 (Proved: the geometric rank-five quotient).}
Put $\lambda_*=-256/3125$ and
$U=\mathbb P^1\setminus\{0,\infty,\lambda_*\}$.  The smooth projective
models of
\[
 C_\lambda:\quad y^2=x^5+\lambda^3x+\lambda^4
 \tag{F.7}
\]
give a polarized rank-four rational cohomology system $\mathcal H$ on
$U$.  Its primitive exterior square is rationally isomorphic to
the rank-five quotient of the system with matrices $A,B$ in
Theorem F.2.5 by its invariant first coordinate line.
On the orbifold $\mathcal B'$ obtained by taking roots of orders
five at infinity and eight at zero, the extended Hodge bundle is
\[
 \mathcal E=\mathcal T^3\oplus\mathcal T,\qquad
 \mathcal T=\mathcal O_{\mathcal B'}(2D_\infty-3D_0),\qquad
 \deg\mathcal T=\frac1{40}.
 \tag{F.8}
\]
Here $D_\infty,D_0$ are the reduced root divisors; the extension at
$\lambda_*$ is the nilpotent canonical extension.

\begin{proof}
The discriminant of the degree-five polynomial in (F.7) is exactly
$\lambda^{15}(3125\lambda+256)$.  At $\lambda_*$ its unique double
root is $r=64/625$ and the coefficient of $(x-r)^2$ is
$10r^3\ne0$.  Its other roots are simple.  Thus the fibre has one
ordinary node and irreducible genus-one normalization.  The local
model is a smoothing of a nonseparating node; the vanishing circle
is a primitive nonseparating homology class.  Cutting and regluing
an annular neighborhood proves that its monodromy is a single
Dehn twist, and hence a primitive rank-one transvection on $H^1$.
In particular its logarithm $N_H$ has square zero and rank one.

At zero, the substitutions $\lambda=s^8$, $x=s^6X$, $y=s^{15}Y$
give
\[
 Y^2=X^5+X+s^2.
\]
This is smooth at $s=0$.  The generator $s\mapsto\zeta_8s$ at
fixed $(x,y)$ acts by $(X,Y)\mapsto(iX,\zeta_8Y)$, so its
eigenvalues on the two holomorphic differentials are
$\zeta_8,\zeta_8^3$.  The full rational characteristic polynomial
is $g(T)=T^4+1$.
At infinity put $\lambda=s^{-5}$, $x=s^{-4}X$, $y=s^{-10}Y$.
The resulting smooth family is
\[
 Y^2=X^5+sX+1.
\]
The deck generator acts by $(X,Y)\mapsto(\zeta_5^4X,Y)$,
with holomorphic eigenvalues $\zeta_5^4,\zeta_5^3$.  Its full
characteristic polynomial is $f(T)=T^4+T^3+T^2+T+1$.
Inverses have the same two rational characteristic polynomials.

We identify the representation, rather than merely its spectra.
Choose positive peripheral generators $T_1,T_2$ at infinity and
zero, in that order, with $T_1T_2$ inverse to nodal monodromy.
Set $R=T_1^{-1}$ and $S=T_2$.  Then $R-S$ has rank one and their
characteristic polynomials are $f,g$.  The following elementary
rank-one lemma fixes their simultaneous rational conjugacy class.
If $u$ spans $\operatorname{im}(R-S)$, it is cyclic for $R$:
the span of its iterates is stable under both $R$ and $S$, and
the irreducibility of $f$ makes that span all of $\mathbb Q^4$.
In the rational basis $u,Ru,R^2u,R^3u$, the matrix of $R$ is
the companion matrix of $f$, and $S=R-u\varphi$ for a covector
$\varphi$.  The determinant identity
\[
 g(T)=f(T)+\varphi\operatorname{adj}(TI-R)u
 \tag{F.9}
\]
determines $\varphi$ uniquely: the four coordinates of
$\operatorname{adj}(TI-R)u$ are a basis of the polynomials of
degree at most three, as is seen from their successively distinct
leading degrees in the cyclic basis.  It follows that any two
such pairs are simultaneously rationally conjugate.  Apply the
lemma to the explicit companion pair
\[
 C_f=\begin{pmatrix}
 0&0&0&-1\\1&0&0&-1\\0&1&0&-1\\0&0&1&-1
 \end{pmatrix},\qquad
 C_g=\begin{pmatrix}
 0&0&0&-1\\1&0&0&0\\0&1&0&0\\0&0&1&0
 \end{pmatrix}.
 \tag{F.10}
\]
Their difference has rank one and their characteristic polynomials
are $f,g$.  Thus $(T_1,T_2)$ is rationally conjugate to
$(C_f^{-1},C_g)$, with the chosen peripheral conventions.

Use the lexicographic basis in $\bigwedge^2\mathbb Q^4$.
Its common invariant vector is $(1,1,1,1,1,1)^t$.
Let $Q$ be the identity matrix of size six with first column
replaced by this vector.  It is unimodular.  Direct multiplication
gives the exact identities
\[
 \left(Q^{-1}\bigwedge^2(C_f^{-1})Q\right)_{2:6,2:6}
   =A_{2:6,2:6},\qquad
 \left(Q^{-1}\bigwedge^2(C_g)Q\right)_{2:6,2:6}
   =B_{2:6,2:6}.
 \tag{F.11}
\]
They are checked over $\mathbb Z$ in
\texttt{fourfolds\_periods.py}, not inferred from eigenvalues.
The invariant line is the inverse of the polarization, up to a
nonzero rational scalar.  Projection along that line identifies
the quotient with the primitive exterior square $\mathcal W$.

The differentials $\omega_1=dx/y$ and $\omega_2=x\,dx/y$ span
the extended Hodge bundle everywhere.  At the node they are
independent dualizing differentials: $\omega_1$ has nonzero
opposite residues on the normalization and
$(x-r)\omega_1$ is a nonzero holomorphic differential there.
Here the extension assertion can be checked on the nodal
annulus, without invoking an additional extension theorem.
In coordinates $uv=t$ a relative dualizing differential has
a convergent Laurent expansion
$\sum_n a_n(t)u^n\,du/u$.  Integration on the vanishing circle
is $2\pi i a_0(t)$.  Integration across the annulus has the
single logarithmic term $a_0(t)\log t$; integration of every
nonconstant Laurent monomial and of the complementary fixed
arcs gives a holomorphic remainder.  Thus the period vector
has the form $\exp(sN_H)$ times a holomorphic vector,
$s=\log t/(2\pi i)$.  Untwisting by $\exp(-sN_H)$ is exactly
the nilpotent canonical extension used here.  The two independent
dualizing limits remain a basis of its Hodge subbundle.
Their isotropic pairing, zero away from zero, remains zero
at the limit by holomorphic continuation, so that limit is
Lagrangian.
At zero their expressions in the root coordinates are
$s^{-9}dX/Y$ and $s^{-3}X\,dX/Y$; at infinity they are
$s^6dX/Y$ and $s^2X\,dX/Y$.
Thus their divisor lines are
$\mathcal O(6D_\infty-9D_0)$ and
$\mathcal O(2D_\infty-3D_0)$, proving (F.8).
The nonzero residue of $\omega_1$ also shows
\[
 N_H\ell_\infty\ne0,
 \qquad \ell=\mathcal T^3\subset\mathcal E,
 \tag{F.12}
\]
where $\ell_\infty$ denotes its limit at the nodal cusp.
Indeed integration on the vanishing circle tends to
$2\pi i$ times that residue; the transvection sends the cohomology
class of a differential to that period times its vanishing
cohomology vector.
\end{proof}

\paragraph{Theorem F.3.6 (Proved: a global nonpolarized period family).}
The integral rank-six system with cohomology monodromies $A,B$
admits a holomorphic subbundle $\mathcal F$ of rank three over
$U$ satisfying $\mathcal F\cap\overline{\mathcal F}=0$.
It extends holomorphically on the finite root charts.  At the
cusp an integral marking puts its periods in the form
\[
 Z(t)=\frac{\log t}{2\pi i}I_3+C(t),
 \tag{F.13}
\]
where $C$ is holomorphic at zero and bounded there.

\begin{proof}
Write the exact sequence supplied by (F.11) as
\[
 0\longrightarrow\mathbb Q e\longrightarrow\mathcal V
   \longrightarrow\mathcal W\longrightarrow0.
 \tag{F.14}
\]
Here $e$ is the first coordinate vector.  The nilpotent canonical
extension at the cusp and the ordinary extensions on finite root
charts preserve this sequence.  This follows directly by using
the frames $\exp(-sN)$: each residue has only eigenvalue zero,
and the same exponential construction respects the inclusion and
the quotient.  At finite roots the statement is ordinary exactness
of bundles with trivialized local monodromy.  Denote these extended
bundles by $\mathcal H_D,\mathcal V_D,\mathcal W_D$.

Let $\omega\in\bigwedge^2\mathcal H_D$ be the real, flat
nondegenerate bivector inverse to the polarization, and let
$p:\bigwedge^2\mathcal H_D\to\mathcal W_D$ be primitive projection.
Define the holomorphic rank-three bundle
\[
 \mathcal P=p(\ell\wedge\mathcal H_D).
 \tag{F.15}
\]
Projection is injective on $\ell\wedge\mathcal H_D$: every
nonzero bivector there is decomposable whereas $\omega$ is not.
This argument applies also at the cusp.  The polarization gives
$\mathcal H_D/\mathcal E\simeq\mathcal E^*$, including there,
since $\mathcal E$ is Lagrangian.  Hence
\[
 0\longrightarrow\mathcal T^4\longrightarrow\mathcal P
 \longrightarrow\mathcal O\oplus\mathcal T^2\longrightarrow0.
 \tag{F.16}
\]
After dualizing, the end terms are
$\mathcal O\oplus\mathcal T^{-2}$ and $\mathcal T^{-4}$.

We compute the obstruction to lifting $\mathcal P$ in (F.14).
For the coarse projection $\rho:\mathcal B'\to\mathbb P^1$,
the elementary root-chart calculation of invariant sections gives
\[
 \rho_*\mathcal T^k=
 \mathcal O_{\mathbb P^1}
 \left(\left\lfloor\frac{2k}{5}\right\rfloor+
       \left\lfloor\frac{-3k}{8}\right\rfloor\right).
 \tag{F.17}
\]
Higher direct images vanish, because taking invariants under a
finite group in characteristic zero is exact.  For $k=0,-2,-4$
the displayed degrees are respectively $0,-1,-1$.  Their first
cohomology vanishes.  The dual of (F.16) therefore gives
$H^1(\mathcal B',\mathcal P^*)=0$.
The pullback of (F.14) to $\mathcal P$ consequently splits as
holomorphic bundles: choose a holomorphic lift
$\sigma:\mathcal P\to\mathcal V_D$ of its inclusion in
$\mathcal W_D$.  This is an actual extension-class calculation
on a compact orbifold, not an approximation or an Oka principle.
To specify the choice concretely, take a finite root-coordinate
cover with local bundle lifts $\sigma_i$ and write their
differences as $\xi_{ij}e$, where
$\xi_{ij}$ is a Cech cocycle in $\mathcal P^*$.
Choose any solution $b_i-b_j=\xi_{ij}$ of this explicitly
vanishing cohomology class and glue the maps
$\sigma_i-b_ie$.  Every such solution is allowed in the
construction that follows.

There is also a canonical holomorphic functional
$\alpha:\mathcal P\to\mathcal O$.
For $v\in\mathcal P$, its unique preimage in
$\ell\wedge\mathcal H_D$ is
\[
 j(v)=v+\alpha(v)\omega
 \tag{F.18}
\]
under the primitive decomposition.
At a smooth point, the positive Hodge metric implies that
$\ell$ and $\bar\ell$ form a nondegenerate symplectic pair.
Their common wedge line is $\ell\wedge\bar\ell$.
The bivector $\omega$ does not belong to
$\ell\wedge\mathcal H+\bar\ell\wedge\mathcal H$, since its
component on the two-dimensional symplectic orthogonal complement
is nonzero.  Projection $p$ therefore preserves the intersection,
and
\[
 \mathcal P\cap\bar{\mathcal P}
      =p(\ell\wedge\bar\ell).
 \tag{F.19}
\]
This is the complexification of a real line.  If $u$ spans
$\ell$, then $i u\wedge\bar u$ is real and its coefficient on
$\omega$ is nonzero, by positivity.  Thus $\alpha$ is real and
nonzero on that real intersection line.  There is a unique smooth
real section $n$ of it such that $\alpha(n)=1$.

Since the quotient of $\sigma(n)$ is real, there is a unique
real smooth function $h$ on $U$ such that
\[
 \sigma(n)-\overline{\sigma(n)}=2ih e.
 \tag{F.20}
\]
For $c\in\mathbb R$ set
\[
 \mathcal F_c=(\sigma+ic\alpha e)(\mathcal P).
 \tag{F.21}
\]
This is a holomorphic rank-three subbundle projecting
isomorphically to $\mathcal P$.  An intersection with its complex
conjugate must project to (F.19).  On its real generator the two
lifts differ by $2i(h+c)e$.  It follows, including arbitrary
complex multiples of that generator, that
\[
 \mathcal F_c\cap\overline{\mathcal F_c}=0
       \quad\Longleftrightarrow\quad h+c\ne0.
 \tag{F.22}
\]

We verify the required cusp property before choosing $c$.
In a complex symplectic basis $a,b,d_1,d_2$ of the limiting
$\mathcal H_D$, relation (F.12) permits us to arrange
$N_Hb=a$, $N_Ha=N_Hd_1=N_Hd_2=0$ and $\ell_\infty=\mathbb Cb$.
Then
\[
 \mathcal P_\infty=
 \langle k=p(b\wedge a),\ b\wedge d_1,\ b\wedge d_2\rangle,
 \quad
 N_W(b\wedge d_i)=a\wedge d_i,
 \quad N_Wk=0.
 \tag{F.23}
\]
The line generated by $k$ is
$\ker N_W/\operatorname{im}N_W$; all other kernel directions
come from the displayed image.  The exact matrices give
$N_V^2=0$, $\operatorname{rank}N_V=3$ and
$\operatorname{rank}N_W=2$.
Every lift of an image vector $a\wedge d_i$ is killed by $N_V$,
because that vector has a lift in $\operatorname{im}N_V$ and
$N_Ve=0$.  Rank three now forces $N_V$ to take any lift of $k$
to a nonzero multiple of $e$.  Therefore
$N_V:\mathcal F_{c,\infty}\to\operatorname{im}N_V$ is an
isomorphism for every $c$.  In particular its domain is a
complement to $\ker N_V$.

The primitive index-one cusp has a real integral basis
$a_i,b_i$ with $N_Vb_i=a_i$, by Proposition F.2.1, also for the
dual or inverse convention.  In this basis the untwisted
subbundle near zero is the graph $b+C_c(t)a$, with $C_c$
holomorphic and bounded.  Restoring the flat frame by
$\exp(sN_V)$, $s=\log t/(2\pi i)$, gives (F.13).
Thus $\operatorname{Im}Z=\operatorname{Im}s\,I_3+
\operatorname{Im}C_c$ is nonsingular on a sufficiently small
punctured disc.  Apply this first with $c=0$.
Equation (F.22) shows that $h$ is nonzero and has constant sign
on that connected punctured disc.

At the two finite root points, all ingredients in (F.19)--(F.20)
are smooth and $\alpha(n)=1$ remains valid, so $h$ extends and
is bounded.  On the compact complement of the cusp disc it has
a finite bound $K$.  Choose $c$ with the same sign as $h$ near
the cusp by the exact rule
$c=\operatorname{sign}(h(t_0))(K+2)$, where $t_0$ is any point
of that punctured disc and $K$ is the supremum of $|h|$ on
the compact complement.  Then $h+c$ is nonzero both on that
disc and on its complement.  Equations (F.21)--(F.22) prove
global real transversality, and the preceding cusp calculation
proves the bounded correction for this same $c$.

Finally, the fourth power of the order-eight action on
$\mathcal H$ is $-I$.  It acts trivially on
$\mathcal W,\mathcal P,\alpha$ and $\mathcal V$.
Hence these bundles, the lift and the constructed periods descend
from $\mathcal B'$ to the orbifold $\mathcal B$ with orders five
and four.  Indeed any equivariant bundle map between bundles with
trivial action of the extra order-two stabilizer descends on its
root chart.  This gives exactly the finite orders required for
the rank-six datum.
\end{proof}

\paragraph{Theorem F.3.7 (Proved: compact holomorphic filling).}
The smooth eight-manifold of Theorem F.2.6 carries an integrable
complex structure.  There is a proper holomorphic map to
$\mathbb P^1$ whose regular fibres are complex three-tori,
whose two finite reduced fibres are the free affine quotients
of orders five and four specified in Theorem F.2.5, and whose
remaining fibre is the index-one Kuhn degeneration.

\begin{proof}
We first fix the dual convention.  The preceding construction
has cohomology monodromies $A,B$; Theorem F.2.5 uses homology
monodromies $A^t,B^t$.  All three special values are real.
Choose a positive real basepoint and paths along the positive
real circle to infinity and to zero.  Complex conjugation fixes
these paths and reverses the respective peripheral meridians.
Pull back the preceding period family by conjugation of the
base and conjugate its complex fibres.  Its Hodge subbundle is
holomorphic, with local expression
$\overline{\mathcal F(\bar\lambda)}$, and its cohomology
monodromies are $A^{-1},B^{-1}$.  Its homology monodromies are
therefore exactly $A^t,B^t$.  The third monodromy is fixed by
the product relation, so this also fixes its cusp marking.
The primitive index-one condition and the bounded holomorphic
cusp correction are preserved, after replacing the real
angular basis by its negative if necessary.

For a holomorphic Hodge subbundle $\mathcal F$ with
$\mathcal F\cap\bar{\mathcal F}=0$, restriction of integral
homology functionals embeds the homology lattice discretely in
$\mathcal F^*$.  Real transversality makes its real span the
whole real vector space.  The quotient
$\mathcal F^*/\mathcal V_{\mathbb Z}^*$ is consequently a
holomorphic family of compact complex three-tori over $U$,
with its zero section.  This also proves local properness,
by a smoothly varying fundamental parallelepiped.

On the root disc at the finite point of order $m_i$, the family
extends smoothly and its deck map has linear part $M_i$.
The invariant vector $v_i$ of Theorem F.2.5 gives a holomorphic
torsion section $v_i/m_i$.  Use the diagonal affine deck action
\[
 (z,s)\longmapsto(M_i z+v_i/m_i,\zeta_{m_i}s).
 \tag{F.24}
\]
Here $M_i z$ denotes the induced complex-linear map
$\mathcal F_s^*\to\mathcal F_{\zeta_{m_i}s}^*$ from the marked
lattice action, not multiplication of a complex three-vector
by a six-by-six matrix.  Likewise $v_i/m_i$ means its evaluation
as a point of the target torus.
It has order $m_i$, since the $m_i$th translation is integral.
Every nonidentity power is fixed-point-free on the central
torus: the invariant integral covector $\ell$ takes $v_i$ to
$1$ or $-1$, so a fixed point would require the noninteger
$\pm j/m_i$ to vanish in $\mathbb R/\mathbb Z$.
Away from the central fibre the base rotation is already free.
The quotient is therefore smooth and proper over the coarse
disc, with precisely the required multiple fibre.

Its marked punctured family is holomorphically identified with
the original torus family by logarithmic translation.  On the
root cover use $r(s)=(\log s/(2\pi i))v_i$.  Then
$r(\zeta_{m_i}s)-M_i r(s)=v_i/m_i$, modulo the integral lattice;
the change of branch of $\log s$ is itself an integral vector.
Thus translation by $r$ conjugates the linear deck action to
(F.24).  This fixes the affine clutching, including its sign.

At the cusp use the integral marking of (F.13), the fan (E.4)
and the action (E.5) with $C_3=C$.
Its maximal charts are $\mathbb C^4$ with map
$t=z_0z_1z_2z_3$, and its transition maps are the toric changes
of coordinates multiplied by the holomorphic units
$\exp(2\pi i C(t)\lambda)$.
The proof of Theorem E.1.2 used only boundedness of $C$,
integrality of the standard cusp, and unimodularity of that
fan.  It therefore proves smoothness, proper discontinuity
and properness for this model without a symmetry condition
on $C$.  Glue by the displayed integral marking and use no
additional cusp translation.

The three punctured overlap maps are holomorphic isomorphisms
with the same marked torus family, so their cocycle condition
is literal.  Proper local models and a compact complementary
base give a Hausdorff compact total space.  Every chart is
smooth.  The affine vectors, cusp fan and clutching are exactly
those in Theorems F.2.5--F.2.6.  Varying their bounded corrections
and varying the fibre complex structures along a smooth path
does not change these marked smooth fillings: on the finite
root models the underlying real torus family is the fixed
local-system bundle; on the cusp use the stratum-preserving
parameter vector fields and uniform properness estimate of
Theorem E.1.3.  Thus the resulting underlying smooth manifold
is the one whose integral topology was computed there.
\end{proof}

\paragraph{Theorem F.3.8 (Proved: the unpolarized period domain has LSAP).}
Let $D_3$ be the component of
\[
 \{F\in\operatorname{Gr}(3,\mathbb C^6):F\cap\overline F=0\}
\]
containing the standard complex structure, so that
$D_3\cong GL^+(6,\mathbb R)/GL(3,\mathbb C)$.
At every point, $D_3$
is dominated by a spanning tube of nine complex lines.  In particular,
it has the local spray approximation property (LSAP) and is Oka--1.
These assertions concern the unpolarized domain.

\begin{proof}
Use the graph chart $F_Z=\operatorname{im}\binom Z{I_3}$.
The open set
\[
 \Omega=\{Z\in M_3(\mathbb C):\det(\operatorname{Im}Z)>0\}
 \tag{F.O1}
\]
lies in $D_3$ and contains $iI_3$.  For a real column $w$, write
$K_w x=w\times x$, put $q_w=w^tw$, and define
\[
 A_w=q_wI_3-ww^t+K_w.
 \tag{F.O2}
\]
Take the six vectors
\[
 e_1,\ e_2,\ e_3,\ e_1+e_2,\ e_1+e_3,\ e_2+e_3
\]
and the corresponding six matrices $A_w$, followed by the three
matrices $K_{e_1},K_{e_2},K_{e_3}$.  Call these nine matrices
$A_1,\ldots,A_9$.  Their symmetric parts are a basis of the
six-dimensional space of real symmetric matrices: the first three
give all diagonal matrices and the last three add the three distinct
off-diagonal entries.  The final three skew matrices form a basis
of the skew matrices.  Thus the nine directions are linearly
independent over $\mathbb R$ and also over $\mathbb C$.
For reference, the determinant of their nine row-major coordinate
columns is $-16$.

Each $A_w$ is real normal, is zero on $\mathbb Rw$, and has
eigenvalues $q_w\pm i\sqrt{q_w}$ on $w^\perp\otimes\mathbb C$.
For every real $s$ its two transverse squared singular values are
\[
 (1+q_ws)^2+q_ws^2
 =q_w(q_w+1)\left(s+\frac1{q_w+1}\right)^2
       +\frac1{q_w+1}\ge\frac13,
 \tag{F.O3}
\]
since $q_w\in\{1,2\}$; the remaining squared singular value is one.
For a pure skew direction $K_{e_j}$ all singular values of
$I_3+sK_{e_j}$ are at least one.  Their determinants are positive
at $s=0$ and never vanish, so they are positive for all real $s$.

Let $\Lambda=\bigcup_{j=1}^9(iI_3+\mathbb C A_j)$ and let
$B_{1/4}$ be the open Frobenius ball of radius $1/4$ in
$M_3(\mathbb C)$.  Then
\[
 T=\Lambda+B_{1/4}\subset\Omega.
 \tag{F.O4}
\]
Indeed, the imaginary part of $iI_3+tA_j+E$ is
$I_3+(\operatorname{Im}t)A_j+\operatorname{Im}E$.
Its unperturbed least singular value is at least $1/\sqrt3$,
whereas $\|\operatorname{Im}E\|_{\rm op}<1/4$.
It therefore stays invertible throughout the straight perturbation
from $E=0$, preserving the positive determinant.
The set $T$ is open and connected and is a union of affine
translates of the spanning tree $\Lambda$.  Thus it is a spanning
tube in the precise sense of \cite[Definition 2.1]{ffAlarconForstneric25},
not merely a neighborhood of one entire line.  Its inclusion in
$D_3$ is a local biholomorphism at $iI_3$.

The real group $GL^+(6,\mathbb R)$ acts holomorphically and
transitively on $D_3$.  Transporting this tube gives the same
domination at every point.  The observation preceding
\cite[Proposition 7.13]{ffAlarconForstneric25} gives LSAP from
this domination, and its part (b), or Theorem 2.2 of that paper,
gives Oka--1.  The exact program
\texttt{code/fourfolds\_oka\_tubes.py} checks all nine integer
matrices, their linear independence and the symbolic identity
(F.O3); the displayed uniform bound proves (F.O4) for all complex
parameters and is not a finite positivity test.
\end{proof}

\paragraph{Theorem F.3.9 (Proved: holomorphic periods on every open base).}
Let $B$ be an open connected Riemann surface and
$\rho:\pi_1(B)\to GL^+(6,\mathbb R)$ any representation.
Its associated flat $D_3$-bundle admits a holomorphic section.
More generally, every continuous section which is holomorphic
near a compact Runge set can be deformed through sections to a
holomorphic section, with arbitrary uniform approximation there
and any prescribed finite-order interpolation at finitely many
points of that set.  For an integral representation, this gives
a holomorphic family of complex three-tori with the prescribed
lattice local system and a zero section over $B$.

\begin{proof}
The flat bundle
$\mathscr D=\widetilde B\times_\rho D_3\to B$ is a holomorphic
fibre bundle: its transition maps are constant real linear
transformations, acting holomorphically on the Grassmannian.
Its fibre has LSAP by Theorem F.3.8, so the bundle map is an
Oka--1 map by \cite[Proposition 7.13(d)]{ffAlarconForstneric25}.
The lifting property in Definition 7.7 of that paper, applied
to the identity map of $B$, gives precisely the stated section
deformation and approximation, including the finite jets.
This uses the section theorem, rather than an inference from
the Oka--1 property for maps to the fibre alone.

There is a continuous section to start with.  An open Riemann
surface has the homotopy type of a one-dimensional CW complex;
a bundle with connected fibre admits a section over its vertices
and edges, and the homotopy lifting property transports this
section to $B$.  Connectedness holds for the chosen component
$D_3$.  Taking the Runge set empty therefore proves existence.

For an integral representation let $\mathbb V_\mathbb Z$ be its
lattice local system and let $\mathcal F\subset
\mathbb V_\mathbb Z\otimes\mathcal O_B$ be the resulting
holomorphic rank-three subbundle.  Real transversality says that
projection identifies the real vector bundle with the underlying
real bundle of
$\mathcal Q=(\mathbb V_\mathbb Z\otimes\mathcal O_B)/\mathcal F$.
Hence the image of $\mathbb V_\mathbb Z$ is a full lattice in
every fibre of $\mathcal Q$.  Locally it is uniformly discrete
and acts by holomorphic translations.  The quotient
$\mathcal Q/\mathbb V_\mathbb Z\to B$ is consequently a smooth
proper holomorphic family of compact complex three-tori,
with its evident zero section and the required marking.
\end{proof}

\paragraph{Scope of F.3.8--F.3.9.}
Taking $B=\mathbb P^1\setminus\{\hbox{singular values}\}$
settles existence of nondegenerate holomorphic periods on the
smooth locus for every oriented rank-six datum.  It does not
impose good reduction after the finite root changes or the bounded
correction to the logarithmic cusp period.  Those are conditions
at points missing from $B$, whereas the interpolation theorem
controls only finite jets at points lying in $B$.  Exact agreement
with a prescribed germ on an annulus cannot be requested either:
the identity theorem would then determine its analytic continuation.
Thus this theorem advances the unrestricted F.3 problem but does
not decide all its compactification conditions.

\paragraph{Lemma F.3.10 (Proved: retaining isolated extension points).}
Let $R$ be an open Riemann surface, $V\to R$ a holomorphic vector
bundle, and $\mathscr U\subset\operatorname{Gr}_r(V)$ an open
subset whose projection onto $R$ is surjective.  Let $S\subset R$
be finite.  Suppose that over $R\setminus S$ this projection is
a holomorphic fibre bundle with an LSAP fibre.  Let $K\subset R$
be a compact Runge set containing $S$ in its interior.
Every continuous section of $\mathscr U$ which is holomorphic
near $K$ can be deformed through sections to a holomorphic
section, with arbitrary uniform approximation on $K$ and
interpolation of any given finite jets at finitely many points
of its interior.  No fibre-bundle or Serre-fibration assertion
is required at the points of $S$.

\begin{proof}
Here are the details of the relative version of the LSAP argument
in \cite[Lemma 5.1, Theorem 6.1 and Proposition 7.13(d)]
{ffAlarconForstneric25}.  They also specify why isolated extension
points cause no new analytic step.  Enlarge $K$ slightly inside
the domain of holomorphicity, retaining the Runge property.
Choose a strongly subharmonic Morse exhaustion of $R$ starting
with this compact set.  All subsequent boundary arcs, attached
discs and critical one-handles can and will be chosen disjoint
from $S$.  On them $\mathscr U$ is an ordinary holomorphic
bundle with LSAP fibre.

The global spray on the old compact set can be taken to be
vertical even though the fibres at $S$ differ.  Indeed, let $f$
be the current holomorphic section on a Stein neighborhood $W$
of that set.  Put $A$ equal to the finite interpolation set,
including $S$, and choose an ideal sheaf $\mathcal I$ vanishing
to one more than the prescribed jet orders at its points.
Cartan's theorem A supplies finitely many holomorphic sections
$T_1,\ldots,T_N$ of
$\operatorname{End}(V)|_W\otimes\mathcal I$ which generate
$\operatorname{End}(V)$ on the compact attaching arcs, after
shrinking $W$ if needed.  The formula
\[
 F(z,a)=\exp\left(\sum_{j=1}^N a_jT_j(z)\right)f(z)
 \tag{F.O5}
\]
defines a holomorphic spray of sections of the Grassmannian
bundle.  Its differential in $a$ surjects onto the vertical
tangent space at every attaching point, since
$\operatorname{End}(V_z)\to
\operatorname{Hom}(f(z),V_z/f(z))$ is onto.
For $a$ sufficiently small it has range in $\mathscr U$ on
a neighborhood of the old compact set, by openness and
compactness.  Every parameter value preserves all the chosen
jets.  Domination at $A$ is neither asserted nor used.

For a noncritical extension, subdivide the new annular regions
into small discs lying in bundle trivializations over
$R\setminus S$.  Restrict (F.O5) to a thin attaching disc.
LSAP of the fibre extends this spray across the attached disc,
approximating the old spray as closely as desired on the
intersection.  The spray gluing used in the proof of
\cite[Lemma 5.1, equations (5.6)--(5.8)]
{ffAlarconForstneric25} now applies verbatim to vertical sprays:
the inverse-function step is needed only on the attaching
intersection, where (F.O5) is dominating.  It produces
reparametrizations of the two sprays and a glued section.
The old part remains in $\mathscr U$ because it is evaluated
in (F.O5); the new part remains in $\mathscr U$ because the
LSAP spray takes values in the bundle fibre.  The fixed jets
survive the reparametrizations, being independent of the
spray parameters on the old side.

For a critical extension, first extend over the attached real
arc using the given continuous section and the homotopy already
constructed on the old compact set.  Apply Mergelyan
approximation with the same finite jets to a neighborhood of
the union of the old set and this arc, and then apply the
noncritical step.  To see explicitly that this approximation
is by sections, trivialize $V$ holomorphically on the open
Riemann surface: its underlying complex bundle is trivial over
a graph spine, and the frame bundle therefore has a continuous
section.  Its fibre $GL(\operatorname{rank}V,\mathbb C)$ has
the dominating spray $(g,A)\mapsto g\exp A$, so
\cite[Theorem 1.2]{ffForstnericPrezelj02} gives a holomorphic
frame.  The section is then a map to the
fixed Grassmannian, to which
\cite[Theorem 2.13]{ffAlarconForstneric25} applies.  The base
coordinate is kept unchanged.  Close approximation on the
compact arc and old set, followed by shrinking its
neighborhood, preserves membership in the open subset
$\mathscr U$.  This is the critical case of Theorem 6.1 of
that paper, with no approximation performed at a singular
point of the projection: here the projection is a submersion
everywhere, as $\mathscr U$ is open in a Grassmannian bundle.

Choose successive errors summably small on the preceding
compact sets and smaller than their distances to the complement
of $\mathscr U$.  The sections converge uniformly on every
compact set to a holomorphic section with range in
$\mathscr U$.  The finite jets persist.  The close homotopies
on successive stages concatenate in the usual exhaustion
construction.  All new attachment regions lie outside $S$;
the old section remains holomorphic on a full neighborhood of
each point of $S$ at every stage.  This proves the lemma.
\end{proof}

\paragraph{Theorem F.3.11 (Proved: the bounded-cusp period decision).}
Let
$B^*=\mathbb P^1\setminus\{p_0,p_1,\ldots,p_k\}$,
where $k\ge1$, and let
$\rho:\pi_1(B^*)\to GL^+(6,\mathbb R)$ be a representation.
Assume that the monodromy $M_0$ at $p_0$ satisfies
$(M_0-I)^2=0$ and $\operatorname{rank}(M_0-I)=3$.
Then there is a holomorphic period family with this monodromy
and nondegenerate real transversality everywhere on $B^*$,
equivalently a nondegenerate imaginary part in each local
normalized period chart.
In a marking with $M_0=\left(\begin{smallmatrix}I&I\\0&I\end{smallmatrix}\right)$
its cusp period is
\[
 Z(t)=\frac{\log t}{2\pi i}I_3+C(t),
 \qquad C(t)\text{ holomorphic at }t=0.
 \tag{F.O6}
\]
Any prescribed finite jet of $C$ at zero may be imposed.
For an integral primitive cusp the marking can be integral.
Consequently every topological survivor of F.2 admits the
period family requested literally in F.3.  No good-reduction
condition at the other, still punctured, points is included
in this assertion.

\begin{proof}
Write $N=M_0-I$ and choose a real basis with
$N=\left(\begin{smallmatrix}0&I\\0&0\end{smallmatrix}\right)$.
Its integral version is Theorem F.2.1.  Choose the component
of the period domain selected by the positive logarithmic
model in (F.O6).  This choice matters only when an orientation
has been fixed independently of the marking; every element
of $GL^+(6,\mathbb R)$ preserves either component.

Retain the cusp as a point of the \emph{open} surface
\[
 B=\mathbb P^1\setminus\{p_1,\ldots,p_k\}.
\]
The flat complex vector bundle on $B^*$ has a holomorphic
extension $V_D$ across $p_0$: in a local coordinate $t$ use
the single-valued frame obtained by the change
$\exp(-sN)$, where $s=(\log t)/(2\pi i)$.
Changing the logarithm by $2\pi i$ is cancelled by $M_0$.
This specifies the extension directly, without a growth
assumption on a period section.

Over $B^*$ let $\mathscr U$ be the flat bundle of real-transverse
planes in the chosen component.  At $p_0$ adjoin only the
planes transverse to $\ker N_\mathbb C$, an affine space of
matrices $C\in M_3(\mathbb C)$.  This gives an \emph{open}
subset $\mathscr U\subset\operatorname{Gr}_3(V_D)$.
Indeed, in the chart $F_C=\operatorname{im}\binom C{I_3}$
the actual flat period is
$\exp(sN)F_C=F_{sI_3+C}$.  For $t\ne0$ membership is exactly
\[
 \det\bigl(y(t)I_3+\operatorname{Im}C\bigr)>0,
 \qquad y(t)=-\frac{\log|t|}{2\pi}.
 \tag{F.O7}
\]
For every fixed $C_0$, and uniformly for $C$ in a sufficiently
small bounded neighborhood of $C_0$, the matrix in (F.O7) is
invertible with positive determinant when $|t|$ is sufficiently
small.  This follows, for example, from
$y(t)>\|\operatorname{Im}C\|_{\rm op}$ and a straight-line
homotopy to $y(t)I_3$.  Thus a full neighborhood of
$(0,C_0)$ lies in the adjoined set, proving openness.
The projection is consequently a holomorphic submersion,
although its fibre at $p_0$ is $\mathbb C^9$ rather than $D_3$.

Start with a polynomial $C_0(t)$ realizing the desired finite
jet.  Formula (F.O7) makes it a holomorphic section on a small
closed cusp disc $K$, which is Runge in $B$.
This section extends continuously over $B$.  To justify the
relative extension, $D_3$ is connected and simply connected:
it retracts onto $SO(6)/U(3)$, and
$\pi_1(U(3))\to\pi_1(SO(6))$ is onto, since a unitary rotation
in one complex coordinate is a generator of the real
two-plane rotation loop.  A cellular decomposition of the
surface outside the cusp disc now extends its boundary
section over edges by connectedness and over two-cells by
simple connectivity.  The flat bundle is trivial over each
cell.  This argument uses no fibration property at $p_0$.

Over $B\setminus\{p_0\}$ the fibre is the LSAP manifold
$D_3$ from Theorem F.3.8.  Lemma F.3.10, with $S=\{p_0\}$,
produces a holomorphic section of $\mathscr U$ over all of
$B$, with the chosen jet at $p_0$.  Since its value at $p_0$
lies in the affine chart just adjoined, it has a holomorphic
matrix $C(t)$ there.  Undoing the canonical frame gives
exactly (F.O6), with a bounded holomorphic correction.
Everywhere else membership in $\mathscr U$ is precisely
nondegeneracy of the complex-torus periods.
More explicitly, choose three members of a real lattice basis
whose projected periods are complex-linearly independent, and
normalize them to $I_3$.  The remaining three periods form a
matrix $Z$, and real transversality is exactly
$\det\operatorname{Im}Z\ne0$.  This normalization persists
locally.  A single Schubert chart is not imposed globally,
and none is needed to define a marked torus period family.

For F.2 survivors, Theorem F.2.7 excludes $k=3$ and
Proposition F.2.15 proves orientation preservation of the
finite matrices for $k=2$.  The cusp determinant is also one.
Thus all survivors satisfy the hypotheses.  This is a
universal existence decision and does not require completion
of their simultaneous integral classification.
\end{proof}

\paragraph{Finite good reduction remains a distinct requirement.}
Theorem F.3.11 resolves the period-existence clause of F.3 as
written: its base is the punctured sphere, and its specified
end condition is the bounded cusp correction.  It does not
claim that every such family extends at the other punctures
after the finite root changes.  Requiring that stronger
simultaneous extension would retain every puncture and leave
a compact base, to which the Runge exhaustion used above does
not apply.  The specific compact fourfold in F.4 already has
all its finite extensions by Theorems F.3.5--F.3.7; its
existence does not depend on this stronger unresolved general
compactification question.

\subsection{The oriented smooth type}

\paragraph{Theorem F.4.1 (Proved: a complex structure on the stated connected sum).}
The compact complex fourfold of Theorem F.3.7 is three-connected,
has Euler number six, and is orientation-preservingly diffeomorphic to
\[
 3\mathbb H\mathbb P^2\ \#\ (-\mathbb H\mathbb P^2).
 \tag{F.25}
\]
Its integral middle lattice is $I_{3,1}$ and, in a diagonal integral
basis, $p_1/2=(1,1,1,1)$.  In particular the conclusion is an exact
smooth diffeomorphism, with no remaining homotopy-sphere ambiguity.

\begin{proof}
Theorem F.2.6 gives three-connectivity and free integral homology
with middle rank four.  The complex orientation and Theorem F.1.2
give signature two, the odd unimodular form $I_{3,1}$, and
$c_2^2=2$.  Since $H^2=0$, we have $c_1=0$ and
$p_1=-2c_2$.  The reduction of $c_2$ modulo two is $w_4$.
The Wu formula, with $w_1=w_2=0$, says that $w_4$ is characteristic
for the middle form.  Thus
$v=p_1/2=-c_2$ is characteristic and $v^2=2$.
It is primitive: a nontrivial integral divisor would have square
dividing two.  This excludes every divisor of absolute value
at least two.

We verify the orbit of this vector in $L=I_{3,1}$.
Its orthogonal complement $K=v^\perp$ is even, indefinite of
signature $(2,1)$, with determinant of absolute value two.
Unimodularity of $L$ and primitivity of $v$ give $w\in L$ with
$(w,v)=1$.  The vector $k=w-v/2$ represents the nonzero element
of $K^*/K\simeq\mathbb Z/2$.  Its discriminant square is
\[
 k^2=w^2-\frac12\equiv\frac12\pmod {2\mathbb Z},
 \tag{F.26}
\]
because characteristicity makes $w^2$ odd.
Consequently $K$ has the signature and discriminant form of
$U\oplus\langle2\rangle$, hence the same genus by
\cite[Corollary 1.9.4]{ffNikulin}.
Nikulin's Theorem 1.14.2 applies: the rank is three, the
odd-primary discriminant lengths are zero, and the two-primary
length is one, so its additional equality-case condition at two
does not arise.  The genus has a single isomorphism class.
Thus $K\simeq U\oplus\langle2\rangle$.
The extension from $\mathbb Zv\oplus K$ to the odd unimodular
lattice $L$ has index two.  Its gluing subgroup is the unique
nontrivial graph between the two discriminant groups of order
two.  Every isomorphism of complements therefore extends, taking
$v$ to a chosen characteristic vector of square two.
The vector $(1,1,1,1)$ is such a vector in $I_{3,1}$.
This proves the asserted orbit, including the class of $p_1$.

Let $M=3\mathbb H\mathbb P^2\#(-\mathbb H\mathbb P^2)$ with
the displayed orientation.  Its intersection form is $I_{3,1}$,
and its standard generators have $p_1=2\sum_i x_i$, as used
and derived in Theorem F.1.3.  Therefore the two manifolds have
isomorphic middle forms and first Pontryagin classes.
These determine Wall's full normal invariant in dimension eight.
Indeed the normal rank-four bundle of a representative embedded
$S^4$ has Euler number its self-intersection and Pontryagin
number the evaluation of $p_1(TM)$, since $p_1(TS^4)=0$.
The clutching group is
$\pi_3\operatorname{Spin}(4)=\mathbb Z\oplus\mathbb Z$;
with compatible generators $(a,b)$ its two characteristic
numbers are $e=a-b$ and $p_1=-2(a+b)$.
This pair is injective, so no additional unstable normal
invariant is omitted.

Wall's classification theorem in the precise formulation
\cite[Recollection 2.2, Definition 2.5 and Theorem 2.6]{ffBHS}
now gives an oriented diffeomorphism after connected sum with
some $\Sigma\in\Theta_8$.  We remove this ambiguity using the
exact inertia calculation
\cite[Theorem 1.4 and its proof in Section 9]{ffNagy2026}:
for a closed three-connected eight-manifold its inertia group
is zero exactly when $p_1$ is divisible by eight; otherwise,
since $\Theta_8=\mathbb Z/2$, it is all of $\Theta_8$.
Here $p_1$ has divisibility two.  Thus $M\#\Sigma$ is oriented
diffeomorphic to $M$ for both possibilities for $\Sigma$.
This proves (F.25).
\end{proof}

The application of Nagy's theorem uses the published 2026 paper
by Csaba Nagy, whose theorem and proof were opened in this run;
it is not an inference from the earlier announcement with
Crowley.  Mandal's almost-complex criterion alone gives no
integrability statement.  Here integrability comes from the
periods and the compact charts of Theorems F.3.5--F.3.7, while
the cited classification results identify their smooth manifold.


\paragraph{Conclusion F.2.4 (Proved: resolution of the classification problem).}
The previously incomplete linear classification and all-twist
integral topology computation are completed by F.2.31 and
F.2.35--39, with the fixed cusp filling and full marking
specified as required by F.2.3.  Every linear tuple has an
exact canonical representative and a returned unimodular
conjugator; the finite bound depends on the supplied input.
The unbounded trace family proves that no fixed global
height cutoff can replace this procedure.  All actual affine
orders, numerators, clutch vectors and cusp markings are
included in the general fundamental-group and integral
homology algorithms.  The complete constructive survivor
list, its integral tests and the common oriented smooth
type are given in F.2.30.

This conclusion does not claim an executed enumeration of
the infinite general class set.  The exhaustive finite
involution catalogue, the general comparison tests, the
three full topology integrations and the unified wrapper
checks have the exact executed scopes stated in F.2.35--39.
The earlier obstruction record is retained verbatim in
\path{records/fourfolds_F2_4_superseded.tex}; its finite-order
filter and homology-level obstructions were superseded by
the complete canonical and chain constructions.

\paragraph{Conclusion F.3.3 (Proved: the previously open period decision).}
Theorem F.3.11 now gives a positive answer for every topological
survivor of F.2: there is a holomorphic period family on the
punctured sphere, with nondegenerate imaginary part in every
local normalized period chart and a bounded holomorphic
correction at its primitive maximal cusp.  The proof is universal
and does not require that the remaining integral classification
be finite.  The proposed polarized abelian-surface source is
excluded by Theorem F.3.4; the universal construction instead
uses the explicit unpolarized LSAP domain and the relative
section principle in F.3.8--F.3.11.  Simultaneous good reduction
at every finite puncture is a stronger compactification condition
than the literal F.3 statement and is not inferred from this
argument.  It is already proved for the explicit compact
fourfold of F.4.
